% =====================================================================
%  Harald Høffding, DANSKE FILOSOFER.  Kjøbenhavn og Kristiania:
%  Gyldendalske Boghandel, Nordisk Forlag, 1909 (MDCCCCIX).  206 pp.
%  TRANSCRIPTION (Danish) of the 1909 print, complete.
%
%  SOURCE: the Royal Library's (KB) scan, ~/bibliotek/Høffding, Harald/
%  danske-filosoffer.pdf (SCANS.tsv).  PAGE MAP (pagemap.py): printed page =
%  PDF page - 15; p. 1 = PDF 16 ... p. 206 = PDF 221.  Front matter
%  (unnumbered): half-title PDF 10, title PDF 12, contents PDF 14; these
%  carry "% --- front matter ---" comments and no \opage.  Not
%  transcribed: blank leaves, the publisher's pages after p. 206, folios,
%  running heads, sheet signatures („Harald Høffding: Danske Filosofer. 1“),
%  library stamps, pencil marks and underlinings in this copy.  Every page
%  turn is marked "% --- p. N ---" + \opage{N} at the first word of the
%  printed page; a word divided over a page turn ends in %.
%
%  CONVENTIONS: letterspacing (the print's emphasis, chiefly names) ->
%  \emph; italic -> \textit; printed dashes -> ---, number ranges -> --;
%  quotation marks „…“ as printed.  Chapters begin on a new page under a
%  numbered head in capitals („I. LUDVIG HOLBERG“, then „2.“-„17.“ as printed); the Indledning has no
%  head.  Footnotes at the page foot, marked *), **) on each page;
%  \footnote[k] gives the note's rank on its printed page and a comment
%  "% footnote *) on p. N" precedes each.  A note that runs over to the
%  next page is not split.
%
%  METHOD (the e-book method, 2026-10-02; TRANSCRIPTION-PLAYBOOK §00 point
%  0).  No one typed the text, except the Indledning (pp. 1-2), which the
%  SAGA e-book lacks: it is the scan's text layer read against the image.
%  For pp. 3-206 the base text -- words, punctuation, paragraphs, italics,
%  the footnote texts -- is the SAGA e-book (Danske_filosoffer.epub), which
%  sets the 115 footnotes as endnotes; .parts/ebook/build.py aligned it
%  letter by letter with the scan's text layer, split per page by type
%  size into text and footnote lines, took every page turn from the layer
%  and every letterspaced word from the layer's spacing (names, verified
%  on a sample at the image; short words excluded).  Of 392 letter
%  disagreements 287 were settled by rule (the layer's glyph faults where
%  the e-book's word is attested elsewhere in the repository, letterspaced
%  capitals, specks, running heads, stamps, blocks the two streams place
%  differently); 105 letter, 91 punctuation and 20 paragraph disagreements
%  were decided at the image from crop sheets (decisions.tsv).  The print
%  corrects the e-book at: p. 14 Forgaaen; p. 17 Regering; p. 19 til;
%  p. 33 Fornuft; p. 45 Schein; p. 47 Oehlenschläger (twice); p. 61
%  nogetsomhelst; p. 78 sin; p. 79 til; p. 85 heldede; p. 97 forbigaaet;
%  p. 114 Bestaaen; p. 119 hans; p. 134 religionsfilosofiske; p. 137 für,
%  regnet; p. 178 dens; p. 193 externo; p. 90 n. den; and 27 places of
%  punctuation, chiefly commas the e-book adds („bøje, alt“, „sin,
%  frejdige“) or points it gives as commas.  PRINTED AS IS: p. 3 „at at“; p. 37
%  „Bamærkninger“ (corrected in pencil in this copy).
%
%  RESIDUAL RISKS: (1) a misprint that the e-book and the layer both
%  normalise is invisible to the method; (2) italics are the e-book's,
%  which the layer cannot witness; (3) letterspacing on words of fewer
%  than four letters is not recorded unless the word stands in a spaced
%  name; (4) p. 195, where the layer's line positions are unusable, had
%  no paragraph-indent check.
% =====================================================================
\documentclass[12pt,a4paper,openany]{book}
\usepackage[utf8]{inputenc}
\usepackage[T1]{fontenc}
\usepackage{textalpha}
\usepackage{libertinus}
\usepackage{libertinust1math}
\usepackage{graphicx}    % for \reflectbox (the old Danish ɔ: id-est mark)
\DeclareUnicodeCharacter{0254}{\reflectbox{c}}  % ɔ (U+0254) = reversed c
\usepackage[danish]{babel}
\usepackage[protrusion=true,expansion=true]{microtype}
\usepackage{setspace}
\usepackage[margin=1.2in]{geometry}
\usepackage{fancyhdr}
\setlength{\headheight}{28pt}
\addtolength{\topmargin}{-13.5pt}
\usepackage{hyperref}
\hypersetup{pdftitle={Danske Filosofer}, pdfauthor={Harald Høffding}, hidelinks}
\pagestyle{fancy}
\fancyhf{}
\fancyhead[LE,RO]{\thepage}
\fancyhead[RE]{\textit{Harald Høffding: Danske Filosofer}}
\fancyhead[LO]{\textit{1909}}
\setstretch{1.3}
\usepackage{marginnote}
\newcommand{\opage}[1]{\marginnote{\footnotesize\textit{[#1]}}}
% Footnotes as the print marks them, per page: *), **), ***), †).
\renewcommand{\thefootnote}{\ifcase\value{footnote}\or*)\or**)\or***)\or†)\fi}
\newcommand{\tocline}[3]{\noindent\makebox[2.2em][r]{#1}\hspace{0.6em}#2\dotfill\ #3\par}

\begin{document}

% --- front matter: half-title (PDF 10) ---
\begin{center}\vspace*{6em}DANSKE FILOSOFER\end{center}
\clearpage

% --- front matter: title (PDF 12) ---
\begin{center}
{\large HARALD HØFFDING}\\[6pt]
\rule{2cm}{0.4pt}\\[12pt]
{\Huge DANSKE}\\[6pt]
{\Huge FILOSOFER}\\[8em]
% (the publisher's device)
{\small GYLDENDALSKE BOGHANDEL}\\
{\small NORDISK FORLAG}\\[2pt]
\rule{0.6cm}{0.4pt}\\
{\scriptsize MDCCCCIX}
\end{center}
\clearpage

% --- front matter: contents (PDF 14) ---
\begin{center}{\large INDHOLD}\end{center}
\medskip
\tocline{}{Indledning}{1}
\tocline{1.}{Ludvig Holberg}{3}
\tocline{2.}{Sneedorf og Wolfianerne}{16}
\tocline{3.}{Danske Kantianere og deres Modstandere}{25}
\tocline{4.}{Henrik Steffens}{40}
\tocline{5.}{Hans Christian Ørsted}{49}
\tocline{6.}{Grundtvig og H. C. Ørsted}{57}
\tocline{7.}{Baggesen og Grundtvig}{66}
\tocline{8.}{Anders Sandøe Ørsted}{73}
\tocline{9.}{F. G. Howitz}{81}
\tocline{10.}{Niels Treschow}{89}
\tocline{11.}{F. C. Sibbern}{97}
\tocline{12.}{Poul Møller}{118}
\tocline{13.}{Heiberg og Martensen}{128}
\tocline{14.}{Søren Kierkegaard}{147}
\tocline{15.}{Frederik Dreier}{175}
\tocline{16.}{Rasmus Nielsen}{184}
\tocline{17.}{Hans Brøchner}{196}
\clearpage

% ===== INDLEDNING (pp. 1--2): not in the e-book; the scan's text layer, read at the image =====
% --- p. 1 ---
\setcounter{footnote}{0}%
\opage{1}ET lille Land som vort ligger aabent for Indflydelser udefra. De
aandelige Strømme fra de store Kulturlande gribe bestemmende ind i dets
Tankeliv, og dette kan ikke faa den faste Sammenhæng gennem Tiderne, som er
muligt, naar et Land kan frembringe en stor Mangfoldighed af tænkende Aander,
der aldrig lade Udviklingen stanse og ved de forskellige Typer, de frembyde,
kunne bryde Paavirkningerne udefra, saa at Sammenhængen i Folkets Tankeliv
bedre kan bevares.

Alligevel er der blevet tænkt alvorligt og selvstændigt herhjemme om de
Spørgsmaal, det er Filosofiens Kald at drøfte, --- om Videnskabens
Forudsætninger, om dens Resultaters Betydning for Livsanskuelsen, om
Grundlaget for etisk Vurdering, om den menneskelige Aands Natur og Love, om
Forholdet mellem Religion og Humanitet. De ydre Strømninger ere blevne
tilegnede paa ejendommelig Maade, og vor Nationalkarakter har lagt sig for
Dagen i den Maade, paa hvilken Problemerne ere blevne behandlede. Ved
Behandlingen af saadanne Spørgsmaal som de anførte, der ligge paa Grænserne
mellem den strenge Videnskab og Livet, maa Tænkernes Personligheder
nødvendigvis medvirke i højere Grad end ved andre Undersøgelser. Og hvor
Personligheden virker efter sin hele Ejendommelighed, gøre ogsaa de nationale
Egenskaber sig gældende. De ere jo ikke Lag eller Stykker, der udefra ere
føjede til de Enkeltes andre Egenskaber, men bestemme fra først af disse og
klinge med, hvergang Sindets Strenge tone.

At dansk Tænkning har sin Ejendommelighed, kan ses ved Sammenligning med
svensk Tænkning. Sveriges Aandsliv
% --- p. 2 ---
\setcounter{footnote}{0}%
\opage{2}har modtaget de samme Strømninger fra Udlandet som Danmarks, men de ere paa
Tankens Omraade blevne tilegnede og udviklede paa anden Maade og i anden
Retning end hos os. Spekulation og Mystik have spillet en langt større Rolle i
Sverige end i Danmark, og i Sammenhæng dermed har en Trang til systematisk
Afslutning gjort sig gældende, der aldrig har fundet fast Jordbund hos os. Den
danske Tænkning har væsentligt været baaret af psykologisk og etisk Interesse,
og der har atter og atter gjort sig en besindig Kritik gældende overfor
Systemdannelserne.%
% footnote *) on p. 2
\footnote[1]{Se min Afhandling om svensk og dansk Filosofi i Samlingen
\textit{Norden} (1902), saavel som min Afhandling \textit{Filosofien i Sverige}
(Det Letterstedtske Tidsskrift 1879). --- Om de filosofiske Bevægelser i
Udlandet, der faa Indflydelse paa dansk Filosofi, vil man finde nærmere
Oplysning i mit Værk \textit{Den nyere Filosofis Historie} ved Hjælp af
Registret. --- Biografier giver jeg som Regel ikke, især fordi jeg kan henvise
til de fortrinlige Artikler af min Kollega Professor \emph{Kroman} i
„Biografisk Leksikon“.} Hos vore ypperste tænkende Aander (Holberg,
Ørstederne, Sibbern, Kierkegaard) gør denne Ejendommelighed sig gældende, og
den afføder Tanker og Værker, der vel ikke have grebet ind i den store
internationale Diskussion om Erkendelsens og Livets Spørgsmaal, men som dog
have deres Værdi --- og det ikke blot for os. En Karakteristik af danske
Filosofer vil kunne lære os Mangt og Meget til Orientering overfor hine
Spørgsmaal, og ofte vil det vise sig, at vore Tænkere meget vel vilde have
kunnet gøre sig gældende paa en større Skueplads, om Forholdene havde anvist
dem en saadan. Tankens Himmel er, ligesom den fysiske Himmel, tilgængelig for
Alle, ikke blot for dem, der fødtes i et stort Folk. Denne Bog vil derfor kunne
kaste Lys over de filosofiske Problemer, ikke blot tjene til Karakteristik af
dansk Aandsliv saaledes som det ytrer sig paa Tænkningens Omraade.

\begin{center}\rule{2cm}{0.4pt}\end{center}
\clearpage

% ===== TEXT (pp. 3--206) =====
% --- p. 3 ---
\setcounter{footnote}{0}%
\opage{3}

\begin{center}\large I. LUDVIG HOLBERG\end{center}

Dansk Filosofi begynder med selve vor Literaturs egentlige Grundlægger. Han er
den Første herhjemme, der paa selvstændig Maade tager Stilling til Tankens
Problemer, samtidigt med, at han fører de Tanker og Løsningsforsøg, der vare
opstaaede i Udlandet, ind i dansk Aandsliv.

Holberg staar vel i første Linie for os som Digter, men det vilde være urigtigt
at se ham alene under dette Synspunkt og tyde hans Stemninger under hans Livs
Afvexlinger blot ud fra den Skæbne, hans Digtervirksomhed mødte. Forskningen
har en Hovedinteresse hos ham og udfyldte en væsentlig Del af hans Liv ligefra
hans unge Dage til hans høje Alderdom. Han mente som hans Erasmus Montanus, at
Studier er til Pryd i Lykkens Dage og til Trøst i Ulykkens. Intet forefaldt ham
saa tungt, at det ikke blev lettere ved Studeringer, --- har han selv sagt. Det
er karakteristisk, at en lignende Ytring tillægges Leibniz, der ogsaa havde en
alsidig Interesse for Læsning og Forsken. Men Leibniz havde et lysere og
lettere Blik paa Livet og Menneskene end Holberg, hvis tunge Natur ikke lod ham
nøjes med det glasklare Prisme, til hvilket Leibniz havde udslebet sin
Livsanskuelse. Han kunde ikke ved spekulative Betragtninger føres til at at %
% PRINTED AS IS: „at at“ (p. 3)
overse Livets Disharmonier eller at forsones med dem. For Holberg var det Onde
i Verden en Nød, der ikke var saa let at knække. Menneskelivet led særligt
under en stor Modsætning mellem Skin og Virkelighed. „Skyggen tages for
Legemet“ --- dette Udtryk kommer atter og atter igen hos Holberg, og han skrev
sine Komedier, sin „Niels Klim“ og sine „Moralske Tanker“
% --- p. 4 ---
\setcounter{footnote}{0}%
\opage{4}for at modarbejde denne stadige Forvexling. Han vil ikke prædike
Moral, men sønderrive Illusioner og vække til Selverkendelse.

Naar han saaledes tillægger sig et udtrykkeligt moralsk Formaal med alle sine
Værker, er dette kun at forstaa som en Tydning, han bagefter har anlagt paa sin
Forfattervirksomhed. Der har naturligvis rørt sig hos ham en uvilkaarlig Trang
til at udfolde sin Aandskraft paa Digtningens, Historieforskningens og
Tænkningens Omraade, og først da han i sin Alderdom saa tilbage paa sin rige og
mangeartede Virksomhed, har han i en moralsk Hensigt fundet en noget kunstig og
filistrøs Forklaring.%
% footnote *) on p. 4
\footnote[1]{Smlgn. min Afhandling om Holberg i \textit{Mindre Arbejder,} anden
Række p. 150---153.} Men hans Tankearbejde var en Side af hans Virksomhed, han
med særlig Glæde saa tilbage paa. Ligesom han mente at kunne „sige uden
Skrupler, at han næsten er den første Danske, der i Stedet for magre og konfuse
Krøniker har givet os nogenlunde vel indrettede Historier“, saaledes roste han
sig ogsaa af, at han havde „sat det moralske Studium istand igen, der lige til
denne Tid laa som begravet her i Norden“.

Men Holberg var ikke blot Moralfilosof. Han har sysselsat sig med den nye
Videnskab og den nye Verdensanskuelse, der havde udviklet sig paa Renaissancens
Tid, og hvis mest iøjnefaldende og tillige mest indgribende Resultater vare
Jordens forandrede Stilling i Verdensbilledet og Verdens Uafsluttelighed. Og
derhos var hans Tanke optagen af Verdens og Livets store Modsætninger, som
kulminerede i det Ondes Problem, der netop i Holbergs Ungdom drøftedes af saa
store Tænkere som Bayle og Leibniz.

Hans Tanke yndede at stille Modsætninger skarpt overfor hinanden, forat
Problemet ret kunde springe frem ved det stærke Sammenstød,--- Modsætningen f.
Ex. mellem den uendelige Horizont, Videnskaben havde aabnet, og Menneskelivet
med dets Stræben og dets Ønsker, --- Modsætningerne indenfor Menneskenaturen,
--- Brydningen af det Gode
% --- p. 5 ---
\setcounter{footnote}{0}%
\opage{5}og det Onde. Han yndede ved Behandlingen af den Art Spørgsmaal
„paradoxe Meninger“, hvormed han forstod saadanne, „som med gode Grunde
kuldkaste de Meninger, der have taget Borgerskab“. Det var ikke Hang til
Spidsfindighed eller til Tankeleg, der ledede ham herved: Paradoxerne skulde
have „Rimelighed“, d. v. s. være virkeligt begrundede, og det vare de, naar det
var udvidet Erfaring, der førte til dem. Han saa, at naar man stiller sig paa
Erfaringens brede Grund, --- naar man grunder sine Domme paa Erfaring, ---
kommer man til at tvivle om Meget, der hidtil stod som selvfølgeligt. Det er
intet Under, at man faar faste Meninger, naar man holder sig til de givne
Overleveringer og lader dem bestemme, hvad man erfarer, og hvorledes man tyder
det, man har erfaret.

Paa dette Punkt ligger især Holbergs Betydning som Tænker --- i Vægten paa saa
vid Erfaring som muligt og paa de Problemer, Erfaringen til hver given Tid
frembyder. Begge disse Tendenser gøre afsluttende Systematisering vanskelig, om
ikke umulig. Naar Holberg i Fortalen til \textit{Moralske Tanker} taler om at
gøre sit System færdigt,%
% footnote *) on p. 5
\footnote[1]{„Jeg anseer mig selv, i Henseende til min Alder og Svaghed, som en
reisefærdig Mand, der maa forfatte sit Systema, førend han tager Afskeed og
begiver sig paa Reisen.“} saa mente han derved sikkert kun en Afrunding eller
Afslutning for sit eget personlige Vedkommende. Selv den mest aabne og skarpe
Tænker stanser ved en Grænse, men han behøver ikke at betragte denne Grænse som
al Tænkens og Erfarings Grænse.

Hermed stemmer det, at naar Holberg skal sige, hvilke Filosofer, han skylder
mest, nævner han \emph{John} \emph{Locke}, Erfaringsfilosofiens Hovedmand, og
\emph{Pierre} \emph{Bayle}, den store Kritiker og Problematiker. Han skylder
idethele engelsk og fransk Aandsliv mest. Intetsteds befandt han sig paa sine
Rejser i Udlandet saa godt som i Oxford og Paris. Han glædede sig over, at han
var kommen til at studere ved de franske og de engelske Akademier og kaldte
særligt England „Filosofiens rette Skole“.

Naar Holberg efter sine Rejser, efter den Erfaring, han
% --- p. 6 ---
\setcounter{footnote}{0}%
\opage{6}havde indsamlet, og efterat hans Tanke var skærpet ved hans ihærdige
Syslen med Tidens Problemer forsøgte at filosofere i den danske studerende
Verden, rejste der sig fra forskellige Sider store Vanskeligheder for ham.

Det danske Sprog var kun lidet tildannet til videregaaende Tænkning; det havde
snevre Grænser for sin Udbredelse --- og indenfor disse Grænser var det ikke
agtet af sine Egne. „Mit Navn,“ siger Holberg, „kunde maaske have været
nogenledes bekendt udenlands, hvis jeg ikke havde skrevet i det danske Sprog,
der er saa indskrænket, at man ofte i Dannemark selv maa søge efter den, der
forstaar det.“ Des større Taknemmelighed maa vi føle overfor vor Literaturs
Grundlægger, at han ikke foragtede dette Sprog, men gjorde, hvad han kunde, for
at dyrke det op til at blive et Redskab for Tanken. Havde han for sin
Berømmelses Skyld skrevet paa et fremmed Sprog, vilde et dansk Tankeliv ikke
have udviklet sig. Hvad han begyndte, fortsattes af hans Efterfølgere, især af
Eilschow, Sneedorff og Kraft.

Men paa dette endnu lidet udviklede Sprog kunde man ikke udtale sig saa frit
som i de andre Lande, især i England. „Her i vort Norden,“ siger Holberg,
„plages vi med skarpe Censurer, hvorved al Munterhed og Geist dæmpes hos en
Skribent.“ Naar en Bog forsynes med „Imprimatur“, er det gerne et Tegn paa, at
der Intet er ved den. Det var især fra teologisk Side, at Censuren
inspireredes; det mærkede Holberg især ved Niels Klim. „Jeg tog mig især iagt
for, at Geistligheden ikke skulde faa Leilighed til at sætte sig i Harnisk mod
mig, siden jeg veed af Erfarenhed, at disse Folk ere mest uforsonlige.“ „Kun de
Kontraparter er jeg bange for, der iføre sig Religionens Harnisk.“ Og i sine
seneste Aar blev han endnu mere forsigtig: „Jeg sætter mine Ord og Skrifter paa
Skruer mere end tilforn, saasom jeg har bemærket, at de paa en Tid have været
nøjere examineret.“ --- Holbergs religiøse Meninger vare ikke ortodoxe. Derom
er der ingen Tvivl, skønt det kan være vanskeligt at afgøre, hvor Alvoren hører
op i hans Udtalelser, og hvor Forsigtigheden begynder. Man har overdrevet baade
Alvoren og Forsigtig%
% --- p. 7 ---
\setcounter{footnote}{0}%
\opage{7}heden, --- baade fundet mere Ortodoxi og mere Taktik hos ham, end der
er. Han har haft sin Teologi, som han ansaa for fornuftig og stemmende med hans
Erfaring. Den gaar ikke meget videre end Rousseaus og Kants, men han har ment,
at den ved ret Fortolkning kunde siges indeholdt i de bibelske og dogmatiske
Forestillinger.

De fleste Studerende (to Trediedele) ved Københavns Universitet paa Holbergs
Tid studerede Teologi, og den skolastiske Filosofi, som endnu dyrkedes
herhjemme, betragtedes som et vigtigt Middel til at værne den rette Lære. Ved
formel Logik skulde Tanken øves i Distinktioner og Syllogismer for at kunne
sætte Dogmatiken i System. I \textit{Erasmus Montanus} fremstiller Holberg en
ung dialektisk Fægter af den Art, de stadige Disputereøvelser maatte uddanne.
Hans formelle Virtuositet driver ham til at paatage sig at bevise hvadsomhelst
det skal være. „Derfor,“ siger Erasmus Montanus, „har jeg studeret, og derfor
har jeg lagt mit Hoved i Blød, at jeg kan sige, hvad jeg vil, og forsvare det
for hele Verden.“ I de Scener, hvor Erasmus disputerer med Folk paa Bjerget,
lader Holberg ham dog ikke bruge den Logik, der lærtes paa Universitetet.
Erasmus bevæger sig gerne frem gennem positive Slutninger i anden Figur,%
% footnote *) on p. 7
\footnote[1]{A er B og C er B, altsaa er A C! Nille kan ikke flyve, en Sten kan
ikke flyve, altsaa er Nille en Sten!} som Skolastiken meget vel vidste var
ugyldige. Men Holberg anvender saadanne Sofismer, som allerede Aristofanes og
Platon havde tillagt deres Modstandere. Istedetfor at vise det formelle
Pindehuggeris Unyttighed har han opnaaet en komisk Virkning gennem den Maade,
paa hvilken den Respekt, man paa Bjærget trods Alt havde for den københavnske
Lærdom, fik det rene Nonsens til at staa som bevist Sandhed.

Blandt de Emner, paa hvilke den logiske Skarpsindighed anvendes, nævner
Esrasmus først og fremmest de teologiske, f. Ex. det Spørgsmaal, om Englene ere
skabte før Menneskene. Men der nævnes ogsaa saadanne Spørgsmaal som det, om
Jorden er rund eller oval og om Solens og Stjernernes Afstand og Størrelse, og
det er paa et Spørgsmaal af denne
% --- p. 8 ---
\setcounter{footnote}{0}%
\opage{8}sidste Art, at Erasmus bliver fældet af den hjemlige Inkvisition. Det
er nu karakteristisk, at Spørgsmaal af denne sidste Art netop interesserede
Holberg selv i høj Grad og førte ham ind paa nogle af hans dristigste Tanker,
saaledes som det skal vises i det Følgende. Her havde han selv „Bjærget“ og „de
Lærde paa Bjerget“ imod sig, og han vilde være kommen i lignende Konflikt som
Erasmus Montanus, hvis han her aabant havde draget Konsekvenserne. Men i
Komedien var det ham Hovedsagen at stille nyttige Studier overfor
Spilfægteriet, og han har da været saa ubarmhjertig at lade Erasmus falde netop
der, hvor han var i sin Ret. Lieutenanten i femte Akt, Stykkets Moralist, har,
ligesom Holberg, især studeret „latinske Autores, Naturens Ret og moralske
Sager“, der af Montanus betragtes som Lapperi. Stykkets Moral er udtalt i
Lieutenantens Ord: „Det første Bud i Filosofien er at kende sig selv.....
Desmere vil man synes, der staar tilbage at lære. Men I gør Filosofi til en
Fægtekunst.“

Istedetfor de Spidsfindigheder eller de Grublerier over Guds og Englenes Natur,
man dengang kaldte Metafysik, vilde Holberg vende sig til Emner, der kunde give
positivt Udbytte for Filosoferingen. Istedetfor utallige uvisse Ting vilde han
have nogle faa visse. --- Lad os da fremdrage de Hovedspørgsmaal, der især
sysselsatte Holbergs Tænkning.

\begin{center}\rule{2cm}{0.4pt}\end{center}

Ofte erklærer Holberg i sin senere Tid, at Moralfilosofien er hans egentlige
Studium. Hovedværket er her \textit{Moralske Tanker} (1744). Det er ikke en
systematisk Bog, men en Række Smaaafhandlinger, som lejlighedsvis strejfe ind
paa mindre videnskabelige Emner, idet der ikke blot tales om Poesi og om
Studeringer, men ogsaa om Grundene til, at Forfatteren ikke har giftet sig, og
om det interessante Spørgsmaal, hvorvidt en slem Mave er ligesaa slem som en
slem Kone. Men den Hovedtanke, til hvilken der atter og atter vendes tilbage,
er, at der gives en Moral, som er grundet paa Naturens Lys. Ogsaa i sine
Epistler, den anden Hovedkilde
% --- p. 9 ---
\setcounter{footnote}{0}%
\opage{9}til hans Filosofi, hævder han denne Sætning. „Det naturlige Morale er
et Studium, det teologiske et andet. Naar man tager sig for at skrive udi det
første, udfordres, at man grunder Lærdommen paa Naturens Lys.“ Men den
naturlige Moral er ikke, som man paastod, „et hedensk Morale“. Ogsaa „Christi
Morale“ grunder sig paa Naturens Lys; Forskellen er blot, at samme „Morale“
udføres i større Fuldkommenhed.

Holbergs Etik er religiøs, forsaavidt som han forbinder den med „den naturlige
Religion“. Men han vil have Moralen stillet foran de specielle Dogmer og
Mysterier. „Man maa ved Christi Morale og en sund Filosofi lægge Grundvold, før
Christendommens Dogmata og Hemmeligheder læres.“ „Børn maa gøres til Mennesker,
før de blive Christne.“ En menneskelig Konfirmation maa gaa forud for den
kristelige; ti hvis En lærer Teologi, før han lærer at blive Menneske, bliver
han aldrig Menneske. „Man begynder med at raabe: Tro! før man viser, hvad man
bør tro og ikke tro. Det er at gaa til Værks ligesom hin Dommer, der begyndte
med Exekutionen og siden lod Sagen examinere.“ Man bør bibringe Menneskene en
kritisk Evne, at de ikke skulle tage Skyggen for Legemet.

For Holberg var Reformationens Grundvold „den kristelige Frihed i at ransage,
før man tror“. „Det staar i et Menneskes Magt at ransage, men ikke at tro.“
Vildfarelse efter nøje Examen er mere at undskylde end Ortodoxi uden
foregaaende Efterforskning. Strid er bedre end brutal Samdrægtighed. Holberg
selv læste hellere Kættere end Ortodoxe, fordi disse Sidste mere svare med
Skældsord end med Grunde. Og det var hans Overbevisning, at man maatte opgive
Troslærens Udenværker forat kunne forsvare Hovedfæstningen desbedre. For sit
Vedkommende var han villig til at underskrive alle fundamentale Troens
Artikler, der findes i Konfessionen; men han vilde dog hellere indskrænke end
forøge deres Antal, ti der behøves ikke nær saa mange Trosartikler, som
gemenligen foregives at være ganske nødvendige. --- Desværre siger Holberg
ikke, hvilke han vil stryge.

Han vil ikke have Undersøgelsen stanset derved, at man
% --- p. 10 ---
\setcounter{footnote}{0}%
\opage{10}tager sin Tilflugt til „Guds Motiver og forbeholdne Aarsager“ eller
„Fornuftens Ufuldkommenhed og Guds uransagelige Veje“. Det er Tegn paa
Dovenskab. En saadan Beraabelse er kun tilstedelig, naar man af sin yderste
Magt har stræbt at udfinde Aarsagen og søgt at forene de forekommende
Vanskeligheder med Fornuften.

Alligevel er der, som allerede antydet, i Holbergs moralske Katekismus
\textit{(Epist. 46)} baade et religiøst og et rationelt Element. Den er grundet
paa den naturlige Religion. De tre første Artikler handle om Gud og
Udødelighed. Den fjerde Artikel fordrer, at ingen Lærdom maa antages, før den
gøres bevislig; ti ellers kan man ligesaavel forfalde til falsk som til sand
Tro, og Fornuften, Guds ypperste Gave, vilde være til ingen Nytte. Den femte
Artikel fordrer Forkastelse af alle Lærdomme, der stride mod Guds hellige
Egenskaber, og den sjette lyder saaledes: Had og forfølg Ingen, der mod sin
Villie farer vild i Troen!

Af de to Elementer ligger det rationelle dybest, og logisk set burde den fjerde
Artikel være sat foran de tre forudgaaende. Meget tydelig er en Udtalelse fra
„Moralske Tanker“ om, at den rette Metode vilde være at forklare de Unge
Naturens Hovedbud, fremstille, hvad Billighed, Retfærdighed og Godhed ere, ---
at de ere Dyder hos Guds Skabninger og \textit{derfor} konsekvent maa være
Guddommens egne Egenskaber. --- Altsaa de Egenskaber, der tillægges Gud, maa
først kendes fra den menneskelige Verden. Det er da en Cirkelslutning, naar man
udleder Moralen af Teologien, begrunder den menneskelige Retfærdighed ved den
guddommelige. Det Etiske faar ifølge Holberg Prioriteten for det Religiøse;
dette forstaas ved hint, før det Omvendte kan blive Tilfældet.

„Naturens Lys“ finder Holberg aabenbaret i visse Grundsætninger, der for ham
ere umiddelbart indlysende. „Det store Naturens Bud“: Gør ikke mod Andre, hvad
du ikke vil, at de skulle gøre mod dig! --- vil, mener han, indrømmes af
Enhver, som hører det forklaret, --- selvom Locke har Ret i, at det ikke kan
siges at være medfødt. Det er lige%
% --- p. 11 ---
\setcounter{footnote}{0}%
\opage{11}saa klart som, at en Ting ikke paa én Gang kan være og ikke være. ---
Moralen har altsaa sin egen Kilde, selvom den staar i nøje Sammenhæng med den
naturlige Religion. Om hin Grundsætning virkeligt er selvindlysende, er et
andet Spørgsmaal. Det Problem, hvorledes vi komme til at opstille de
Grundsætninger, paa hvilke al etisk Bevisførelse beror, kommer Holberg ikke ind paa.

Hvad de enkelte Egenskaber og Handlinger angaar, indskærper Holberg, at al
Fuldkommenhed beror paa en bestemt Begrænsning: „Al Dyd bestaar udi
Mediocritet, og saasnart den gaar ud over Grænserne, bliver den metamorfoseret
til Last.“ „Det danske Ordsprog: Frygt Gud og følg Landevejen! er helt vel
grundet; ti de, som leder altid efter Stier og ville forkorte Vejene, komme
oftest senest frem. Og det er om saadanne konstige Vandringsmænd, man her udi
Landet siger, at de gaa Per Gantes Genvej.“ --- Desværre er der maaske noget
Dansk i disse Regler. Men man vilde dog grovt misforstaa Holberg, hvis man
troede, at han anbefaler en Middelvejsmoral. Det gør han ligesaalidt, som
Aristoteles gjorde det i sin Lære om Dyden som den rette Midte.%
% footnote *) on p. 11
\footnote[1]{Se herom min \textit{Etik} XI, 9 (kfr. IV, 2). (Tredie Udgave p.
221 fr. p. 89).} Hvad der er for meget eller for lidet, bestemmes ved den
Enkeltes Natur og kan være forskelligt fra Individ til Individ. Hvad der for
den Ene er en Sti, kan være en Landevej for den Anden. Det var ikke Holbergs
Mening at udelukke nye Veje. Hans Sammenligninger og Exempler vise klart, hvad
han mener. I Sygdoms Tilfælde skal en Patient have én Pille, en anden to eller
tre. En Cholericus maa dømmes anderledes end den, der er født med en
„temperered Complexion“. En, der ingen Galde har, kan ikke berømmes for
Koldsindighed. At ride paa et spagt Æsel er ingen Merite, men at regere en vild
og modig Hest fortjener Applausum og Hænders Klappen. Udi Misgerningers
Revselse bør Straffen rettes efter Misdæderens Konstitution; en Misgerning kan
være lige stor og dog ikke lige strafbar hos forskellige Personer.

% --- p. 12 ---
\setcounter{footnote}{0}%
\opage{12}Det er altsaa en Individualisering af Moralen, Holberg kræver. Der
skal ved etiske Vurderinger og Fordringer tages Hensyn til de indre personlige
Betingelser hos de Enkelte, og Alle kunne ikke skæres over én Kam. Vel uddyber
Holberg ikke det Problem, som her dukker op, navnligt hvorvidt Hensynet til de
enkelte Personers individuelle Dispositioner kan forenes med de Krav,
Samfundets Bestaaen stiller. Men Problemet om den Enkelte træder her for første
Gang op i dansk Tænkning. --- Det vil ved nøjere Undersøgelse vise sig at hænge
sammen med Spørgsmaalet om etiske Grundsætningers Bevislighed og Gennemførlighed.%
% footnote *) on p. 12
\footnote[1]{Smlgn. min Afhandling om \textit{Forholdsloven i Etiken.} (Etiske
Undersøgelser. 1891, p. 42---83).}

\begin{center}\rule{2cm}{0.4pt}\end{center}

Holberg har ikke blot følt Trang til at tænke over etiske Spørgsmaal. Han er
tillige den Første herhjemme, der har grublet over den Indflydelse, den af
Kopernikus, Galilei, Descartes og Newton grundlagte nye Verdensopfattelse
maatte faa paa selve Livsanskuelsen. Han betoner her den Udvidelse,
Verdensbilledet har faaet ved de nye Opdagelser og Synsmaader. Den nye
Astronomi har udvidet vore Forestillinger og givet os et større Begreb om Gud.
Mere end nogensinde staa nu Spekulationer om Gud som urimelige. Han erklærer
sig enig med Grækeren Simonides i, at jo mere man tænker over saadanne
Spørgsmaal, des vanskeligere blive de. Det nytter ikke med „meenige Theologi“
at sige, at Simonides taler ud af hedensk Blindhed, men at Aabenbaringen har
bragt Oplysning; ti de billedlige Udtryk, Bibelen overvejende bruger, giver
intet ret Begreb, ja lede til „grove og kiødelige Concepter“. Holberg finder,
at Naturkyndige have højere Tanker om det guddommelige Væsen end den gemeene
Almue, ja end mange Gejstlige selv, idet de ikke som disse gøre Forsøg paa at
udtrykke det Ubegribelige. Jo dybere vi kommer ind i Undersøgelse af, hvad
Evighed, Uendelighed, Alvidenhed, Allestedsnærværelse („Alværenhed“) og Almagt
er, des mere drukne vi i dem og se stedse mindre.

% --- p. 13 ---
\setcounter{footnote}{0}%
\opage{13}Efter det nye Verdensbillede staar derhos Jorden som et Sandskorn
overfor det uendelige Univers; den er ikke længere Midtpunktet i Verden, og
selv den Sol, den drejer sig om, er kun en Stjerne blandt utallige andre
Stjerner. Menneskene ere som Orme, der krybe om paa en liden Tue og dog bilde
sig ind, at Alt er til for deres Skyld. Ogsaa fra denne Side viser det sig,
hvor udsigtsløse metafysiske og teologiske Spekulationer ere.

I den naturvidenskabelige og den filosofiske Verden stod endnu paa Holbergs Tid
Striden mellem Descartes's og Newton's Lære om Himmellegemernes Bevægelser og
de Love, der ligge til Grund for dem. Descartes antog, at Universet bestod af
Systemer, der hvert for sig var i hvirvlende Bevægelse. De Kloder, der hørte
til samme Hvirvel, dreves rundt ved den omgivende Masses Bevægelse. At Kloderne
ikke i Kraft af Inertiens Lov fore ud i Universet bort fra deres Systems
Midtpunkt, forklaredes af, at den omgivende Masse, i hvilken de ligesom
svømmede, stedse førte dem rundt om Midtpunktet. Newton forkastede denne
Hvirvelteori og forklarede Klodernes Forbliven i deres Baner ved Tyngden eller
Tiltrækningen, der stadigt modvirkede Tendensen til at fare ud i en lige Linie.
Holberg synes at have været meget optagen af denne Diskussion. „Efter den
løselige Kundskab, som han derom haver“, holder han sig helst til den
kartesianske Opfattelse, fordi den forekommer ham simpel og anskuelig, hvorimod
Newton forekommer ham at beraabe sig paa en mystisk Kraft, naar han taler om
Tiltrækning.

Men naar Universet bestaar af mange Verdener i stadig Hvirvelbevægelse,
hvorledes rimer dette sig da med den bibelske Fortælling om Verdens Skabelse?
Dette Punkt havde i Holbergs yngre Dage voldet ham megen Tvivl og Uro. Han
lagde sig dog Sagen til Rette paa den Maade, at den mosaiske Skabelseshistorie
kun fortæller om Tilblivelsen af det Hvirvelsystem, til hvilket Jorden hører.
Der staar jo Intet om, at det var den første eller den eneste Verden, der
% --- p. 14 ---
\setcounter{footnote}{0}%
\opage{14}blev skabt! Maaske finder der en stadig Opstaaen og Forgaaen af
Verdener Sted!

Det er den dristigste Tanke, Holberg antyder. Og der er med denne Tanke givet
en stor Baggrund for Holbergs Verdensanskuelse, som hans egen Forsigtighed og
Tidens Smaalighed hindrede ham i at drage bestemtere frem, men som dog mærkes i
den Aandsfrihed, hans Tænkning idethele vidner om. Hans Forestillinger om de
afgørende Spørgsmaal have udfoldet sig under Indflydelse af en stor Horizont,
der vilde være traadt tydeligere frem, hvis han ikke selv for meget i Ordets
uheldigere Betydning havde frygtet Gud og fulgt Landevejen.

Endnu staar dog et Punkt tilbage som det, i hvilket Holberg paa den mest
afgørende Maade afveg fra den ortodokse Tankegang, --- nemlig hans Stilling til
det Ondes Problem, et Problem, der kort før hans Optræden var blevet drøftet af
Bayle og Leibniz. Pierre Bayle havde hævdet Umuligheden at forklare det Ondes
Opstaaen, naar man gik ud fra Forestillingen om en algod og almægtig Gud som
Ophav til Alt, --- og erklæret, at den eneste konsekvente Antagelse, naar man
stillede sig paa den naturlige Fornufts Standpunkt med Bortseen fra
Aabenbaringen, var Forudsætningen om to oprindelige Principer, et godt og et
ondt, i stadig Brydning og Kamp, saaledes som Manikæerne havde lært. Holberg
kunde dog ligesaalidt gaa ind paa denne dualistiske Opfattelse, som han kunde
følge Leibniz i hans optimistiske Spekulationer. Men en af hans Venner, som han
betegner ved B. S. (man formoder, det er Rektor Bernhard Schnabel i Roskilde),
havde vist ham en Udvej: dersom Gud fra Begyndelsen af staar overfor en kaotisk
Materie, der ikke fuldkomment er modtagelig for de Former, han vil indpræge i
den, men gør en vis Modstand, forstaas det, at ikke blot det Fuldkomne
existerer, men ogsaa det Ufuldkomne. Man undgaar da den manikæiske Antagelse af
to aktive Principer, en god og en ond Gud, og fastholder Troen paa en god Gud,
hvis Hensigters Udførelse hæmmes ved en stridig Materie. En saadan Synsmaade
(af lignende Art som den, Rousseau
% --- p. 15 ---
\setcounter{footnote}{0}%
\opage{15}og Voltaire, og senere Stuart Mill og William James havnede i) vilde
--- mener Holberg --- have afvæbnet Bayle, hvis ikke Aabenbaringen havde lært
en Skabelse af Intet. Paa dette Punkt er hans Akkomodation og Skinunderkastelse
overfor Kirkelæren aldeles tydelig. Med hans Lære om vort Verdenssystems
Opstaaen som Afløser af et tidligere, vilde den angivne Løsning af det Ondes
Problem have passet fortrinligt.

Iøvrigt har Holberg ogsaa et andet Argument mod Bayle. De forskellige Væsener
maa maales med forskellig Maalestok. Et Menneskes Fuldkommenhed er en anden end
en Flues og end en Engels; ethvert Væsen er fuldkomment i sin Art. Denne
Tankegang stemmer godt med Holbergs Individualiseren af Etiken, og den
forudsætter, at Maalestokken hentes fra selve de enkelte Væseners Natur hvert
for sig. Den stiller et lignende Problem som Etikens Individualisering, og i
skarpere Form.

Holberg har ligesaalidt her som paa flere andre Punkter udførligt og konsekvent
udviklet sine Tanker. Men han har gjort et grundlæggende Arbejde og gjort fri
Tankebevægelse mulig herhjemme. For første Gang drøftedes i hans Skrifter store
Problemer paa Dansk og i dansk Aand.

\begin{center}\rule{2cm}{0.4pt}\end{center}

% --- p. 16 ---
\setcounter{footnote}{0}%
\opage{16}

\begin{center}\large 2. SNEEDORFF OG WOLFIANERNE\end{center}

Holbergs Skikkelse fylder saa meget i det danske Aandsliv i det attende
Aarhundrede, at der synes at være et tomt Rum efter ham. I selvstændig Tænkning
kom Ingen i Slutningen af Aarhundredet ham nær. Men han havde vakt Lyst til at
tænke og vist Muligheden af at benytte det danske Sprog dertil. Hos de Bedste
blandt de Yngre rørte der sig en Skamfølelse over, at der endnu ikke existerede
en dansk videnskabelig Literatur, medens andre europæiske Folk forlængst vare
begyndte at bruge deres Modersmaal i Videnskaben.

Den, der som Tænker nærmest virkede i Holbergs Aand, var \emph{Jens}
\emph{Schielderup} \emph{Sneedorff}. Ligesom Holberg var han paavirket af
engelsk og fransk Literatur. Især Shaftesbury og Montesquieu have øvet
Indflydelse paa ham. Han var født 1724, studerede i Göttingen, hvis Universitet
blev besøgt af mange Udlændinge, og hvor Filosofi, Statsvidenskab og
Naturvidenskab blomstrede. Senere virkede han (1751---1761) som Professor i
Rets- og Statslære ved det nyoprettede Akademi i Sorø, der i sin korte
Blomstringstid blev Vuggen for dansk videnskabelig Literatur. Han døde allerede 1764.

Sneedorff priser Holbergs store Fortjenester af vort Sprog og vor Literatur.
„Men vi vare ikke værd at have en saadan Medborger, dersom vi af Dovenskab
eller Nederdrægtighed vilde tro, at der var Intet for os at gøre.“ Kampen for
Modersmaalet og for sand, naturlig Tænken gaar Haand i Haand hos Sneedorff.
Bort med barbariske Kunstord og aka%
% --- p. 17 ---
\setcounter{footnote}{0}%
\opage{17}demisk Spilfægteri, --- med død Hukommelse og død Abstraktion!

Ifølge Sproghistorikerne blev Grundvolden til det nyere danske Sprog lagt af
Sneedorff og hans Samtidige. ---

Sneedorffs Tænkning angaar dels Rets- og Statsforhold, dels almindelige
filosofiske Synspunkter.

I første Henseende er Hovedskriftet \textit{Den borgerlige Regering} (1757).
Vægten lægges her paa, at en god Regering maa følge Grundsætninger, der bygge
paa Erfaring og stemme med Folkets Natur; kun en saadan Regering kan vinde
Folkets Tillid. Skrevne Love, ikke mindst Grundlove, ere unødvendige, ja kunne
være til Skade paa Grund af deres Uforanderlighed, der hindrer den stadige
Vexelvirkning med Erfaringen. Den danske Enevælde var i Sneedorffs Øjne det
eneste absolute Monarki, der var grundet paa Folkets Tillid, og han protesterer
derfor imod, at den skulde kunne kaldes Despoti. --- Men forat Erfaringer kan
gøres, og forat Folkets Natur kan lægge sig for Dagen, maa den offentlige
Mening udfolde sig frit, og dertil hører igen Oplysning saavel som Frihed i
alle borgerlige Forhold for alle Stænder, ikke mindst for den hidtil saa
tilsidesatte Bondestand. Friheden gør Mennesker virksomme og fornuftige, og
uden Frihed ingen Ære! ---

Enevælden anerkendes her, fordi den er oplyst og folkelig, ikke længere, fordi
den er af Guds Naade. En Vending i Tankegang og Begrundelse, der maatte faa
videregaaende Konsekvenser!

Sneedorff var en af dem, der paa de almindelige Ideers Omraade virkede til at
forberede de store Landboreformer i Danmark i Slutningen af det attende
Aarhundrede. Medens Holberg endnu havde ment, at Bonden maa holdes til Arbejde
ved tvingende Love, hævdede Sneedorff og hans Meningsfæller, at baade Hensynet
til den personlige Frihed og til det almene Vel, --- Hensyn, der i Grunden
falde sammen, --- krævede Forandringer i Bøndernes Stilling.%
% footnote *) on p. 17
\footnote[1]{Smlgn. E. \emph{Holm}: \textit{Om det Syn paa Kongemagt, Folk og
borgerlig Frihed, der udviklede sig i den dansk-norske Stat i Midten af det
attende Aarhundrede.} (Universitetsprogram 1883). p. 107---114.}

% --- p. 18 ---
\setcounter{footnote}{0}%
\opage{18}Sneedorffs almindelige filosofiske Standpunkt lægger sig allerede for
Dagen i hans Rets- og Statslære, men han udvikler den nærmere, tildels i
polemiske Former, i sit Tidsskrift \textit{Den patriotiske Tilskuer} (især i
første Bind 1761).

Hvor det gælder Erfaring og Beskrivelse af Erfaring, eller hvor det gælder at
finde Udtryk for Følelseslivet, der kan man ifølge Sneedorff ikke anvende
strenge Definitioner og Bevismetoder. Det kommer væsentligt an paa at skaffe
sig saa fyldigt et Indhold som muligt, og en logisk eller matematisk
Formulering vilde her kun virke indsnevrende, ja, i Grunden bero paa en
Illusion. Paa Erfaringens Grund naa vi kun til Sandsynligheder
(„Rimeligheder“), ikke til Nødvendigheder. Og vi maa være varsomme med at finde
Modsigelser, ti disse ville kunne falde bort med udvidet Erfaring, der kan lære
os, at der gives langt flere Muligheder, end vi tænkte os.

Det er karakteristisk, at Sneedorff i sin Omtale af Erfaringens Forhold til
logiske Modsigelser fremhæver den modsatte Side af den, Holberg dvælede ved.
Holberg lagde Vægt paa, at jo større Erfaring vi faa, des flere Modsigelser og
Vanskeligheder ville vi finde. Det var hans Sag at finde og stille Problemer.
Sneedorff dvæler især ved udvidet Erfarings Virken til at faa de foreløbigt
fundne Modsigelser til at falde bort. Det hænger sammen med den lysere Tone,
der er over hans Filosoferen, og med hans Haab til Fremskridtet, men ogsaa med
en Svækkelse i Sansen for Problemernes --- de stadigt tilbagevendende
Problemers --- Skarphed.

Men dogmatisk er Sneedorff ikke. Som han har et klart Blik for det Hypotetiske
i de Sætninger, vi uddrage af Erfaringen, Hovedkilden til vor Erkendelse,
saaledes har han ogsaa Øje for den Rolle, Analogien spiller paa Tankens
Grænseomraade, ja, allerede, naar vi ville slutte os til, hvad der foregaar i
andre Væseners Indre. Og han hævder, at kun de enkelte Ting ere virkelige,
medens Arter og Slægtsbetegnelser skyldes vor Tænkning. Der ytrer sig i denne sidste
% --- p. 19 ---
\setcounter{footnote}{0}%
\opage{19}Sætning en Sans for det Individuelle, som vi allerede have truffet
hos Holberg og ville genfinde hos senere danske Tænkere (Treschow, A. S.
Ørsted, Kierkegaard). Sneedorff protesterer mod, at Filosofien ene skulde være
Metafysik og Logik. Studiet af Sjælens Virkninger og Bevægelser kræver Finhed
og Højhed i Tanken ikke mindre end den rene Metafysik, og Retslære og Æstetik
kræve nu deres Ret, efterat Teologi, Filologi, Matematik og Metafysik saa længe
have ført Ordet.

\begin{center}\rule{2cm}{0.4pt}\end{center}

Holberg og Sneedorff høre til den Art Tænkere, som bestræbe sig for at bygge
paa Erfaring, med fuld Bevidsthed om, at man ad denne Vej ingen absolut
Afslutning kan naa. De træde op imod den Tænkning, der gaar ud fra visse
dogmatisk opstillede Principer, ud fra hvilke man saa bevæger sig frem gennem
strenge Slutninger. Erfaringen faar her en underordnet Betydning. For en stor
Del rettes Sneedorffs Udviklinger imod \emph{Jens} \emph{Kraft} (1720---1765),
Professor i Matematik i Sorø og Forfatter til en hel Række filosofiske
Lærebøger. Kraft havde i sine \textit{Kritiske Breve til Videnskabernes
Fremvæxt og Smagens Forbedring} (1761) forsvaret den matematiske (deduktive)
Metode i Videnskaben. Han fremhævede den abstrakte Tænknings og Matematikens
Uundværlighed, dersom Erkendelsen skal være grundig og dersom den skal anvendes
i Praxis, og han frygter for, at „de skønne Videnskabers“ Overvægt vil føre til
Magelighed og overfladisk Aandrighed. Paa Sneedorffs Fremhæven af
Erfaringsmetoden gaar han derimod ikke ind.

Nogle Aar iforvejen (1751---1753) havde Kraft udgivet en systematisk
Fremstilling af Metafysiken efter den tyske Filosof Wolffs Metode. Ved
Metafysik forstaar han „Læren om de første og nødvendigste Grundlærdomme, som
behøves til den menneskelige Kundskab“. Den første Del er Ontologien, Læren om
de alleralmindeligste Sandheder, der kunne reduceres til to: at det
Selvmodsigende er umuligt, og
% --- p. 20 ---
\setcounter{footnote}{0}%
\opage{20}at enhver Ting har sin Aarsag. Den anden Del er Kosmologien, Læren
„om de naturlige Tings Kræfter og første fysiske Beskaffenheder“, Nøglen til
Fysiken. Han skiller sig her fra Wolff ved at antage Newtons Hypotese og
forkaste Descartes', Bogen betegner derved et Vendepunkt i dansk Videnskab. Den
tredie Del er Psykologien. Her skelnede Wolff mellem den rationelle
(spekulative) og den experimentale (empiriske) Psykologi. Kraft sammensmelter
dem begge, idet han bemærker: „Jeg har ikke fundet iblandt alle de Theoretiske
Decouverter om Sjælen uden faa af Vigtighed, hvorover jeg har holdt for, at man
med Føye og Ret kunde føre dem ind med det Experimentale i et samlet Værk.“
Kraft har her altsaa selv villet lade Erfaringen komme til sin Ret. Den fjerde
og sidste Del er „Natur-Lærdommen om Gud, eller den naturlige Theologie.“ Her
bevises Guds Existens dels ved, at der maa være en første Begyndelse til
Forandringerne i Verden, da disse ellers ikke vilde have deres tilstrækkelige
Grund, og Spørgsmaalet om Oprindelsen ellers stedse vilde komme igen uden at
blive besvaret, dels ved, at det ligger i Begrebet om Gud som det uendelige og
fuldkomne Væsen, at han maa existere, da han har sin Grund i sig selv, ikke er
afhængig af noget Andet, der kunde udelukke hans Existens.

Krafts Værk er det eneste fuldendte System, som er skrevet paa Dansk. Det
bygger, afset fra de enkelte fremhævede Punkter, ganske paa Wolffs Filosofi,
det første tyske System. Chr. Wolff var en Discipel af Leibniz, og hans
Filosofi er paa én Gang en Popularisering og en Systematiseren af Mesterens
Tanker. Han har fortrængt den gamle Skolastik paa de tyske Universiteter,
ligesom Descartes havde fortrængt den i Frankrig og Holland, og Locke i
England. Men hans teologiske Kolleger i Halle frygtede Konsekvenserne af den
nye Filosofi, dels fordi den med saa stor Iver hævdede „den tilstrækkelige
Grunds Princip“, som syntes dem at føre til Fatalisme, dels fordi den lærte en
rent menneskelig Morals Mulighed. I Danmark var den ogsaa ilde set, medens
Pietismen herskede. En af de første Danske, der paavirkedes af den, Henrik
Stampe, skal af den Grund have bot „halvt in%
% --- p. 21 ---
\setcounter{footnote}{0}%
\opage{21}cognito, idetmindste for Hofpræst Bluhme“ i Wolffs Hus. Wolfianismen
blev først indført ved Universitetet i København af Jurister (Stampe, Kofod
Anker) og Teologer (Gunnerus). Men den første rent filosofiske Forfatter
herhjemme af denne Retning var \emph{Frederik} \emph{Christian} \emph{Eilschow}
(1725---1750).

Ligesom Wolff i Tyskland skabte en filosofisk Terminologi, der skulde erstatte
den latinske, og ogsaa for en Del gjorde det, saaledes var det ogsaa en af
Eilschows Bestræbelser at erstatte de mange fremmede videnskabelige Ord med
danske, og disse maatte da først laves. Eilschow har ofte været heldig med sine
Forslag i denne Henseende. Ord som Enkelthed, Sandsynlighed, Sædelighed,
Personlighed, Selvbevidsthed, Sindsbevægelse, Overflade skyldes ham. Han
hævder, at Videnskaberne bør fremstilles i Folkets eget Sprog. Overfor dem, der
mente, at der ingen Grund var til at forlade Latinen, da praktiske Folk dog
ikke havde Tid til Læsning af videnskabelig Art, beraaber han sig paa sin egen
Erfaring. Der er Mange, som uden at høre til de Lærde, købe og læse
videnskabelige Bøger, hvis Sprog de forstaa; mange Timer, ja hele Nætter
anvende de paa at fordybe sig i dem.%
% footnote *) on p. 21
\footnote[1]{\emph{Eilschow}: \textit{Cogitationes de scientiis vernacula
lingva docendis} Hafniæ 1747. p. 7.}

Eilschow priser Englænderne, fordi de lægge saa stor Vægt paa Erfaringen, og
paa Grund af de mange dybsindige Tanker, der stamme fra England. Men han
slutter sig overvejende til Wolff, hvis Filosofi han vil gøre mere tilgængelig,
ogsaa for Damer,%
% footnote **) on p. 21
\footnote[2]{\textit{Forsøg til en Fruentimmerphilosophie.} Kbhvn. 1749.} og
hvis Værk han tillige vil fortsætte.%
% footnote ***) on p. 21
\footnote[3]{\textit{Philosophiske Breve over adskillige nyttige og vigtige
Ting.} Kbhvn. 1748.} Wolff havde jo sat Metafysiken i System, idet han har
inddelt den i Grundlære (Ontologi), Læren om Verden (Kosmologi), Læren om
Sjælen (Psykologi) og den naturlige Teologi. Eilschow har ikke selv behandlet
alle disse Fag; en saadan Behandling blev, som allerede omtalt, given af Kraft
nogle Aar senere. Det er den naturlige Teologi, Eilschow
% --- p. 22 ---
\setcounter{footnote}{0}%
\opage{22}især beskæftiger sig med. Det er ud fra Begrebet Mulighed, Eilschow
vil føre Bevis for Guds Tilværelse. Ingen endelig Tings Mulighed har sin Grund
i sig selv; men alle Muligheder have deres sidste Grund i Guds Mulighed --- og
denne har sin Grund i sig selv, ligesom alle Tal forudsætter Ettallet. Dette er
dog ikke at forstaa saaledes, at Tingenes Muligheder skulde skyldes Guds
Villie. De staa uafhængigt af hans Villie i hans evige Tanke; kun
Virkeliggørelsen skyldes Villien. Mod de evige Muligheder formaar Villien Intet
--- den guddommelige ligesaa lidt som den menneskelige.

Nu ligger det Ondes --- Smertens og Syndens --- Kilde i den Begrænsning, alt
Endeligt som saadant er underlagt, og som følger af, at de endelige Væsener
staa i Sammenhæng med hverandre. Denne Sammenhæng, der i sig selv er et Gode,
medfører Skranker og dermed Mulighed for Ufuldkommenhed. Det Onde kan da ikke
skrives paa Guds Regning, da han ligesaa lidt har skabt Mulighederne, som han
har skabt sin egen Forstand.%
% footnote *) on p. 22
\footnote[1]{\textit{Philosophiske Breve} p. 150. --- En lignende Tankegang
findes i \emph{Jens} \emph{Krafts} \textit{Naturlige Theologie.} p. 13---16.
--- Eilschow udtaler dog ogsaa, at vi ikke kunne slutte fra de
Ufuldkommenheder, der findes i en Del af Verden, til det Heles Ufuldkommenhed.}

Ligesaa lidt som Leibniz og Wolff har Eilschow besvaret den Indvending, som
Bayle rejste, og som Holberg gentog: hvorfor har da Gud skabt en Verden, om
hvilken han vidste, at den indeholdt Muligheden for Smerte og Synd, en
Mulighed, hverken guddommelig eller menneskelig Villie kunde ophæve? Holberg er
konsekventere end Eilschow. Forstand og Villie, Mulighed og Virkelighed stride
her i den Grad imod hinanden, at de egentligt maa tilhøre to forskellige
Væsener. Den Gud, der gerne vil, men ikke kan, er ikke almægtig; han har en
Skranke --- hvad enten man nu med Holberg søger denne i en evig Materie, eller
med Manikæerne i et evigt ondt Væsen, en Ahriman. --- Eilschow mangler Holbergs
Sans for Problemernes Alvor.%
% footnote **) on p. 22
\footnote[2]{Et rent filosofisk Standpunkt overfor det Ondes Problem er angivet
i min \textit{Religionsfilosofi} §§ 88---89 og min \textit{Etik} cap. VI.}

% --- p. 23 ---
\setcounter{footnote}{0}%
\opage{23}Ved Spørgsmaalet om Forholdet mellem Fornuft og Aabenbaring gør
Eilschow ogsaa Brug af Distinktionen mellem Mulighed og Virkelighed. Gud kan
aabenbare Noget, man ikke kan udlede af Fornuften og af de Muligheder, den
godtgør; men han kan ikke aabenbare Noget, der kuldkaster den naturlige
Kundskab, modsiger Sansernes klareste Efterretninger og omstøder alt det,
hvorpaa alle Sandheders Vished beror.

Ligesom Sneedorff anerkendte Enevælden, fordi den virkede i Oplysningens og
Folkelighedens Tjeneste, og ikke fordi den udsprang af Guds Naade, saaledes
anerkendte Eilschow og de andre Wolfianere Aabenbaringen i Tillid til, at den
ikke modsiger Fornuften og de evige Muligheder. Spørgsmaalet maatte blive ---
og blev det ogsaa --- paa hvilken Side Uretten ligger, naar der faktisk viser
sig en Modsigelse. Der kom den Tid henimod Aarhundredets Slutning, da man ikke
var i Tvivl om, at der var en Modsigelse tilstede, og heller ikke var i Tvivl
om, paa hvilken Side Uretten var.

Foreløbigt var Eilschow sikker paa, at han havde Midler i Hænde til at gendrive
Fritænkerne, d. v. s. dem, der negte det guddommelige Forsyn.%
% footnote *) on p. 23
\footnote[1]{Ordet Fritænker stammer fra England (free-thinker) og forekommer
først omkring Aar 1700, brugt om og af Mænd som Toland og Collins. Tidligere,
fra det 16. Aarhundrede af, brugtes i Frankrig Benævnelser som Deister og
Naturalister.} De forsynde sig imod Metafysiken og gendrives ved Hjælp af den
omtalte Lære om Mulighederne og disses sidste Grund. Uden Metafysik farer man
let vild, netop i saadanne Spørgsmaal som det Ondes Oprindelse og om Forsynet.
Dog mener Eilschow, at Fritænkerne egentligt blot ønske Frihed til at synde, og
at de negte Sandheden for at slippe for at lyde den. Der er vel Nogle, som f.
Ex. Spinoza, der have ført et smukt Levned --- men det er inkonsekvent, da de
ikke vente Løn eller Straf efter Døden. ---

Venligere stemt end overfor Fritænkerne er Eilschow overfor Dyrene. Han er en
af de første Filosofer i den nyere Tid,
% --- p. 24 ---
\setcounter{footnote}{0}%
\opage{24}der hævder, at vi have Pligter imod Dyrene. Forat man skal have Pligt
mod et Væsen, er det ikke nødvendigt, at det har Fornuft, naar det kun har Evne
til at lide. Der bestaar et Samfund mellem dem og os, idet vi tage deres
Tjeneste og til Gengæld have Pligter overfor dem. ---

Wolffs systematiske Filosofi raadede ved Universitetet til ind i det nye
Aarhundrede. Den sidste Repræsentant var \emph{Børge} \emph{Riisbrigh} (f.
1731, d. 1809), der var Professor i Filosofi fra 1767 til 1803. Hans
Undervisning roses for dens formelle Klarhed, men han var ikke meget produktiv.
Det betydeligste Arbejde fra hans Haand er en Afhandling (fra 1802---03) om
Kants Filosofi, et Værk, der gør ham Ære, og som senere hen vil blive omtalt.

\begin{center}\rule{2cm}{0.4pt}\end{center}

% --- p. 25 ---
\setcounter{footnote}{0}%
\opage{25}

\begin{center}\large 3. DANSKE KANTIANERE OG DERES MODSTANDERE\end{center}

Det Wolfiske System blev i Tyskland fortrængt ved et af de mest dybtgaaende
Tankearbejder, Filosofiens Historie kender --- ved \emph{Kant's} kritiske
Filosofi. Kant var bleven vakt ved Humes Kritik af Aarsagssætningen, den
Sætning, paa hvilken al Naturvidenskab og al spekulativ Filosofi og Teologi
byggede. Samtidigt stod Newton's matematisk udviklede Naturvidenskab for ham
som Idealet af videnskabelig Erkendelse. Naar han nu spurgte, hvorledes saadan
Videnskab var mulig trods Humes Paavisning af, at Aarsagssætningen hverken var
selvindlysende eller bevislig, saa fandt han Løsningen i, at der er visse
Former, den menneskelige Aand efter sin Natur stedse arbejder med, og som den
ogsaa da maa gøre Brug af, naar den skal fastslaa, hvad der har objektiv
Gyldighed, --- hvad der er Sandhed og Virkelighed. Derved blev det paa en Gang
muligt at anerkende Videnskaben i dens Selvstændighed og at se dens Grænser;
disse maa ligge der, hvor den menneskelige Aands Anskuelses- og Tankeformer
ikke længere finde noget givet Stof at anvendes paa. Videnskaben kan derfor
ikke gaa længere, end Erfaringens Mulighed gaar. Naar denne Mulighed glipper,
kan --- og maa --- Mennesket have en personlig Tro til Støtte for sin
Fastholden ved Samvittighedens Bud; men en spekulativ Teologi er umulig.

Samvittighedens Lov var ifølge Kant en ren Fornuftlov: ligesom den menneskelige
Aand i sin Videnskab stræber at bringe Iagttagelserne ind under Anskuelsens og
Forstandens
% --- p. 26 ---
\setcounter{footnote}{0}%
\opage{26}Former og Love, --- under Rummets og Tidens, Størrelsesbegrebets og
Aarsagsbegrebets Former, --- saaledes kan den kun da overbevise sig om en
Handlings Moralitet, naar den kan paavise, at Handlingen kan anerkendes ud fra
en almindelig, for Alle gældende Lov; den Grundsætning, jeg følger i min
Handling, maa være en saadan, at Enhver i mit Sted kunde følge den.

Som Kant fra den udviklede Videnskab søger tilbage til selve den arbejdende
Tankes inderste Natur, saaledes søgte han paa det etisk-religiøse Omraade
tilbage til Samvittighedens Selvlovgivning. Ud fra de Principer, han saaledes
vandt, bedømte han de spekulative Systemer og de teologiske Dogmer. Hans store
Sans for de oprindelige Kræfter førte ham fremdeles til paa det æstetiske
Omraade at pege paa Geniet som Kunstens uvilkaarligt fremsprudlende Kilde. I
Tanken, i Samvittigheden og i Geniet ytrer sig for Kant en og samme Grundkraft
i forskellige Former. Ved denne store Ide danner Kant paa én Gang Afslutningen
af det attende Aarhundredes Aandsarbejde og peger langt ud over dets Horizont.

Det var en ung, lovende Forsker, der især vakte Opmærksomhed for Kants Filosofi
i Danmark. I Foraaret 1793 holdt \emph{Christian} \emph{Hornemann} (f. 1769)
Forelæsninger over Kants Fornuftkritik. De blev afbrudte ved Forfatterens Død
samme Aar, og de store Forhaabninger, man havde stillet til ham, faldt til
Jorden. Med stor Begejstring havde han optaget Kants Hovedtanker i sig, og et
Studieophold hos Reinhold i Jena havde udviklet hans filosofiske Interesse. I
sine Forelæsninger (der ere trykte i hans \textit{Efterladte Skrifter} 1795)
gjorde han klart Rede for den kritiske Filosofis Forhold til tidligere
Opfattelser og begyndte derefter at gennemgaa Kants Hovedværk; men han naaede
ikke længere end til Slutningen af Læren om Rum og Tid. --- Hornemanns
Hovedopgave synes at have været den at vise, hvorledes det religiøse Problem
vilde stille sig under Forudsætning af den nye Opfattelse.

Den rent teoretiske Side ved Kants Filosofi blev kun lidet
% --- p. 27 ---
\setcounter{footnote}{0}%
\opage{27}fremdragen herhjemme, og ligesaa lidt den æstetiske. Det var den
moral- og religionsfilosofiske Side, der især blev paaagtet. Den betydeligste
Forkæmper for Kants Etik og Retslære var den Mand, der senere blev Danmarks
største Jurist, \emph{Anders} \emph{Sandøe} \emph{Ørsted} (f. 1778); han vandt
i Aaret 1798 Universitetets Guldmedaille for en Prisafhandling \textit{Om
Sammenhængen mellem Dyds- og Retslærens Princip,} en for sin Tid fortrinlig
Fremstilling af Kants Etik og Retsfilosofi og en Redegørelse for de Drøftelser,
til hvilke denne gav Anledning i Samtiden, saavel i Udlandet som herhjemme. I
den staaende Strid mellem „Bibelpartiet“ (Palæologerne) og „Fornuftpartiet“
(Neologerne) --- det første repræsenteret af Biskop Balle, det andet af
Horrebow (Udgiver af Tidsskriftet „Jesus og Fornuften“) tog Ørsted Del med en
længere Afhandling (i „Philosophisk Repertorium“ 1797---1798), der var rettet
mod Horrebow's Populær- og Halvfilosofi. Ud fra et lignende Standpunkt som
Lessing's i „Erziehung des Menschengeschlechts“ og Fichte's i „Kritik aller
Offenbarung“ peger han ud over de to stridende Partier. --- Ørsted blev --- som
vi i en senere Sammenhæng skulle se --- ikke staaende ved det strengt
kantianske Standpunkt, han i disse Afhandlinger havde indtaget. Men Beundringen
og Taknemmeligheden overfor hans store Lærer fulgte ham i hele hans Liv, og han
har i sin høje Alderdom udtalt den paa en smuk Maade.

Det var, som Ørsted bemærker i sin Selvbiografi%
% footnote *) on p. 27
\footnote[1]{\textit{Af mit Livs og min Tids Historie.} I (1851). p. 22. --- En
ret mærkelig og uafhængig Stilling overfor Kants Etik indtog den daværende
Professor i Moralfilosofi \emph{Anders} \emph{Gamborg}, der i Tilslutning til
Hutcheson og Adam Smith hævdede en af Religion uafhængig, af Menneskekærlighed
udsprungen Moral. I sit første Skrift \textit{(Forskel mellem Dyd og gode
Handlinger.} 1783) kender han endnu ikke Kants Skrifter. I et senere Skrift
\textit{(Jesu Moral.} 1799) erklærer han, at han hverken er Eudæmonist eller
Kantianer; men han glæder sig dog over, at han i sin Hævden af Sandhedspligtens
Ubetingethed stemmer overens med Kant, skønt deres Begrundelse er forskellig.},
ikke saa meget de unge Teologer som de unge Jurister, der beskæftigede sig med
Filosofi. Og i de Spørgsmaal, der be%
% --- p. 28 ---
\setcounter{footnote}{0}%
\opage{28}skæftigede det næste Tiaar af det attende Aarhundrede, vilde den
kraftige og frie Aand, der karakteriserer Kants og Fichtes etiske og
retsfilosofiske Synsmaader, have kunnet virke overordentligt frugtbringende.
Men der opstod ikke her, som i Tyskland, en Række af selvstændige og energiske
Tænkere og Reformatorer, som dem, de store Reformer i Preussen efter 1806
skyldtes. „Kønigsberger-Idealismen“, hvem ifølge Gneist Preussens Genoprejsning
for saa stor en Del skyldtes, satte ikke blivende Virkninger i Folkets Liv
herhjemme og afløstes tidligt af Romantikens mere passive og drømmende Retning.
For en Del kom dette af ydre politiske Forhold; for en Del vel ogsaa af, at vi
tolv Aar før Aarhundredets Slutning havde gennemført en af de største Reformer,
vort Lands Historie kender, en Reform, ved hvilken iøvrigt --- som paa sit Sted
berørt --- ogsaa filosofiske Ideer havde været medvirkende. Foruden Ørsted,
hvis filosofiske og juridiske Ungdomsstudier hørte med til Forberedelse af hans
lange Virksomhed som Dommer og Lovgiver, maa dog nævnes en Kantianer, hvis
Skrifter ved deres Klarhed og Energi staa langt over det Meste af, hvad der
frembragtes her hjemme i den over alle Bredder strømmende offentlige Diskussion
i det attende Aarhundredes sidste Tiaar, inden den strenge Trykkefrihedsordning
af 1799 stansede den frie Drøftelse af offentlige Sager. Præsten \emph{Michael}
\emph{Gottlieb} \emph{Birchner} (f. 1756, d. 1798) hævdede i sit Skrift
\textit{Om Trykkefriheden og dens Love} (1797), at selv i et absolut Monarki
kan fuldstændig Trykkefrihed meget vel herske. --- Det var den oplyste
Absolutismes Ideal --- fri Drøftelse af alle Tanker og dog absolut Lydighed i
alle Handlinger --- som her gennemførtes i sin yderste Konsekvens. Kant selv
havde iøvrigt i sin mærkelige lille Afhandling \textit{Was ist
Aufkl}\textit{ä}\textit{rung?} (1784) udtalt, at kun en oplyst Enehersker, ikke
en Fristat, turde sige: Ræsonner saa meget I vil, --- men adlyd! Men det var
saa ogsaa Kants Mening, at en saadan Enehersker havde Verden kun havt i
Friedrich II af Preussen. Birckners Kantianisme havde ført ham til at fortsætte
en Tankegang, som allerede Sneedorff og hans Slægtled havde
% --- p. 29 ---
\setcounter{footnote}{0}%
\opage{29}bevæget sig i. Ogsaa paa andre Omraader fremtræder Birckner som
udpræget Kantianer. Endnu stærkere end Kant og uden den Sans for det Ophøjede,
der præger Mesterens Fremstilling, hævder han Fornuftens Eneret i det etiske
Liv. Al Følelse --- med Undtagelse af den rene Følelse af Agtelse, Moralloven
indgyder --- var for ham egentligt Egenkærlighed. Selv Fædrelandskærlighed og
Menneskekærlighed vilde han derfor ikke anerkende som de rette etiske Motiver;
de kunne ikke kræves og paabydes, men maa udvikle sig uvilkaarligt, og det er
ikke den umiddelbare Agtelse for det Rette, der kundgør sig i dem, men den
Enkeltes individuelle Betagethed.%
% footnote *) on p. 29
\footnote[1]{Se hans Afhandlinger \textit{Om Kierlighed til Fædrelandet} og
\textit{Om Christendommens højeste Moralprincip.} (Efterladte Skrifter. Udgivne
af A. S. Ørsted. Kbhvn. 1800). --- Om de etiske og psykologiske Misforstaaelser
i Kants og Birckners Opfattelse se min \textit{Etik} XII, 1.} ---

Blandt de danske Kantianere maa endnu \emph{Baggesen} nævnes. Hos ham var det
ikke den strenge, rationelle Side af Kants Filosofi, der vakte Interessen, men
netop den Karakter af Ophøjethed, som Kants Tanker kunne frembyde, især naar de
bevæger sig paa Erfaringens Grænselinier. Kant tilfredsstillede for en Tid den
unge Digters sværmeriske Søgen, skønt han indrømmede, at den rent teoretiske
Entusiasme var ham fremmed. Han filosoferede, som han selv siger i et Brev til
Reinhold,%
% footnote **) on p. 29
\footnote[2]{\textit{Aus Jens Baggesens Briefwechsel mit Reinhold und Jacobi.}
I. p. 6.} af etisk og religiøs Interesse, og hans Tankearbejde gik egentligt
altid ud paa at skaffe Symboler for det, der bevægede sig i hans Følelse og
hans Fantasi. Filosofien var for ham „Gedankenspiel“: „Da jeg ikke arbejder i
Filosofien og paa Filosofien, er min Filosoferen --- hvad enten den bestaar i
Læsning af filosofiske Værker eller i Samtale med filosofiske Venner --- kun
Leg, men den kæreste af alle Lege.“%
% footnote ***) on p. 29
\footnote[3]{Anf. Skr. II. p. 56. --- Af Interesse er det at høre
\textit{Reinholds} Dom om Baggesen til tredie Mand: „Sein darstellender Genius
ist sein Plagegeist. In den Augenblicken, wo er ihn in Ruhe lässt, ist Baggesen
im Denken Philosoph und im Handeln verständiger Mann von der ersten Grösse, ---
aber das sind auch nur Augenblicke --- sonst ist er Engel oder Teufel, und
befindet sich im Himmel oder in der Hölle --- leider! öfters da als dort. Als
Korrespondent \textit{(wenn} er schreibt) und als auf einige Tage besuchender
Gast hat er wohl Niemanden über sich.“ Reinhold til Erhard \textit{2.} Aug.
1796. (Denkwürdigkeiten des Philosophen und Arztes Erhard. Herausg. von
Varnhagen von Ense. 1830. p. 426).} Poesien og Filosofien stredes om denne
forunderlige Aand. Og han stod ikke blot paa Grænsen mellem to Aandsomraader,
men ogsaa paa
% --- p. 30 ---
\setcounter{footnote}{0}%
\opage{30}Grænsen mellem to Aarhundreder. Vi ville møde ham senere i det
mærkelige Sammenstød med Grundtvig, der lovede saa Meget, men blev til saa
Lidet, fordi hos Baggesen Legen med Symboler fik Overtaget over Tankens Alvor.

\begin{center}\rule{2cm}{0.4pt}\end{center}

Kantianismen blev mødt med ikke ringe Modstand herhjemme fra forskellige Sider,
ikke mindst fra den selvtilfredse „Oplysnings“ Side. Og literære Ordførere som
Rahbek ærgrede sig over de unge Kantianeres Overmod og de nye Kunstord, hvormed
de slog om sig. „Oplysningen“ var iøvrigt enig med Ortodoxien om, at det kunde
ikke være saa vanskeligt, som Kant mente, at faa Tilværelsens Gaader løste.
Hvortil de mange Forberedelser, naar man havde sin Bibel og sin Katekismus
eller sin „sunde Menneskeforstand“ at ty til?

Men den nye Tankeretning fandt dog ogsaa virkeligt filosofiske Modstandere, og
nogle af disse have endog med Klarhed paavist Mangler og uløste Spørgsmaal i
den kritiske Filosofi. Tyge Rothe og Johannes Boye stødte sammen med
Kriticismen under deres Arbejde paa at udvikle deres egne Ideer, medens
Riisbrigh, Niels Schow og Treschow direkte undersøgte dens Grundlag og Rækkevidde.

\emph{Tyge} Rothes (f. 1731, d. 1795) Forfattervirksomhed falder i to Perioder,
en historisk-politisk og en filosofisk. Til den første Periode hører hans
Skrifter om Kristendommens Indflydelse paa Folkenes Tilstand i Europa og om
Nordens Historie, som gav ham Lejlighed til at virke for Bondestandens
Frigørelse. Til den anden Periode hører hans filosofiske Hovedskrift
\textit{Philosophies Ideer til Kundskab om vor}
% --- p. 31 ---
\setcounter{footnote}{0}%
\opage{31}\textit{Art og til Glæde over den} (1788---1789). Det er her ikke,
som man har sagt og gentaget, en af de almindelige Ordførere for Oplysning og
Lyksalighedslære, der har Ordet. Rothe er paavirket af Mænd som Herder og
Jacobi. Han tænkte paa at nedlægge sin Pen, da Herder's „Ideen“ udkom, men
fattede dog Mod til at gaa i hans Fodspor. Han er klar over, at hverken
Menneskets Adel eller dets Lykke kan bedømmes blot efter dets Oplysning. Den
frie Humanitet, der arbejder sig frem til sjælelig Harmoni og til Sympati med
alt Menneskeligt, er for ham Menneskets Bestemmelse. Samfundslivet udvikler sig
af Hjerternes fri og uforstyrrede Trang, ikke af Nød og Elendighed, heller ikke
af ren Fornuft. Og ligesom sine to tyske Mestre finder han i den umiddelbare
Følelse et Organ for det Højeste: ved en Naturaabenbaring, ikke gennem
overleverede Troslærdomme eller ad den rene Fornufts Vej, opdage vi vort eget
Væsen som foraarsaget ved et evigt Væsen, og dermed er da selve Aarsagsloven
given i sin Uundgaaelighed. Al Filosofi er derfor Teologi. Den positive
Religions Lærdomme fortolkes efter Naturaabenbaringens Lærdomme, som ogsaa
Erfaringen bekræfter; Newtons Paavisning af Naturens Lovmæssighed fører tilbage
til Naturens Udspring af et eneste Væsens Villie; han har givet os den
Fysikoteologi, der saa herligt karakteriserer vor oplyste Tidsalder.

Kants Lære hører (ligesom Spinozas, Berkeleys og Humes) til „de mørke
Filosofisystemer“. Den rokker ved den Erfaringsfilosofi, vi ene kunne bygge paa
med Tryghed. --- Det er aabenbart dels Kants Drøftelse af Aarsagssætningen som
ikke umiddelbart given eller indlysende, dels hans Kritik af Beviserne for Guds
Tilværelse, der frastøder Rothe. Han er vis paa, at naar Kant faar nøje
bestemt, hvor langt den strenge Viden gaar, da ville Menneskene føle, at de
ikke kunne nøjes med det, og arbejde sig fra den mørke Labyrint ud i den
umiddelbare Forvisnings klare Dagslys.

Rothe repræsenterer et Standpunkt, der ellers ikke kommer frem herhjemme i
Slutningen af det attende Aarhundrede. Han mangler Herders Originalitet og
Jacobis mær%
% --- p. 32 ---
\setcounter{footnote}{0}%
\opage{32}kelige Forening af Kritik og Patos, og hans Sentimentalitet og
Kunstleri stikker stærkt af mod deres Hævden af det Uvilkaarliges og
Umiddelbares Ret i Aandslivet. Men det var ikke uden Betydning, at ogsaa en
Tone af denne Art, var den end svag og undertiden falsk, kom til at klinge med
blandt Tidens mange Røster herhjemme. ---

Rothe havde ikke blot modtaget Paavirkning af tyske Forfattere, men ogsaa af
fransk Literatur. Især var det Genferen Charles Bonnet, Naturforskeren og Filosofen,%
% footnote *) on p. 32
\footnote[1]{\emph{Bonnet} hørte til de Udlændinge, der understøttedes af den
danske Regering under Frederik V. Hans endnu lærerige psykologiske Skrift
\textit{Essai analytique sur des facultés de l'ame} (især af Interesse for
Opmærksomhedens og og Villiens Psykologi) er trykt i København 1760 og
dediceret til Frederik V. Skønt Republikaner priste han de Danske lykkelige,
fordi de havde en saadan Monark: Ce republicain envieroit le sort de l'Iheureux
Danois, si un citoyen de Genève pouvoit envier quelque chose!} der begejstrede
ham og var hans Vejleder i Naturbetragtningen. Han hylder Bonnet's
Præformationslære, ifølge hvilken hver Art var skabt for sig og saaledes, at
Spirerne til alle senere Individer indeholdtes i Artens første Repræsentanter.
Alle Væsener danne en stor Kæde, men Leddene i denne Naturkæde forbliver gennem
Tiderne specifik („artsligen“) forskellige, og ingen Forandring af Arterne
eller Overgang fra den ene Art til den anden gaar for sig. De oprindelige Anlæg
i Naturen kunne ikke ophæves af nogen mekanisk Kraft. Mennesket er fuldkomment
i sit første Anlæg, men ydre Forhold, store Jordomvæltninger have grebet
hæmmende ind og forandret Anlægget, saa at en lang og møjsommelig Udvikling
blev nødvendig. Paa ethvert Trin i Udviklingen gælder det dog, at et Væsen
vinder den Tilfredsstillelse, det er det muligt at naa. Gennem denne
Betragtning hævder Tilværelsens Harmoni sig for Rothes Tanke. Han afviser
udtrykkeligt den Udvej, som Eilschow og Kraft (i Tilslutning til Leibniz og
Wolff) havde grebet: at skelne mellem Mulighederne for den guddommelige
Forstand og Virkeligheden, som den udspringer af den guddommelige Villie. „Kun
én kunde Guds Plan være; den Plan var den eneste iværksættelige, thi
% --- p. 33 ---
\setcounter{footnote}{0}%
\opage{33}bort med det Taabesnak, at Gud pønsede, og vejede og valgte!“ „Hvad
der kunde være i Guds Skabningsrige, det er der. Hvad der kunde bestaa, det
bestaaer.“ Hvorfor Paavirkninger (eller som Tyge Rothe siger, „Paastød“) hæmme
Anlæggenes Udfoldelse, det kan man ifølge Rothe ikke give nogen Forklaring af.
For vor Skyld forandres de evige Love ikke; de ere heller ikke givne for vor
Skyld. Men denne Tanke formørker ikke Rothes lyse Blik paa Tingene. ---
Optimismen er stigende --- fra Holberg gennem Eilschow og Kraft til Tyge Rothe. ---

Medens Rothe opfatter Udviklingen som blot Udfoldelse af oprindelige Spirer,
staar \emph{Johannes} \emph{Boye} (f. 1756, d. 1830) som Udviklingsfilosof i
mere moderne Betydning af Ordet. I sit mærkelige Skrift \textit{Statens Ven}
(1792 ff.) lægger han den Tanke til Grund, at enhver med Kraft begavet Ting,
ogsaa Mennesket, maa kæmpe med Modstand og Vold for at naa sit Maal. Ingen
sover sig god, men ethvert Skridt fremad mod Fuldkommenhed maa tilkæmpes. Kun
gennem Strid og Arbejde hjælpes Mennesket frem til det, han er. Der føres en
Kamp af Alle mod Alle. Tapperhed var en Dyd, før Fromhed blev det, og fysisk
Overmagt raadede, før aandelig Overmagt kunde øves. Naturen opdrager sine Børn
strengt --- men derved opdrager den Helte. De Dovne og Svage fortrænges i den
Kræfternes Strid, gennem hvilken Orden og Harmoni vindes. Denne livsalige Strid
er Midlet, som Forsynet bruger til Verdensplanens Virkeliggørelse.

Ret og Moral afhænger af Kulturens Udvikling, og det er Historiens Aarbøger,
ikke spekulativ Fornuft, der kan lære os Noget om deres Udspring. Mennesket har
udviklet sig fra en dyrisk Tilstand og fra Brutalitet gennem alle Kulturgrader
til den Højde, det nu indtager. Og Staten er det vigtigste Hjælpemiddel for
Kulturen. Statslivet grunder sig paa Forstaaelse af den Nødvendighed at give
Afkald paa nogle af sine Rettigheder for at bevare de andre. Gennem
Samfundslivet og ved Sprogets Hjælp er Mennesket hørt op at være Dyr; uden dem
levede vi endnu blandt Dyrene i Sko%
% --- p. 34 ---
\setcounter{footnote}{0}%
\opage{34}ven og aad Olden med dem. Men Udviklingen stanser ikke ved en
Selvbegrænsning i Forventning om Gensidighed. Grundtrangen i os sker først
Fyldest i Sympati og ædel Selvfølelse. Mennesket føler sig da som Del af et
Hele; hans Tanke skuer ud over større Kredse, og hans Interesse udvides langt
ud over hans eget Jeg. Paa dette Standpunkt er Dyden sin egen Belønning. Selv
Opofrelsen bliver kun mulig ved Trangen til at bruge den store, af Sympati
ledede Kraft. ---

Ud fra dette Standpunkt bekæmper Boye den kritiske Filosofis Forsøg paa at
bygge Moral og Ret paa den rene Fornuft. Denne Filosofi stanser, siger han, ved
en uforsonet Modsætning mellem den rene Fornuft og de psykologiske Motiver. „En
god Villie“ maa dog være en Villie, der er god imod Nogen, og skulde
Menneskekærlighed være forkastelig, fordi Mennesket i den følger en dyb Trang?
Ideerne om det Gode og Rette udvikle sig af menneskelige Interesser og deres
Kamp. ---

Boye's Tankegang, som her er givet i sammentrængt Form, kommer ikke til sin Ret
i hans Bøger og misforstaas derfor let af overfladiske Læsere. Det er en mandig
og fri Anskuelse, der rører sig hos ham, og den havde ikke fortjent den
Forglemmelse, der er bleven dens Lod, men som forklares ved Tidsaanden, der
først var letvindt optimistisk og snart efter sværmende romantisk. Dog manglede
Boye tilstrækkelig historisk Sans og tilstrækkeligt Materiale til at give sine
Ideer et grundigt Underlag.

A. S. Ørsted, der i sin Prisafhandling besvarede Boyes Angreb paa Kants Etik,
anerkender hans Skrift som „aandrigt og smagfuldt“. Han er enig med Boye i, at
den rimeligste Hypotese er den, at Menneskeheden i sin Barndom ikke har staaet
højt over Dyrene, en Antagelse, som ogsaa Kant havde fremsat (i sin Afhandling
„Muthmaaslicher Anfang des Menschengeschlechts“ 1786). Men han indskærper
Forskellen mellem den historiske Udvikling af Erkendelsen eller af
Samvittigheden og den filosofiske Paavisning af de Forudsætninger, paa hvilke
al Erkendelse og al Samvittighed byg%
% --- p. 35 ---
\setcounter{footnote}{0}%
\opage{35}ger, --- eller mellem Erkendelsens Begyndelse og dens Grundlag. ---
Ørsted har Ret i, at Kants Modstandere stadigt oversaa (og tildels endnu
overse) denne Forskel; men Kant har unægteligt selv bidraget Meget til
Misforstaaelsen, især ved ikke at give den psykologiske og historiske
Undersøgelse, hvad der tilkommer den. ---

\emph{Børge} \emph{Riisbrigh} fandt i sin høje Alderdom Tid og Lyst til at
drøfte den Filosofi, som fortrængte den Lære, han i hele sin lange
Universitetsvirksomhed havde forkyndt. Og han gør det med en Skarpsindighed og
en Uhildethed, der gør ham Ære. Hans Undersøgelse fremkom i Videnskabernes
Selskabs Skrifter 1803%
% footnote *) on p. 35
\footnote[1]{Dens Titel er: „Undersøgelse om og hvorvidt deres Mening er
grundet, som holde for, at sand Filosofi og et fuldstændigt Begreb om sand
Filosofi først i vore Tider ere blevne til.“ --- Filosofien hørte ikke til de
Fag, der lige fra det danske Videnskabernes Selskabs Stiftelse (1742) vare
repræsenterede i det. Først under A. P. Bernstorff's Præsidenttid blev (1795)
„de egentlige filosofiske Videnskaber eller den saakaldte spekulative Filosofi,
der i senere Tid har faaet saa megen Vigtighed“, føjede til de matematiske,
fysiske og historiske Fag. C. \emph{Molbech}: \textit{Det kongelige danske
Videnskabernes Selskabs Historie i dets første Aarhundrede.} p. 283.}, samme
Aar, som han trak sig tilbage fra sin Lærervirksomhed, og den er endnu af
Interesse. Han viser først, at man allerede før Kant har været opmærksom paa,
at hvad vor Erkendelse forudsætter og anvender i sin Undersøgelse af
Objekterne, det kan i selve Anvendelsens Øjeblik ikke være kommet fra selve
Objekterne, men maa have været i Sindet tilforn, hvad enten det i sin første
Oprindelse stammer fra Erfaring eller ikke. Der gives altsaa en
„Forerfaringskundskab“ (Kundskab a priori). Men Kant har først søgt Kilden til
disse Grundsætninger i selve det menneskelige Sind og vist deres rent formale
Karakter, samt givet en Fortegnelse over dem. Han har derved lagt Grunden til
en ny Videnskab indenfor Filosofien, den videnskabelige Fornuftkritik. Men han
har ikke, som nogle af hans Efterfølgere, paastaaet, at man før ham hverken
havde sand Filosofi eller vidste, hvad Filosofi var. Og han kastede
% --- p. 36 ---
\setcounter{footnote}{0}%
\opage{36}ingenlunde Vrag paa Erfaringen; han antog baade en Erfarings- og en
Forerfaringsfilosofi.

Riisbrigh ønsker som Sidestykke til Kants Undersøgelse af den rene Fornuft ---
en Undersøgelse af den empiriske Fornuft, og han mener, at man da vilde komme
til at se dennes store Betydning. Man vilde da igen paaskønne, hvad Sokrates,
Aristoteles, Baco og Locke havde lært. Paa Videnskabens, ja paa hele den
menneskelige Kulturs Vegne maa man advare mod at undervurdere
Erfaringsvidenskaben. Vor Forsken maa dog beskæftige sig med andre Objekter end
dens egne Former. Og paa den anden Side ere disse Former selv noget Givet:
Sindet har ligesaa lidt givet sig selv sine væsentlige Former, som det har
frembragt Formerne i Sanseverdenen! ---

I samme Bind af Videnskabernes Selskabs Skrifter, som bragte Riisbrighs
Afhandling, findes en Afhandling af Filologen og Arkæologen \emph{Niels}
\emph{Schou} (f. 1754, d. 1830), der ogsaa frembyder ikke ringe Interesse%
% footnote *) on p. 36
\footnote[1]{Dens Titel er: „Den nuherskende Skepticisme eller Kritik i
Filosofien kan, naar den grunder sig paa den menneskelige Fornuft og rigtigt
anvendes, ej være farlig, hverken for videnskabelig Cultur eller Moralitet.“}.
Forfatteren finder Kants Betydning især i den Bevægelse i Tankerne, han havde
vakt, hvorimod han finder den Kritik, G. E. Schultze havde anlagt paa Kants
Filosofi (især i „Kritik der theoretischen Philosophie“ 1801) aldeles
afgørende. Det er især Kritiken af Kants underlige Begreb „Ding an sich“, det
drejer sig om. Schou hævder, at om den absolute Realitet, hvis Tilværelse vi
erkender, indeholdes i Subjektet eller i Objektet, eller i hvilket Forhold den
staar til begge, vide vi ikke; vi spore kun dens Virkninger. Det overskrider
den menneskelige Fornufts Grænse at give en positiv Bestemmelse. Selv de mest
ophøjede filosofiske Fremstillinger --- som Spinozas --- ere her
Antropomorfismer. Filosofiens Opgave er at bestemme de Grænser, vor Viden
(„vort Vide“) stedse underligger, --- at drøfte Erfaringsvidenskabens Resultater i
% --- p. 37 ---
\setcounter{footnote}{0}%
\opage{37}dens Betydning for de sidste Spørgsmaal, --- og at arbejde hen til en
Harmoni mellem Handlen og Viden („Handle og Vide“). Ud over det Hypotetiske
komme vi ikke.

Allerede flere Aar før Riisbrighs og Schous Afhandlinger fremkom havde
\emph{Niels} \emph{Treschow} (f. 1751, d. 1833) i Kristiania, hvor han da var
Rektor, holdt Forelæsninger over Kants Filosofi%
% footnote *) on p. 37
\footnote[1]{\textit{Forelæsninger over den Kant'ske Philosophie.} Kiøbenhavn
1798.} og givet en kritisk Vurdering af denne, der gik dybere end noget af de
hidtil omtalte Arbejder om dette Emne. I en klar og livlig Fremstilling giver
Treschow en Oversigt over Kants Filosofi og knytter dertil sine kritiske
Bamærkninger. Han finder Udledningen af Forstandens rene Former
utilfredsstillende hos Kant og paaviser navnligt, at Kant ikke med Rette kunde
beraabe sig paa den gængse Logiks Inddeling af Dommene, naar han selv forandrer
denne. Og selvom det var lykkedes Kant at finde alle de Former, Erkendelsen gør
Brug af, saa maatte det dog endnu paavises, hvad det er i det Givne, i
„Stoffet“, der muliggør og foranlediger Anvendelsen af en bestemt speciel Form
i det enkelte specielle Tilfælde. Hvad gør f. Eks., at jeg her ser en Firkant,
hist en Trekant? Det kan ikke udledes af Rummet som en almindelig
Anskuelsesform. Hvad gør, at jeg her anvender Størrelsesbegrebet, hist
Aarsagsbegrebet? Det kan ikke udledes af Forstandens almindelige Evne til at
forbinde Iagttagelser i visse Former (Kategorier). Der maa i Tingene selv ligge
Noget til Grund for de specielle Forbindelsesmaader. Allerede vore Fornemmelser
maa lære os, hvad Naturen selv har forbundet, og hvad den har adskilt.

Der savnes ifølge Treschow i Kants Fremstilling af Tankens Former
(Kategorierne) saa vigtige Begreber som Lighed (Identitet) og Forskel,
Samstemmen og Modstrid, som dog anvendes i al Tænkning, og som ligesaavel
anvendes paa det Givne som Størrelses- og Aarsagsbegrebet. Og de nævnte
Begreber ere endog af ganske særlig Betydning: Forstandens Arbejde bestaar
egentligt blot i Sammenligning af
% --- p. 38 ---
\setcounter{footnote}{0}%
\opage{38}det i Erfaringen Givne. Det er efter Tingenes Lighed eller Ulighed,
deres Samstemmen eller Modstrid, at Forstanden afgør, hvilke Forbindelser
mellem dem der have Gyldighed, og hvilke ikke. --- Treschow har her peget paa
et vigtigt Punkt i Erkendelsesteorien. Kant havde ikke set, at de formelle
logiske Principer ligesaavel som de matematiske Principer og Aarsagsprincipet
anvendes paa det givne Stof, saa at ogsaa deres Gyldighed maa undersøges af
Fornuftkritiken.

Det er Erfaringsfilosofiens Standpunkt, Treschow hævder overfor Kant. Han
lovpriser Britterne, der her ikke alene selv ere komne Sandheden nærmest og
have gjort betydelig Fremgang i Kundskab om den menneskelige Natur, men ogsaa
ere blevne Lærere for Andre. Han hævder John Locke's Standpunkt overfor Leibniz
og Kant og bebrejder Tyskerne, at de ikke have villet lære af Englænderne.
Locke havde Ret i, at enhver Form forudsætter et Stof, og at der ingen Begreber
kan dannes, naar der ikke forud er givet visse Forestillinger.

Men Treschow misforstaar Kant, naar han siger, at ifølge Kant danner Forstanden
sine Begreber, „førend nogen Erfarenhed kommer til“. Det er jo, som A. S.
Ørsted paaviste overfor Boye, Kants udtrykkelige Lære, at Erkendelse
\textit{begynder med} Erfaring, selv om ikke al Erkendelse
\textit{grunder}\textit{sig paa} Erfaring. Selve Erfaringen var ifølge Kant
sammensat af givet Stof (Fornemmelserne) og uvilkaarligt anvendte Former
(Anskuelsesformer og Kategorier).

Med den Begrænsning, Kant fastslaar for vor Erkendelse, er Treschow heller ikke
tilfreds, saalidt som med den Maade, paa hvilken Kant vil have den Trang
tilfredsstillet, som Videnskaben ikke kan stille. „Kant selv kan ikke negte, at
der ligger i vor Natur en ubeskrivelig Længsel efter at vide mere, end han
holder for muligt at opnaa, ja, han tilstaar, hvad der er endnu mere, at vor
Tilværelses Øjemed nødvendig fordrer en Kundskab, vi ej paa nogen af de angivne
Veje [Sanseiagttagelsens og Forstandens] ere i Stand til at erlange. Agter han
da vel at bringe os tilbage under den blinde Troes og Hierarchiets Aag, eller
er det med Flid, han lige%
% --- p. 39 ---
\setcounter{footnote}{0}%
\opage{39}som i et Heltedigt forvikler Knuden for at drive Opmærksomheden til
det Højeste og fortjene des større Beundring ved at løse den?“ --- Treschow
mistvivler om Kants Kunst og vil ikke bygge paa den nye Grundvold, Kant vilde
lægge for Livsanskuelsen. Det er betænkeligt at give den praktiske Fornuft og
Troen Ret til at overskride Erkendelsens Grænser. Hvorledes kan det bevises, at
visse Trosforestillinger angive den eneste Maade, paa hvilken de etiske Ideer
kunne hævdes? Og opgiver man al Bevisførelse, appellerer man blot til den
personlige Interesse, saa aabner man Døren for Trosstridigheder og
Forfølgelser. Om det, der blot skal have praktisk Betydning, vil Enhver --- og
ikke mindst Magthaverne --- mene at have en Stemme. --- Treschows
Betænkeligheder ere ikke ugrundede. Et andet Spørgsmaal er, om Kant ikke har
Ret i, at det Afgørende for enhver Livsanskuelse tilsidst er en personlig, af
Livserfaring fremgaaet Trang. Med denne maa enhver Tro, hvad enten den støtter
sig til Videnskab eller ikke, kunne stemme. Det Store hos Kant er netop den
Kunst, hvormed han strammer Modsætningen mellem Tankens Evne og Livets Trang.
Han saa kun ikke, at Livets Trang ikke én Gang for alle kan godtgøres at have
fundet sin Tilfredsstillelse i den „naturlige Religions“ saalidt som i nogen
anden Religions Dogmer. ---

Treschow, som vi her kun have mødt som Kritiker, ville vi senere møde som en
ejendommelig og selvstændig Tænker, der i flere Retninger har øvet en betydelig
Indflydelse paa dansk Tankeliv.

\begin{center}\rule{2cm}{0.4pt}\end{center}

% --- p. 40 ---
\setcounter{footnote}{0}%
\opage{40}

\begin{center}\large 4. HENRIK STEFFENS\end{center}

Det synes at være en Strid, der gentager sig ved hvert Aarhundredes Slutning,
naar det nye Aarhundrede begynder. I Weimar var man ikke i Tvivl om, naar det
nittende Aarhundrede begyndte; det blev fejret af Goethe, Schiller, Schelling
og Steffens, da Klokken slog tolv om Aftenen den 31. December 1800. Denne lille
Kreds karakteriserer Overgangen mellem de to Aarhundreder. Det attende
Aarhundrede havde været Kritikens og Rationalismens Tid. Men Mænd som Rousseau,
Lessing og Herder, Kant, Goethe og Schiller pegede ud over deres Tid ved deres
Sans for det Oprindelige, det Naturlige, det Geniale, --- for de dybt liggende
Kræfter, hvis Væsen aldrig ganske udtømmes i de Former, en Tidsalder eller et
Aarhundrede formaar at give det. Disse Mænd repræsentere den nye humanistiske
Renaissances Værk og Kamp. Og ligesaa højt, som de staa over det attende
Aarhundredes gængse Rationalisme, ligesaa højt staa de over den saakaldte
Romantik, der betegner det nittende Aarhundredes første Reaktion imod sin
Forgænger. Det var de nævnte Mænds Ideer, der nu i Romantikens Filosofi og
Poesi førtes ud i den spekulative Tankes og den sværmende Fantasis Former.
Romantiken bygger paa deres Aandsværk, samtidigt med, at den betegner en vis
Reaktion imod det. Saaledes betegne de to yngste Deltagere i hin Nytaarsfest,
Schelling og Steffens, en Reaktion mod Kants mandige Hævden af Tankens
Selvstændighed og Selvbegrænsning, samtidigt med at de ville bygge paa det
Grundlag, han havde lagt.

% --- p. 41 ---
\setcounter{footnote}{0}%
\opage{41}Medens Frankrig og England stredes om Verdensherredømmet, fordybede
man sig i Norden i Ideernes og Fantasiens Verden. Schiller havde anslaaet Tonen
i sit Digt ved det nye Aarhundredes Begyndelse:

Freiheit ist nur in dem Reich der Träume,

und das Schöne blüht nur im Gesang.

Hvad Danmark angaar, havde der i de sidste Tiaar af det attende Aarhundrede
været ført en ivrig Diskussion om offentlige Forhold, og der havde udviklet sig
en levende Sans for politisk Frihed. Sneedorff, Tyge Rothe, Birckner og
Johannes Boye vidne derom. Men det Haab, især Birckner havde havt, at Enevælde
og absolut Trykkefrihed kunde forenes, brast, og den almindelige europæiske
Reaktion trængte ogsaa ind over vort Lands Grænser. Det var et Tryk fra
Ruslands Side, der gav det afgørende Stød til Fremkomsten af
Trykkefrihedsforordningen af 1799, som stansede den frie Diskussion.%
% footnote *) on p. 41
\footnote[1]{Se herom: Edvard Holm: \textit{Den offentlige Mening og
Statsmagten i den dansk-norske Stat i Slutningen af det attende Aarhundrede.}
(Universitetsprogram 1888). p. 175 f.} Mærkværdigt hurtigt faldt Sindene til Ro%
% footnote **) on p. 41
\footnote[2]{De Norske tog Sagen paa en mere energisk Maade end de Danske, og
de fik ogsaa snart en stor politisk Opgave. „For Norge fik Sæden, der var
udstrøet i Slutningen af det attende Aarhundrede, Spirekraft, da Tanken om
Nødvendigheden af en fri Forfatning kom til at gaa sammen med Spørgsmaalet om
at føre et selvstændigt politisk Liv. De fleste af de betydeligere
Ejdsvoldsmænd havde som unge Mennesker i København i Halvfemserne faaet de
første vækkende Indtryk i politisk Henseende.“ E. \emph{Holm}: Anf. Skr. p. 183
f. --- Grundtvig udtalte engang i en Samtale med Sibbern, at Bevægelsen i det
sidste Tiaar af det attende Aarhundrede ikke fik Lov til „at løbe Livet af
sig“. Den kom en Menneskealder efter op igen. (\emph{Sibbern}:
\textit{Dikaiosyne.} 1843. p. 9).} --- dels i Filisteriets, dels i Romantikens
Verden. De, der søgte ud over Døgnets Kval eller søgte Trøst for de Ulykker,
der snart brød ind over Landet, fandt, hvad de søgte, i Spekulationens og
Fantasiens Verden.

Og dog, --- det var for de Danske, der hverken havde haft nogen Kant eller
nogen Goethe, en stor aandelig Skat, der blev aabnet med den nye aandelige
Strømning, som kaldes Romantiken. Vor Literatur, vort Tankeliv og vor Poesi blev
% --- p. 42 ---
\setcounter{footnote}{0}%
\opage{42}anden Gang genfødt. Tiden efter Holberg havde, hvad Originalitet
angaar, været en Bølgedal. Nu gik det opad igen, og her som paa Holbergs Tid
blev det Nye optaget og bearbejdet paa selvstændig Maade. Folket satte sit eget
Præg paa, hvad det tilegnede sig.

Det var \emph{Henrik} \emph{Steffens}, der først forkyndte Romantikens
Evangelium i Danmark. Født i Stavanger (1773), altsaa Normand ligesom Holberg
og Treschow, levede han dog sin Barndom i Danmark og studerede i København og i
Kiel. Den store aandelige Bevægelse i Tyskland drog ham til sig, og efter et
rigt Samliv med Tidens største Aander paa Digtningens, Naturforskningens og
Filosofiens Omraade vendte han i Aaret 1802 hjem til København, hvor han i
Vinteren 1802---1803 holdt sine berømte Forelæsninger, der tildels foreligge
trykte \textit{(Indledning til filosofiske Forelæsninger.} Kbhvn. 1803).

\begin{center}\rule{2cm}{0.4pt}\end{center}

Ordet „romantisk“ (eller, som man tidligere sagde „romanisk“) synes fra England
at være kommet til Frankrig og Tyskland. Det kommer af „Roman“, en eventyrlig
Fortælling af den Art, der først udviklede sig hos de romanske Folk. Først
brugtes det i dadlende Betydning: vild, fantastisk, regelløs, --- senere, især
fra Rousseau's Tid, i rosende Betydning. Ordet findes brugt om Alperne ---
først i den første, siden i den anden Betydning.%
% footnote *) on p. 42
\footnote[1]{Smlg. herom \emph{Ludwig} \emph{Friedländer's} Exkurs i anden Del
af hans \textit{Sittengeschichte Roms. 5.} Udg. II. p. 245---247.} Det betegner
Modsætningen til det Klare, Bestemte, Overskuelige og Forstaaelige, og bruges
først om Landskaber og om Menneskekarakterer. Det er blevet det staaende Udtryk
for Grundstemningen og Tankegangen i Tyskland og i Norden i den første
Menneskealder af det nittende Aarhundrede.

Romantisk er det Uendelige, det, der fører ud over al Begrænsning, saa at
Sindets Higen ganske kan udfolde sig. Det
% --- p. 43 ---
\setcounter{footnote}{0}%
\opage{43}Romantiske danner Modsætningen til det Klassiske i Kunsten og til det
Rationelle i Tænkningen.

Romantisk er det Fjerne, der vækker Vemod og Længsel, særligt det Fjerne i
Tiden, de gamle Tider, der længst ere svundne. Romantikens Tidsalder levede ---
eller mente at leve --- i Erindringen om det Store og Værdifulde, der var
forbigangent, medens det attende Aarhundrede levede i Glæden over, hvad der var
kommet, eller i Haabet om, hvad der skulde komme. Oldtid og Middelalder stode
for Steffens i mystisk Glans som Tider, i hvilke de Modsætninger, der nu delte
Livet, ikke vare fremtraadte, --- „i hvilke Videnskab, Kunst og Handling i
uafladelig Forening formaaede at frembringe det Guddommeligste og Højeste“, og
han spurgte: „Er ikke det Skønneste, Herligste som gennemtrænger vor Tid Sansen
for hin gamle, forsvundne?“

Hvad man længtes tilbage til, var især, som det ses af de anførte Ytringer, den
samlede, koncentrerede Maade, hvorpaa Livet, det aandelige som det materielle,
kunde føres, før Arbejdsdelingen havde grebet splittende ind. Allerede Rousseau
og Schiller havde energisk fremdraget Arbejdsdelingens mislige Følger, og de
ere ikke mindst paa dette Punkt Romantikens Forgængere. Steffens udtaler (i sin
syvende Forelæsning), at hvad der især udmærkede hin svundne Tid, var Enheden
af det, vi nu maa adskille som Ord og Handling. „I levende Handling udprægedes
Alt, hvad vi nu søge ved Ord, alle Naturens Revolutioner og Udviklinger, alle
Religionens Mysterier. Livet selv, frisk som det udsprang fra Guddommen, havde
umiddelbart den Betydning, vi nu kun ane. Derfor var alle Forestillinger
levende. Intet Spor til den adsplittende Død, som i nyere Tider har gjort
Videnskaber og Poesi og, ved dem, Religionen saa monstrøse.“

Denne Adsplittelse skal nu ophøre, Arbejdsdelingen skal ophæves, og Religion,
Poesi og Videnskab igen blive Et! --- Det var Romantikens Evangelium.

\begin{center}\rule{2cm}{0.4pt}\end{center}

% --- p. 44 ---
\setcounter{footnote}{0}%
\opage{44}For at Filosofien kan træde i det nye Evangeliums Tjeneste, maa
Metoden forandres. Steffens fordrer af sine Tilhørere --- foruden visse
almindelige Kundskaber og Interesse for at løse Tilværelsens Gaade --- en Evne
til Helhedsanskuelse, som lader alle Enkeltheder staa som Led i et absolut
Hele. Han vil give „et levende System, ikke en død Klassifikation“, --- vil
„aabne et betydningsfuldere Blik over Livet og Tilværelsen end det, den
almindelige Erfaring i det daglige Liv, indknebet af endelig Trang, leder os
til. Af det dybere Blik over Naturen, Historien og vort inderste Gemyt,
udspringer det filosofiske Problem som af sin egentlige Kilde“.

Men Helheden kan ikke naas ved Opdyngen af Enkeltheder, ved et Aggregat af
Fakta, saaledes som Naturforskere og Historikere i Reglen give os det. Bringe
de en Forbindelse tilveje mellem de sondrede Dele, bliver det kun en ganske
udvortes, mekanisk Forbindelse. Thi det Enkelte forstaas kun i Sammenhæng med
det Hele, og derfor danner en Helhedsanskuelse en nødvendig Forudsætning for
enhver Syslen med Enkeltheder. Og hos alle store, grundlæggende Forskere rører
der sig da ogsaa en saadan Anskuelse, selvom det kun er i Instinktets eller
Anelsens Form. --- Ved Steffens' egne Studier besjæledes han --- som han siger
i sin Selvbiografi --- af en dunkel Anelse om et altomfattende Liv, længe kun
som fri Fantasi og mere af digterisk end videnskabelig Betydning.

Det var Romantikens skæbnesvangre Fejl, at den var saa sikker paa sin
Forudsætnings Gyldighed, at den ikke anvendte de rette Midler til at arbejde
sig Skridt for Skridt hen imod det store Maal, der var stillet ved
Forudsætningen om Tilværelsens Enhed. Den Svælgede i Forestillingen om dette
Maal istedetfor at besjæles af den til Fordybelse i Enkelthederne. Da den
mekaniske Naturvidenskab, der i to hundrede Aar havde paavist, hvorledes
mangfoldige Forandringer i Naturen kunde forklares som Bevægelser, fremkaldte
ved Tryk, Stød eller Tiltrækning, ikke syntes at fyldestgøre det romantiske
Ideal, blev den betragtet som en Vildfarelse.
% --- p. 45 ---
\setcounter{footnote}{0}%
\opage{45}Videnskaben, klager Steffens, har lukket Naturen for os; Nøglen til
Naturens Hemmeligheder kan kun findes i vor egen Aands Indre. Naturen maa
betragtes som Aandens Forhistorie og læses ovenfra nedad. --- Istedetfor den
Analogi, til hvilken Naturvidenskaben hidtil havde støttet sig, --- Analogien
mellem Forandring og Bevægelse, --- vil Steffens altsaa sætte en anden:
Analogien mellem Forandringer i Naturen og psykologiske Processer. Naturen skal
ses som en Række af Trin, ad hvilke den slumrende Aand successivt vaagner til
Bevidsthed. Romantiken har sin Berettigelse overfor en materialistisk
Dogmatisme (eller, som man har kaldt den, „en mekanisk Mytologi“), men den
Analogi, den vil sætte istedetfor Naturvidenskabernes frugtbare mekaniske
Analogi, kan kun føre til Poesi, ikke til Videnskab. ---

Man kunde tro, at naar Naturen skal være Aandens Forhistorie, maatte der være
Tale om en real Udvikling, en Udvikling i Tiden, ligesom i den moderne
Udviklingslære. Men Steffens erklærer det udtrykkelig for en falsk Hypotese at
antage en virkelig, faktisk Overgang fra en Livsform til en anden. Han
anerkender, at det er en Anelse om Naturens Enhed, der rører sig hos
Naturforskere, som hylde en saadan Antagelse; men her har denne Anelse ledet
vild. Den tilsyneladende Vorden af nye Former er kun et Skin, der fremkommer
ved Relativiteten af alle Modsætninger. Som Steffens siger i et senere tysk
Skrift: „Die Stufenfolge der Formen ist nur für den Schein“%
% footnote *) on p. 45
\footnote[1]{\textit{Grundzüge der philosophischen Naturwissenschaft.} Berlin
1806. p. 187 (kfr. 170).}. Tidsforskellen var for den romantiske Bevidsthed for
udvortes til at have afgørende Betydning. Man følte sig i den Grad ind i
Naturen, at man var lige hjemme i alle Former og saa enhver af dem som evigt
Udtryk for det Ene, der aabenbarer sig i Alt.

Alligevel lægger Steffens stor Vægt paa Forskellen mellem de forskellige
Naturformer, og han finder den især i den forskellige Grad af Individualitet,
de forskellige Væsener besidde. Der ytrer sig i Naturen en individualiserende
Tendens, der naar sit Højdepunkt i Mennesket. Hos det og i det
% --- p. 46 ---
\setcounter{footnote}{0}%
\opage{46}aabenbarer Naturens midtpunktsøgende, inderliggørende Tendens sig.
Mennesket har en hel Verden overfor sig; men derfor aabenbarer sig ogsaa en hel
Verden i hans Indre. Aandens Forhistorie er her afsluttet --- paa selve det
Punkt, hvorfra vi maa gaa ud for at forstaa de andre Naturvæsener. Det sidste
Afsnit i det tyske Skrift, der ligger til Grund ved de danske Forelæsninger%
% footnote *) on p. 46
\footnote[1]{\textit{Beiträge zur inneren Naturgeschichte der Erde.} Freyberg.
1801. p. 275 ff. --- Skriftet er decideret til Goethe og „med hellig Sky
nedlagt i den højere Poesis delfiske Tempel“.}, har til Overskrift: „Durch die
ganze Organisation sucht die Natur nichts als die individuelleste Bildung“.
Gennem Reproduktion, Irritabilitet og Sensibilitet nærmer den organiske Natur
sig --- formedelst stigende Individualisering --- Intelligensernes Rige, og
Alt, hvad dér viser sig, ligger allerede som slumrende Anlæg i den bevidstløse
Natur. --- I sin Selvbiografi \textit{(Was}\textit{ich erlebte.} Tiende Bog)
siger han endog: „Hvad jeg kalder Naturfilosofi, er ikke Andet end den
Overbevisning, at der... i Alt, hvad der er Genstand for Forskning, maa
anerkendes et Ejendommeligt, Noget, som udvikler sig ud af sig selv.“

\begin{center}\rule{2cm}{0.4pt}\end{center}

Steffens kom ikke til at fortsætte sin Virksomhed herhjemme. Man vilde ikke
have den Sværmer til Lærer for de Unge. Han drog til Tyskland og blev Professor
i Halle, senere i Breslau og tilsidst i Berlin, hvor han døde 1845.

Den romantiske Enhed af Poesi, Religion og Filosofi blev ham senere hen ikke
nok. Han saa, at dette Ungdomsideal kun langsomt og ad spredte Veje kunde
realiseres, og den Trang til at fordybe sig i store Emner med hele sin Sjæl,
der i hans Ungdom havde ført ham til Naturspekulation, førte ham senere til at
hylde den lutherske Lære. Hans Sjæl fandt Hvile i den religiøse Tro, og
Religionen blev ham nu det eneste Lyspunkt i Historiens Kaos. (Was ich erlebte.
Ottende Bog)%
% footnote **) on p. 46
\footnote[2]{Smlgn. den udførlige Karakteristik, jeg har givet af Henrik
Steffens i mine \textit{Mindre Arbejder.} Anden Række. p. 162---180. --- Til
sin smukke Udgave af Henrik Steffens' Forelæsninger (1905) har Hr. B. T.
\emph{Dahl} knyttet en interessant Afhandling: \textit{Nogle Sider af Henrik
Steffens' ydre og indre Liv. ---}}.

% --- p. 47 ---
\setcounter{footnote}{0}%
\opage{47}Ved sin Tænkning har han øvet en betydelig Indflydelse især gennem
\emph{Sibbern}, der kom i et nært Forhold til ham i Halle og i Breslau. Det er
især Steffens' stærke Hævden af det Individuelle%
% footnote *) on p. 47
\footnote[1]{Det var ogsaa paa dette Punkt, at Steffens fik Indflydelse paa
Schleiermacher: Smlgn. \textit{Den nyere Filosofis Historie. 2.} Udg. II. p.
193 f.}, der spores i senere dansk Filosoferen, --- især i Linien Sibbern-Poul
Møller-Kierkegaard.

Til Poesiens Historie hører hans Indflydelse paa \emph{Oehlenschläger}, som han
selv smukt har udtrykt saaledes: „Jeg gav ham ham selv.“ Ogsaa paa
\emph{Grundtvig} øvede han stor Indflydelse, vistnok mest ved den anden, ikke
udgivne Række af Forelæsninger, der behandlede Poesien og kom ind paa
historisk-filosofiske Betragtninger. Den fælles Modsætning til Rationalismen
var et Foreningspunkt; og Grundtvig udtalte (i et Brev til Ingemann 1824), at
hans egen poetisk-historiske Retning var en Fortsættelse af Steffens' fra 1803.
Dog følte han en stor Modsætning mellem dem og erklærede samtidigt, at de vare
himmelvidt adskilte. Hvad der især voldte Adskillelsen, var, at Steffens
aldeles ikke vilde vide af Grundtvigs Ideer om „Nordens og især Danmarks
aandelige Vilkaar og dybe historiske Vigtighed“; saadanne Tanker var i
Steffens' Øjne „Judaisme“ --- men for Grundtvig var Danmark netop „Historiens
Palæstina, det Folkeland, hvor Kristendommen skal naa sin højeste jordiske
Forklaring“%
% footnote **) on p. 47
\footnote[2]{I det anførte Brev til Ingemann (se Grundtvigs og Ingemanns
Brevveksling 1821---1859. p. 35) omtaler Grundtvig et „Ultimatum“, han sendte
Steffens (efter dennes Besøg hos ham 1824). Dette er nu trykt hos \emph{Fr}.
\emph{Nielsen}: \textit{Grundtvigs religiøse Udvikling.} p. 239---243.}.

\begin{center}\rule{2cm}{0.4pt}\end{center}

To unge danske Mænd var der, som ingen Vækkelse behøvede af Romantikens
Apostel: \emph{Brødrene} \emph{Ørsted}. I Stilhed havde de selv tilegnet sig,
hvad den nye Retning kun%
% --- p. 48 ---
\setcounter{footnote}{0}%
\opage{48}de bringe af Værdifuldt. Selv udgaaede fra det Bedste og Dybeste i
det attende Aarhundredes Aandsudvikling, stode de ved deres store Selvarbejde
friere overfor den nyere Retning, idethele to af de alsidigste Aander, Danmark
har frembragt. Vi ville senere vende tilbage til dem. Her skal kun anføres en
Udtalelse af H. C. Ørsted om Steffens i et Brev til Oehlenschläger (1. Novbr.
1807), der var bleven vred og uretfærdig mod sin tidligere Ven:

„Hvad Steffens angaar, saa ved jeg ret godt, at han vil indklemme Naturen i et
System, og det i et meget snevert. Gud bevare mig fremdeles for at blive
Steffensianer...... Med en sjelden Mangfoldighed, skønt ikke Dybde, i
Kundskaberne forener han mangfoldige lykkelige Overblik. Naar man ikke tager
disse for uimodsigelige Resultater af den dybeste Filosofi, men for det, de
ere: Lynglimt til Vejledning, saa lærer man Meget af ham. Saaledes har jeg
gjort, og kan uden Fare vedblive at gøre.“

\begin{center}\rule{2cm}{0.4pt}\end{center}

% --- p. 49 ---
\setcounter{footnote}{0}%
\opage{49}

\begin{center}\large 5. HANS CHRISTIAN ØRSTED\end{center}

En gærende Bevægelse, som Romantiken var, maa godtgøre sin Berettigelse ved den
Udvikling, og især ved det Arbejde, den fører til. Ingen var saa vel forberedt
til at modtage de vækkende Paavirkninger af den nye Strømning, uden dog at
hildes i den, --- som Hans Christian Ørsted.

Født i Rudkøbing (1777), hvor der var meget liden Adgang til Undervisning,
maatte han ved Selvstudium og ved Samarbejde med sin et Aar yngre Broder skaffe
sig de Kundskaber, han behøvede for at blive Student. Derefter fulgte nogle
rige Ungdomsaar i København. Hans alsidige Interesse viste sig ved, at han ---
alt imedens han dyrkede sit farmaceutiske Studium --- vandt Prisen for en
æstetisk og en medicinsk Afhandling og ivrigt studerede Kants Filosofi, en
Interesse, han delte med Broderen. Men medens denne mest interesserede sig for
Kants Etik og Religionsfilosofi, var det Kants Opfattelse af Naturerkendelsen
og dens Forudsætninger, der især interesserede ham. I „Filosofisk Repertorium“
(1798) udgav han en Afhandling om \textit{Naturfilosofiens første Grunde,} i
hvilken han først giver en Oversigt over Kants Skrift „Metaphysische
Anfangsgründe der Naturwissenschaft“ og derpaa udvikler sin egen Opfattelse,
især gennem en Kritik af Atomteorien, som endnu vil have Interesse. Mærkeligt
er det dog --- men karakteristisk for Tidsalderen --- at skønt det nævnte
Skrift af Kant allerede er mere dogmatisk, end Kants egen Erkendelsesteori
tillod, mente dog den unge Forfatter at maatte gaa et Skridt videre. Kants
Metode bestod i stedse at spørge, hvilke Forudsæt%
% --- p. 50 ---
\setcounter{footnote}{0}%
\opage{50}ninger videnskabelig Erfaring hviler paa, og det var ad denne Vej han
i sit nævnte Skrift mente at kunne slaa fast, af der kun kunde gives to
bevægende Grundkræfter i den materielle Verden, en tiltrækkende og en
frastødende. Et materielt Punkt kan, mente han, kun virke paa et andet
materielt Punkt i en lige Linie mellem de to Punkter, og her er da kun to
Bevægelser mulige, en, ved hvilke de to Punkter nærme sig til hinanden, en
anden, ved hvilken de fjerne sig fra hinanden. Hertil slutter Ørsted sig
ganske, men han mener, at Kants Tankegang kun fører til en Hypotese, fordi den
bestaar i en Slutning tilbage fra Erfaringen. Ørsted vil derimod udlede alle
Naturlove „af vor Kendeevnes Natur, hvilken Kant saa fortræffeligen har
udviklet i sin „Kritik der reinen Vernunft“ “. Men han overser, at Kant i dette
sit Hovedværk netop havde indskærpet, at det kun var Erfaringens
Grundforudsætninger, der fandt deres Forklaring ved selve den menneskelige
Erkendelses Natur, medens de specielle Naturlove maatte udfindes Haand i Haand
med selve Erfaringen: Kants „Metaphysische Anfangsgründe der Naturwissenschaft“
afvige forsaavidt fra Hovedværket.

Da Ørsted skrev den nævnte Afhandling, i hvilken han vilde være mere kantiansk
end Kant selv, var allerede Schellings første naturfilosofiske Skrifter
udkomne. Men overfor dem gør Ørsted bestemt Front. Han udtaler: „Disse Bøger
fortjene vistnok Opmærksomhed for de skønne og store Ideer, som man finder i
dem, men de miste Meget af deres Værdi formedelst den ikke ret strenge Metode,
hvorved Forf. blander Erfaringssætninger ind uden at adskille dem
tilstrækkeligen fra de aprioriske Sætninger, --- især da de Erfaringssætninger,
han anfører, ofte er grundfalske.“

Faa Aar efter denne Afhandlings Fremkomst kom Ørsteds Vandreaar saavel som hans
første experimentale Arbejder. Allerede før Rejsen havde han begyndt sin
Lærervirksomhed ved Universitetet (1800), der skulde blive af saa stor
Betydning for dansk Videnskab og Aandsliv, saavel som for vor materielle
Kultur. Vi ville ikke her følge hans Udvikling i det Enkelte, men især holde os
til den Bog, i hvilken han
% --- p. 51 ---
\setcounter{footnote}{0}%
\opage{51}paa sine gamle Dage samlede sin Opfattelse af Naturvidenskaben og
dens Betydning for Livs- og Verdensanskuelsen: \textit{Aanden i Naturen.} Denne
i vor Literatur enestaaende Bog udkom i Aaret 1850, men den bestaar af
Afhandlinger, der ere skrevne til forskellige Tider, og den vidner om
Sammenhængen i Ørsteds Liv og Tænkning fra Ungdommens Dage til Alderdommens
Tilbageblik paa Arbejdet%
% footnote *) on p. 51
\footnote[1]{I det af V. Østergaard udgivne Værk „Vort Folk i det nittende
Aarhundrede“ (1897) har jeg givet en Karakteristik af Brødrene Ørsted af mere
populær Art end den, jeg i nærværende Skrift giver af de to Mænd.}.

\begin{center}\rule{2cm}{0.4pt}\end{center}

Medens Steffens klagede over, at Videnskaben skjulte Naturens Aand for os,
vilde Ørsted netop gennem Videnskaben trænge ind i Naturens Aand. Det var hans
Overbevisning, at det var den samme Fornuft, der ytrer sig i den ydre Natur og
i vor egen Tænken og Følen. Det Bestaaende i Naturen er Lovene. I et Vandfald
er det Bestaaende jo ikke de enkelte legemlige Dele, men den Lov, der binder
dem, deres Bevægelser og disses Virkninger --- Draabeadspredelsen,
Skumdannelsen, den ved Styrtningen, Brusningen og Skumningen foraarsagede Lyd
--- sammen til en eneste Naturhandling. Vi mærke ikke blot en Mangfoldighed,
men en Samfoldighed, der udtrykker en Naturtanke, en Ide. En Tings Væsen er
Indbegrebet af de samvirkende Naturlove, i Kraft af hvilke den bestaar.
Naturlovene ere i Naturen, hvad Tankerne ere i vor Aand.

Kantianeren er her bleven Romantiker. Kant havde kun fundet en Analogi mellem
vore Tanker og deres Overgang i hverandre gennem Slutning paa den ene Side,
Naturprocesserne paa den anden Side. Ørsted gør, ligesom Schelling og Steffens,
denne Analogi til Identitet. Men sin „dynamiske“ Naturopfattelse, ifølge
hvilken Stoffet ikke er en for sig bestaaende død Væren, men
Virksomhedsytringer, der bestemmes og begrænses ved Naturlovene, havde han
allerede vundet gennem Studiet af Kant.

Det religiøse Udtryk for „Aanden i Naturen“ er Gud. Ti
% --- p. 52 ---
\setcounter{footnote}{0}%
\opage{52}det Grundvirksomme og det Ordnende i Naturen er ikke to forskellige
Ting, men en uendelig, levende Fornuft. Guds Villie er Religionens Udtryk for
Tilværelsens evige Love. Og, som Ørsted et andet Sted udtrykker sig: „den
evige, levende Fornuft kalde vi Gud, naar vi betragte den i dens
Selvbevidsthed, dens Personlighed.“ Modsætningen mellem Gud og Verden gælder
egentligt kun, naar man ved Verden forstaar det i Endeligheden, der fjerner os
fra Gud. I de højeste Øjeblikke falder Modsætningen bort, idet det, man kalder
Verden, staar som en Guddomsvirkning.

\begin{center}\rule{2cm}{0.4pt}\end{center}

Naar Ørsted opfattede Stoffet som afledet af Kraften, --- naar han i sit Skrift
„Den almindelige Naturlæres Aand og Væsen“ (1811, optrykt 1847) lærte, at
Stoffet ikke er Andet end det formedelst Naturens Grundkræfter opfyldte Rum, og
at det, der giver Tingene deres uforanderlige Særegenheder, er de Love,
hvorefter de frembringes, --- saa har han dermed udtalt en Anskuelse, der i
Fysikens nyere Historie har faaet stigende Bekræftelse. Fysiken har mere og
mere udviklet sig til en Lære om Naturkræfter. Med den anførte Udtalelse af
Ørsted fra 1811 kan sammenstilles en Udtalelse af Faraday fra 1846: „Hverken i
det, man kalder det tomme Rum, eller i Legemerne kan jeg opdage Andet end
Kræfterne og de Linier, hvorefter de virke.“ Professor Christiansen ser i
Ørsteds Tanker paa dette Punkt „en Forudanelse af, hvad en fjern Fremtid skulde
bringe“.%
% footnote *) on p. 52
\footnote[1]{Chr. \emph{Christiansen}: \textit{H. C. Ørsted som Naturfilosof.}
(Det danske Videnskabernes Selskabs Oversigter. 1903). p. 488.}

\begin{center}\rule{2cm}{0.4pt}\end{center}

Trods den skarpe Dom, Ørsted i sin Tid fældede om Schellings og Steffens'
Naturfilosofi, kom han dog ogsaa selv undertiden ind paa fantastiske Analogier
og Kombinationer. Saaledes naar han sammenstiller Verdenshjørnerne med Ilt,
Brint, Kulstof og Kvælstof og finder Overensstemmelse mellem de elektriske
Figurer og de organiske Former. Men der
% --- p. 53 ---
\setcounter{footnote}{0}%
\opage{53}var én Tanke, som den romantiske Naturfilosofi havde udtalt, og som
ledede Ørsted paa frugtbar Maade under hans speciellere Undersøgelser, --- det
var Tanken om Naturkræfternes Slægtskab og Sammenhæng. Denne Tanke forfulgte
han i det Enkelte. Den førte ham til at opfatte Lysstraalen som en Række af
smaa elektriske Gnister. Han har dermed antydet den senere af Maxwell og Hertz
udviklede elektromagnetiske Lysteori.%
% footnote *) on p. 53
\footnote[1]{Smlgn. Chr. \emph{Christiansen}: \textit{Den elektromagnetiske
Lysteori.} (Det danske Videnskabernes Selskabs Skrifter. 1889). p. 189---193.})
--- Af særlig Betydning for hans Forsken blev dog Ideen om Naturkræfternes
indre Sammenhæng derved, at den Skridt for Skridt ledte ham til at foretage de
Forsøg, hvorved han (Juli 1820) opdagede, at en elektrisk Strøm paavirker
Magnetnaalens Bevægelse. Med Stolthed kunde han erklære, at denne Opdagelse
ikke skyldtes et Tilfælde, men var et Resultat af lang Tids Eftertanke. Han
havde paa dette Punkt maattet berigtige de Forudsætninger, han i sin Ungdom
havde optaget efter Kant. Ifølge den Maade, paa hvilken Kant havde
systematiseret Newtons Lære, skulde der kun kunne findes to bevægende Kræfter,
Tiltrækning og Frastøden, virkende ad Forbindelseslinien mellem to materielle
Punkter. Men hvis en elektrisk Strøm skulde have Indflydelse paa en Magnet,
maatte Bevægelsen netop gaa paa tvers af Forbindelseslinien. Hin
naturfilosofiske Tanke kunde han da kun gennemføre paa Trods af Newtons og
Kants Autoritet.%
% footnote **) on p. 53
\footnote[2]{Smlgn. Chr. \emph{Christiansen}: \textit{Nogle Bemærkninger om
Naturvidenskabernes Betydning for vor Tid.} (Universitetsprogram 1905). p. 41.
--- Som bekendt forfulgte Ørsted ikke sin Opdagelse videre. Det var
\emph{Ampère}, der først i den af Ørsted paaviste Kendsgerning saa et enkelt
Exempel paa en almindelig Lov, som han gennemførte teoretisk og experimentalt.
Filosofien har ikke Skylden for, at Ørsted ikke gik videre; thi Ampère var
ogsaa Filosof og har en smuk Plads i Filosofiens Historie. Se \textit{Den nyere
Filosofis Historie. 2.} Udg. II. p. 304---308.}

\begin{center}\rule{2cm}{0.4pt}\end{center}

Som det ses af det Foregaaende var der for Ørsted en inderlig Sammenhæng mellem
hans naturvidenskabelige, filoso%
% --- p. 54 ---
\setcounter{footnote}{0}%
\opage{54}fiske og religiøse Forestillingskreds. Det var ogsaa en af hans Livs
Hovedtanker, at Naturvidenskaben ikke skulde staa isoleret i det menneskelige
Liv, men gribe ind i Tænkning, Tro og Handling. Med Udgivelsen af „Aanden i
Naturen“ tilsigtede han, som han aabent udtalte, en Forandring i den
menneskelige Verdensanskuelse. Han opfattede selve Videnskabsdyrkningen som
Religion: den søger jo det, som bestaar under alle Forandringer, --- „det
Enhedsbaand, der gør, at Tingene i alle deres mangfoldige Fordelinger og
Adskillelser dog ikke adsplittes.“

I sin Biografi af H. C. Ørsted siger \emph{Carsten} \emph{Hauch} om ham: „Det
var hans Overbevisning, at der i en grundig Opfattelse af de virkelige
Verdensforhold laa en dybere Filosofi end den, der findes i de fleste
filosofiske Systemer, og en ægtere Poesi end den, der findes i mange Digte; ja,
han mente endog, at selv vore religiøse Anskuelser derved kunde berigtiges.“

I det store, prægtige Brev til Oehlenschläger af 1. og 6. November 1807, i
hvilket han læser sin Ven Texten, fordi han vilde løbe for tidligt fra Skole,
udtaler han sig særligt om Forholdet mellem Videnskab og Poesi. Tænkeren, mener
han, vil først ved Slutningen af sin Vandring naa det Punkt, hvor Kunstens
Verden begynder, medens omvendt Digteren, naar han har formet sine Anelser med
fuld Klarhed, nærmer sig Videnskaben. Han havde ikke Romantikens Sky for
Videnskaben og Virkeligheden. Og i sit Digt „Luftskibet“ gjorde han selv et
Forsøg paa at vise, hvorledes Naturvidenskaben kan træde i Poesiens Tjeneste,
--- et Forsøg, som dog netop kan tjene til at vise, at det er Menneskestræben
og Menneskeskæbne, der er Sjælen i al Poesi: Digtet handler i Grunden om
Montgolfier's, Luftskibets Opfinders, Skæbne, ikke om Naturfænomener. Megen
Poesi er der dog ikke deri. Poesien var i Ørsteds egen Sjæl, i hans Tænken,
hans Begejstring og hans Arbejde, men han formaaede ikke at lægge den ind i sit Digt.

Han bebrejdede Digterne, at de gjorde Overtroen poetisk, istedetfor at finde
det Poetiske i selve Naturen. Og ved Overtro forstod han den Indbildning, at
Noget kan virke i Natu%
% --- p. 55 ---
\setcounter{footnote}{0}%
\opage{55}ren anderledes end efter dens Love. Hans Hensigt med Udgivelsen af
„Aanden i Naturen“ var, som han skriver i et Brev (8. Februar 1850), „at vise
vor Tidsalders ofte overpoetiske og overtroende Mennesker, at Forstanden, som
de saa hyppigt miskende, har det store Kald at værne om Samvittighedens og
Troens Helligdom, hvoraf jeg rigtignok vil have endel gammel Overtroens og
Fordommens Surdejg udfejet.“ Allerede tidligere --- i sin Tale paa
Kristendommens Tusindaarsfest i Danmark (1826) --- havde han da ogsaa forsvaret
det attende Aarhundredes „dristigt granskende Mænd“, over hvem „visse af vor
Tids stolte og haarde Mænd i Kristendommens Navn bryde Staven.“ Han delte ikke
Romantikens Dom over den nærmest foregaaende Tid. Der var jo ogsaa i hans egen
Udvikling en indre Sammenhæng mellem, hvad han havde taget med sig fra det
foregaaende Aarhundrede, og hvad han havde lært i det nye.

\begin{center}\rule{2cm}{0.4pt}\end{center}

Tankegangen i første Del af „Aanden i Naturen“ førte til en Polemik med Biskop
Mynster, som især angik to Punkter. --- Mynster hævdede, at Udviklingen maa
være begyndt med en oprindelig Fuldkommenhedstilstand. Men Ørsted lærer, at Alt
i Verden begynder fra noget Uudviklet og gennemløber en Række af Trin. Hvorfor
det er saaledes, kan han ikke sige, men Erfaringen viser det. Teologernes Lære
om, at Naturen er bleven slettere ved Menneskets Synd, kan ikke forsvares.
Naturlovene have ikke forandret sig, og der har været Lidelse og Død i Verden,
før Mennesket blev til. „Jeg henstiller til Dogmatikerne at overveje, hvorvidt
deres Videnskabs Lære om Syndigheden i alle Maader maa betragtes som
ubestrideligt rigtig, eller kunde vinde ved en ny Bearbejdelse.“ --- Det andet
Punkt var, hvorvidt vore Ønsker og vor Trang kan bestemme, hvad der er Sandhed.
Herimod udtalte Ørsted sig bestemt og i smukke, inderlige Ord: „Maatte vi ikke
skamme os i vort Inderste, naar vi grebe os selv i at ville have en anden
Sandhed end den virke%
% --- p. 56 ---
\setcounter{footnote}{0}%
\opage{56}lige?.... Nej, lader os give Sandheden Æren! Med den er det Gode
uopløseligt forbundet. Den fulde Sandhed bringer selv sin Trøst med sig.“

I denne Udtalelse kundgør sig ikke blot Ørsteds store Sandhedskærlighed, men
ogsaa hans lyse Blik paa Livet. Han var en stor Optimist. En Sikkerhed for, at
Sandheden altid bringer Trøst, kan Ingen skaffe, og desuden kan der være en
lang Vej til den fulde Sandhed!%
% footnote *) on p. 56
\footnote[1]{Smlgn. mine Bemærkninger om Ørsteds Udtalelse i min
\textit{Religionsfilosofi} (§ 115).} Men Ytringen er karakteristisk for H. C.
Ørsteds lykkelige Natur. Frederika Bremer, den svenske Forfatterinde, gik i
Rette med ham for hans Optimisme. Han lover ikke at glemme Livets mørke Side.
„De tvang mig,“ skriver han (1849) til hende, „til at underkaste mine Tanker
over Ulykkerne, Smerten, det Hæslige, en ny Overvejelse, som ikke har været mig
ufrugtbar, og som jeg endnu fortsætter. Under min Glæde over det Heles Harmoni
skal jeg stræbe at indrømme Smerten og Sorgen deres lovlige Rettigheder!“

Saaledes skrev han to Aar før sin Død. Hans Overvejelser have neppe ændret hans
Blik paa Livet. Hans eget Liv havde i bedste Forstand været lykkeligt. Hans
Ungdoms Begejstring satte Frugt i en stor Opdagelse og i Fremstillingen af en
betydningsfuld Verdensanskuelse, saavel som i praktisk Virken for hans
Videnskabs Indflydelse paa Livet (ved Stiftelsen af Selskabet for Naturlærens
Udbredelse og ved Grundlæggelsen af polyteknisk Læreanstalt). Der var Harmoni i
hans Personlighed som i hans Tanke. Ungdommens Begejstring satte Frugt i
Manddommens Værk, og Alderdommen tabte ikke Forstaaelsen af Ungdommens Tanker.

\begin{center}\rule{2cm}{0.4pt}\end{center}

% --- p. 57 ---
\setcounter{footnote}{0}%
\opage{57}

\begin{center}\large 6. GRUNDTVIG OG H. C. ØRSTED.\end{center}

Hos H. C. Ørsted var det ene af de Elementer, Romantiken vilde sammensmelte til
Enhed, blevet udskilt og hævdet i sin Selvstændighed, --- Videnskaben. Og
Forholdet mellem Videnskaben og de andre Elementer i Romantikens Enhed var for
ham snarest det, at den skulde ligge til Grund, eller at vi gennem den skulde
komme til Poesien og Religionen.

I modsat Retning bevægede en anden af vort Folks store Aander sig ---
Grundtvig. For ham blev Religionen Hovedsagen, og Poesi og Videnskab skulde
være den underordnede, uselvstændige i Forhold til den. Intet Under, at der kom
et skarpt Sammenstød mellem de to Mænd. Der førtes en betydningsfuld Strid
imellem dem i Aarene 1814 og 1815.

Man kunde undre sig over, at Grundtvig overhovedet har haft Noget at gøre med
Filosofien. Og dog har den ogsaa grebet ind i hans Udviklingsgang og har endog
afgivet Motiver, der vedblive at virke, efterat han havde naat sit definitive
Standpunkt. Vor mest nationale Retning har faaet vigtige Impulser fra tysk Tænkning.

Steffens' Forelæsninger --- eller, som Grundtvig siger i „Kirkespejlet“, „den
forskrækkelige Støj, som min halvtyske Fætter, Henrik Steffens gjorde paa
Ehlers' Kollegium, saa Regensen nær var falden ned“ --- vakte ham især i
poetiskhistorisk Henseende, som allerede omtalt i det Foregaaende. Sansen for
det Store og Oprindelige, Romantikens bedste Side, blev styrket hos ham, næret
som den allerede var bleven ved Tidens store Begivenheder og ved Fordybelsen i de
% --- p. 58 ---
\setcounter{footnote}{0}%
\opage{58}gamle Sagaer. 22 Aar gammel gik han (1805) som Huslærer til
Langeland, og her var det, han med Schiller, Fichte og Schelling grublede over
Livets Gaader. Især har Fichte havt stor Indflydelse paa ham. Han grebes af den
Maade, paa hvilken Fichte indskærpede Nødvendigheden af en i Villien fæstnet,
af Menneskets Indre fremgaaende Overbevisning i Modsætning til gold Reflexion
og flagrende Forestillinger („Die Bestimmung des Menschen“. 1800). Naar Fichte
i „Grundzüge des gegenwärtigen Zeitalters“ (1806). tog til Orde mod Tidens
Aandløshed, den døde Lærdom og den selvtilfredse Oplysning, og ikke mindst naar
han fremhævede det mundtlige Ords Betydning i Modsætning til det skrevne, var
det Tanker, der bar Frugt gennem hele Grundtvigs Liv og Virksomhed. Og
utvivlsomt blev han ogsaa betagen af den Maade, hvorpaa Fichte senere (i „Reden
an die deutsche Nation“ 1808) begejstret hævdede Nationalitetens uforgængelige
Betydning og særligt henviste til de tyske og skandinaviske Folk som dem, der
fra urgammel Tid havde bot paa samme Jord, --- saavel som af Fichtes
Opdragelsesideer, der krævede en Udvikling for de Unge ad den frie Tilegnelses
og Selvvirksomheds Vej. Da Fichte døde, hyldede Grundtvig --- der dengang dog
var naat til sit udpræget positivt-religiøse Standpunkt --- ham i et smukt
Mindedigt.%
% footnote *) on p. 58
\footnote[1]{Smlgn. min Afhandling: \textit{Fichtes Taler til den tyske Nation
og de danske Nordslesvigeres nationale Kamp} (Mindre Arbejder. Anden Række. p.
192---200). (Om Fichtes Indflydelse paa dansk Aandsliv se særligt p. 195 f.).}

Poesi og Filosofi vare ikke nok for den unge Mand, hvis Sind var i saa stærk
Bevægelse. Han tager Afsked med Filosofen i en mærkelig religionsfilosofisk
Afhandling fra Aaret 1807 (trykt i „Theologisk Maanedsskrift“): \textit{Om
Religion og Liturgie.} Det er en religionsfilosofisk Afhandling --- ti man kan
nu engang ikke tage Afsked med Filosofien uden gennem Filosofi.

Der spørges først: Trænger vi til Religion? --- Vi leve ved et aabent Hav.
Poesien søger at flyve derover, men forsvinder i Skyerne, der belyses af
Glansen fra det Fjerne. Filoso%
% --- p. 59 ---
\setcounter{footnote}{0}%
\opage{59}fien undersøger, hvorfra denne Glans kommer, --- bygger Skibe til
Fart, men stanses af Bølgerne og maa enten blive Poesi eller lade sig nøje med
at forstaa Jorden og Livet der i den endelige Verden. Filosofien er Tidens
kloge Datter, men kender ikke sin Moders Fødsel. Den kan paavise Livets
Disharmonier, men ikke vise, hvorledes de udsprang af en oprindelig Harmoni,
eller hvorledes de kunne opløses i Harmoni. Ved al sin Spekulation kommer den
dog ikke ud over en Hypotese. Saa maa vi da enten opgive at vinde Kundskab om
Forholdet mellem Tid og Evighed --- men det vilde føre tilbage til Dyriskhed,
--- eller Kundskaben maa komme ovenfra.

Spørges der dernæst: Hvad er Religion? --- saa vidner Historien, at den er en
guddommelig Kraft, der hæver Mennesket over alle Jordens Omgivelser til det
Højeste, Oversanselige, og ved hvilken der stiftes et Samfund mellem det
Endelige og det Evige. I Kristendommen fremtræder disse Træk tydeligt og med
Fuldkommenhed. I den ere Poesi og Filosofi, der i de hedenske Religioner stedse
laa i Strid, forbundne til Enhed, og de kunne derfor begge hver paa sin Maade
udtrykke Religionen i Endelighedens Verden%
% footnote *) on p. 59
\footnote[1]{Maanedsskriftets Udgiver (P. E. Müller) erklærer i en Note, at han
„ikke ganske forstaar dette“. Jeg forstaar det saaledes: Religionen er det
Fundamentale, men Poesien vidner om den i sine Billeder, Filosofien i sine
Spørgsmaal. Denne Tydning stemmer med Afhandlingens Tankegang og med Grundtvigs
daværende Standpunkt.}. Poesien faar sin Realitet, og Tanken kan arbejde med
sine Tilnærmelser. Men selve Enheden af Poesi og Filosofi er et Mysterium, som
man fornegter, baade naar man holder sig til Poesien, og naar man blot holder
sig til Filosofien. Og det er i Liturgien --- særligt i Sakramenterne --- at
dette Mysterium finder sit Udtryk. --- Og dermed har den unge Forfatter da
banet sig Vej over i Teologien og til sine Forslag om Kirketjenestens Ordning.

Skønt denne Afhandling ligger forud for Grundtvigs religiøse Gennembrud er den
af stor Interesse som Vidnesbyrd om en Stræben efter at stille det religiøse
Problem i skarpere
% --- p. 60 ---
\setcounter{footnote}{0}%
\opage{60}Form, end Rationalismen og Romantiken gjorde det muligt. Religionens
Forskel fra Poesi og Videnskab betones skarpt --- og dog fastholdes endnu
Muligheden af en Harmoni mellem alle tre Omraader, kun saaledes, at det er det
religiøse Element, der bestemmer Klangfarven. Uklart staar det hen, hvorledes
der kan siges at være vundet „Kundskab“ om Forholdet mellem det Evige og det
Endelige, naar Resultatet bliver et Mysterium. Fremdeles: hvorledes kan
Historien vise os guddommelige Kræfter virkende? Historien kan gennem sine
Undersøgelser fremstille, hvad der er sket, og hvilke Kræfter, der have virket;
men at disse Kræfter vare guddommelige, kan den ikke fastslaa. Og endeligt:
skulde det virkeligt staa saaledes til med os, at vi kun have den Udvej at
henfalde til Dyriskhed, hvis vi ikke kunne vinde fuld Kundskab om Forholdet
mellem Tid og Evighed? ---

Disse kritiske Betænkeligheder gjorde sig ikke gældende for den unge Teolog,
der nu gik sine indre aandelige Kampe imøde. Han forargedes over en
rationalistisk Historieskrivers Omtale af Korstogene og førtes ind i en Periode
af Sværmeri og Kamplyst, indtil han gennemrystedes af det personlige Spørgsmaal
om hans egen Sjæls Frelse og bragtes til Vanvidets Rand. I et Brev til Molbech
(9. Marts 1811) erklærede han kort efter: Min hele Tilstand vilde være den
uforklarligste Gaade, dersom Satan ej var medvirkende! ---

Efter at have gennemgaaet denne Krisis optraadte han saa (især i
„Verdenskrøniken“ 1812) med Bibelen i Haanden overfor den slappe Tidsalder,
dens Literatur og dens Videnskab. Da Molbech havde angrebet hans Bog og dens
Omtale af forskellige Personligheder som uhistorisk og fordrede, at Historien
skulde være uafhængig af Teologien, svarede Grundtvig: „Skal en Kristen skrive
Historie, da maa han skrive den bibelsk, kalde al Lærdom, som strider mod
Bibelen, ligefrem uden videre Bevis Løgn og Vildfarelse, lade Bibelen dømme
Mænd og Tider, ti ene den har Ret dertil; den er Kristi Statholder paa Jorden.“%
% footnote *) on p. 60
\footnote[1]{\textit{Chr. Molbech og N. F. S. Grundtvig.} En Brevveksling.
Udgivet af L. Schrøder. 1888. p. 123 (smgln. p. 101).}

% --- p. 61 ---
\setcounter{footnote}{0}%
\opage{61}Der er nu ikke mere Tale om Religionen som en højere, om end mystisk
Enhed af Poesi og Filosofi. I Bibelen taler Religionen med Myndighed, og al
anden aandelig Stræben maa underordne sig den.

Medens \emph{Molbech} forsvarede Historien overfor Profetens haarde Domme,
traadte H. C. \emph{Ørsted} i Skranken paa Filosofiens og Naturvidenskabens Vegne%
% footnote *) on p. 61
\footnote[1]{\textit{Imod den store Anklager.} Af H. C. \emph{Ørsted},
Professor. 1814. (Motto: Du skal ikke sige falsk Vidnesbyrd imod din Næste).
--- Grundtvig's Svar førte Titelen: \textit{Imod den lille Anklager, det er
Prof. H. C. Ørsted,} med Bevis for, at Schellings Philosophie er ukristelig,
ugudelig og løgnagtig. Ved N. F. S. \emph{Grundtvig}, Præst 1815. (Motto:
Stoppes skal deres Munde, som tale Løgn).}.

I Sammenstødet mellem Ørsted og Grundtvig mødtes to modsatte Tankeverdener ---
Videnskabens med dens stadige Udviklingsstræben og Bibeltroens med dens Stræben
efter at bøje alt Aandeligt ind under en enkelt, en Gang for alle, given
Forestillingskreds. Og det var to modsatte Naturer, der mødtes. Den Ene --- en
af de sundeste Sjæle, dansk Aandsliv kender --- gik sin frejdige Gang gennem
Forskningens Verden, sikker paa, at intet Sandt og Stort, ingen virkelig Værdi
vilde mistes ved den stadige Søgen, men at Livet netop vilde blive rigere ved
den. Den Anden havde i en svar Krise faaet fæstet sit Blik paa et eneste Punkt
i Aandens og Historiens Verden, og det, hvori han havde fundet sit Tilhold,
holdt han nu op for Tidsalderen som det Eneste og det Højeste, --- som det, der
skulde dømme Alt og Alle, men ikke selv dømmes af Nogen.

Striden mellem dem gjaldt Videnskabens Selvstændighed. Grundtvig erklærer, at
det ikke er „selve“ Videnskaben, han kæmper imod, --- kun dens Selvstændighed.
Al Videnskab skal styres af en kristelig Ide. Ørsted svarer hertil; at en
uselvstændig Videnskab, --- en Videnskab, der tager nogetsomhelst andet Hensyn
end det, dens Emne kræver, --- ikke er Videnskab, og at vi jo dog lære Naturen
bedre at kende gennem Naturen selv end gennem Bibelen! --- Der stredes ogsaa
om, hvorvidt Schellings Filosofi er kristelig eller ikke;
% --- p. 62 ---
\setcounter{footnote}{0}%
\opage{62}dette har nutildags mindre Interesse og vilde føre os for vidt.
Derimod har det sin Interesse at se noget nærmere, hvorledes Grundtvig paa
dette Tidspunkt opfattede Forholdet mellem Religion og Videnskab.

Filosofien som selvstændig Tænkning kan --- mener Grundtvig--- have sin
Betydning i en ukristelig Tid. Saaledes havde Kant og Fichte en stor Gerning at
udøve véd deres strenge Idealitet og deres Fordring om, at Sandheden om
Mennesket, dets Kaar og dets Bestemmelse skal findes ad Tankens Vej. Naar
Tidsalderen ikke fulgte disse Filosofer, var det, fordi de fordrede for Meget
af den slappe Tid, der kun havde forladt Kristendommen forat undgaa den Kamp,
der følger med al højere Stræben. Fra Kristendommens Synspunkt er en
selvstændig Filosofi en Uting; her skal Fornuften underkaste sig det skrevne
Ord, idet den erkender dets guddommelige Kraft til at overvinde Verden, og da
er en kristelig Logik en ganske herlig Ting, naar det gælder at stoppe
Negternes og Spotternes Mund. --- Hvori denne „kristelige“ Logik er forskellig
fra den sædvanlige Logik, --- om den har andre Grundsætninger end denne, ---
faar man desværre ligesaalidt at vide her, som senere hos Martensen, da han
proklamerede den „genfødte“ Erkendelse som Teologiens Organ.

Det gaar ikke andre Videnskaber bedre end Filosofien. Grundtvig paastaar, at
hvor Kemi, Astronomi og Matematik ret blomstre, er det aandelige Liv sin
Undergang nær. Det er tiltrækkende for vantro Sjæle at forlyste sig ved dem, da
de bestyrke Forestillingen om, at Mennesket kun er et Legeme, der opløses ved
Døden, og da Syslen med de naturlige Aarsager til de store Himmelfænomener let
fører til at glemme Gud. I Historieforskningen, hvis Selvstændighed Grundtvig
havde negtet overfor Molbech, finder han ingen saadan Fare. Dog gælder det om
alle Videnskaber: naar de ville være selvstændige, er Kristendommen dem i
Vejen, og de ere Kristendommen i Vejen.

Det er dog ikke blot fra et religiøst Synspunkt, Grundtvig er betænkelig ved
Naturvidenskaben. „Jeg anser,“ siger han, „
% --- p. 63 ---
\setcounter{footnote}{0}%
\opage{63}Fysikens lidenskabelige Studium og højere Vurdering, især den
experimentales, der ender sig i Mekaniken, for et Tegn paa Videnskabelighedens
nære Grav.“ --- Her kundgør sig den romantiske Opposition mod den nyere
Naturvidenskab, saaledes som denne havde udviklet sig fra Galilei's Tid. Et nyt
Videnskabsbegreb skulde --- som vi have set hos Steffens --- træde i Stedet,
men blev ligesaalidt formuleret, som den „kristelige Logiks“ Principer.

Hovedsagen er dog for Grundtvig det religiøse Synspunkt. Det er jo dog
Religionen, der skal lære os Menneskets Vilkaar og Bestemmelse og Vejen til at
naa denne at kende, og som skal styre vor Attraa, Tanke og Gerning. Den maa da
ogsaa adspørges, før vi søge nogen anden Kundskab, og maa staa os levende for
Øje, medens vi sysle med Videnskaberne. For den Troende er Videnskaben enten
noget Afskyeligt eller ogsaa Religionens Tjener. Og ved Religionen tænkte
Grundtvig paa „den gammeldags, enfoldige Tro“, hvis Indhold findes i Bibelen,
forstaaet aldeles ordret. ---

Det var ikke blot Grundtvigs personlige Trang, der saaledes førte ham til at
fordre Videnskaben underordnet under Religionen. Herved medvirkede ogsaa hans
store Interesse for det Folkelige. Folkets aandelige Enhed vilde jo være
sikret, naar de Ulærde fandt deres aandelige Næring i Bibelen, tagen efter dens
Bogstav, og de Lærde bøjede deres Videnskab under ganske den samme Autoritet.
Det var Fortsættelsen af denne Tankegang, der senere (1825) førte ham til den
Opdagelse, at ikke Bibelen, men de tre Trosartikler vare den egentlige
Grundvold. Bibelen kan jo fortolkes paa forskellig Maade, og den kræver Lærdom
for ret at forstaas. Det mente Grundtvig ikke var nødvendigt for
Trosartiklernes Vedkommende, og kun ved, at de ligge til Grund, kunne de Ulærde
blive uafhængige af de Lærdes Vidnesbyrd.%
% footnote *) on p. 63
\footnote[1]{\textit{Kirkespejlet.} p. 343.} ---

Grundtvig har Ret i, at de Lærde ikke skulle være Autoriteter. De Lærdes Ret
til at høres beror paa deres Grunde og deres Beviser, og om disse slaa til, kan
ogsaa kun afgøres ved Grunde og Beviser. Tankens Verden staar aaben for
% --- p. 64 ---
\setcounter{footnote}{0}%
\opage{64}Enhver, og de Lærde, der ligesaavel ere en Del af Folket som de
Ulærde, have den Pligt at værne om Tankens Verden som en Del af Folkets
aandelige Goder; men denne Pligt kunne de kun øve, naar deres Arbejde kan gøres
i fuld Frihed og Selvstændighed. Det beror naturligvis paa de Ulærde selv,
hvorvidt de ville benytte den Vejledning, de Lærde kunne øve.

Sætningen om, at de Ulærde ikke skulle være afhængige af de Lærdes Vidnesbyrd,
betyder egentligt et Tilbagetog fra Standpunktet i Striden med Ørsted, hvor han
omvendt vilde gøre de Lærde afhængige af de Ulærdes Vidnesbyrd. Og dog havde
Grundtvig ogsaa senere ondt ved at forstaa, at man kan dyrke Videnskaben for
dens egen Skyld, og at den netop derved kan blive af stor menneskelig Værdi. I
Vinteren 1843 ---44 holdt han Foredrag over græske og nordiske Myter og udtalte
da, at han ikke begreb, hvorledes man kan interessere sig for „det uendelige
Rum og det kopernikanske System“, naar man ikke gik ud fra den Antagelse, at
Dyrekredsen og Mælkevejen, Sol, Maane og alle Stjerner ere til for Menneskets
Skyld. Han indvendte desuden mod Astronomien, at den var bleven for matematisk;
den var bleven Sort paa Hvidt istedetfor Guld paa Blaat!%
% footnote *) on p. 64
\footnote[1]{\textit{Brage-Snak om} græske og nordiske Myter og Oldsagn. 2.
Udg. p. 1; 165 f. --- Grundtvig skal stadigt have fastholdt Troen paa, at Solen
bevægede sig om Jorden. Smlgn. \emph{Sofus} \emph{Høgsbro}: \textit{Mit Forhold
til Grundtvig, Tscherning og Monrad.} 1902. p. 13. --- Grundtvig staar endnu
overfor Kopernikanismen, som Luther og Melanchton stode, da deres unge
begejstrede astronomiske Kollega Rheticus, Kopernikus' første Discipel, bragte
den nye Lære ind paa Wittenberg Højskole. Se \textit{Den nyere Filosofis
Historie.} 2. Udg. I. p. 103. --- Iøvrigt er det ikke Kopernikus selv, der har
bevirket en saa afgørende Forandring i Verdensanskuelsen; en saadan kom først
ved Giordano Brunos konsekvente Udvidelse af Kopernikanismen. Anf. Skr. p.
117---124.} ---

Fra Romantikens spekulative Ideer kom han i sine senere Aar bort. Paa sin baade
kostelige og geniale Maade tog han derimod i de nævnte Foredrag midt under
Hegelianismens Herredømme herhjemme Ordet for en Erfaringsfilosofi efter
engelsk Mønster. Odin, siger han, betyder Nordens Aand, og
% --- p. 65 ---
\setcounter{footnote}{0}%
\opage{65}Odin var, hvad de tyske Filosofer med en meget fornem Mine kalde en
Empiriker. Han vandrede viden om, og han saa sig langt omkring fra Hlidskjalf;
han havde desuden Hugin og Munin, og han raadspurgte Mimer. Det er meget
klogere med den gamle Odin og Engelskmanden at tage Livet og Verden, som de
ere, og at gøre os dem saa nyttige som muligt, end med den unge Filosofi og
Tyskerne at spilde sin Tid og sine Kræfter paa at skabe Livet og hele Verden om
efter sit Hoved, og saa springe i Flint, fordi det ikke vil lykkes, fordi det
Umulige ikke vil lade sig gøre.%
% footnote *) on p. 65
\footnote[1]{\textit{Brage-Snak.} p. 201---205.}

Det er en sund Realisme, der her lovprises; om det er den „kristelige Logik“,
Grundtvig i sin Ungdom profeterede, er et andet Spørgsmaal. Der ligger deri
ingen Løsning af det Spørgsmaal, han drøftede med H. C. Ørsted. Og med al den
Betydning, den fra ham udgaaende aandelige Retning har havt for vort Folk,
staar dog det Spørgsmaal stadigt aabent og vil maaske endog antage skarpere
Former end nogensinde før: hvorledes den aandelige Enhed kan bevares i Folket
trods Modsætningen mellem de „Ulærdes“ Tro og de „Lærdes“ selvstændige Videnskab.

\begin{center}\rule{2cm}{0.4pt}\end{center}

% --- p. 66 ---
\setcounter{footnote}{0}%
\opage{66}

\begin{center}\large 7. BAGGESEN OG GRUNDTVIG\end{center}

Faa Aar efter Sammenstødet mellem H. C. Ørsted og Grundtvig kom det
ejendommelige Sammenstød mellem Baggesen og Grundtvig. De to Digtere stødte
sammen i det store Spørgsmaal om Poesiens Forhold til Livet, --- eller rettere
i Spørgsmaalet om Værdimaalet for det Liv, Poesien skal udtrykke. Det trak op
til en Kamp i poetisk Form mellem store Tanker --- men uheldigvis trak det
over. Vor Literatur fik kun Heroldernes Signaler. Maaske kan Poesien heller
ikke give Mere.

Vi forlod Baggesen som sværmerisk Kantianer. Sværmeriet var stærkere end
Kriticismen, --- og han fandt snart, hvad han --- foreløbigt, som altid,---
søgte, i Jacobi's Følelsesstandpunkt. I et langt Brev til Reinhold (af 20.
Novbr. 1801) gjorde han Rede for sit nye Syn. „Tro, min kære Reinhold, er
højere end Viden --- er det Ufornuftige i Fornuften, Fornuftens egentlige
Kraft, uden hvilken dens Evne vilde være en blind Evne til tomme Former.“ Den
betegner en Afslutning af al Reflexion og Analyse og kan ikke selv analyseres.
Den er en Kraft, der ikke kan deles, tælles, maales eller begribes, --- men det
er den, der bærer Livet. I den aabenbarer sig det Hellige. Religionen skal ikke
blot tone i vor Viden, men i vor hele Sjæl. Man maa ikke borttage den Gysen,
den vækker, --- ikke berøve Stjernehimlen den Nat, i hvilken den ene er synlig.

Han ved meget godt, at og hvorfor han ikke duer til videnskabelig Folosof: „For
det Første har jeg mere Indbildningskraft end Forstand, mere filosofisk Fantasi
end filosoferende Fornuft... For det Andet er mit Sprog for billedrigt. For
% --- p. 67 ---
\setcounter{footnote}{0}%
\opage{67}det Tredie er jeg maaske allerede af Naturen altfor gennemtrængt af
det Sande, til at jeg skulde anstrenge mig for at søge Sandheden ad kunstig
Vej. For det Fjerde tror jeg ikke paa noget fuldendt System af menneskelig
Erkendelse, fordi jeg tror, at den menneskelige Aand aldrig kan staa stille, ej
heller skal staa stille.“

Disse sidste Ord gjorde han selv til Sandhed. Han trængte til at se sig mere om
i Verden, end Jacobi's Tro opfordrede til. Hans Reflexion arbejdede paa at
forbinde Natur, Moral og Religion i en stor Helhed. Jacobi havde overfor
Lessing indrømmet, at hans Trosstandpunkt kun kunde naas ved et Spring, et
salto mortale. I Modsætning til dette „den blinde Tros salto mortale“ vilde
Baggesen nu sætte „den seende Tros salto vitale“, --- og hans Bestræbelser, der
udfylde en meget stor Del af hans senere Aar, kende vi fra det efter hans Død
af hans Sønner udgivne \textit{Philosophischer Nachlass.} Han vil give en
spekulativ Teologi, der udvikler Begrebet om Gud som „alle Sandheders Sandhed,
paa én Gang aabenbaret i vor Samvittighed, i Stjernehimlens Glans og i det nye
Testamente“.%
% footnote *) on p. 67
\footnote[1]{\textit{Philosophischer Nachlass} (1858 ff.) I. p. 409. 465 p.}
Han betragtede sig nu som Filosof og vilde gerne have været filosofisk
Professor i København, da her i Aaret 1813 var en Post ledig. Han skriver endog
til En af sine Sønner i Aaret 1811: „Ikke Poesi, men Filosofi har fra min
Barndom været mit Fag og Genstand for mine Bestræbelser“, og endnu bestemtere
udtaler han sig i denne Retning i et Brev til en af sine Venner ti Aar efter.%
% footnote **) on p. 67
\footnote[2]{Disse to Breve findes ikke i Baggesens Biografi, men citeres i Kr.
\emph{Arentzen's} V\textit{ærk Baggesen og Oehlenschläger.} VI. p. 34. ---
Dette Værk har været mig til meget stor Nytte, især ved Studiet af Forholdet
mellem Baggesen og Grundtvig. Det er Frugt af mange Aars taalmodigt og af
inderlig Begejstring for vor Literaturs Stormænd baaret Arbejde. Ubehjælpsomt
og sammenstykket i sin Form er det dog en Skatkiste for dem, der beskæftige sig
med denne Tids Literatur. En mærkværdig Modsætning er der mellem dets Form og
dets livlige og entusiastiske Forfatter, der havde en sjelden Evne til at vække
og begejstre de Unge, hvis Lærer i Dansk han var, og som ofte bleve hans
Venner. Den Brug, jeg ved nærværende Arbejde har gjort af hans Bog, har fornyet
min taknemmelige Erindring om min Lærer og Ven.} --- Baggesens psykolo%
% --- p. 68 ---
\setcounter{footnote}{0}%
\opage{68}gisk-teologiske Reflexioner beholdt stedse en individuel Karakter,
hvad han ogsaa vedkendte sig: „enhver Snegl maa have sit eget Hus!“ Der findes
i hans „Philosophischer Nachlass“ mange aandfulde Reflexioner, mange Vidnesbyrd
om alvorlig Eftertanke og Tilløb til religiøs Spekulation. Han bevæger sig, som
Sibbern har sagt om ham, paa „den med Skarpsindighed, Vid og Lune sig udtalende
filosoferende Reflexions Omraade“.%
% footnote *) on p. 68
\footnote[1]{I „Zeitschrift fur Philosophie“ 1859 findes en intesessant
Anmeldelse af Baggesens „Philos. Nachlass“ ved \emph{Ulrici}. --- En teologisk
Bedømmelse af Baggesens Forsøg paa Kristendomsfilosofi giver \emph{Martensen}:
\textit{Af mit Levned.} III. p. 162---172. Martensen anerkender ham som
„kristelig Tænker“, men beklager dog, at han ikke ret er trængt ind i
Treenighedslæren, og at han negter, at vi kunne erkende Gud, som han er i sig
selv. Saa langt var nemlig den „kristelige Tænken“ naaet paa Martensens Tid.}

I Baggesens Poesi fik denne Baggrund af alvorlig Eftertanke af og til
dybsindige Udtryk, --- saaledes i „Gengangeren“ (1807), hvor det fordres, at
Digteren skal erkende sit guddommelige Kald

I evigt Hang til det Uendelige

og i uendelig og evig Lyst

til Sammenstemning af det Endelige

med Evighedens Grundklang i hans Bryst.

Men ofte legede han med Digtningen, som han legede med Tanken --- og det var
dette, der fremkaldte Sammenstødet mellem ham og Grundtvig. I Striden mod
Oehlenschläger sympatiserede Grundtvig med Baggesen og blev angreben af dennes
Modstandere. Han mente, ligesom Baggesen, at den store Digter var gaaet tilbage
og ikke længere opfattede sit Kald med Alvor. Og begge toge Ordet for Kunstens
Sammenhæng med Livet og dets Gaader. Men Grundtvig fandt, at hans Forbundsfælle
drev sin Kritik og Polemik for meget som en Leg. I sine „Rimelige Strøtanker
ved Jens Baggesens Grav“ (1815) siger Grundtvig om ham:

% --- p. 69 ---
\setcounter{footnote}{0}%
\opage{69}Det var for ham da al Ulykken,

for meget havde han paa Nakken

og altfor Lidt i Ryggen!

Havde han --- mente Grundtvig --- havt en Kirkemur (som den han, Grundtvig,
selv havde fundet) i Ryggen, vilde han have været bedre faren overfor sine
Modstandere.

Og i Baggesens bekendte Gaade („Det evige Sindbillede, ikke blot Literaturens
Oprindelse, men Filosofiens Ende. En Gaade som Prisopgave for aandeligt
Levende“) kunde Grundtvig ikke se andet end Tanke- og Ordleg med det Højeste.
Det Sindbillede, Baggesen mente, synes (ifølge et Brev til En af hans Sønner)
at have været Korset som det, der er dannet af to Linier og derfor udtrykker
den store Dobbelthed eller Modsætning, der aabenbarer sig i Alt og hindrer
enhver ad en eneste Linie forløbende Tankegang i at udtømme Tilværelsen. Korset
paa Golgata er et enkelt, stort Exempel: Gud selv aabenbares i Lidelsen. Disse
Sindrigheder brød Grundtvig sig ikke om. For ham var Gaaden én Gang for alle
gættet i Bibelen, og Baggesens Gækkeri var ham en Forargelse. Han tilraabte ham:

Du est dig selv en Gaade,

det har jeg længe vidst,

og til den ret at raade

du har kun stakket Frist! ---

Hvad er din Gaade? Kun et Tidsfordriv?

Du har dog nu ej megen Tid at spilde!

En Skygge af dit hele Skjaldeliv?

Maaske, men sandelig, da var det ilde!

Et Billede af Livet i det Hele?

Bestaar da det i sine brudte Dele?

I Baggesens Svar staa de stolte Linier:

Vel kender jeg til Bogen,

hvorpaa du peger hen,

maaske saa godt som Nogen,

og Løsningen i den;

% --- p. 70 ---
\setcounter{footnote}{0}%
\opage{70}men skal af den og øses,

hvad slukker Tankens Tørst,

maa gamle Seglet løses,

og Brønden aabnes først.

Omsonst af den du venter

Allivets Kildevand,

hvis ej din Aand det henter

dybt med Fornuftens Spand.

Baggesens „seende Tro“ kunde ikke lade sig nøje med Henvisningen til „Bogen“.
Han vilde jo sammenarbejde Bibelordet med Naturerkendelsen og med
Samvittighedens Udsagn. Hans Tanke kunde ikke gaa frem ad én Linie. Han mente
vel ikke, at det Hele kun bestod i brudte Dele, men at det er det brudte Lys,
vi kende, og at Eftertanken maa samle Delene, uden at kunne naa den fulde
Enhed. Det var Grundtanken i hans filosofiske Reflexioner, og det var Tanken
med hans Gaade, her dog skjult og forstyrret med alskens Ordleg.

Kampen kom nu til at staa om det Spørgsmaal: staar den Livsanskuelse højest,
der støtter sig til Tanken, medens den dog tager fra de religiøse
Overleveringer, hvad den kan forene med Naturkundskaben og med Samvittigheden,
--- eller den, der ene holder sig til den overleverede Tro i fuld Underkastelse?

Baggesen mente at staa højest, fordi han vidste og skuede, hvad hans Modstander
troede, og fordi han let kunde sætte sig paa dennes Standpunkt, medens det
Omvendte ikke var Tilfældet:

Fra Stadet, hvor jeg staar, jeg vel i dit mig ned kan sætte; men
Opgangstrinnene fra dit til mit er ej saa lette!

Grundtvig svarer med en Henvisning til Dommedag:

Om virkelig du staar paa Tempeltind,

og om jeg op til dig har mange Trin,

det klarer sig engang vel paa det Sidste.

% --- p. 71 ---
\setcounter{footnote}{0}%
\opage{71}Men ser jeg ret ved Ordets Lampeskin;

jeg mod din Højde ikke gør et Trin,

og aldrig skal din Modenhed mig friste.

Aldeles afset fra Baggesens og Grundtvigs specielle Standpunkter er det et
betydningsfuldt Spørgsmaal, der kom op imellem dem: hvad er højt og hvad er
lavt i Aandens Verden? hvilken Maalestok have vi her? --- Det er naturligvis
ingen Rangstrid mellem Personer, det gælder, men det gælder at finde Værdiernes
Skala. Grundtvig henviser til Bibelen og til Verdensdommen, og Baggesen bygger
paa sin religiøse Spekulation. Men skulde der ikke være en rent menneskelig
Maalestok for Livets Værdier?

Det saa et Øjeblik ud, som om Kampen mellem de to store Digtere skulde komme
ind paa dette Spørgsmaal for Alvor. Baggesen udæsker til fortsat Kamp, ti „i
Aandeverdenen bør der være Krig!“ --- og han vil ingen Fred, før Alt er klaret.
Og Grundtvig byder ham velkommen til Kampen („Velkommen du, som kender Aanders
Lov!“) og betegper nærmere, hvad der kæmpes om og hvorledes der skal kæmpes.
Det er ikke Digterkransen, de kæmpe om, ---

Nej, hvad er Lys og Aand, og hvad er Mørke,

ja, hvad er Sandhed, hvad er Tant og Løgn;

hvad man skal tro, hvem man skal trolig dyrke, ---

er Kundskab hellig, eller er den søgn? ---

Og han vil lægge sin Bibel bort. Naar det gælder om, hvad og hvorfor man skal
tro, kan man jo heller ikke begynde med at beraabe sig paa sin Trosbekendelse,
som den anden Part ikke deler. Her udtaler han de store og smukke Ord, der egne
sig til Motto ved al Strid om Livets største Spørgsmaal.

Ej spørge her vi stort om Konfessioner,

men søge Grunden til de brudte Toner

og spejde Lys med aandelig Fornuft!

% --- p. 72 ---
\setcounter{footnote}{0}%
\opage{72}Dette varslede godt. Men det blev ved Varslet. Kampsignalet blev
givet, men ikke fulgt. Baggesen begyndte igen sin aandelige og legemlige
Omflakken, og Grundtvig fordybede sig i historisk Gransken og i Udforming af
sine kirkelige Anskuelser, hvor snart Garantierne for, hvad man tror, kom til
at spille en større Rolle end Aarsagerne, der føre til Tro. Konfessionen
traadte i Forgrunden, den aandelige Fornuft i Baggrunden. Først senere blev i
dansk Aandsliv den Drøftelse optagen, til hvilken de poetiske Herolder havde
opfordret hinanden, uden selv at gøre Alvor af at kæmpe Striden til Ende.

\begin{center}\rule{2cm}{0.4pt}\end{center}

% --- p. 73 ---
\setcounter{footnote}{0}%
\opage{73}

\begin{center}\large 8. ANDERS SANDØE ØRSTED\end{center}

Af de to Dioskurer, der herhjemme mødte det nye Aarhundrede med størst
Modenhed, have vi allerede omtalt den ene, Naturforskeren. Juristen, Anders
Sandøe (1778---1860) var et Aar yngre. Hans Tankeliv frembyder den samme Enhed
i sin Udvikling som Broderens. Hans Christian var den lykkeligste, idethele en
af de lykkeligste Skikkelser i dansk Aandsliv. Anders Sandøe var tungere; hans
Tanke bevægede sig frem gennem grundig Prøvelse af alle Enkeltheder og
Muligheder. --- Det er her ikke som Skaberen af dansk Retsvidenskab, han skal
omtales, men som den filosofiske Tænker.

Han levede i Kredsen af de Bedste herhjemme. Oehlenschläger skyldte de to
Brødre, at han var saa vel forberedt til at modtage Steffens' Paavirken.
Oehlenschlägers Søster blev Ørsteds Hustru, og i det Ørstedske Hus bleve de
Værker, den unge Digter sendte hjem fra Rejsen, forelæste, før de udkom. Med
sin Broder levede Ørsted i et inderligt og paa store fælles Interesser bygget
Samkvem, og Mænd som Mynster og Sibbern hørte med til Kredsen. Han blev vraget
af Universitetet, --- en af dettes Synder, som dog vendtes til det Gode, da han
derved til Gavn for sin Videnskab og sin Udvikling førtes ud i praktisk
Virksomhed som Dommer og Lovgiver. Senere beklædte han høje Embeder, som
Generalprokurør, Kancellideputeret og som Minister.

Som tidligere omtalt var han i sine unge Dage en ivrig Kantianer. Snart fik
Fichte stor Indflydelse paa ham, især ved sit Forsøg paa at grundlægge en af
Etiken uafhængig Retslære. Ørsted var en personlig Ven af Fichte, der boede
% --- p. 74 ---
\setcounter{footnote}{0}%
\opage{74}i hans Hus under sit Besøg i København i Sommeren 1807. Skønt han
ikke fortsatte sine filosofiske Studier, optagen, som han blev, af praktisk
Arbejde, bevarede han dog de Grundtanker, han havde tilegnet sig, hele sit Liv
igennem. Spekulative Anlæg mente han ikke at have, og han saa med Mistro paa
den tyske Spekulations (Schellings og Hegels) Indtrængen her i Landet. Derimod
mente han at besidde Evner i Retning af logisk og etisk Tænkning. Om Kants
Indflydelse paa ham og hans Samtid har han i sin Alderdom aflagt et smukt
Vidnesbyrd:%
% footnote *) on p. 74
\footnote[1]{\textit{Af mit Livs og min Tids Historie.} 1851. p. 34.}

„Skønt jeg forlængst er kommen tilbage fra det Afguderi med Kants Lære, hvormed
jeg indtraadte paa Forfatterbanen, saa har dog dens rene og kraftfulde Moral
haft en Indvirkning paa mit Sind, som aldrig har tabt sig. Ogsaa er jeg
overbevist om, at det Samme er Tilfældet med en stor Del af den Slægt, som til
samme Tid fik sin Manddoms Dannelse. Men selv paa den senere Slægt, for hvem
den Kantiske Lære kan synes efterhaanden at være bleven fremmed, har den
sikkert middelbart havt en uberegnelig Virkning; ti skønt hans Moral i den
Skikkelse, hvori han overleverede den, har, og det langt mere end vedbørligt,
tabt sin Anseelse, saa har dog den Grundtone, som gik igennem samme, længe
bevaret sig under den forandrede Skikkelse, som Moralen senere har
modtaget..... Det er af Mange blevet erkendt, at Kant og Fichte, hvilken Sidste
i Begejstring for ren, al Selvkærlighed forsmaaende Sædelære i fuldt Maal
sluttede sig til Kant, have virket meget til det Opsving, som Preussen tog i de
Aar, som fulgte nærmest efter den ulykkelige Krig i 1806 og 1807. Da en af de
Statsmænd, som havde havt Del i, at lede Preussens Genfødelse, og som stedse
var blevne den Aand tro, hvorfra den gik ud (Statsminister v. Schön) i
Anledning af sit Embedsjubilæum modtog en skøn og varm Hyldest af sine
Medborgere, vidnede han, at dersom han havde udrettet noget Gavnligt for sit
Fædreland, saa skyldte han det til den Tænkemaade, som hans Lærer Kant havde
indpodet ham,
% --- p. 75 ---
\setcounter{footnote}{0}%
\opage{75}og at han derfor maatte føre denne Tak, man havde bragt ham, tilbage
til den Kilde, hvorfra han kun var en liden Bæk. Uden iøvrigt at kunne gøre
nogen Sammenligning, gribes jeg af en lignende Følelse ved at tænke over, hvad
Kant har været mig og mange Andre i mit Fædreland.“ ---

Hvad der giver A. S. Ørsted en interessant Plads i den danske Tænknings
Historie, er den Betydning, som Forholdet mellem Tanken og Erfaringen, mellem
almindelige Sætninger paa den ene Side og Virkelighedens mangfoldige og
skiftende Indhold paa den anden Side, har havt i hans Udvikling --- baade som
Filosof og som Jurist. Det var hans praktiske Arbejde som Dommer og Lovgiver i
Forbindelse med Studiet af forskellige Landes Lovgivning og med Indflydelsen af
Filosofer som Fichte og Treschow, der bragte ham til at lægge Vægt paa
Erfaringen og til at se, at Tanken hverken i filosofisk Systematisering eller i
Lovgivning kan udtømme Livets Rigdom. Dens Betydning er at orientere i Livet,
ikke at skabe Livet. --- Jeg fremhæver nogle af de Hovedpunkter, som i denne
Henseende ere af Interesse.

\begin{center}\rule{2cm}{0.4pt}\end{center}

Det var efter Ørsteds eget Udsagn, i hans Skrift \textit{om
Trykkefrihedslovgivningen,} der udkom 1801, at hans ændrede Tankegang først
udtaltes. Det havde fra først af været hans Opfattelse, at Retsordenens
Grundlag ligger i Nødvendigheden af, at den Enkeltes Frihed sikres saaledes, at
den ikke blot beror paa de Andres gode Villie. Der kræves da en Magt; som
sætter Grænser for Overgreb. Men de Midler, ved hvilke denne Begrænsning skal
naas, maa bestemmes ved Folkets hele moralske Tilstand til den givne Tid, og
ingen Lov kan formulere Regler, som uden videre kunne anvendes paa det enkelte
Tilfælde. Dommeren maa ved et Skøn afgøre, om dette hører ind under Loven: „Det
rene Faktum, saaledes som det viser sig for Sanserne, lader sig ikke
umiddelbart henføre under et Begreb, hvortil Loven kan knytte en Straf eller en
anden Retsvirkning, men Domstolene maa først ved
% --- p. 76 ---
\setcounter{footnote}{0}%
\opage{76}en paa alle Forholdenes Overvejelse grundet Akt af Dømmekraften
erkende, at det givne Faktum har den retlige Egenskab, som Lovbuddet forudsætter.“%
% footnote *) on p. 76
\footnote[1]{\textit{Af mit Livs og min Tids Historie.} I. p. 91. (Ørsted giver
her selv et Resumé af det nævnte Skrift om Trykkefrihedslovgivningen).} I en
Afhandling fra samme Aar \textit{om Regeringens Ret til at ophæve eller
forandre Stiftelser, som private Mænd have oprettet,} udtaler Ørsted, at man
ved at undersøge hvilketsomhelst Forhold af moralsk Natur vil overbevise sig
om, at almindelige Grundsætninger ikke ere nok til at ordne det, naar man ikke
understøttes af en sund praktisk Takt og en redelig Villie.%
% footnote **) on p. 76
\footnote[2]{\textit{Eunomia.} Afhandlinger henhørende til Moralphilosophien,
Statsphilosophien og den dansk-norske Lovkyndighed. I. p. 37.} Og senere kommer
han, som vi skulle se, udførligere tilbage til dette Emne.

En absolut Ejendomsret gives der ikke. Ejendomsretten grunder sig paa, at den
Enkelte maa have Midler til fri Udvikling, til at naa sit fornuftige Maal. Men
ved denne Begrundelse sættes der straks en Grænse. Ejeren kan ikke have Ret til
at træffe evigt gældende Bestemmelser for Anvendelsen af sin Ejendom. Det vilde
kunne føre til, at de Døde kom til at herske over de Levende. Testamenter til
Stiftelser kunne da kun anerkendes, forsaavidt de i de enkelte Tilfælde findes
at stemme med det almindelige Vel og den almindelige Villie.%
% footnote ***) on p. 76
\footnote[3]{Smlgn. \textit{Af mit Livs og min Tids Historie} I. p. 118, hvor
Ørsted henviser til, at en Forordning af 1781 erklærede alle
Villiestilkendegivelser for ugyldige, der kunde hindre Landets Velfærdssag,
Jordfællesskabets Ophævelse.}

Som man ser, søger Ørsted at sætte en erfaringsmæssig Begrundelse af Moral og
Ret istedetfor den rene Fornuft. Det er Hensynet til, hvad i hvert enkelt
Tilfælde „det almindelige Vel og den almindelige Villie“ kræver, der bliver det
Afgørende. Udtrykkene kunne kun forstaas saaledes, at det er Samfundets mod
„det almindelige Bedste“ stræbende Villie, hvorover der --- af Regeringen ---
skal skønnes i de enkelte Tilfælde. Udtrykkene minde baade om Bentham og om
Rousseau, hvis „volonté générale“ egentligt kundgør sig i Kants kategoriske
Imperativ. Hine Udtryk ere karakteri%
% --- p. 77 ---
\setcounter{footnote}{0}%
\opage{77}stiske for en Kantianer, der har indset Nødvendigheden af
Erfaringshensyn i Etik og Retslære.

I en Afhandling fra 1807 \textit{om Religion og Stat} (Eunomia I) hævder han,
at alle de Betingelser for et Folks Velfærd, der kunne formuleres som en
almindelig Ordning, bør drages ind under Retsordenen. Altsaa tjener Retten ikke
blot (som Ørsted i tidligere Arbejder havde antaget) til at sikre den Enes
Frihed overfor de Andres, men faar en mere positiv Betydning. Udelukket af
Statsøjemedet er kun, hvad der bedre iværksættes ved Enkeltes frie, til sig
selv overladte Omsorg. Men en aldeles betryggende retslig Mekanisme er ikke
opnaaelig. Samvittighedsfuld Overvejelse af de foreliggende Forhold maa søge at
anvende de almindelige Love paa hensigtsmæssig Maade.

Ved de Forandringer, der saaledes vare komne ind i Ørsteds Opfattelse af Moral
og Ret, havde han brudt saavel med den gamle Naturret som med Kants og Fichtes
Teorier. Men derimod opgav han ikke Muligheden af en Retsfilosofi. Det blev saa
Opgaven for en saadan at undersøge, hvorledes Retsforholdene skulle ordnes,
naar der tages Hensyn saavel til at sikre de Enkeltes Frihed og Personlighed
som Samfundets fremskridende Velbefindende, dette Ord taget i omfattende
Betydning, saa at ogsaa de aandelige Livsgoder indgaa derunder. Sammenlignende
Retslære saavel som Nationaløkonomi faa stor Betydning for en saadan
Retsfilosofi, der vil blive frugtbarere end den gamle Naturret med dens
abstrakte Principer og dens kunstige Deduktioner. Forarbejder til en saadan
Videnskab finder Ørsted i Montesquieu's, Filangieri's og Benthams Værker. Selv
mener han i sit Værk om sit Livs og sin Tids Historie at give Bidrag dertil.%
% footnote *) on p. 77
\footnote[1]{\textit{Af mit Livs og min Tids Historie.} I. p. 145---150.}

Det var i de foregaaende Arbejder især Retslæren, Ørsted havde behandlet. For
Etikens Vedkommende har den mærkelige Afhandling \textit{Om Teori og Praxis i
Sædelæren} (1811) (optrykt i Eunomia I) en lignende Betydning, som Afhandlingen
om Religion og Stat havde for Retslæren. Ogsaa her hævdes,
% --- p. 78 ---
\setcounter{footnote}{0}%
\opage{78}at der er et Mellemrum mellem det almindelige Begreb og det enkelte
Tilfælde, der skal gaa ind under det, et Mellemrum, som kun kan udfyldes af en
sund Dømmekraft.

Nogle Aar iforvejen havde \emph{Niels} \emph{Treschow}, som nu var Professor
ved Universitetet, udgivet en interessant Afhandling%
% footnote *) on p. 78
\footnote[1]{Det danske Videnskabernes Selskabs Skrifter. Femte Bind. 1810.},
hvis Titel var: \textit{Gives der noget Begreb eller nogen Ide om enslige} [ɔ:
enkelte, individuelle] \textit{Ting?} I denne mod Schelling rettede
Undersøgelse hævder Forfatteren, at de enkelte Væsener ere Mere end flygtige
Billeder af det Absolute. Intet alment Begreb udtømmer deres Natur. Det er ikke
Arterne, men Individerne, der ere det egentligt Reale, og der findes neppe
faste Grænser mellem Arterne. Det er vor Forstands Begrænsning, der gør, at vi
ikke kunne trænge ind til det Ejendommelige, Enkelte og Individuelle.
Almenbegreberne ere kun Hjælpemidler til Oversigt.

Til denne Afhandling henviser Ørsted, som dog allerede var kommen ind i en
lignende Tankegang ad sin egen Vej. Det er paa det etiske Omraade, han hævder
det Individuelles Betydning. Ingen almindelige moralske Bud kunne udtømme, hvad
Forholdene kræve i det enkelte Tilfælde, og navnligt kunne de ikke lade den
Enkeltes Ejendommelighed komme til dens fulde Ret. Det Aandelige antager en
særegen Skikkelse i hvert enkelt Menneske. Den, der vilde ophæve sin
Ejendommelighed og søge at forvandle sig til et abstrakt Fornuftvæsen, vilde
arbejde paa at oprykke Livets Rod og gøre sig selv til et marvløst
Skyggebillede. Det Bidrag til Verdensplanens Udførelse, den Enkelte skal give,
kan han kun give i den for hans aandelige Natur ejendommelige Form. Foruden de
Opgaver, han har fælles med Andre, har han sin særegne Opgave, sin specielle
Etik. Ikke blot har hver Enkelt sit Kald, men den almindelige Kamp mod Daarskab
og Ondskab maa Enhver føre paa sin Maade. Hvis den Sagtmodige vilde paatage sig
den Djærves, og denne hins Rolle, kunne de begge være visse paa at tale og
handle forgæves. Desværre have vi her intet Middel til Konstruktion eller
Udmaaling. Det er ikke engang muligt at henføre en bestemt, ved Iagt%
% --- p. 79 ---
\setcounter{footnote}{0}%
\opage{79}tagelse given Individualitet til et Begreb: hvor længe jeg end bliver
ved at føje Kendemærke til Kendemærke, vil jeg ikke kunne udtømme den uendelige
Mængde af Nuancer og Grader. Individualiteten sætter en Grænse for vor Stræben
efter at gennemskue Sammenhængen mellem Enhed og Mangfoldighed.

Ingen Moralteori kan opstille Begreber, der ere tilstrækkelige til at løse alle
de Opgaver, som Livets mangfoldige Forviklinger medføre. Derved bortfalder ikke
Betydningen af at finde Ideer og Typer; men det videnskabelige Arbejde maa her
tilsidst gaa over til en Kunst. Videnskaben kan forberede denne Kunst, men ikke
frembringe den. To forskellige Drifter virke hos os: til at finde Enhed i Livet
som i Naturen --- og til at lade Livet udfolde sig frit og ejendommeligt. Disse
to modstridende Drifter egne sig godt for et Væsen, hvis Bestemmelse er en
stadigt fremskridende Udvikling, ikke en fast Besiddelse. ---

Hos Holberg have vi allerede mødt en lignende Tankegang, og vi ville møde den
senere hos andre danske Filosofer. Afhandlingen om Teori og Praxis i Sædelæren
er ikke blot et interessant Bidrag til Etikens Teori, men ogsaa et
karakteristisk dansk Bidrag. Hvad Ørsted angaar, udtaler han 40 Aar senere i
sin Selvbiografi, at han endnu stadigt ganske vedkender sig Afhandlingens
Indhold. ---

\begin{center}\rule{2cm}{0.4pt}\end{center}

Ørsteds Forfattervirksomhed stansedes for en lang Tid i Aaret 1826. Han havde i
Juridisk Tidsskrift skrevet et Indlæg i den bekendte kirkelige Strid mellem
Grundtvig og H. N. Clausen, under Titel: \textit{Behøver den danske
Kirkeforfatning en omgribende Forandring?} Ørsted svarer Nej. Han hævder
overfor Clausen, at det ikke var nødvendigt at ændre de symbolske Bøger, ---
ligesaalidt som det var berettiget (hvad Grundtvig krævede) at udelukke Enhver
fra et kirkeligt eller teologisk Læreembede, der ikke tog Bibelen og de
symbolske Bøger ordret. Han gjorde her en lignende Tanke%
% --- p. 80 ---
\setcounter{footnote}{0}%
\opage{80}gang gældende, som han 1798 havde hævdet overfor Bibelog
Fornuftpartiet, og en Tankegang, der var analog den, der laa til Grund for
Afhandlingen om Regeringens Ret til at ophæve eller ændre private Stiftelser.
Han mente, at Fortolkningen maatte skifte med de skiftende Tider; bogstavelig
Tilslutning var ikke nødvendig, naar kun det Væsentlige anerkendtes. Som
Punkter, hvor en bogstavelig Opfattelse ikke var nødvendig, nævner Ørsted:
Syndefaldet, Treenigheden, Inspirationen og de evige Helvedesstraffe.

Kong Frederik VI fandt, at en høj juridisk Embedsmand ikke burde tage Del i
Diskussionen om en Sag, i hvis Behandling for Domstolene han maaske kunde komme
til at tage Del. Ørsted maatte love at ophøre med sin Forfattervirksomhed.%
% footnote *) on p. 80
\footnote[1]{Se om denne Sag L. Koch: \textit{Anders Sandøe Ørsted.} 1896. p.
85---96.} ---

Om Ørsteds senere politiske Virksomhed henviser jeg til L. Kochs smukke
Biografi, samt til min Karakteristik af Brødrene Ørsted (i „Vort Folk“).

I Trediverne var han „Kongens og Folkets Mand“ og som saadan Præsident i
Stænderforsamlingerne. Vendepunktet i hans Forhold til sine Landsmænd kom, som
han selv har udtalt%
% footnote **) on p. 80
\footnote[2]{\textit{Af mit Livs og min Tids Historie.} IV. p. XLII.}, ved hans
Optræden (i Viborg Stænder 1842) i Sagen om det danske Sprogs Ret i Slesvigs
Stænder. Han, der saa stærkt havde indskærpet Erfaringens og Individualitetens
Betydning og Ret, saa ikke Betydningen af den vaagnende Nationalitetsfølelse og
af de nationale Forskelligheder. Modsætningerne --- de religiøse, nationale og
sociale --- havde skærpet sig i Løbet af Aarhundredet, og det var intet Under,
at selv højt begavede Mænd, der vare opvoxede under de mere harmoniske Forhold
ved Aarhundredets Begyndelse, ikke kunde tage det rette Stade i disse
Modsætningers Strid. Ørsteds tragiske Skæbne som historisk Personlighed har dog
ikke hindret Anerkendelsen af hans Forskergerning. I Bissens smukke Statue er
det den grublende Tænker, der er fremstillet.

\begin{center}\rule{2cm}{0.4pt}\end{center}

% --- p. 81 ---
\setcounter{footnote}{0}%
\opage{81}

\begin{center}\large 9. F. G. HOWITZ\end{center}

Paa en enkelt Undtagelse (Joh. Boye) nær vare siden Holbergs og Sneedorffs Tid
de, der filosoferede herhjemme, aldeles overvejende paavirkede af tysk
filosofisk Literatur, og især var det Tilfældet i det nittende Aarhundredes
første Tiaar, da den tyske Romantik førte Ordet. Dertil bidrog ogsaa, at der i
engelsk og fransk Filosofi herskede Ebbe --- efter Hume i England og efter
Rousseau og Diderot i Frankrig. Som et forfriskende Pust kommer da Lægen Franz
Gothard Howitz's Optræden i Aaret 1824. Howitz (f. 1789, d. 1826) var Professor
i Retsmedicin og skrev i Anledning af en Sag om en anklaget Kvindes
Tilregnelighed en Afhandling \textit{(Om Afsindighed og Tilregnelse,} et Bidrag
til Psychologien og Retslæren), der vakte en Strid om „Villiens Frihed“, i
hvilken A. S. Ørsted, Mynster, Sibbern, J. L. Heiberg o. fl. toge Del. Howitz's
Modstandere stode alle paa et kantiansk eller tysk-spekulativt Standpunkt, og
han fandt i Stridens Løb gentagne Gange Lejlighed til at henvise til engelsk og
fransk Tænkning, saavel som til Spinoza, hvis Standpunkt han nærmest gør til
sit. Han gør opmærksom paa, hvor Meget der stadigt er at lære af hine ældre
Tænkere, og han mener, at selvom man ikke kan negte den tyske Forsken
Betydning, saa har tysk Literatur i Begyndelsen af det nittende Aarhundrede
taget en Retning henimod det Dunkle og Fantastiske, der har skadet ikke blot
Filosofien, men ogsaa andre Videnskaber. Vor Nationalkarakter er desuden
forskellig fra den tyske; ligesom den hollandske kan den siges at danne en Art
Overgang mellem den tyske og den engelske. Men fremfor
% --- p. 82 ---
\setcounter{footnote}{0}%
\opage{82}Alt gælder det, at lade Erfaringen komme til sin Ret overfor
Spekulationen, og han henviser til, at der er Tegn i selve Tyskland til en
Vending i den Retning, ogsaa paa det filosofiske Omraade.%
% footnote *) on p. 82
\footnote[1]{\textit{Determinismen eller Hume imod Kant.} Et philosophisk
Forsvar for Afhandlingen om Afsindighed og Tilregnelse. 1824. Fortale. --- Det
er Filosofen Beneke, der sigtes til i de sidst anførte Ord. Smlgn. anf. Skr. p.
76 og \textit{Ultimatum} (1825). p. 50. --- I „Determinismen eller Hume mod
Kant“ giver Howitz en Oversættelse af Hume's fortrinlige Afhandling om Frihed
og Nødvendighed.}

Naar Striden mellem Howitz og hans Modstandere kom til at dreje sig om „Frihed“
og „Nødvendighed“, var det dog ikke denne Side af Sagen, Howitz havde lagt
Hovedvægten paa i sin første Afhandling. Hans Hovedformaal var at indskærpe, at
der gives Mellemformer og Overgange mellem Afsindighed og almindelig
menneskelig Fornuftighed, og at der ikke i Naturen er den skarpe Forskel,
Juristerne plejede at antage. Hans Tankegang her er egentlig en Fortsættelse af
Treschows og A. S. Ørsteds. Mennesket danner sig, siger han, skarpt afsondrede
Begreber, opretter Definitioner og foretager Inddelinger, men i den virkelige
Verden smelte Tingenes Grænser sammen, og der er en Kædeforbindelse, der
tilvejebringes ved idelige Overgange. De systematiske Naturforskere finde det
ofte vanskeligt at fastsætte Grænserne mellem forskellige Arter eller Grupper,
som f. Ex. mellem Planter og Dyr; men Vanskelighederne maa blive endnu større,
hvor vi --- som hvor det drejer sig om Menneskets psykologiske Natur --- ikke
have helt gode, haandgribelige Kendemærker at holde os til. Der er mange Grader
mellem Frihed og Ufrihed, mellem menneskelig forstandig Normaltilstand og
Afsindighed, mellem Tilregnelighed og Utilregnelighed. Disse Modsætninger gælde
kun, naar man tager Tilfælde, der ere Extremer i forskellige Retninger, men
ikke for de mellemliggende Tilfælde. Der er saaledes mange Grader af
Afsindighed, især naar man tager Hensyn til, at Afsindighed ikke altid bestaar
i virkelig Mangel paa Fornuft. Der kan være sygelige Tilstande, hvor Fornuften
Intet fejler, men hindres i at
% --- p. 83 ---
\setcounter{footnote}{0}%
\opage{83}øve sin Indflydelse; --- fixe Affekter gives der ligesaa vel som fixe
Ideer. Der kan være pludseligt frembrydende Tilskyndelser, som staa i den
skarpeste Modsætning til Individets hele Karakter, og som det ikke kan hæmme,
uagtet det kan være sig sin Tilstand klart bevidst. Af denne Art er den af
Pinel beskrevne mania sine delirio og den saakaldte folie raisonnante, hvor
Patienten med stor Kløgt véd at skjule sin Tilstand og besvare alle
Indvendinger, der rettes mod hans Forestillinger og Planer. Og som der er mange
Arter af Sindssygdom, er der ogsaa mange Overgange mellem Afsindighed og
almindelig Menneskefornuft. De lyse Mellemrum vidne derom; ligeledes de lavere
Grader af alle Arter af Afsindighed og den Afsindighed, der er begrænset til
ganske enkelte Omraader, samt Begyndelsesstadiet og Rekonvalescensen for den
enkelte Sindssygdoms Vedkommende. I sædvanlige Menneskers Halvsøvn, i de med
halv Bundethed fortsatte Morgendrømmerier, der især for nervesvage Mennesker
ere saa tillokkende, have vi ogsaa Overgangstilstande; i denne Morgenskumring
kunne Skrækkebilleder fra Drømmen overføres paa virkeligt tilstedeværende
Personer og føre til Angreb paa disse. Endeligt frembyde saakaldte normale
Menneskers Lidenskaber og Daarskaber ofte Træk, der gøre det vanskeligt at
sætte en skarp Grænse mellem det Sunde og det Syge.

Og som Afsindigheden har sine Grader, har Friheden det ogsaa, naar man tager
dette Ord i den simple Betydning af Evne til at gøre, hvad man vil eller
ønsker, uafhængigt af fremmed Bestemmelse og fremmed Villie. I denne Betydning
har Friheden en ubestridt Virkelighed --- men er i de enkelte Tilfælde tilstede
i højst forskellig Grad. Det vil komme an paa, i hvilken Grad Mennesket kan
hæve sig over enkelte, øjeblikkelige Tilskyndelser. Som overalt i Naturen Kraft
staar mod Kraft, saaledes ogsaa her: et levende Væsens Selvvirksomhed er
forskellig fra den ydre, mekaniske Nødvendighed; overfor den dyriske Attraa
staar den fornuftige Betænksomhed; overfor de egoistiske Drifter staa de
Følelser, der føre til Arbejde for det Heles Vel. Styrkeforholdet mellem
% --- p. 84 ---
\setcounter{footnote}{0}%
\opage{84}de forskellige Kræfter er højst forskelligt i de enkelte Tilfælde og
hos de enkelte Individer.

Howitz fandt, at den gængse juridiske Opfattelse, han optraadte imod, led under
en falsk Psykologi, der i den sidste Menneskealder var bleven bestyrket ved
Kant's Filosofi. Det var især Følelseslivet, der ikke kom til Sin Ret. For
Kants strenge Fornuftmoral stod al Følelse, al Lyst og Ulyst, al Glæde og Sorg,
al Medfølelse og Æresfølelse som „sanselig“ eller „egoistisk“. Howitz
indskærper, at Følelser kunne være af højst forskellig Værdi, og at man ikke
uden videre kan sidestille Sympatiens Glæde med Ganens Nydelser. Enhver Følelse
hænger sammen med en Trang eller en Drift, og Individets Udvikling beror paa
Forholdet mellem de forskellige Drifter. Hvor ophøjet og ædel en Handling end
er, saa bliver den dog kun mulig ved, at en Følelse eller Drift sejrer over
andre. Spinoza har Ret i, at en Lidenskab kun kan fortrænges ved en anden
Lidenskab. Følelseslivet har sine egne Love, og Fornuften --- Evnen til at
dømme om Tingenes Forhold --- er kun tjenende; den afgør, hvilke Midler der ere
nødvendige, og hvilke Veje der kan gaas, men den afgør ikke, hvad der bliver
det sidste Formaal. Fornuften er Styrmanden med sit Ror og sit Søkort, men
Følelsen er Vinden, der fylder Sejlet. Lysten driver Værket, Interessen
udvikler, ja, skaber Geniet, og den Rolle, et Menneske skal spille i Verden,
beror mere paa hans herskende Karakter og Temperament end paa hans Tænkeevnes
logiske Fuldkommenhed.

Gennem Følelses- og Driftslivet komme vi ifølge Howitz tilbage til Menneskets
Naturgrund. Der er en Sammenhæng mellem alle sjælelige Livsytringer, selvom de
ere af højst forskellig Værdi. Denne Sammenhæng betegnes ved, at
Lyksalighedsdriften rører sig paa alle Trin; Forskellen mellem Trinnene beror
paa Forskellen mellem de Genstande, ved hvilke Lyksalighedsdriften
tilfredsstilles. --- Paa dette Punkt blev Howitz, ikke uden egen Skyld,
misforstaaet af sine Modstandere, der, skønt de toge Afstand fra Kants
Psykologi, ved Lyksalighedsdrift blot tænkte paa Egoisme eller Trang til
sanselig Nydelse. Men hans Hovedbestræbelse var at hævde
% --- p. 85 ---
\setcounter{footnote}{0}%
\opage{85}den sjælelige Naturs Lovmæssighed. Og medens Romantikerne, som f. Ex.
Steffens, vilde læse Naturen ovenfra nedad, vilde Howitz gaa den omvendte Vej,%
% footnote *) on p. 85
\footnote[1]{Allerede \emph{Johannes} \emph{Boye} havde i „Statens Ven“ og
\emph{Treschow} i sin „Historiens Filosofi“ (1811) lært en stor
Udviklingsproces fra de lavere organiske Former til de højeste. Mere
videnskabeligt havde et Par Aar før Howitz's Skrifters Fremkomst Statsøkonomen
\emph{Chr}. \emph{Olufsen} i en Afhandling Om \textit{Menneskets Rolle i den
fysiske Verden} (1822) fremsat den Anskuelse, at „Livskraften“ i Naturen fra de
Lag og Former, hvis Tid er forbi, arbejder sig frem til nye Former, der bedre
passe til de forandrede Betingelser. Skønt Olufsen ikke synes at antage en
Kontinuitet mellem selve de forskellige organiske Former, har hans Fremstilling
dog Interesse for Udviklingstankens Historie. --- Howitz har neppe kendt
Olufsens Afhandling, men hans Opfattelse vilde kunne føre i samme Retning.}
betragte Mennesket som Naturvæsen, delagtigt i de Egenskaber, alle organiske
Væsener besidde, og stige op fra de simpleste, mest elementære Handlinger til
de mere sammensatte og udviklede Villiesytringer, stedse søgende de bestemte
virkende Kræfter eller Motiver.

Det er i Virkeligheden en sund og dygtig Erfaringspsykologi, Howitz gør
gældende, paavirket som han er gennem Studiet af Spinoza og Hume. Han var sine
Modstandere fuldstændigt jævnbyrdig, skønt nogle af disse, især Mynster, paa en
saare lidet tiltalende og endnu mindre begrundet Maade vilde slaa sig til
Ridder paa ham. --- Med stor Klarhed udtaler han sig f. Ex. om Forholdet mellem
det Aandelige og det Materielle i Anledning af, at man havde bebrejdet ham, at
han heldede til Materialismen. Han holder sig til Erfaringen, der viser
Mennesket udstyret baade med legemlige og aandelige Egenskaber. Kun ved
Abstraktion kan man skille det Aandelige fra det Materielle, eller omvendt. De
to Arter af Egenskaber stige og falde med hinanden; men det er ikke rigtigt at
sige, at der her er et Aarsagsforhold tilstede, da vi ikke kunne udlede det
Aandelige af det Materielle, saalidt som det Materielle af det Aandelige. Vi
have her et Væsen med to forskellige Egenskaber; men fordi to Egenskaber ere
forskellige, behøve de ikke at høre til to forskellige Verdener. Det
% --- p. 86 ---
\setcounter{footnote}{0}%
\opage{86}beror maaske paa vor Opfattelsesevne, at vi maa gøre en saadan
Forskel mellem Egenskaberne. Jeg finder, siger Howitz, ingen tilfredsstillende
filosofisk Anskuelse om denne Sag undtagen Identitetssystemet, der erklærer
Legeme og Aand for i Grunden at være den samme Substans, kun at denne
fremstiller sig paa en anden Maade for den indvortes end for den udvortes Sans.
--- Howitz henviser ikke blot til Spinoza, men ogsaa til Treschow, der havde
udtalt en lignende Opfattelse, som vi i næste Afsnit ville omtale. Ogsaa
Sibbern hyldede en saadan Opfattelse.

\begin{center}\rule{2cm}{0.4pt}\end{center}

Det var, som allerede bemærket, ikke Howitz's Hensigt med hans Optræden at lære
Determinismen, men at godtgøre, at Spørgsmaalet om Tilregnelighed var
uafhængigt af spekulative Teorier. Dette vilde han opnaa ad to Veje. For det
Første antog han, at naar der i Erfaringen foreligger en Mængde Overgange
mellem Afsindighed og aandelig Sundhed og mellem Svaghed overfor øjeblikkelige
Tilskyndelser og Modtagelighed for Eftertanke, saa vil man ikke kunne paapege
nogen absolut Grænse mellem Frihed og Ufrihed. Han undrer sig over, at Ørsted
og Sibbern, de af hans Modstandere, for hvilke han har størst Respekt, kunne
anerkende hine mange Grader og dog mene at være Indeterminister. --- For det
Andet hævder han, at en Jurist som Ørsted, der grunder Straffens Berettigelse
paa Nødvendigheden af Afskrækkelse, egentligt ikke behøver og ikke kan bruge
Antagelsen af en aarsagsfri Villen. Det kommer an paa, om Individets Villie kan
paavirkes af Hensyn til Straffen; Straffens Berettigelse og Nødvendighed ligger
i, at der skal tilvejebringes en Modvægt, et Kontramotiv overfor de Motiver,
der kunne lede til Forbrydelse. Straffen er et Forebyggelsesmiddel. Det kommer
da ikke an paa, om Individet i Handlingens Øjeblik havde kunnet modstaa
Tilskyndelsen; Straffen gaar ud paa at fremkalde de Motiver, der have vist sig
ikke at være tilstede.
% --- p. 87 ---
\setcounter{footnote}{0}%
\opage{87}Kun den, der ikke kan paavirkes af Hensyn til Straf, vilde det være
meningsløst at straffe.

Naar Howitz erklærer sig mod Indeterminismen, er det ikke mindst, fordi den gør
Straf og Tilegnelse umulige. Frihed --- i den Betydning, i hvilken Howitz
antager dens Mulighed --- er Modtagelighed for Motiver og for Eftertanke. Han
negter kun Frihed i Betydning af Evne til under samme indre og ydre Betingelser
lige godt at kunne fatte to forskellige Beslutninger (bilateral Vilkaarlighed).
Men han gør opmærksom paa, at de, der forsvare Indeterminismen, bruge Ordet
Frihed i tre eller fire Betydninger, hvad der ikke gavner Klarhed. Særligt er
Forvexlingen af Aarsagsfrihed og Villiesstyrke hyppig.%
% footnote *) on p. 87
\footnote[1]{I min \textit{Etik} (Kap. 5) skelner jeg mellem sex forskellige
Betydninger af Ordet Frihed.}

Ogsaa Ordet Nødvendighed har forskellige Betydninger. Det betyder ikke det
Samme som Tvang. Determinisme (eller, som Howitz noget uheldigt ogsaa kalder
den, Fatalisme) er ikke det Samme som Tyrketro. Tyrketroen lader jo, ligesom
Indeterminismen, det, der sker, være uafhængigt af alle forudgaaende Aarsager,
medens Determinismen hævder, at vi gennem Forstaaelse af Aarsagerne og
Anvendelse af denne Forstaaelse kunne forebygge det Onde og fremhjælpe det
Gode. Kun derved bliver Opdragelse og Karakterudvikling mulige. Naar
Menneskekundskab er saa vanskelig at opnaa, er det ikke, fordi Aarsagsloven
ikke gælder i det menneskelige Liv, men fordi her saa mange forskellige
Aarsager virke sammen.

\begin{center}\rule{2cm}{0.4pt}\end{center}

I sin etiske Opfattelse slutter Howitz sig væsentligt til Hutcheson og Hume,
idet han i de sympatiske Drifter ser Grundlaget for det Moralske. Disse Drifter
udformes og klares ved Forstaaelsen af, hvad Samfundslivets Bestaaen kræver.
Derfor hører der en udviklet Erkendelse til, for at der kan handles moralsk.
Men baade med de sympatiske
% --- p. 88 ---
\setcounter{footnote}{0}%
\opage{88}Drifters og med Erkendelsestrangens Tilfredsstillelse er der
forbunden Lystfølelse, og forsaavidt kan Trangen til Lykke siges at være
Grundtrangen. Hvis en Moder ikke følte større Tilfredsstillelse ved --- trods
egen Træthed og Savn --- at pleje sit Barn end ved at forlade det, saa vilde
Moderkærligheden være umulig.

Som allerede bemærket falder hos Howitz Lyksalighedsdrift og Egoisme ikke
sammen. Lyksalighedsdriften kan antage Former, der ere indbyrdes højst
forskellige, ja modstridende. Saaledes er Menneskets Kærlighed til det Gode
eller Almennyttige og hans Tilbøjelighed til at stræbe ud over sin endelige
Tilværelse en egen Art Bevæggrunde, forskellig fra de sanselige Tilskyndelser. ---

De Anskuelser, Howitz gjorde gældende paa det psykologiske og det etiske
Omraade, fandt ingen Tilslutning i vor filosofiske Literatur. Han stod saa godt
som ene med dem. Blandt hans Modstandere var Sibbern den, der stod ham nærmest,
og Sibberns senere Udviklingsgang lod dette Slægtskab træde stærkere frem.

Efter den filosofiske Holmgang trak Howitz sig tilbage til sit eget Omraade,
hvor det desværre ikke længe blev ham forundt at fortsætte sin Virksomhed.

\begin{center}\rule{2cm}{0.4pt}\end{center}

% --- p. 89 ---
\setcounter{footnote}{0}%
\opage{89}

\begin{center}\large 10. NIELS TRESCHOW\end{center}

Flere Gange i det Foregaaende er Treschow bleven nævnt; han var den, der mest
indgaaende fremstillede og vurderede Kants Filosofi, og baade A. S. Ørsted og
Howitz henvise til hans Arbejder. Det er nu Stedet til at give en
sammenhængende Karakteristik af ham.

Treschow blev født ved Drammen i Aaret 1751 og blev Student i København endnu
ikke 15 Aar gammel. Hans Trang til Selvtænkning blev tidligt vakt. „Jeg
fulgte,“ siger han i sin Selvbiografi, „Descartes' Regel, uden at kende denne
Filosof, at begynde min videnskabelige Bane med Tvivl om Alt, hvad jeg havde
lært. Disse Tvivl haabede jeg dog omsider at faa opløste --- ti jeg saa, at
lærde Mænd ikke mere tvivlede.“ Efter meget omfattende Studier, der i
Filosofien foruden den herskende Wolfiske Skole ogsaa omfattede det attende
Aarhundredes engelske Filosofer, drog han igen til Norge som Lærer ved Skolen i
Trondhjem. Senere blev han Rektor i Helsingør og kom derfra igen til Norge som
Rektor i Kristiania. Her holdt han 1796---1797 sine Forelæsninger over Kants
Filosofi. Det var vel nærmest disse, der gjorde, at han 1803, da Riisbrigh tog
sin Afsked, uden Ansøgning blev Professor i Filosofi i København. Regeringen
vilde ikke have Steffens, denne romantiske Sværmer, til Lærer for de Unge. Det
Valg, der blev gjort, var særdeles heldigt. Treschow var en livlig og
selvstændig Mand, der øvede stor Indflydelse paa de Studerende%
% footnote *) on p. 89
\footnote[1]{\emph{Carsten} \emph{Hauch} har i sine Ungdomserindringer givet en
Karakteristik af ham som Lærer.} og blev savnet, da han 1813 drog til Kri%
% --- p. 90 ---
\setcounter{footnote}{0}%
\opage{90}stianias nye Universitet. Han holdt ogsaa Foredrag for større Kredse
med stor Tilslutning. --- I Norge blev han senere Statsraad (Kultusminister) og
levede sine sidste Aar tilbagetrukken og beskæftiget med filosofiske Arbejder
indtil sin Død i Aaret 1833.

\begin{center}\rule{2cm}{0.4pt}\end{center}

Treschow har forfattet den første danske Erfaringspsykologi: \textit{Om den
menneskelige Natur, især fra dens aandelige Side (1812).} Denne Bog, som endnu
har sin Interesse, vidner om fin Iagttagelsesevne og om varm Interesse for
Sjælelivet og dets Udvikling. Skønt han endnu forsaavidt er det attende
Aarhundredes Barn, som han anser Forestillingslivet som det Fundamentale, saa
at alle Drifter skulle have deres Oprindelse af „dunkle Forestillinger“, saa
stiller han dog Følelse og Villie som selvstændige Sider af Sjælelivet overfor
Forestillingen. Afsnittet om Følelseslivet er især rigt paa gode Iagttagelser
og Bemærkninger, og her igen særligt Afsnittet om Sympati, Beundring,
Entusiasme og lignende Følelser.

Med Hensyn til Forholdet mellem Sjæl og Legeme lægger Treschow Spinozas
Opfattelse til Grund. „Mennesket kan vel betragtes som et sammensat Væsen, men
ikke som sammensat af Sjæl og Legeme, ti begge ere kun modsatte Sider af det
Samme, saavidt det nemlig kan være Genstand for den indvortes og den udvortes Sans.“%
% footnote *) on p. 90
\footnote[1]{I sit „filosofiske Testamente“ \textit{(Om Gud, Ide- og
Sandseverdenen.} Kristiania 1831. II. p. 51) siger Treschow: „Sjælens Tanker og
Følelser maatte, hvis det var muligt ved de udvortes Sanser at bemærke dem,
selv forekomme os som visse Arter af Bevægelse; hvad de ere mere, lærer den
indvortes Sans os, paa samme Maade, som vi ved Synet og Hørelsen opdage
Egenskaber, hvorom vi ved Lugten og Smagen ej engang faa nogen Anelse, og
omvendt.“} Materialisterne antage kun det Legemliges, Spiritualisterne kun det
Aandeliges Existens. Begge have Uret overfor Erfaringen. Hverken om Materien
eller om Sjælen have vi umiddelbar Erfaring; det er i begge Tilfælde ved
Slutning, vi fra de skiftende Fænomener naa til Tanken om, hvad der ligger til
Grund, og det er tilsidst et og samme Princip, i hvilket Grunden maa søges.

Ved denne Synsmaade var Treschow istand til at give den
% --- p. 91 ---
\setcounter{footnote}{0}%
\opage{91}naturvidenskabelige Opfattelse dens fulde Ret, uden nogen romantisk
Omtydning. Psykologien kunde bevare sin Selvstændighed som udgaaende fra et
særeget Synspunkt for den samme Tilværelse, som Naturvidenskaben, særligt
Fysiologien betragtede fra sit Synspunkt. Det er til selve de
naturvidenskabelige Grundsætninger, Treschow støtter sig: Aarsag og Virkning
kunne ikke være forskellige i Væsen; en Bevægelse kan kun fremkalde en anden
Bevægelse; synes der at opstaa Mere i Virkningen, end man kan udlede af
Aarsagen, f. Ex. Fornemmelse foruden Bevægelse, saa maa Forklaringen ligge i,
at vi kun ufuldstændigt og ensidigt have opfattet Aarsagen.

Det var denne Opfattelse, Howitz sluttede sig til. Den ligger ogsaa til Grund i
Sibberns Psykologi, saavel som for min Fremstilling af Psykologien. Der er god
dansk Tradition for den. Treschow synes paa dette Punkt at være paavirket af
Hartley, foruden af et direkte Studium af Spinoza; men hans egen Tænkning har
gjort Udslaget. ---

\begin{center}\rule{2cm}{0.4pt}\end{center}

Maaske det mærkeligste af Treschows Arbejder er den allerede foreløbigt omtalte
Afhandling (i Videnskabernes Selskabs Skrifter 1810), som drøfter det
Spørgsmaal, om der kan dannes Begreber om individuelle Genstande.
\textit{(Gives der noget Begreb eller nogen Ide om enslige Ting?)} Den er
rettet mod den spekulative Filosofis Tilbøjelighed til at bevæge sig i almene
Ideer og betragte de enkelte, individuelle Fænomener som ufuldkomne Udtryk for
dem. Treschow fremhæver, at der indeholdes Mere i de individuelle Genstande end
hvad der kan indgaa i et Almenbegreb. Den specielle, konkrete Form eller Type,
som det enkelte Individ frembyder, maa det ligesaavel være en Opgave at
erkende, som det, der er fælles for mange Individer. I den virkelige Verden
gives der kun Individer, og vore Almenbegreber ere egentligt kun Hjælpemidler
til at lære dem at kende. Men da et enkelt Individ i hele sin særlige
Ejendommelighed frembyder et uudtøm%
% --- p. 92 ---
\setcounter{footnote}{0}%
\opage{92}meligt Indhold, saa er det en meget vanskelig Opgave at finde den
Ide, der udtrykker, hvad der ligger til Grund for dets skiftende Tilstande.%
% footnote *) on p. 92
\footnote[1]{Naar Treschow her taler om „en Ide, der ej i den Henseende er
almindelig, at den omfatter flere Ting, men forsaavidt den udtrykker alle
mulige Tilstande af den samme“, mener han det Samme, som jeg i min Psykologi
(VB, 9) har kaldt en typisk Individualforestilling. --- Smlgn. hvad H. C.
Ørsted kaldte „Samfoldighed“.} Det er især paa det organiske Livs og Historiens
Omraade, at dette Problem stilles. Og der er et dertil svarende praktisk
Problem: vor Værdi („vort Væsens hele Fortræffelighed“) beror paa
Individualitet, og Formaalet for vor Stræben maa være at opnaa den højeste Grad
af Individualitet.

Naar Individerne ere det egentligt Virkelige, ikke blotte Midler eller Former
for noget Alment, bliver det et Spørgsmaal, hvilken Gyldighed der kan tillægges
Artsbegreberne. Skulde Naturen virkeligt have dannet faste Grænser mellem
Arterne? Ere Artsbegreberne Mere end Hjælpemidler, vi anvende til at skaffe os
Oversigt? Er det ikke vor Forstands Ufuldkommenhed, der gør saadanne
Hjælpemidler nødvendige, fordi vi ikke formaa at trænge ind i det Enkeltes hele
Ejendommelighed? ---

Omtrent samtidigt med disse Tanker fremkom \textit{Lamarck's} Kritik af
Artsbegrebet \textit{(Philosophie zoologique. 1809).} Artsbegrebet er for
Lamarck et kunstigt Middel, man ikke maa henregne til „selve Naturens Love og
Virksomheder“. Treschow har ikke kendt Lamarck, hvorimod han henviser til ældre
franske Naturforskere (Maillet, Buffon), hvis Ideer allerede gik i samme Retning. ---

Naar Artsbegrebet kritiseres paa denne Maade, ligger det ikke fjernt at antage,
at Forskellen mellem de forskellige Arter ere opstaaede ved successiv
Udvikling. Og denne Tanke spiller ogsaa en stor Rolle hos Treschow. Den
fremtræder i hans \textit{Elementer til Historiens Filosofi} (1811), som han
først fremstillede i Forelæsninger i Vinteren 1806---1807. Ogsaa paa dette
Punkt traadte Treschow i Modsætning til Romantiken. Naar denne talte om
Udvikling, var det kun symbolsk:
% --- p. 93 ---
\setcounter{footnote}{0}%
\opage{93}det er i Formernes Spil, i Naturens eller Naturaandens Fantasi, at
Overgangene finde Sted; det Evige er paa én Gang udtrykt i uendeligt mange
Former; men en historisk Overgang, en Udvikling i Tiden under bestemte ydre
Betingelser antages ikke. Steffens benegter jo saaledes udtrykkeligt en real
Udvikling i Naturen. Treschow derimod forsøger paa sin Vis at give en naturlig
Skabelseshistorie, idet han fortolker første Mosebog med samme Dristighed som
Holberg. Skabelsen er en lang Proces, fremskridende fra Trin til Trin. En
„umiddelbar Skabelse“ er et tomt Ord. „Fra Stoffer som synes raa, fordi man
deri endnu ingen Form bliver var, sker en umærkelig Overgang til mere bestemte
Skikkelser... Naturens Gang i de Afvigelser eller Udartninger, hvorved de
forrige Slægter og Arter forvandles til ganske nye, er ej blot krybende, men
saa langsom tillige, at hverken egen Erfarenhed eller Historiens altfor lidet
Omfang sætter os istand til at følge dens næsten ubemærkelige Spor. Efter
Aartusinders Forløb give de mindst tvetydige Levninger af Dyr og Mennesker, der
kunne anses for Stammefædre til de nærværende, skønt temmeligt forskellige, dog
Anledning til Indvendinger og Tvivl.“ --- Mennesket har først efterhaanden
faaet de Egenskaber, der nu karakterisere det. Menneskets Naturhistorie er en
Del af hele Naturens Historie. Det har i Aarhundreder ført et dyrisk Liv, endog
af højst ufuldkommen Art; dets Udvikling skyldtes en Vexelvirkning mellem indre
og ydre Betingelser. Før vi kunde opløfte vore Hoveder til Himlen og føle os
beslægtede med den, maatte ikke blot vor Organisme udformes, men ogsaa
Jordoverfladen, Plante- og Dyreverdenen undergaa Forandringer.

„Have vi maaske,“ spørger Treschow, „Aarsag til at skamme os ved saadan
Oprindelse? Dette Spørgsmaal kan besvares ganske kort. Vi behøve ej at gaa saa
langt tilbage i Tiden forat finde stærkere Grunde til vor Ydmygelse. Hvor mange
Aar er det vel siden, at Enhver af os har staaet langt under Dyrene, ja selv
under den ringeste Plantes Liv? Er det ej vigtigere at tænke paa, hvad vi nu
ere og engang ganske vist
% --- p. 94 ---
\setcounter{footnote}{0}%
\opage{94}skal blive, end hvad vi vare, somom den stolte Ceder med Haan vilde
se ned paa den unge Spire, der oprinder ved dens Rødder, eller vi til det
Foster, hvis Hjerte neppe har begyndt at slaa i Moders Liv? Overalt har det
Levende samme Rod: guddommelig er vor Herkomst, evig og uendelig vor
Bestemmelse.“ ---

Treschow har ikke Ideen om Kampen for Tilværelsen, der derimod, som ovenfor
omtalt, findes hos Joh. Boye og kommer frem endel Aar efter hos Chr. Olufsen i
en Afhandling, som paa en interessant Maade staar imellem den romantiske og den
darwinistiske Opfattelse.

Som det allerede ses af Treschows anførte Udtalelser lægger han Vægt paa
Analogien mellem det enkelte Individ og hele Slægtens Udvikling. Baade fysisk
og aandeligt gennemgaar den Enkelte Slægtens Udvikling. Han strækker denne
Analogi meget vidt og finder saaledes, at Slægten ligesom Individet har havt
sin Fostertilstand, sin barnlige Uskyldighedstilstand, sin Drengealder o. s. v.
Tilsidst vil den ogsaa faa sin Alderdom. --- En saadan Analogi er naturligvis
vilkaarlig, og den svigter paa de afgørende Punkter. Kun billedligt kan Slægten
opfattes som et enkelt Individ, allerede af den Grund, at der i Historien
samtidigt foregaa forskellige Processer, der ikke kunne indordnes i en eneste
fremskridende Række. Man maa da i hvert Tilfælde --- som Sibbern senere gjorde
--- lægge Vægt paa, at Udviklingen foregaar sporadisk ud fra mange forskellige
Punkter paa én Gang. Men derved skydes Problemet selv længere ud, idet det
store Spørgsmaal bliver, om og hvorledes de mange Processer kunne forenes i én
samlet Strøm. ---

\begin{center}\rule{2cm}{0.4pt}\end{center}

Bag de specielle Teorier ligger hos Treschow en religiøsspekulativ Enhedslære.
Enheden i Alt er for ham Tænkningens sidste Konsekvens, --- det Punkt, vi naa
til, naar vi Skridt for Skridt have fulgt Rækken af Grunde og Følger, men hvor
den Sætning, at der er en Grund til Alt, ikke læn%
% --- p. 95 ---
\setcounter{footnote}{0}%
\opage{95}gere kan anvendes, fordi den kun gælder for det Mangfoldige og Forskellige.

Overfor Kant havde Treschow hævdet, at vi kun kunne forbinde i vor Tanke, hvad
der allerede er forbundet i Naturen. Det er de i Erfaringen fremtrædende
Modsigelser og Vanskeligheder, der drive os frem til Ideen om en højeste Enhed.
Det er Nødvendigheden af at forbinde Tingene, der godtgør denne Ides Gyldighed.
Fra Sanseverdenen komme vi saaledes til Ideverdenen, og fra denne til Gud. Men
det er ikke tre forskellige Verdener, vi her have for os. „Naturens Kræfter ere
fra Guddommens egne ej forskellige, undtagen forsaavidt denne i sig selv
betragtet er uendelig, men i hine aabenbarer sig paa endelig Maade, og i sin
Uendelighed alene som en uophørligt tiltagende Størrelse. I denne Mening kan
man sige, at Materien har dannet sig selv, efterhaanden er gaaet frem fra de
mindre fuldkomne til den mest fuldkomne menneskelige Form. Man kan ligeledes
med Sandhed sige, at Gud, efterat have frembragt Materien selv, eller ved dens
Egenskaber gjort sit eget Væsen saa anskueligt som det kunde blive, deraf igen
har frembragt alle Ting. Disse Talemaader, som de Fleste holde for himmelvidt
forskellige, betyde dog egentligt det Samme og give kun tvende menneskelige
Forestillingsarter tilkende.“ --- Ligesom H. C. Ørsted var Treschow overbevist
om, at hvad Filosofien kalder den højeste Enhed er det Samme, som Religionen
kalder Gud. En umiddelbar Aabenbaring antog han ikke. End ikke i sit
Alderdomsarbejde \textit{Om Gud, Ide- og Sanseverdenen} (1813), hvor han
udvikler sine Anskuelser paa en noget dogmatisk Maade og tager mere Hensyn end
før til teologiske Forestillinger, forlader han sit rent filosofiske
Standpunkt, der jo ogsaa havde en saadan Dybde og Bredde, at det kunde lade
naturvidenskabelige Hypoteser og psykologiske Iagttagelser frit udfolde sig,
uden at behøve at frygte for at komme i Modstrid med dem. ---

Efter Treschows Bortgang fra København blev det det æstetiske og teologiske
Spørgsmaal, der for lang Tid optog Interessen. Men Treschows Eftermand ved
Universitetet, som
% --- p. 96 ---
\setcounter{footnote}{0}%
\opage{96}vi nu skulle tale om, arbejdede væsentligt i samme Aand som han,
skønt han ogsaa var stærkt paavirket af Romantikens Aand og kun ved den stadige
Stræben efter Sundhed og Sandhed, der prægede hans lange Liv, arbejdede sig ud
af dens Ensidigheder.

\begin{center}\rule{2cm}{0.4pt}\end{center}

% --- p. 97 ---
\setcounter{footnote}{0}%
\opage{97}

\begin{center}\large 11. FREDERIK CHRISTIAN SIBBERN\end{center}

Da Treschow gik til Norge 1813, blev Sibbern (f. 1785, d. 1872) hans Eftermand.
Baggesen havde tænkt sig at blive Professor i Filosofi ved denne Lejlighed og
udtaler sig endel Aar senere meget harmfuld over at være forbigaaet for en
„uvittig Dreng“.%
% footnote *) on p. 97
\footnote[1]{Se Kr. \emph{Arentzen}: \textit{Baggesen og Oehlenschläger.} VI p.
35. --- At Sibbern optraadte paa Oehlenschlägers Side i den store Strid, stemte
naturligvis ikke Baggesen blidere.} Denne Dreng skulde dog voxe til Mand og
blive en af de ejendommeligste Skikkelser i dansk Tænknings Historie.

Sibberns Stilling i vort Aandsliv er betegnet ved, at han, der i sin Ungdom var
stærkt greben af Romantikens Aand, baade i filosofisk, religiøs og poetisk
Henseende, gennem utrætteligt Arbejde med Tanken og med hele sin Personlighed
aabnede sig for nye Synsmaader og fortsatte sin Udvikling helt op i
Oldingealderen, hvor den for de Fleste en Gang for alle er afsluttet. Det var
især hans psykologiske Sans og hans varme Sympati, der aabnede hans Øjne for
Livets Virkeligheder i deres Mangfoldighed og Ejendommelighed hver for sig.
Istedetfor Romantikens Spekulation traadte en paa Erfaringspsykologi og
Naturvidenskab grundet Verdensanskuelse; istedetfor den Ortodoxi, han af Lede
ved Rationalismen havde sluttet sig til, udviklede der sig en fri religiøs
Livsanskuelse; istedetfor Romantikens Konservatisme traadte en stadigt voxende
Sympati med Folkets, især „Lavmandsfolkets“ Fremstræben, og han var En af de
Første, der her
% --- p. 98 ---
\setcounter{footnote}{0}%
\opage{98}hjemme havde Sans for det sociale Spørgsmaal og udtalte sig for en
radikal Løsning af det. Fælles for alle Stadier i hans Udvikling var den
levende Sans for alt Menneskeligt, Aabenheden og den kostelige Naivitet. I hans
ydre Optræden var der en vis Ubehjælpsomhed, og der kom desuden tidligt en
pedantisk og stiv Form i hans Tale og Skrift, en Skal, der for Mange skjulte Kærnen.%
% footnote *) on p. 98
\footnote[1]{„Der har for mig,“ siger han i et Brev, „altid været et stort
Skridt fra det, der stod for mit indre Syn, til at komme i Gang med at faa det
frem for mig paa en tilfredsstillende Maade.“ \textit{Breve til og fra
Sibbern.} II. p. 230.} Det Livfulde i hans Skrifter hindredes --- især i de
senere Udgaver --- af en stærk Tilbøjelighed til Inddelinger og Rubriceringer.

Hans egne Oplevelser og hans eget Temperament blev i flere Henseender afgørende
for den psykologiske Opfattelse, der spillede saa stor en Rolle i hans
Filosofi. I sin Ungdom havde han en Gæringsperiode, i hvilken en haabløs
Kærlighed var det stærkest virkende Element. Han lærte her af egen Erfaring at
kende Brydningen mellem opstigende Begær og beundrende Glæde. I første Del af
\textit{Gabrielis' Breve} har han skildret denne Krise, ofte med de samme Ord,
der findes i hans egne Breve fra samme Tid. I hans Sind brødes stadigt
Bitterhed og Nedstemthed med frejdig Optimisme, og han har herved følt
Modsætningen mellem de Tider, hvor Aandsforladthed og Ufrugtbarhed raade, og de
Tider hvor „Musen kommer i Besøg“.

Sin største Indflydelse paa de Studerende øvede han vistnok i 20'erne og
30'erne. For dem, der kom i Berøring med ham i Slutningen af hans 55-aarige
Universitetsvirksomhed, var det ofte vanskeligt at trænge ind til Kærnen bag
den noget stivnede Optræden. Og dog var Kærnen stadigt levende. Endnu 86 Aar
gammel udtalte han, at han følte Sommerlighed hos sig, og at man vel ikke kunde
lægge en Alen til sin Væxt, men dog maaske en Tomme, naar man holder sig rank!%
% footnote **) on p. 98
\footnote[2]{I mine \textit{Mindre Arbejder} (første Række) har jeg givet en
udførligere Karakteristik af Sibberns Personlighed. Ogsaa hvad hans Lære
angaar, kan jeg henvise til denne Fremstilling; den i det følgende givne
Fremstilling er paa nogle Punkter udførligere, paa andre Punkter mere
kortfattet. I hin Karakteristik er særligt Paavirkningen af Steffens fremhævet.
Dertil maa føjes Treschows Indvirkning, som ikke har været ringe, og som
Sibbern selv gentagne Gange fremhæver. Og begge Paavirkninger gik i Retning af
det Individuelles Betydning i Modsætning til det abstrakt Almene, det Reales
Betydning i Modsætning til det Formale.}

% --- p. 99 ---
\setcounter{footnote}{0}%
\opage{99}Sibbern var først og fremmest Psykolog. Han havde stor Evne til
Iagttagelse og stor Glæde ved at gøre Iagttagelser paa det sjælelige Omraade.
Han sammenligner sin Psykologiseren med en Botaniseren. Og ligesom der for
Botanikeren intet Ukrudt gives, saaledes interesserer Psykologen sig for alt
Menneskeligt, hvad enten det er stort eller ringe, godt eller ondt. Han var
overbevist om en nøje Sammenhæng mellem det Sjælelige og det Legemlige. Det er
en og samme Grundvirken, der fremtræder som sjælelig og som legemlig, et og
samme Liv, der ytrer sig under to Former. Hvis man vilde antage, at der fra
først af kun foregik en legemlig Bevægelsesproces, vilde man ikke kunne
forstaa, hvorledes den i sin Fortsættelse kunde frembringe Tanker. Det
Legemlige kan kun frembringe et Legemligt. Omvendt kan man ikke forstaa,
hvorledes Tanker kunne frembringe legemlige Bevægelser, og det Sjælelige kan
derfor heller ikke være det Oprindelige, men maa have sin Grund i noget Dybere,
der tillige formaar at frembringe Bevægelse. Ud af den samme Grundkilde eller
Grundvirken, ved hvilken de materielle Formationer ske, sker ogsaa den
aandelige Virken.%
% footnote *) on p. 99
\footnote[1]{\textit{Om Forholdet mellem Sjæl og Legeme.} 1849. p. 84---86.
\textit{--- Spekulativ Kosmologi. 1846. p. 41.} --- Hos Treschow og Howitz have
vi allerede truffet en lignende Opfattelse. Sibbern henviser idethele gentagne
Gange til Treschow som sin Forgænger. Se Fortalen til \textit{Psykologisk
Patologi} (1828) og flere Udtalelser i hans Alderdomsværk \textit{Meddelelser
af et Skrift fra Aar 2135. ---} Hvad der stiller Sibbern i en vis Modsætning
til vor Tids Videnskab, er hans Vitalisme, d. v. s. hans Lære om en særegen
Livskraft, om en fra en indre Kilde kommende Virksomhed. Dog sætter han igen
denne Livskraft i indre Forbindelse med den i al uorganisk Natur, saaledes
allerede i Urtaagen, virkende Kraft.}

Sjælelivet er ligesom det virkelige Liv stedt i stadig Skiften og Vorden. Livet
arbejder sig stedse frem til nye Former;
% --- p. 100 ---
\setcounter{footnote}{0}%
\opage{100}ikke blot fra først af, men under hele Livets Gang arbejdes der i
Sjælens dybe Grund; nye Anlæg og Tilskyndelser opstaa. Derom vidner ikke blot
det Geniale, men ogsaa de pludselige Vækkelser, der kunne opstaa af ubevidste
Oprindelser. Det er Livet, der vil frem i os. Der er Meget, som hører med til
Sjælelivet, og som virker, før Bevidstheden vaagner. Sjælelivet er, ligesaalidt
som det legemlige Liv, en Gang for alle formet; Sjæleformationsprocessen gaar
stadigt sin Gang, hvad enten vi umiddelbart mærke det eller ikke.

Sjælen existerer ikke som et Væsen for sig, afset fra Sjælelivet, de sjælelige
Funktioner, ligesaalidt som Tanken existerer afset fra Tankegangen. Ordet „i
sig selv“ kan kun betyde: „i sine Grundforhold“. Fichte forkastede med Rette
Forestillingen om en Sjæl, naar man derved forstaar en Ting, et Noget, der
\textit{har} Liv, eller \textit{har} Bevidsthed. Det er i Grunden en Slags
Materialisme, der ytrer sig i en saadan Forestilling; man gør „Sjælen“ til et
Væsen, der kunde være at træffe etsteds i Rummet, f. Ex. et eller andet Sted i
Hjernen.%
% footnote *) on p. 100
\footnote[1]{Smlgn. Afhandlingen \textit{Om den menneskelige Sjæls Subsistens i
Liv og Død.} (Filosofisk Arkiv og Repertorium. 1829).} Den Leven eller Virken,
i hvilken Sjælen bestaar, kan stadigt undergaa Forandringer, og Aarsagerne til
disse Forandringer behøve, som sagt, ikke at findes i det bevidste Sjæleliv.
Hele Livet igennem kan Sjælen voxe ved Indflydelser fra den dybe Naturgrund,
der bærer den. --- Sibberns Opfattelse minder her om, hvad især William James i
den nyeste Tid har gjort gældende, at der kan ske nye Indskud i
Bevidsthedslivet fra Noget, som ligger under Bevidsthedens Tærskel (det Subliminale).

Men den Udvikling, der saaledes foregaar, tager ikke sin Udgang fra et enkelt
Punkt i Sjælen for derfra at brede sig til det øvrige Sjæleliv. Der rører sig
paa én Gang forskellige, ofte indbyrdes stridende Anlæg og Tilskyndelser, Ideer
og Drifter. Udviklingen foregaar \textit{sporadisk,} og det kommer da an paa,
om der kan komme et Hele ud deraf. Baade det „Lavere“, det, der bør fortrænges,
og det „Højere“, det, der bør fremmes, opkommer først som Glimt, Ansporelse, Beta%
% --- p. 101 ---
\setcounter{footnote}{0}%
\opage{101}gelse. Og deraf følger da Røre i Sindet, Gæring, og ofte en svar
indre Kamp eller Debat. Det er herved ikke rent individuelle Kræfter, der ere i
Bevægelse; ti den Leven og Virken, hvori Sjælens Væsen bestaar, er en Del af
den universelle Virken, der rører sig i alt Liv, i den hele Natur. Og baade om
det rent Individuelle og om det Universelle gælder det, at det fremtræder
sporadisk og ofte stridende, før Enhed og Harmoni kan naas.

„Jo dybere Livet aabner sig i et Gemyt, desto snarere kan det komme til at
gennemgaa en kaotisk Elementernes Blanding og Røring i sit Indre, og den
Verden, der skal opbygge sig i det, maa det nu være med til at skabe, som om
det skulde genskabes fra Grunden af. Og da heri ej blot Livets Gehalt
overhovedet rører sig efter sin Væsenhed --- dog i sine forskellige
Grundvirkninger som fra forskellige Sider og som i stridige Elementer, der
først søge Harmonien, --- men nu ogsaa Jegheden rører sig efter sin Jeghed, saa
er det ganske naturligt, at i den indvortes Kamp en vis sygelig Hildethed i sig
selv, med vexlende Selvtillid og Mistvivlelse, let kan faa Overhaand i visse
Livets Perioder som naturlig kritisk Overgangstilstand. Denne maa da gennemgaas
med Tillid og Selvopmandelse, men og med resignerende Oppebien og Hengivenhed i
den altomfattende Alvirken, som just her viser sin Magt i Individet, og hvis
Kar man just her føler sig at være.'“%
% footnote *) on p. 101
\footnote[1]{\textit{Psykologisk Pathologie.} 1828. p. 223. --- Dette Skrift
udgør anden Del af den første Udgave af Sibberns Psykologi. Det blev optrykt
1885 (Hundredaaret efter hans Fødsel) under Titel: „Læren om de menneskelige
Følelser og Lidenskaber.“ Ordet „Pathologie“ betyder nemlig her ikke det syge
Sjæleliv, men idethele det i Følelser, Drifter og Sindsbevægelser fremtrædende
Sjæleliv. --- Det er den fyldigste og livligste psykologiske Fremstilling,
Sibbern har givet.}

Denne Udtalelse viser Sibberns Ejendommelighed: personlig Erfaring har aabnet
hans Blik for en væsentlig Side af det sjælelige Livs Kaar, og den lægges da
til Grund for en almindelig Teori; tillige er det ikke blot Psykologen, men den
praktiske Livsfilosof, der fører Ordet.

% --- p. 102 ---
\setcounter{footnote}{0}%
\opage{102}Medens han i den anførte Udtalelse skelner mellem den universelle
Virken og Jegheden, dvæler han andre Steder netop ved saadanne Tilstande, hvor
begge gaa i Et, som i den instinktive og geniale Virken, eller i den fulde
Hengivelse til og Opgaaen i en Sag. Her kan der ikke tales om Motiver, medens
man nok kan tale om Følelsesbevægelser. Kun hvor man ikke gaar op i umiddelbar
Virken, --- naar altsaa Handlingen ikke „forstaar sig af sig selv“, --- kun der
opstaar Bevidsthed om Motiver som noget fra Villien Forskelligt. Den fulde
Indlevethed er da ikke tilstede.

Følelseslivet skildrer Sibbern bedst. Han har klart Blik for Kontrasternes
Betydning, for Modsætningen mellem Begær og Beundring, for forskellige
Følelsers Brydning og den nøje Forbindelse, de netop formedelst Brydningen
kunde indgaa. I sidstnævnte Henseende skelner han mellem Følelsers „Blanding“
og deres „Blandethed“. Naar to Genstande (eller en Genstands to forskellige
Sider) paa én Gang sætte Sjælen i to modsatte Følelsesbevægelser, sker der
enten en udvortes Blanding af de to Følelser, idet de ligesom ligge ved Siden
af hinanden, eller Sindet svinger imellem dem; i begge Tilfælde ere de
stridende Elementer endnu uforsonede. Anderledes, naar Modsætningen opløses til
Enhed. Da opstaar en blandet Følelse, idet de to Elementer virke saaledes i Et,
at de tilsammen frembringe en eneste Totalfølelse. Saaledes naar Frejdigheden
vindes ved Overvindelse af Frygten, eller Modet og Lysten forhøjes ved selve
den Modstand, der mødes. Det er ikke den blotte Kontrast, der virker; det er en
ny Følelse, til hvilken begge de kontrasterende Følelser have givet deres Bidrag.

\begin{center}\rule{2cm}{0.4pt}\end{center}

Hvad Sibbern havde fremhævet som ejendommeligt for den sjælelige Udvikling,
genfandt han i hele Tilværelsen. Ideen om den sporadiske Udvikling er en
Hovedtanke hos ham og af afgørende Betydning for hans hele Livsanskuelse, eller
for hans hele Stemning ved Livet.

% --- p. 103 ---
\setcounter{footnote}{0}%
\opage{103}Sibbern udtaler sig mod ethvert Forsøg paa apriorisk Konstruktion af
Naturen, bl. A. mod Kants Forsøg i denne Retning.%
% footnote *) on p. 103
\footnote[1]{Særligt indvender Sibbern mod Kant, at han opfatter Tyngden
(Tiltrækningen) som en Grundkraft, skønt den Mulighed ikke er udelukket, at den
kan være Resultatet af en Proces, udspringe af noget langt mere Kompliceret,
end der ses i det simple Fænomen. \textit{Spekulativ Kosmologi.} 1846. p. 36.
--- Dette Skrift indeholder i de store Træk Sibberns almindelige
Verdensanskuelse. Sin Opfattelse af Filosofien har han fremstillet i sit Skrift
\textit{Om Filosofiens Begreb, Natur og Væsen.} 1843.} Vi ere endnu, siger han,
altfor langt fra at være trængte saa vidt ind i Naturens Studium, at en
Bestræbelse for, blot efter de Grundideer, hvortil vi nu ere komne, at ville
konstruere sig et Naturens System a priori, skulde kunne Andet end være højst
misligt. Vil man naa til en Filosofi, maa denne udvikle sig under stadig
Sammenligning (Konferens) med det Givne. Den er Maalet, ikke Begyndelsen.
Selvom det, der virker i den hele Tilværelse, ogsaa virker i vor Tanke, saa kan
dog denne vor Tanke kun udvikle sig gennem Vexelvirkning og Brydning med de
andre Elementer i Tilværelsen. Vor Tanke er kun et af de mange
Tilværelseselementer. Her viser sig netop \textit{det Sporadiske.}

Erfaringen vidner paa alle Omraader om, at Alt i Naturen fra først af opstaar
sporadisk. I den Urtaage, hvoraf vort Planetsystem har udviklet sig, og som
maaske selv er et tidligere Klodesystems sidste rent opløste Residuum, udgik
der Virkninger fra forskellige Dele, og gennem denne Vexelvirkning skred
Udviklingen frem. Ved Krystallisationen, f. Ex. ved Isdannelsen, skyde Naale
sig frem fra forskellige Punkter for tilsidst at danne en fast Sammenhæng.
Under Fosterudviklingen begynder Organisationsprocessen, især Knokkeldannelsen,
ud fra forskellige Steder og mødes saa forat slutte sig sammen til en Helhed. I
Sjælelivet bevæge sig, som vi allerede have set, mange forskellige Drifter,
Tilskyndelser og Tanker, om ikke kaotisk, saa dog ad forskellige Veje, indtil
den rette harmoniske Opgaaen til Totalitet bliver mulig. Paa lignende Maade
gaar det med Menneskene
% --- p. 104 ---
\setcounter{footnote}{0}%
\opage{104}i Samfundet, og med de forskellige Samfund indenfor Menneskeheden.

Gennem den ved det Sporadiske foranledigede Debat eller Kamp af Alt imod Alt
skrider nu Udviklingen frem. Betingelsen for en gennemorganiseret og fyldig
Leven er, at den har udfoldet sig successivt, og at den er opstaaet ved
Harmonisering af sporadisk opstaaede Elementer. Tid og Sondring ere altsaa
væsentlige Forudsætninger for Udviklingen. --- Sibbern træder her i skarp
Modsætning til Romantikens Filosofi, der kun antog en Udvikling i rent logisk
Betydning, og for hvem Udvikling var en Udfoldelse af et absolut Princip uden
Forudsætning af reale Forskelle. Med Fremhævelse af det Sporadiske hænger
Sibberns store Sans for det Individuelle sammen: hvert enkelt Individ er
Udgangspunkt for nye, ejendommelige Processer, er et af Tilværelsens Elementer,
et Indlæg i Livets store Debat. Ved Fremhævelsen af Tidens afgørende Betydning
for alt Existerende viser Sibbern sig som ivrig Udviklingsfilosof. Han anvender
sin Udviklingstanke ogsaa paa Mennesket. Menneskeheden har udviklet sig af
Dyriskheden. En vis Dyreformation har fra en ringe Begyndelse gennem en Række
af Generationer, indenfor hvilke Uddannelsesprocessen stedse førtes videre,
naat det Punkt, da nu dens sidste Affødning blev et Væsen, der saa at sige slog
Aandelighedens Øje op. Den Proces, der saaledes er foregaaet, havde igen sit
Udspring fra det første Røre i Urtaagen, ud fra hvilket der blev skredet frem
til et Punkt, hvor Livets Udvikling kunde begynde.

En første Begyndelse spore vi ikke, og om et Endemaal kunne vi kun tale for den
Del af Tilværelsen, hvori vi nærmest befinde os, og som vi kende til. Og selve
denne Del er endnu sandsynligvis i sit første Stadium, paa Vej til sin rette
Begyndelse! --- Et absolut Endemaal vilde forudsætte en Tid, hvis Uendelighed
svarer til Tilværelsesindholdets Uudtømmelighed.

Sibberns Udviklingslære viser tilbage til Lamarck; men i den Betydning, han
tillægger den store Debat eller Kamp af Alt mod Alt, peger han fremad mod
Darwinismen og Muta%
% --- p. 105 ---
\setcounter{footnote}{0}%
\opage{105}tionslæren. I den danske Literatur har han Boye, Treschow og Olufsen
til sine Forgængere. Ja, allerede hos Holberg viste der sig Antydninger af en
stor Udviklingsproces i Tilværelsen.

\begin{center}\rule{2cm}{0.4pt}\end{center}

Karakteristisk for Sibbern er baade hans Overbevisning om, at der stedse kan
gøres ny Erfaringer, og hans ikke mindre faste Overbevisning om, at der trods
det Spredte og Kaotiske i de opdukkende Fænomener og Processer, vil udvikle sig
Harmoni imellem dem. De lange, urolige Overgangsperioder, der ofte frembyde
Billedet af noget Forvredent og Abnormt, kunne vække Tvivl og sætte
Taalmodigheden paa Prøve. Men Livets store Debat er nu engang Betingelsen for
en kraftig Enhedsdannelse. „Jo rigere Endelighed, des rigere Uendelighed; jo
større Mangfoldighed, desto fyldigere Enhed! Alt sondrer sig for paa desto
intensivere Maade at gaa sammen igen til Enhed!“

Her gaar Sibberns filosofiske Opfattelse i Et med den religiøse. I den
Forudsætning om en indre Sammenhæng i Alt, hvorpaa al Videnskab og al
videnskabelig Stræben hviler, finder han Vidnesbyrd om et ideelt Princip, et
Grundvæsen, en natura naturans, som det i alle Ting inderste Tilstedeværende og
til Grund Liggende. Dette Alvirksomme bærer i sig et Skaberdyb, hvoraf stedse
Nyt vælder frem, hvad der især ses, naar Resultatet af en sammensat Proces
indeholder Mere end der følger af Elementernes Samvirken. Ordet Gud betegner
dette allestedsnærværnde og alvirksomme Grundvæsen, forsaavidt det tænkes
centraliseret i sig selv og bliver Genstand for Ærefrygt og Fortrøstning. Der
er paa dette Punkt en Grænse for Sibberns Udviklingsfilosofi, ti han kan ikke
tænke sig, at Gud --- der dog efter hans Opfattelse er et personligt Væsen ---
selv udvikler sig. Det Alvirksomme er evigt og uforanderligt; det er kun
Tilværelsens Verdensside, ikke dens Guddomsside, der udvikler sig. Det er fra
Sibberns eget Standpunkt en Inkonsekvens; Ana%
% --- p. 106 ---
\setcounter{footnote}{0}%
\opage{106}logien med hans Opfattelse af Sjælens Begreb maatte føre til at se
Uholdbarheden af denne Skelnen mellem det Uforanderlige og det Foranderlige.

Hvad Forholdet mellem Filosofi og positiv Religion angaar, er der foregaaet en
interessant Udvikling hos Sibbern. I sin Ungdom var han stærkt betagen af
positiv Kristendom. Han fandt Hvile og Vederkvægelse deri for sin Sjæl, og han
fandt deri en stor Aandsfylde, som Rationalismen ikke havde ladet vederfares
Ret. Men han hævdede allerede paa denne Tid Filosofiens Selvstændighed. For det
Første gjorde han opmærksom paa, at Dogmatiken stedse har bygget paa Filosofi:
hvis man af Dogmatiken, saaledes som den lige fra de første Tider har dannet
sig, vilde uddrage Alt, hvad der er udsprunget af Filosofering, vilde den
styrte sammen. For det Andet hævdede han, at Filosofien ikke kan staa
anderledes overfor Kristendommen, end den staar til alle psykologiske,
biologiske og fysiske Fænomener. Der maa vel skælnes mellem teologiske
Antagelsers Forhold til den kirkelige Overlevering og den reale Mulighed saavel
som den egentlige Beskaffenhed og Betydning af de Kendsgerninger, hine
Antagelser angaa. Filosofien staar overfor Tro og Kristendom, som den staar
overfor et hvilketsomhelst andet Emne.%
% footnote *) on p. 106
\footnote[1]{Se Afhandlingerne \textit{Bidrag til Bestemmelse af Filosofiens
Forhold til Teologien} og \textit{Hvad er Dogmatik?} i Filosofisk Arkiv og
Repertorium (1829---1830).}

Allerede nu gjorde han sine Forbehold overfor den kirkelige Ortodoxi, der i
Modsætning til Rationalismen betragtede Djævletroen og Antagelsen af en evig
Fordømmelse for nødvendige Bestanddele af Kristendommen. Hans senere
filosofiske Forsken og personlige Livserfaring førte ham Skridt for Skridt
videre. --- I sin „Kosmologi“ opfattede han Guds Menneskeblivelse, ikke som en
først paa et enkelt bestemt Tidspunkt foregaaet Begivenhed, men som en stadigt
foregaaende Proces, som paa et enkelt Punkt naaede sit Højeste, sin
Kulmination; --- ligesom Menneskeheden udviklede sig af eller i Dyriskheden
udviklede igen Guddomsheden sig af eller i Menneskeheden. Inkarnationen sker
„fra Evighed
% --- p. 107 ---
\setcounter{footnote}{0}%
\opage{107}af“ --- det vil sige: nu og altid! I denne spekulative Omtydning af
det kirkelige Dogme ytrer sig Sibberns Trang til at faa den nøjeste mulige
Forbindelse mellem sine religiøse Forestillinger og sine filosofiske Ideer. ---
Hans personlige Livserfaring førte ham ogsaa bort fra Ortodoxien. Stedse havde
det jo været den Aandsnæring, han fandt i Kristendommen, der havde draget ham
til den. Og nu fandt han altfor meget Udvortes i den gængse, overleverede
Opfattelse af Kristendommen. Naar den personlige Erfaring og Tilslutning er det
Afgørende, saa maa man lade Enhver gaa de Veje, der egne sig for ham. Man maa
lade Enhver beholde det, hvorved han føler sig fremmet i Henseende til det Ene,
som er fornødent, og lade ham virke frit i sin Retning, naar han kun vil
indrømme os Frihed til det Samme. Man lade Enhver beholde sin Fetisch! Enhed
paa det religiøse Omraade har kun Værdi, naar den kommer som naturlig Virkning
af, at Sandhedserkendelse trænger igennem hos Alle. Anmasselsens Veje, de
udvortes Foranstaltninger, Bogstavvæsenet fører kun til at kue det selvstændige
Liv. Ad mange Veje tilstrømmer der aandelig Vederkvægelse --- fra Humanisme og
Naturalisme paa deres Vis ligesom fra Ortodoxien, og fra Videnskab, Kunst og
fri Menneskekærlighed, saavel som fra Religionen. Naar man dogmatiserer, rykker
man Sandheden løs fra dens levende Rod og gør en Mumie af den, i hvilken
Skikkelse den saa overleveres fra Slægt til Slægt. --- Man kalde ikke dette for
Subjektivisme, ti Alt, hvad har Værdi og senere formes som objektiv Lære, har
havt sit første Udspring i Sjælens Inderste. Netop i de inderligste Bevægelser
i Sindet kunne store Tilværelseskræfter røre sig og stræbe at arbejde sig frem.
Læren om den sporadiske Udvikling finder ogsaa Anvendelse her. Det er Aandens
Vej. Og desuden: hvad der har fundet Anerkendelse i en stor Menneskekreds, f.
Ex. en bestaaende Kirke, bliver derfor ikke mere objektivt end det, der hos den
Enkelte subjektivt gør sig gældende. Det Objektive vil sige det Gyldige, og
dette beror ikke paa ydre Formuleringer og Organisationer.

% --- p. 108 ---
\setcounter{footnote}{0}%
\opage{108}Det var i to Universitetsprogrammer fra 1846 og 1847,%
% footnote *) on p. 108
\footnote[1]{\textit{Om den christelige Ytringsfrihed i kirkelig Henseende.}
1846. --- \textit{Bidrag til at opklare den christelige og kirkelige Frihed.}
1847.} at Sibbern fremsatte disse Tanker, der paa et væsentligt Punkt stemme
overens med den Sætning, Søren Kierkegaard samtidigt proklamerede:
Subjektiviteten er Sandheden! --- Sibbern er utvivlsomt kommen til sit
Standpunkt ad sin egen Vej --- og desuden drager han langt radikalere
Konsekvenser af Sætningen, end Kierkegaard gjorde.

Men det blev ikke Sibberns sidste Ord angaaende Kristendommen. I nøje
Sammenhæng med de Anskuelser om etiske og sociale Spørgsmaal, som udviklede sig
hos ham under hans stadigt fortsatte Tankeliv og formedelst hans levende Sans
for det Menneskelige, gjorde han i sine senere Aar en langt skarpere Forskel
mellem det Forgængelige og det Blivende i Kristendommen. I det Fantasibillede
af et Fremtidssamfund, han har givet i sit mærkelige Skrift \textit{Meddelelser
af Indholdet af et Skrift fra Aaret 2135} (1858---1872) retter han strenge
Bebrejdelser mod Kristendommen. Den har i al sin Tid langt mere fostret og
næret Anmasselse, Herskesyge og Hadoptændthed end Menneskekærlighedens Aand.
Fremtidsmenneskene ville derfor ikke kalde sig Kristne, de ville ikke have Navn
sammen med dem, der i Aarhundreder havde øvet Anmasselser og Grueligheder i
Kristendommens Navn, og det ikke blot i Katolicismen (Eneraadighedskirken) men
ogsaa i Protestantismen (Indsigelseskirkerne). De skelne mellem Kristelighed og
Kristenskab. De læse i Naturens og Livets store Bog, ikke blot i Bibelen, skønt
de ogsaa af denne tage, hvad der kan kvæge Aand og Hjerte. De leve ud af en
Kraft, der rører sig i deres Indre; deres egen Fremtidsstræben og den
guddommelige Virken gaar i Et. Jesus er for dem det Menneske, i hvilken Lysets
og Livets Fylde første Gang traadte fuldt og helt frem i Menneskeheden. ---
Igennem Kristendommen er der vel trods Kristenskabet gaaet en
Menneskekærlighedens Strøm. Men hvorledes er man ikke ofte gaaet frem overfor
Hedningene? De mange Smaa blandt Hedningene, der havde et gudfrygtigt Sind,
% --- p. 109 ---
\setcounter{footnote}{0}%
\opage{109}burde Menneskekærligheden ligesaa vel have undgaaet at forarge, som
de Smaa, der troede paa Kristus. Flerguderiet havde sin gode Grund. Enhver
blive ved sin Tro, naar den er ham til Fortrøstning og til Fremme for det, som
Pligt og Menneskeomsorg kræver! --- Spørges der, om der kan vides Noget om et
andet Liv, svare Fremtidsmenneskene (Nyeuropæerne), at hvad enten der er en
Udødelighed eller ikke for de Enkelte, ville de leve Livet, som de leve det;
dog synes det, at det Alvirkende rører sig i Alle, at føre til Tro paa
Udødelighed. ---

Der er en dyb Sammenhæng mellem Sibberns religiøse Ideer, saaledes som de
efterhaanden udvikle sig, og hans filosofiske Anskuelser. Hans Tro paa, at al
Udvikling foregaar sporadisk, ud fra mange forskellige Punkter, og dog fører
til et harmonisk Resultat, maatte naturligt føre ham til at lægge Vægten paa de
enkelte Personligheder med deres Trang og deres Erfaring, og paa det, der kunde
arbejde sig frem i deres Indre. Derfor taler han saa stærkt mod al Anmasselse,
--- og han vil heller ikke selv være anmassende (hvad hans Bøger ogsaa --- ved
deres ubehjælpsomme Form --- vare langt fra at kunne være). Han siger derfor i
Begyndelsen af sin Fremtidsskildring, at Ingen bør læse den, hvis han mærker,
at „den vil forvirre ham det, hvori han har sit gode og trygge Tilhold“. ---
Neppe har hos nogen dansk Forfatter psykologisk Sans og Menneskelighedsfølelse
været forbunden med saa megen Aandsfrihed som hos gamle Sibbern. I hans sære
Bøger er der paa mere positiv Maade end noget andet Sted henvist til en
menneskelig Løsning af det religiøse Problem.

\begin{center}\rule{2cm}{0.4pt}\end{center}

Den Opfattelse, Sibbern havde af Filosofiens Opgave, Metode og Betydning,
hænger nøje sammen med selve Karakteren af de Anskuelser, han nærede om
psykologiske, kosmologiske og religiøse Forhold. Han har udtalt sig om sin
Opfattelse af Filosofien overhovedet dels i sit Skrift om \textit{Filosofiens
Begreb, Natur og Væsen} (1843), dels i sit Skrift om \textit{Hegels Filosofi
betragtet i Forhold til vor Tid} (1838).

% --- p. 110 ---
\setcounter{footnote}{0}%
\opage{110}For Sibbern er det Filosofiens højeste Opgave at føre til en
altomfattende Verdensanskuelse. Men denne skal bygges paa det i Erfaringen
Givne. Og den første Opgave bliver da den at foretage en Gennemgang eller en
Analyse (en „Explikation“) af det Givne. Her forudsætter Filosofien ikke blot
de andre Videnskaber, men den foretager selv en saadan Gennemgang paa det
psykologiske, logiske, æstetiske, etiske og religiøse Omraade. Først naar denne
Analyse er foretagen, kan den spekulative Konstruktion, der vil udlede Alt af
en Grundidé, begynde sit Værk. Han dadlede Hegel, fordi denne vilde gaa lige
til Konstruktion og udvikle alt Tilværelsens Indhold i en eneste fortløbende
Tankerække. Og dog havde Hegel selv ved at skrive sin „Phänomenologie des
Geistes“ indrømmet, at der behøvedes en Grundlæggelse og en Indledning, før
Spekulationen kunde tage Fart. Det er det givne Tilværelsesindhold, Filosofien
vil se i Ideens Lys, --- og den maa da først overbevise sig om, at dette
Indhold virkeligt er givet. Dernæst kommer det an paa, om Indholdet lader sig
føre tilbage til én Grundtanke. Her maa der paa Grund af Indholdets
Mangfoldighed og Forskellighed føres en lang Debat, før et System kan danne
sig. Et er Grundideen --- et Andet er, hvorledes vi kommer til den! Af en
eneste abstrakt Tanke kan Intet udledes, og Hegels tilsyneladende rent logisk
fremskridende Tankegang har da ogsaa paa hvert Punkt ved Tilsnigelse
(Subreption) optaget et i Erfaring givet Indhold. Men Erfaringen lader sig ikke
ordne i en eneste Række; den frembyder ikke blot Elementer, der kunne stilles i
Under- og Overordningsforhold til hverandre, men ogsaa sideordnede Elementer.

Sibberns Kritik af Hegels Filosofi indeholder mange træffende Bemærkninger og
rammer flere Hovedpunkter i denne spekulative Teori, der i Trediverne og
Fyrrerne gjorde sig ikke lidet gældende og fandt Tilhængere i Mænd som Heiberg
og Martensen. ---

At en spekulativ Verdensanskuelse --- paa Grundlag af Erfaring, Analyse og
Debat --- er mulig, faar efter Sibbern sin Forklaring ved, at hvad der
konstituerer Tilværelsen i
% --- p. 111 ---
\setcounter{footnote}{0}%
\opage{111}dens Inderste, ogsaa konstituerer vor aandelige Tilværelse i dens
Inderste. Han giver altsaa sin Erkendelsesteori (hvad den højeste Art af
Erkendelse angaar) et religiøst-metafysisk Grundlag. I sin „spekulative
Kosmologi“ har han selv gjort et Forsøg, hvis Grundtræk vi kende fra det
Foregaaende. Men den spekulative Teologi, han endnu bebudede, blev ikke færdig,
dels vistnok paa Grund af den afgørende Ændring i hans religiøse Anskuelser,
der indtraadte i Slutningen af Fyrrerne, dels paa Grund af, at det sociale
Spørgsmaal fra den Tid af, som vi skulle se i det Følgende, i saa høj Grad tog
hans Interesse fangen.

En fuldstændig Systematisering af Filosofien var nu ifølge Sibbern ikke heller
mulig. Og de Grunde, han anfører hertil, ere i høj Grad karakteristiske for
hans hele Standpunkt og hans Livsanskuelse.

Erfaringen er spredt og mangfoldig. Verden oplader sig ikke fra først af for os
fra sit Inderste alene, men fra mange Punkter paa én Gang, og desuden skrider
den bestandigt videre. Dette Sporadiske vil aldrig ganske kunne samles til ét System.

Og Forskeren er desuden selv et enkelt af disse Punkter. Filosofien er en af de
mange sporadisk opdukkende Processer i Verden. Hvorledes skulde den da kunne
overskue det Hele? --- Stridighederne mellem Filosoferne ere Tegn paa, at hver
har taget sit Udgangspunkt fra et forskelligt Moment i Tilværelsen, som de
fremhæve med ensidig Styrke og i ensidig Retning. Intetsteds ses det tydeligere
end i de filosofiske Debatters og de religiøse Bevægelsers Sfære, at vort
aandelige Liv endnu befinder sig i den første Dannelsesproces.

Selvom nu til en given Tid et System har samlet Tilværelsesindholdets
Hovedmomenter og paavist deres Udspring af en Grundide, saa maa dog --- paa
Grund af Erfaringens og Livets stadige Fremskriden --- Arbejdet i enhver ny
Tidsalder tages op paany. Ti medens Filosofien systematiseres, vil Livet have
aabnet Mere af sin Fylde, og ny Kundskab og nye Synspunkter ville da blive
nødvendige. Debatten maa be%
% --- p. 112 ---
\setcounter{footnote}{0}%
\opage{112}gynde igen. --- Derfor var Sibbern overbevist om, at det ikke kunde
have sit Forblivende med Hegels Filosofi, og han forsvarede overfor Heiberg
Poul Møller, der en Tid havde været Hegelianer, men senere følte Trang til en
anden Anskuelse, hvorfor Heiberg kaldte ham en Desertør. Sibbern hævdede, at
hvad Poul Møller havde tragtet efter, var et rigere Erfaringsindhold end det,
Hegel havde systematiseret.

Endeligt har Filosofien for Sibbern kun „intermediær“ Betydning, --- Betydning
som et Mellemled. Det Opgør, som den filosofiske Eftertanke foretager, skal
føre til en saa meget des dybere Tilegnelse af Livets reale Magter, Kunst og
Poesi, Kærlighed og Religion, saavel som den virkelige Natur, Erfaringens hele
Rige. Filosofien kan være Tankens, men ikke Livets højeste Form, allerede af
den Grund, at Tanken stedse kun er et af Livets Elementer. Ved Slutningen af
sin Filosofi staar Tænkeren stedse overfor de andre Livselementer.

\begin{center}\rule{2cm}{0.4pt}\end{center}

Paa Sibberns specielle Behandling af Logik, Æstetik og Etik gaar jeg her ikke
ind, saa meget Interessant og Dygtigt den frembyder. Derimod vil jeg til
Afslutning af hans Karakteristik i Korthed omtale hans politiske og sociale
Anskuelser.

Fra Stændertidens første Dage helt op i sin høje Alderdom deltog Sibbern i den
politiske Debat. Det var jo en af hans Filosofis Hovedsætninger, at al
Udvikling gaar frem gennem en Debat af Alt imod Alt. Hans egne Indlæg i den
politiske Strid kom ikke til at spille nogen stor Rolle; men de vare i høj Grad
karakteristiske for ham selv ved deres Aand og Tone, og de vidne om, hvor klart
et Blik og hvor varmt et Hjerte han havde.

Der er i politisk Henseende visse Hovedtanker, han stadigt fastholdt. Han lagde
stor Vægt paa en grundig og alsidig Overvejelse af Sagerne saavel som paa
Muligheden af en energisk Regering. Idealet var en raadgivende Forsamling og en
stærk, af varm Omsorg for Almenvellet, især for „Lavmandsfolkets“ Vel besjælet
Regering. Atter og atter saa han
% --- p. 113 ---
\setcounter{footnote}{0}%
\opage{113}tilbage til den oplyste Enevældes Reformer, især til „de Mænd, der
med oplyst og lutret Villie stode for Bondesagen“. Hvad de Liberale kaldte
Frihed, var i Virkeligheden Del i Magten. Naar man forlangte Garantier, oversaa
man, at alle statsretslige Garantier tilsidst ere moralske. Desuden --- Alt
skal kontrolleres, ogsaa den kontrollerende Repræsentation! --- I 30'erne og
40'erne var Sibbern derfor imod en Konstitution, en „Papirspave“.
\textit{(Patriotiske Intelligensblade.} 1835. --- \textit{Dikaiosyne.} 1843).

Men skulde man have en Konstitution, saa skulde den være saa fri som mulig,
staa aaben for saa Mange som muligt. Han fandt ikke, at der i Danmark var
Betingelser for en Modsætning mellem et aristokratisk og et demokratisk Kammer.
Derfor holdt han paa Junigrundloven og udtalte sig baade mod Ministeriet Ørsted
og mod Forfatningsændringen 1866. \textit{(Om og i Anledning af Kampen mellem
de tvende højeste Statsmyndigheder i Danmark.} 1854. ---
\textit{Samfundsbetragtninger.} 1865---1870. Og en Række Pjecer fra 1865 ---66).

Men allerede meget tidligt var det sociale Spørgsmaal for ham vigtigere end det
politisk-konstitutionelle. I denne Henseende er hans Universitetsprogram fra
1849 \textit{(Nogle Betragtninger over Stat og Kirke)} af Interesse. Han
skelner her skarpt mellem Stat og Folk. Statens Magt og Organisation skulde
ophæve den Krig af Alle mod Alle, der menes at herske i Naturtilstanden. Men
har den ikke snarere gjort Krigen stærkere, hvassere, mere omgribende og
indgribende, dels mellem Staterne eller Folkene indbyrdes, dels indenfor det
enkelte Folk? Det er indenfor Folket Kampen for Erhvervet, endog for det simple
Livsophold, der tager Fart. Konkurrencens Evangelium fører til den mest nervøse
Kamp af Alle mod Alle. Konkurrencen har i høj Grad befordret Vindskibeligheden;
men hvilke ere de Guder, for hvilke der derved er rejst Altre, og hvis Dyrkelse
og Rige er derved bleven fremmet? Det gaar ikke blot ud over dem, der stride
for det Nødvendige, men virker ogsaa nervesvækkende og tærende paa de
Formuende. Alle burde derfor være enige om at
% --- p. 114 ---
\setcounter{footnote}{0}%
\opage{114}sætte en anden organiserende Magt istedetfor den blinde og raa
Konkurrence. Den højeste Retfærdighed er den fordelende Retfærdighed, som gaar
ud paa at lade Alle faa Adgang til den hele Menneskeheds almindelige Eje.
Fødselen bør her ikke paa Forhaand udelukke eller begunstige. Hvad den Enkelte
selv erhverver, skyldes for en stor Del de sociale Forhold, under hvilke han
virker, saa at Selverhvervet ikke kan begrunde en absolut Ejendomsret. Det maa
være Statens Opgave at organisere Arbejdet saaledes, at Alle kan faa Del i
Udbyttet og at tilfældige Forspring udelukkes. I Enevældens Arbejde for
Bondestandens Bestaaen og Udvikling have vi Exempel paa en saadan organiserende
Virken. Kun Skade, at den ikke er kommen Husmandsstanden til Gode! ---
Statsmænd maa i Tide sætte sig ind i Kommunismens Ide, der har Realitet til
Grundlag og Realitet til Formaal. Har den end Mørke bag sig og Mørke for sig,
saa har den dog ikke, som det konstitutionelle Væsen, Blændværk eller Halvhed
for og bag.

Ogsaa her anvender Sibbern sin Lære om Udviklingens sporadiske Karakter. Der er
mange spredte Tilløb og Bestræbelser i Retning af det ene rette Maal; men da
der mangler Sammenhæng imellem dem, kan det se ud, som om det Hele var absurd.
Ligesom Menneskene have vist sig oplagte til, ja opsatte paa at tro det
Urimeligste, saaledes have de ogsaa Tilbøjelighed til at gaa de mest forkerte
Veje; Absurditeten synes at have været Menneskehedens Vugge og Amme. Men tager
man Fosteret paa et tidligt Trin, synes det ogsaa forvredent og unaturligt, og
dog kommer en harmonisk Menneskeskikkelse ud deraf. Og tager man en Dreng i
Lømmelalderen, ser man ofte kun Uvornhed og Balstyrighed, Uro og Skiften, og
dog kan der komme en god og tilfredsstillende Tilværelse ud, naar Metamorfosen
er gennemgaaet. Med disse Analogier trøster Sibbern sig: Menneskeheden er
maaske endnu i Lømmelalderen, --- om den ikke endog først har sin Fødsel i Vente!

Da Sibbern udgav dette Program, havde han allerede begyndt Udarbejdelsen af sin
mærkelige Utopi \textit{Meddelelser af et Skrift fra Aar 2135.} Den første
Fortale er dateret 1846.

% --- p. 115 ---
\setcounter{footnote}{0}%
\opage{115}I dette Skrift kaster Sibbern et kritisk Tilbageblik paa det
nittende Aarhundrede i Europa, specielt i Danmark, særligt hvad de politiske,
økonomiske og religiøse Tilstande angaar. Han gentager tildels sin Kritik af
Konstitutionalisme og Parlamentarisme; det er desuden her, at hans stærkeste
Udtalelser om „Kristenskab“ og Kirkevæsen fremkommer. Vi skulle nu særligt
dvæle ved det socialøkonomiske Fremtidsideal, han lader sin af Menneskelivets
Meningsløsheder, Forvridninger og Lidelser vakte Ynk og Harme diktere sin Fantasi.

Ved Slutningen af det nittende Aarhundrede vansmægtede Sjælene. Al Glæde og
Livsfriskhed var borte. Man stønnede under Regerings- og Parlamentuvæsenet,
under Ejendomskvakleriet, under Loves og Skatters Tryk, saavelsom under
Kirkelærdommenes Anmasselse. Overalt havde ydre Former og stive Dogmer trængt
sig ind mellem Mennesket og det indre Aandens Rige. Staten skulde ophæve
Krigen, men havde kun sat indre Krig istedetfor ydre. Der var mange Mennesker,
som ikke vare optagne i det fælles Liv og ikke nød godt af dets Fordele; de
vare blevne Udskud --- men man havde selv skudt dem ud. Det var Mangelen paa
sand Menneskeomsorg, som havde gjort Straffe- og Fattigvæsenet, de europæiske
Staters partie honteuse, nødvendige og mulige.

Der kom nu en Sløvheds og Søvnens Tilstand over Menneskene. Alt stod stille
eller var i Opløsning. Og samtidigt virkede ogsaa Naturaarsager til den store
Hvile, idet underjordiske Dampe stege op og gjorde al Virksomhed umulig. Da man
saa vaagnede op, var man klar over, at man ikke vilde indføre det gamle Stats-
og Kirkevæsen igen, saalidt som Rets-, Handels- og Ejendomsvæsen. Arbejdet
organiseredes saaledes, at Alle toge Del baade i legemligt Arbejde og i
aandelig Stræben og Leven, og at Sjælene ikke bleve kuede, men udviklede og
berigede gennem Arbejdet. Sjælevelstanden blev det Første og det Vigtigste. Der
var offentlige Forraadshuse, i hvilke alle frembragte Produkter opbevaredes, og
fra hvilke Enhver kunde faa, hvad han behøvede. Ingen tog Mere, end han
behøvede --- ti hvad skulde han med det? ---
% --- p. 116 ---
\setcounter{footnote}{0}%
\opage{116}Med en Mængde naive Enkeltheder udstyres Fremtidsbilledet. Overalt
ytrer sig en mærkelig Aandsfrihed og en stor Sans for det indre Sjæleliv og
dets Trang.

Denne sociale Fremtidsfantasi staar paa en interessant Maade i Modsætning til
Romantikens Fordybelse i Fortiden og dens Idealiseren af det Svundne. Haabet
var traadt i Erindringens Sted, og det hos en Mand, der selv med Varme havde
gennemlevet Romantikens Tid og været besjælet af dens Idealer. Det var foruden
den stadige indre Stræben den friske Sans for det Menneskelige, der førte den
gamle Tænker til at fastholde sin Ungdoms Idealer paa saa fri og levende en
Vis, at de kunde iklæde sig helt andre Former. Han havde begyndt sin Ungdom med
at se tilbage og med Inderlighed fordybe sig i, hvad Fortid og Samtid havde
bragt af Tanke- og Livsformer; og netop fordi han saa trofast arbejdede i
disse, blev det muligt, at hans gamle Øjne kunde se fremad med saa stor Fortrøstning.

Det er naturligvis kun en Henvisning til et Problem, ingen Løsning, Sibbern
giver. Hans Socialisme er utopisk, ikke blot fordi den skildrer Tilstande, der
ganske afvige fra dem, vi kende, men ogsaa fordi han aldeles ikke antyder nogen
naturlig og historisk Udvikling fra Nutidens til den ideale Fremtids Tilstande.
Det er jo nemt nok at lade Alle sove sig bort fra alle Ulemper og sove sig til
de nye Fuldkommenheder. Ved sin Fortælling om den store Dvale har han jo ogsaa
ligefrem indrømmet, at han ingen Overgang kan paavise.

Den utopiske Socialisme staar i en ejendommelig Modsætning til den saakaldte
videnskabelige Socialisme, som netop mener at kunne paavise den nødvendige
historiske Udvikling, der fra den nuværende „kapitalistiske“ Periode skal føre
til Fremtidens Socialisme, men som til Gengæld helst undlader at gaa nærmere
ind paa at skildre Fremtidsordningen i det Enkelte%
% footnote *) on p. 116
\footnote[1]{Smlgn. om de forskellige Hovedformer for Socialisme, min
\textit{Etik3.} XXVI, 7.}. --- Senere hen vil jeg faa Lejlighed til at omtale
en ung, dansk Forfatter, der omtrent samtidigt med Sibbern
% --- p. 117 ---
\setcounter{footnote}{0}%
\opage{117}opstillede et lignende socialt Fremtidsideal, med langt mindre Sans
for personlig Inderlighed og for Følelseslivets Trang, men ogsaa med klar
Forstaaelse af, at Vejen henimod et saadant Ideal ikke gik udenom, men netop
igennem Konkurrencen og Parlamentarismen.

\begin{center}\rule{2cm}{0.4pt}\end{center}

Det er et omfattende Tankearbejde, et Tilbageblik paa Sibberns lange Liv viser
os. Desværre lykkedes det ham ikke i noget enkelt Værk at give sine Ideer en
saa klar og sluttet Form, at de kunde hævde sig en stadig og let tilgængelig
Plads i vor Literatur. Nu vil kun den, der sysler med det danske Aandslivs
Historie, stanse ved hans Bøger. Men det er da en ejendommelig Friskhed,
Hjertelighed og Dybde, der træder Læseren imøde, og som --- trods den sære Form
--- sikrer den gamle Tænker et uvisneligt Minde.

\begin{center}\rule{2cm}{0.4pt}\end{center}

% --- p. 118 ---
\setcounter{footnote}{0}%
\opage{118}

\begin{center}\large 12. POUL MØLLER\end{center}

Det er især Treschow og Sibbern, der have givet den danske Filosofi i det
nittende Aarhundrede dens ejendommelige Præg. Det Fremherskende hos dem var
psykologisk Sans og Interesse, Vægt paa Erfaringskundskab, Fremhæven af
Individualitetens og Individualitetsforskellighedernes Betydning, kritisk
Besindighed overfor Spekulationerne. Den største Skikkelse i dansk Tænknings
Historie, Søren Kierkegaard, vil vise os det samme Præg i en ny og ejendommelig
Skikkelse og tillige saaledes som det udformes ved den skarpeste Kontrast til
det mest spekulative af alle spekulative Systemer, det Hegelske, der havde
faaet stor Indflydelse her i Landet, især hos Æstetikere og Teologer.

Men førend vi tale nærmere om Hegelianismen her i Landet og om dens
lidenskabelige Modstander, maa vi dvæle ved et Mellemled mellem Sibbern og
Kierkegaard, et Mellemled, gennem hvilket Kierkegaard hænger nøje sammen med
tidligere dansk Tænkning, --- et Mellemled, der har sin selvstændige Interesse.

Enhver dansk Læser bliver varm om Hjertet ved at høre Poul Møller nævne, den
danskeste af alle danske Forfattere. Her skulle vi kun beskæftige os med en
enkelt Side af ham; men vi ville der genfinde ham helt og klart, langtfra, som
man undertiden har ment, blot i en tilfældigt fremkaldt Stilling. Hverken som
Digter, Filolog eller Filosof har Poul Møller efterladt sig Andet end
Fragmenter, men i disse har han lagt hele sin Personlighed for Dagen. Det Halve
kan jo undertiden være Mere end det Hele.

% --- p. 119 ---
\setcounter{footnote}{0}%
\opage{119}32 Aar gammel gik Poul Møller som Professor i Filosofi til
Kristiania (1826), og efter at have virket her i fem Aar, blev han Sibberns
Kollega indtil sin Død 1838. --- Paa sit Dødsleje sendte han gennem Sibbern en
Hilsen til „den unge Kierkegaard“. Sammenhængen mellem de tre Mænd faar sit
Udtryk heri.

I sin filosofiske Opfattelse var Poul Møller en Tid stærkt paavirket af Hegels
spekulative System. Sibbern hævdede senere overfor Heiberg, at det fulgte af
Poul Møllers Natur, at han maatte interessere sig for et System som Hegels, men
at det ikke mindre fulgte af hans Natur, at han ikke kunde holde fast ved et
saadant System. I Poul Møllers „Strøtanker“, det vigtigste Bidrag til Kundskab
om hans Filosoferen, finde vi Bekræftelse paa dette Udsagn.

Han saa Spekulationens Betydning deri, at ét Synspunkt konsekvent fastholdes og
gennemføres. I Naturen (Erfaringen) raader Mangfoldigheden og Flersidigheden.
Men et ensidigt System, selvom det skulde være en ren Fiktion, yder mere
virkelig Dannelse, naar man studerer det, end Syslen med en Sammenhobning af
Reflexioner. Man lærer at tænke med organisk Konsekvens. Den, der ikke føler
Trang til at ordne sit Tankesystem, maa enten have meget Geni eller liden
Tænksomhed. Det har sin store Betydning at gøre sig Rede for sin
Verdensanskuelse i dens hele Sammenhæng. Konsekvens udelukker Konfusion, som
fremkommer ved Reminiscenser fra en fremmed Verdensanskuelse. --- Træffende er
her den konsekvente Tænknings Sammenhæng med Personlighedens Enhed fremhævet.
Det er karakteristisk, at det særligt var denne Side, der interesserede Poul
Møller i hans spekulative Periode.

At han ikke blev paa det spekulative og systematiske Stade, havde først og
fremmest sin Grund i den Erkendelse, at Problemerne ikke fandt deres Løsning ad
Spekulationens Vej, men egentlig kun viste sig at ligge paa andre Punkter, end
man havde antaget: „Ved den spekulative Filosofi erholder man ikke saa meget
Svar paa de Spørgsmaal, som den almindelige Menneskeforstand fremstiller, som
man kommer
% --- p. 120 ---
\setcounter{footnote}{0}%
\opage{120}til at erkende, at Spørgsmaalene selv, som grundede paa Begreber,
der ingen Sandhed have, maa bortfalde.“ --- Men dernæst savnede han i den
spekulative Filosofi Sans for Individualitetens Betydning. De forskellige
Individualiteter opdages kun ad Erfaringens Vej, og hvorfor skulde det være en
ufuldkommen Erkendelse? Jeg udvider dog faktisk derved min Kundskab, skønt det
ikke sker ad apriorisk Vej. Og hvis kun den aprioriske Erkendelse er den sande,
kan Individualitet slet ikke erkendes. --- Poul Møller staar her ved Tanker,
som især Treschow havde indskærpet, og som havde spillet en væsentlig Rolle for
A. S. Ørsted og Sibbern.

I sin Afhandling \textit{Om Muligheden af Beviser for Menneskets Udødelighed}
(1837) udvikler Poul Møller disse Tanker noget udførligere.

En fuldendt Viden er uopnaaelig. Erfaringsvidenskaberne blive aldrig færdige
med at sammenfatte de existerende Individer og deres Tilværelsesmaader i et
fuldstændigt Begrebssystem. Og selvom en saadan Afslutning var mulig for et
vist begrænset Erfaringsomraade (Jordkloden, Menneskeslægten), saa kan hin
Sammenfatten dog kun ske ved Hjælp af Almenbegreber, og den fulde Anskuen gaar
tabt. Man maa holde sig til Arter og Slægter --- men hvilket uendeligt Indhold
forsvinder da ikke for os med de enkelte Individer!

Forlader man Erfaringens Vej og holder sig til de mest abstrakte Synspunkter,
kan man derved maaske bringe Tilværelsens almindelige Grundforhold til klar
Bevidsthed, men den hele Uendelighed af Bestemmelser (Kvaliteter) kan ikke
derved udtømmes.

En tredie, højere Erkendelsesmaade, der paa en Gang lader det Individuelle og
det Almene komme til sin Ret, staar ikke til vor Raadighed.

Alligevel mener Poul Møller, at der kan udvikles en Verdensanskuelse, i hvilken
ikke blot Individualitetens Realitet og Betydning hævdes, men ogsaa den
personlige Udødelighed kan godtgøres. Han indrømmer, at en saadan
Verdensanskuelse endnu kun er i sin Vorden; men han ser ikke, at han selv har
erklæret den for umulig, da hin „tredie, højere“
% --- p. 121 ---
\setcounter{footnote}{0}%
\opage{121}Erkendelse ikke er opnaaelig. Han appellerer til en Videnskab, hvis
Umulighed han selv har bevist, og som han iøvrigt omtaler i meget vage
Vendinger. --- Sagen var, at der --- især under Indtrykket af en stor Sorg ---
rørte sig hos ham en personlig Trang til Tro, som gjorde, at han nu stillede
sig i et mere sympatisk Forhold til den positive Religion end før. „Den
kristelige Tradition“ paaberaabes ofte i den nævnte Afhandling, og det
forlanges af den filosofiske Verdensanskuelse, at den skal lade denne Tradition
komme til dens Ret. Hvorvidt dette er muligt, kom han ikke til indgaaende at
undersøge. ---

Bruddet med den spekulative Filosofi har næppe været uden Smerte og Harme,
fordi der havde knyttet sig saa store Forventninger til den. Men hans sunde
Humor vandt igen Overhaand. „En Tidlang,“ siger Kierkegaard, „talte han næsten
med Indignation om Hegel, indtil den sunde humoristiske Sans, der var i ham,
lærte ham at smile, især af Hegelianismen, eller, for endnu tydeligere at
erindre om Poul Møller, ret hjerteligt at le af den. Ti hvo har været forelsket
i Poul Møller og glemt hans Humor; hvo har beundret ham og glemt hans Sundhed;
hvo har kendt ham og glemt hans Latter, der gjorde En godt, selv naar det ikke
blev En ganske tydeligt, hvad det var, han lo af; ti hans Distraktion bragte En
stundom i Vildrede.“ --- Dog fastholdt han Beundringen for Hegels store
Tankearbejde. Men det var nu saa, mente han, at enhver Kæmpe i Filosofien
skulde falde for Dværge. Og han havde ikke stor Respekt for dem, der søgte „at
gaa ud over Hegel“. En af dem sammenligner han med en Rotte, der graver sig ud
gennem Muren i en Domkirke. Mere end et Telt til Brug paa en Rejse kunde
foreløbigt ikke faas istedetfor den stolte Bygning. Et saadant Telt arbejdede
han selv paa lige før sin Død (en Lære om Kategorierne eller Grundbegreberne);
saavidt man kan se, havde det næppe egnet sig til en lang Rejse.

\begin{center}\rule{2cm}{0.4pt}\end{center}

% --- p. 122 ---
\setcounter{footnote}{0}%
\opage{122}Skønt Poul Møller henviste til Traditionen som den, der ogsaa skulde
komme til sin Ret for Tanken, var det dog hans bestemte Overbevisning, at kun
det aandelige Indhold har Værdi, der ikke blot tilegnes, men frembringes paany
i hvert enkelt Individ. I hans æstetiske Opfattelse, saaledes som vi lære den
at kende af hans literære Anmeldelser, skelner han mellem den oprindelige
Kunst, der udspringer af Livet og Naturen, og den afledede Kunst, der
overvejende, om ikke udelukkende, er inspireret af den forudgaaende Kunst, ikke
af det egne umiddelbare Forhold til Virkeligheden. Poesien er en Frugt af
Folkets Liv. „Hvad et Folk og dets Individer have oplevet, afgiver i en vis
Betydning Stoffet for dets højere Liv i Digtekunsten. Men naar en rig poetisk
Literatur har udviklet sig, faar Kunsten efterhaanden mer og mer
Selvstændighed, idet Hovedkilden til nye Digtninger stedse mere bliver den
existerende poetiske Literatur, og den oprindelige Poesi, der med Friskhed og
Originalitet gik frem af det virkelige Liv, bliver et bestandigt sjeldnere
Fænomen..... Naar der hos et uforholdsmæssigt Antal af et Folks Individer
opstaar et Hang til at føre et poetisk Tilskuerliv, da maa en poetisk Periode
snart være forbi, og Malmet hentes da sikkert ikke længer fra en Afgrund, der
kun er tilgængelig for Geniet.“ \textit{(Recension af Extremerne).}

Ligesom Romantikens Filosofi var altsaa ogsaa dens Poesi for Poul Møller kommen
til et kritisk Punkt. Et bestemt Program kunde han ikke give for Fremtiden ---
men for begge Omraaders Vedkommende var han sikker paa, hvor Malmet skulde
hentes. I en levende personlig Erfaring, i nøje Sammenhæng med Livet, skulde
Tænkning og Digtning hente den nye Kraft. Og i enkelte af sine Digte, saavel
som i mangfoldige af sine Tanker, har han selv givet Forbilleder --- desværre
kun lidet udførte. Hvad hans Filosoferen angaar, er der især én Kreds af
Tanker, der i hans senere Aar har beskæftiget ham, og som, saa brudstykkeagtig
den er, dog betegner hans vigtigste Indsats i dansk Tænkning. Ved den maa vi
dvæle lidt udførligere.

\begin{center}\rule{2cm}{0.4pt}\end{center}

% --- p. 123 ---
\setcounter{footnote}{0}%
\opage{123}Naar Poul Møller fordrede en nøjere Sammenhæng mellem Tanke og
Erfaring, end Romantiken havde vidnet om, saa udsprang det for en stor Del af
hans levende Sans for det Oprindelige, Ægte og Friske. Han fordrede fremfor
Alt, at Tanken skulde hænge sammen med Menneskets eget inderste Væsen, have sin
Rod i dette og udtrykke dette. Det er et stort Tab for dansk Literatur, at den
Afhandling \textit{Om Affektation,} som han havde paatænkt, ikke blev skreven.
Kun nogle Strøtanker og et Udkast paa nogle faa Blade foreligge. Poul Møller
havde faaet Øje for, hvad allerede Holberg saa stærkt fremhævede, den hyppige
Forvexling af Skyggen med Legemet.

Det var dog ikke blot af erkendelsesteoretisk og psykologisk Interesse, at han
beskæftigede sig med dette Emne. Det var heller ikke blot af etisk Interesse
for Ærlighed og Redelighed i Tankelivet, for personlig Sandhed.%
% footnote *) on p. 123
\footnote[1]{Udtrykket „personlig Sandhed“ forekommer, saavidt jeg véd, første
Gang hos Poul Møller. Enstydigt dermed bruger han Udtrykket „det personlige
Livs Sandhed“.} Det var tillige, som Vilhelm Andersen har paavist, personlige
Skuffelser og bitre Erfaringer, som voldte Livslede og fremkaldte en skeptisk
og trodsig Stemning, som han synes at have tænkt paa at give Udtryk i en
digterisk Fremstilling, til hvis Hovedfigur han havde valgt Ahasverus. Der
synes at være en indre Sammenhæng mellem Forarbejderne til „Affektation“ og til
„Ahasverus“%
% footnote **) on p. 123
\footnote[2]{\emph{Vilhelm} \emph{Andersen}: \textit{Poul Møller.} 1894. p.
244. --- Jeg kan ikke følge Vilh. Andersen i den Indflydelse, han mener,
Ahasverus i Reflexionerne have havt paa Kierkegaard. Derimod er det for mig
(som jeg allerede har udtalt i min Bog \textit{Søren Kierkegaard som Filosof.}
1892. p. 24---27) utvivlsomt, at Poul Møllers Drøftelse af Affektationen har
havt en stor Betydning for Kierkegaard, skønt han giver Lessing hele Æren. Og
bag Poul Møller staar igen Sibbern. --- Hvad Poul Møller --- se nedenfor ---
kalder „vekslende Affektation“ henleder Tanken paa Kierkegaards „Vexeldrift“ (i
første Del af „Enten---Eller“).} Poul Møller hørte til dem, der helst tænkte
med hele Personen. Det er da sikkert, at den Tankegang hos ham, som nu skal
fremdrages, har været af afgørende Betydning for ham.

% --- p. 124 ---
\setcounter{footnote}{0}%
\opage{124}Affektation er paa én Gang Falskhed og Selvbedrag. Derimod er
Falskhed for sig og Selvbedrag for sig ikke Affektation. Selvbedraget kan være
helt uforskyldt og er da ikke Affektation. Og den, der forstiller sig med fuld
Bevidsthed, er heller ikke affekteret. Til Affektation udfordres en Stræben
efter at være, hvad man ikke kan være; men en saadan Stræben er umulig, naar
man ikke for en Stund indbilder sig, at man er, hvad man ikke er.

„Det Menneske, der har Affektation i sit Liv, bestemmer sig forsaavidt ikke med
fuldkommen moralsk Frihed; dets Handlinger have ej deres Udspring af det sande
Selv, som er dets frie, sædelige Villie. Dets Villie bestemmes af et eller
andet blot naturligt Øjemed, hvorved det bringes til at forestille en fremmed
Person eller paatage sig en falsk Rolle, som ej er det anvist i Livet.“

Jeg forstaar dette saaledes: Der er foruden de skiftende Forestillinger og
Interesser hos ethvert Menneske (og i hver enkelt Livsperiode) et centralt
Indhold, visse dybt liggende Tanker, Formaal og Interesser, --- en Villie, der
ikke behøver at ytre sig umiddelbart i hvert Øjeblik, men stadigt ligger til
Grund. Det er, hvad jeg i min Psykologi (V B, 5) kalder det reale Selv. Kun de
aandelige Livsytringer, der stemme med det, og især de, der udspringe af det,
ere i Sandhed frie og ægte; kun i dem er der personlig Sandhed. Men ofte
skjuler det Centrale sig; det overvældes af ydre Paavirkninger, eller man skyr
den Anstrængelse, der kan være forbunden med at gøre det gældende. --- Det er
et mindre heldigt Udtryk, Poul Møller bruger, naar han taler om sædelig Villie
og moralsk Frihed. Ti det drejer sig aabenbart om et Forhold mellem det Indre
og det Ydre, som er det samme, hvad enten det Indre (det reale Selv) er
sædeligt eller ikke. Hvad det gælder om, er, at Menneskets Indre virkeligt
lægger sig for Dagen i Tanke, Ord og Handling. Deri bestaar personlig Sandhed.
Det bliver et Spørgsmaal for sig, om det Indre, der aabenbarer sig, har moralsk
Værdi eller ikke. ---

Poul Møller tilspidser sin Opfattelse i en Sætning, som han ved, vil blive
betragtet som fantastisk: „Ingen Livsytring har
% --- p. 125 ---
\setcounter{footnote}{0}%
\opage{125}Sandhed, uden deri ligger skabende Selvvirksomhed. Udsigelsen,
Fremstillingen af en Tanke har egentlig kun da Sandhed, naar Tanken i samme
Øjeblik frembringes. Siden er der allerede Mimik deri. --- Men desuagtet er det
menneskeligt at reproducere Tanker; ti kun en Guddom er altid skabende.
Mennesket maa i Tankernes Rige regnes til de drøvtyggende Dyr.“

Skønt han tager et Forbehold, fastholder han dog som Idealet den højeste
Selvvirksomhed i Tanke og Ord. Og han venter sig meget af en Tilnærmelse til
dette Ideal: „Naar hver enkelt Mand, uden at befrygte Dadel for Enfoldighed,
dømte om Tingene saaledes, som de fremstillede sig for ham, da vilde herlige
Fænomener fremstaa.“ Men af Magelighed eller Fejghed lader man ikke sin
Personlighed komme frem; man tror ikke paa sin Persons uendelige Dybde. Man
skammer sig for at staa nøgen og tør derfor ikke forlade de traditionelle
Forestillinger. Man ser ikke, at en Tanke, der udsiges med inderlig
Overbevisning, ikke kan være aldeles urigtig, --- at den indeholder et
Synspunkt, selvom dette Synspunkt med Urette anses for det eneste mulige.---

Man ser, at Ahasverus-Stemningen hos Poul Møller kun angaar det hyppige
Misforhold mellem det Indre og det Ydre. Han er overbevist om, at der er
skjulte Skatte i Menneskenes Indre; det er kun hans Sorg, at de ikke komme frem.

I Lyset af Fordringen om skabende Selvvirksomhed fremtræder Henvisningen til
den kristelige Tradition i Afhandlingen om Udødeligheden som en Modsigelse,
eller dog som et Problem: hvorledes kan en Overlevering indeholde Sandheden,
naar Sandheden væsentlig beror paa Menneskets egen Selvvirksomhed? Dette
Problem tog Poul Møller ikke op. For Kierkegaard blev det netop Hovedproblemet.
Han vilde vise, at netop den inderligste Selvfordybelse, den strengeste
Gennemførelse af den Sætning, at Subjektiviteten er Sandheden, vilde føre over
til det Ideal, der indeholdes i „det fra Fædrene Overleverede“.---

Der kan ifølge Poul Møller være forskellige Grader af Affektation. --- I den
\textit{momentane} Affektation er det af Van%
% --- p. 126 ---
\setcounter{footnote}{0}%
\opage{126}vare, man overskrider sin Personligheds Grænser. Man har for en
Stund beruset sig i Noget, man endnu ikke helt har tilegnet sig. Især i
Ungdommen gør man sig let skyldig i smaa halvt frivillige Uærligheder af denne
Art. Det kan endog være et godt Tegn, da det vidner om, at man ikke egoistisk
indeslutter sig i sig selv. Og en dybere Selvforstaaelse kan blive Følgen, naar
den personlige Usandhed opdages. --- Den \textit{faste} Affektation indtræder,
naar den paatagne Rolle fastholdes, bliver Vane og forfølges med Konsekvens.
Der gøres en uvilkaarlig Modstand mod at blive reven ud af Selvbedraget;
Personligheden forvanskes; man gør sig selv til en Skifting. En ærlig Løgn er
ikke saa slem som denne Selvforvanskning; ti den fuldt bevidste Løgn kan man
snarere aflægge end den uægte Rolle, man har vænnet sig til at spille. --- Den
\textit{vexlende} Affektation er tilstede, hvor enhver indre Kærne mangler. Man
danner sig i hvert Øjeblik en ny Personlighed forat aflægge den i næste
Øjeblik, enten af Lune eller af Hensyn til Omgivelserne. I sin Fuldendelse er
denne Art af Affektation Et med et moralsk Selvmord. ---

Som allerede antydet, ses det ikke, at Poul Møller er kommen ind paa det
Problem, hvorledes man --- for at bruge hans eget Udtryk --- fra at være
Drøvtygger kan blive selvskabende. Mennesket maa ligesaavel optage aandelig
Næring som legemlig, og i begge Henseender maa Næringen svare til Organismen,
kunne gøres til Kød og Blod i denne. Der bliver da Spørgsmaal om, hvorvidt al
Næring, der tilbydes, kan gøre Fyldest i denne Henseende. En lille Antydning,
der viser, i hvilken Retning Poul Møller maaske havde kunnet anvende sin Ide om
personlig Sandhed paa det religiøse Omraade, giver følgende Strøtanke: „De
moderne Afguder, som gøres af Ord, staa deri tilbage for de antike af Træ og
Sten, at de lade Tænkningen mindre Frihed end disse. Enhver antik Kunstner gav
dog en Zeus efter sit eget Ideal. Nu skal enhver Teolog fremsige visse Formler
om det højeste Væsen.“ --- En filosofisk Anvendelse indeholdes i den korte
Strøtanke: „Affektation er Symmetri i filosofiske Systemer.“ Meningen er
aabenbart, at man ofrer Iagttagelsers og Tan%
% --- p. 127 ---
\setcounter{footnote}{0}%
\opage{127}kers oprindelige og umiddelbare Natur forat faa dem til at passe ind
i den systematiske Form; saaledes skulde ifølge Hegel Tanken stedse bevæge sig
frem i treleddede Forhold, og Indholdet lempedes efter en Form, der ikke
dannede sig uvilkaarligt. --- Paa det æstetiske Omraade var den afledede Poesi
Tegn paa Affektation: „Den moderne Poesi søger Livet i Poesien paa enkelte
Steder og udpiller dem. Deraf dens Affektation.“ ---

Det er, trods den fragmentariske Form, i hvilken Poul Møllers Tanker ere os
tilgængelige, dog en karakteristisk Type for dansk Filosoferen, der fremtræder
i dem. Han peger tilbage til Holberg og fremad til Kierkegaard.

\begin{center}\rule{2cm}{0.4pt}\end{center}

% --- p. 128 ---
\setcounter{footnote}{0}%
\opage{128}

\begin{center}\large 13. HEIBERG OG MARTENSEN\end{center}

I det Foregaaende er Hegels' Filosofi allerede flere Gange omtalt, og det er nu
Stedet at dvæle ved den Indflydelse, den vandt i Danmark.

Hegel søgte at forbinde Kants ædruelige Stræben efter Begrundelse med
Schellings og Steffens' romantiske Fordybelse i Naturens og Aandens Indre. Han
vilde anvende en bestemt Metode og gøre Ende paa Romantikernes mystiske Svælgen
og Fantaseren. Det blev hans Grundtanke, at naar Menneskeaanden har
gennemvandret Erfaringens og Historiens lange Vej --- en Vandring, han selv,
paa meget romantisk Vis, beskrev i „Phänomenologie des Geistes“ --- aabner der
sig for den en Kreds af Tanker, i hvilken Alt fra det Højeste til det Ringeste
kan og maa finde sin Plads. Og indenfor denne Tankekreds er der en saa inderlig
Sammenhæng, at man ved at begynde med en enkelt bestemt Tanke, var det endog
med den mest abstrakte af alle, Begrebet Væren overhovedet, med logisk
Nødvendighed kommer over til alle de andre Tanker. Enhver enkelt Tanke udløser
sin Modsætning, og derpaa bliver det Opgaven (en Opgave, som ifølge Hegel altid
kan løses) at finde en Tanke, der som en højere Enhed indbefatter begge
Modsætningsled i sig. Enhver Tanke faar gennem denne Dialektik sin
ejendommelige Plads i hele Tankens Verden. Den maa opgive sin absolute
Selvstændighed og sin udvortes Modsætning til andre Tanker, og forsaavidt gaar
den til Grunde; men dens Ejendommelighed sikres den som Moment indenfor den
højere Enhed.

Men hele denne Tankebevægelse og Tankekamp er nu
% --- p. 129 ---
\setcounter{footnote}{0}%
\opage{129}ifølge Hegel ikke blot Noget, der foregaar i den menneskelige Aand.
Det er Tilværelsens inderste Væsen, som derigennem afslører sig. Ved en
Analogislutning, han ikke er sig bevidst opfatter Hegel Alt, hvad der sker i
Naturens og Historiens Verden, som et Udtryk for en kosmisk Dialektik, gennem
hvilken de enkelte Væsener og Begivenheder vise sig som Led i en stigende Række
af højere Enheder. Hvad der i Erfaringen kan se ud som Undergang, er Overgang
til at være Led i en højere Sammenhæng. Og gennem denne indre Brydning, der kun
for den sanselige Opfattelse fremtræder som en Proces i Rum og Tid, ytrer sig
Verdensaandens, Guddommens Liv. Det Absolute lægger sig for Dagen i de relative
Formers Modsætninger saavel som i den Sammenhæng, Modsætningerne og deres
stadige Ophævelse godtgør.

Hegel deler Romantikens Overbevisning om det aandelige Livs indre Enhed.
Religion, Kunst og Filosofi ere for ham Udtryk for Et og det Samme, nemlig for
hin indre Enhed, der kundgør sig i Sammenhængen og ligger til Grund i selve
Brydningen. Men af disse Aandsformer, i hvilke det Absolute bliver aabenbart,
er Filosofien den højeste, fordi den tænker og begriber, hvor Religion og Kunst
kun forestiller, føler, aner eller anskuer.

\begin{center}\rule{2cm}{0.4pt}\end{center}

Det var \textit{Johan Ludvig Heiberg} (1791---1860), der først indførte Hegels
Filosofi i Danmark. Han havde studeret den hos Hegel selv i Berlin og var
traadt i personligt Forhold til Mesteren. Paa Tilbagerejsen fra Berlin kom der
da et Øjeblik, hvor Systemets Grundtanker med Et stode for ham med fuld Klarhed
og i deres hele Sammenhæng. Og Heiberg har hele sit Liv holdt fast ved disse
Grundtanker og ved Overbevisningen om, at kun ved dem kunde der komme Enhed i
det aandelige Liv, og at særligt Kunst og Religion kun ved at ses i Lyset af
dem kunde komme til at staa som fuldt berettigede Aandsformer.

Det var i den Strid om Villiens Frihed, Howitz havde
% --- p. 130 ---
\setcounter{footnote}{0}%
\opage{130}fremkaldt, at Heiberg først optraadte som filosofisk Forfatter
\textit{(Om den menneskelige Frihed.} Kiel 1824). Senere var det især æstetiske
Spørgsmaal, der gave ham Anledning til at anvende den nye Filosofi, navnligt
til at anvise de forskellige Kunstarter deres gensidige Plads \textit{(Om
Vaudevillen.} 1826. --- \textit{Om Malerkunsten.} 1838). Dog behandlede han
ogsaa spekulativ Teologi \textit{(Treenighedslæren.} 1837). Og han har desuden,
til Brug ved sine Forelæsninger over Logik ved den militære Højskole, ladet
trykke en Fremstilling af Logiken efter Hegels Metode. --- Nogle af de nævnte
Afhandlinger fremkom i et Tidsskrift, han begyndte \textit{(Perseus. Journal
for den spekulative Ide),} men af hvilket kun to Numre udkom.

Den Plads, Heiberg i Tilslutning til Hegel anviste Filosofien indenfor det
aandelige Liv, ses især af hans Skrift \textit{Om Filosofiens Betydning for den
nærværende Tid} (1833).

Hvad Tiden lider under, er for det Første --- mener Heiberg --- en
Overbefolkning af Ideer. Ved Siden af Kunst, Poesi og Religion gøre
naturvidenskabelige og politiske Interesser og Synspunkter sig mere og mere
gældende. Men i hvilken Orden, i hvilket gensidigt Forhold skulle alle disse
Ideer og Interesser stilles til hverandre? Her kan kun en streng Sammenligning
af de forskellige Grundtanker indbyrdes bringe Afgørelsen. Allerede i et Brev
til H. G. Ørsted fra Aaret 1825 har Heiberg fremhævet det som Hegels store
Fortrin fremfor andre Filosofer, at Alt findes i hans System og finder sin
bestemte Plads og den Gyldighed, der tilkommer det. --- Men Tidsalderen
frembyder ikke altid en uordnet Mangfoldighed af Ideer. Der fremtræder ofte en
Ensidighed, idet en enkelt Ide eller Interesse, f. Ex. den politiske eller den
religiøse, sætter sig paa Højsædet som den ene afgørende og bestemmende. Ogsaa
her er Lægemidlet kun at søge i en Undersøgelse af de enkelte Ideers eller
Interessers indbyrdes Forhold. Ingen Ide eller Interesse, der har virkelig
Betydning, tør skydes tilside for andre, og det er tilsidst Sandheden, Livet
som Helhed, der skal anvise hver enkelt dens Plads. --- Ad begge Veje føres vi
saaledes til Filosofien som den gennemførte Lære om Grundtankerne (Kategorier%
% --- p. 131 ---
\setcounter{footnote}{0}%
\opage{131}ne) i deres indbyrdes Forhold. Filosofien fremstiller i Sammenhæng
den Sandhedens Ide eller Ideverden, til hvilke alle de enkelte Interesser og
Bestræbelser tilsidst maa appellere deres Sag.

Men Filosofien bringer dog ikke noget absolut Nyt. Den er stedse kun en
Besindelse paa eller en Erkendelse af, hvad der allerede er givet. Den
indeholdes forsaavidt i den sunde Menneskeforstand. Den er Erkendelsen af den
Sandhed, der er indeholdt i de forskellige menneskelige Erfaringer og
Bestræbelser --- i Natur og Stat som i Kunst og Religion. Filosofien udfolder
sig først, naar det gælder om at besinde sig paa en Udviklingsperiode, der er
afsluttet; ved at omsætte det udviklede Indhold i Tankens Form gør Filosofien
det evigt nærværende.

Og ved at aabne vore Øjne for, hvad vi allerede besidde, bringer Filosofien os
ogsaa Hjælp for en Sygdom i det aandelige Liv, der er ikke mindre betænkelig
end Tankernes kaotiske eller ensidige Karakter. Tidsalderen lider af en
Modsætning mellem Ideal og Virkelighed, mellem det Uendelige og det Endelige.
De endelige og stridende Bestræbelser have skudt det Uendelige og Evige ud i en
fjern Baggrund. Hvad vi skulle lære, og hvad Filosofien kan lære os, det er at
se, hvorledes det Evige og Uendelige netop ligger til Grund for de enkelte og
endelige Bestræbelser, saaledes at det, vi søge, er os givet allerede i selve
vor Søgen, ikke blot som et fjernt og hinsidigt Maal. Hegels Filosofi forsoner,
ligesom Goethes Poesi, Idealet med Virkeligheden, vore Fordringer med det, vi
besidde, vore Ønsker med det, som er opnaaet. ---

Det var Heibergs Program. Vi forstaa nu den Interesse, han havde for de
abstrakte Begrebsudviklinger i den hegelske Logik%
% footnote *) on p. 131
\footnote[1]{Heiberg har foretaget en Ændring i Logiken, der utvivlsomt er En
Berigtigelse af Hegels Fremstilling. Han gør opmærksom paa, at \textit{Vorden}
ikke som hos Hegel kan indtage Plads som en „højere Enhed“, da Vorden betyder
Uro, en udvortes Blanding af at være og ikke at være. Istedetfor Hegels første
Trehed „Væren---Intet---Vorden“ sætter Heiberg da Treheden „Væren (der som rent
abstrakt og indholdsløs er Ikke-Væren)---Vorden---Tilværen“, hvor Tilværen er
Resultatet af den Vorden, til hvilken Væren gaar over, naar det ses, at den
indeholder den Modsigelse at være „Ikke-Væren“. \textit{Perseus} II p. 44. ---
Selvfølgeligt hviler ogsaa Berigtigelsen paa Hegels egne Forudsætninger.}: det
var i denne abstrakte Tankeverden, i hvilken
% --- p. 132 ---
\setcounter{footnote}{0}%
\opage{132}hvert enkelt Livs- og Erfaringsomraade var repræsenteret ved de
Grundbegreber, der udtrykte dets Væsen, --- det var dér, at Kampen mellem de
forskellige Bestræbelser tilsidst udfægtedes. I Hegels System var nu Logiken
kun den ene Trediedel; den anden Del var Naturfilosofien, i hvilken de
abstrakte Ideer fremtræde i Rummets og Tidens Former; den tredie Del var
Aandsfilosofien, i hvilket det i Tid og Rum paa udvortes Maade Sondrede
samledes i Aandens Enhedsformer, tilsidst i selve Tanken, hvorved Systemet i
sin Afslutning vender tilbage til sin Begyndelse og altsaa danner en sluttet Kæde.

Heiberg, der overalt søger at finde Treheder og, hvor han finder Treheder,
søger at give dem rent abstrakte Udtryk, sammenligner Systemets tre Dele med
Helvede, Skærsilden og Paradiset i Dantes Digt. Selv har han da holdt mest af
at opholde sig i Helvede (Logiken) og i Paradiset (Aandsfilosofien); Skærsilden
(Naturfilosofien) har han ikke beskæftiget sig med. Det er maaske et Tegn paa
hans Kritik; ti Naturfilosofien var, som Sibbern med Rette bemærkede i sit i
mange Henseender interessante Skrift om Hegels Filosofi, Systemets partie
honteuse. Derimod har han, som allerede bemærket, arbejdet paa Æstetikens og
Religionsfilosofiens Omraade.

Naar Sibbern i det nævnte Skrift kaldte Heiberg en filosofisk Dilettant,%
% footnote *) on p. 132
\footnote[1]{Til Gengæld fik Sibberns Bøger i „En Sjæl efter Døden“ en
Hædersplads i Biblioteket i Helvede (der her er den evige Kedsommeligheds
Hjem).} var det forsaavidt berettiget, som Heiberg er i den Grad overvældet af
den Hegelske Dialektik, at han ofte anvender den paa en meget udvortes Maade,
uden Hensyn til Emnernes Ejendommelighed. Hos Hegel selv finder man dog en
dybere Sans for de Realiteter, han mente at have givet Tankeudtryk for i sine
Begreber. Mellem Mesteren og
% --- p. 133 ---
\setcounter{footnote}{0}%
\opage{133}Disciplen viser der sig her en ejendommelig Forskel baade paa
Æstetikens og Religionsfilosofiens Omraade.

Heiberg har forsøgt en Systematisering af Æstetiken, endnu før Hegels Æstetik
var udkommen. Hans Interesse gaar her ud paa en Inddeling af Kunstarterne --- i
Analogi med de abstrakte logiske Kategorier, og saaledes at Treheden --- et
Element, dets Modsætning og den højere Enhed af begge --- stedse paany kommer
frem. Skønheden fremtræder først som Formskønhed og Naturskønhed; dernæst som
Kunstskønhed; og endelig som Poesi, der danner en Overgang til Filosofien.
Indenfor Poesien træder igen (i \textit{Afhandlingen om Vaudevillen)} Lyrik og
Epos i Modsætning til hinanden; begges højere Enhed er Dramaet, indenfor
hvilket der ogsaa indbyrdes er Modsætning mellem lyrisk og episk Drama (hint
skildrende Karakteren, dette Situationen) --- uden at her den højere Enhed
angives. Den hele Udvikling tjener nemlig til at give Vaudevillen (som en Art
Situationsdrama) dens bestemte Plads, og Undersøgelsen ender derfor her. Senere
\textit{(i Indledningsforedrag til det logiske Kursus)} stiges der fra
Tragedien gennem Komedien op til det af Calderon repræsenterede lyriske Drama
som Poesiens højeste Form.

Overfor Oehlenschlägers Indvending, at saadanne Inddelinger i faste Arter ere
til Skade for Poesiens frie Udvikling, svarer Heiberg, at man ærer de store
Digtere paa bedste Maade ved at betragte de forskellige Former, de have
anvendt, som noget Fast og Bestemt, ikke som noget Vilkaarligt. Man godtgør den
dybe Fornuft, der aabenbarer sig i det digteriske Geni, og sikrer derved
Poesien dens høje Rang i Livet. Men det var dog en underlig formel
Klassifikation, Heiberg her forfaldt til. Hegel gik en anden Vej, idet han
opfattede Kunsten i dens Sammenhæng med den historiske Udvikling og derved fik
den store Trehed af orientalsk, klassisk og romantisk Kunst. Han bevægede sig i
sin Dialektik overvejende frem gennem historisk Karakteristik, og de abstrakte
Kategorier træde i Baggrunden. Sin Æstetik havde Hegel hentet fra samme Kilde
som sit System idethele; men Heiberg havde først sat sig ind i Systemet og
søgte saa bagefter
% --- p. 134 ---
\setcounter{footnote}{0}%
\opage{134}paa en skematisk Maade at anvende det paa de specielle Omraader.

Noget Lignende gælder hans religionsfilosofiske Udvikling. Ligesom for Hegel,
saaledes have ogsaa for Heiberg Religion og Filosofi det samme Indhold, kun
Formen er forskellig: det Absolute (Gud) fremtræder hist i Forestillingens
(Fantasiens), her i Tankens Form. Ifølge \textit{En Sjæl efter Døden} kan man
komme i Himmerige paa to Maader: enten ved at fordybe sig i Kristi historiske
Skikkelse, eller ved spekulativt at have „fattet Gud“. Kun denne sidste Form er
den fuldkomne Form, og til den maa der derfor atter og atter gaas tilbage,
dersom Religionen skal have Gyldighed. Ti kun det er gyldigt, der kan antage
Tankens Form. Filosofien beviser ikke Guds Tilværelse, der slet ikke kan
bevises, naar den ikke fra først af forudsættes. Men Filosofien gør
Gudsforestillingen, som forudsættes given, til et Gudsbegreb, tænker Gud som
det Absolute, der ikke har nogen Modsætning overfor sig, fordi det selv er den
Enhed, der gør alle Modsætninger og al højere Enhed af Modsætninger mulig. I
den kristelige Treenighedslære finder Heiberg en Overensstemmelse med Hegels
Filosofi: „Den gennem det Hegelske System gaaende stadige Trilogi er selve
Treenighedens Afglans i Tankens, Naturens og Aandens Riger og den absolute
Betingelse for al Filosofi.“

Skal Gud tænkes som uendelig, maa man ved Uendelighed ikke tænke paa en
ubegrænset Udstrækning i Rum og Tid, der kun er en stadig Overgang fra et
Endeligt til et andet, en Sum af Endeligheder; det er hvad Hegel kalder den
slette Uendelighed.%
% footnote *) on p. 134
\footnote[1]{Heiberg mener, at det kopernikanske System har virket til at
befæste Forestillingen om den slette Uendelighed. „Man er blevet ubillig mod
vor Klode, ligesom om man kom Gud nærmere ved at bo i et større og prægtigere
Hus, ligesom ikke vor Erkendelse af Gud var en Erkendelse af det Højeste og
altsaa ikke kunde overgaas af nogen anden Erkendelse i det mindste i Kvalitet.“
\textit{(Perseus} I, p. 27.)--- Smlgn. \emph{Hegels} \textit{Encyclopädie} §
280, hvor Planeten er den højere Enhed af Solen og Maanen (eller Kometen). ---
Hvad der foregaar paa den, maa altsaa være det Vigtigste! Men der er saa mange
Planeter!} Det sande Uendelige er det, der i sig selv
% --- p. 135 ---
\setcounter{footnote}{0}%
\opage{135}har en Fylde, hvis Elementer staa i indre, nødvendigt Forhold til
hverandre. Saaledes er der i Selvbevidstheden et indre Forhold mellem Subjekt
og Objekt; de ere paa intet Punkt udenfor eller ligegyldige for hinanden. Der
er her en i sig selv tilbageløbende Cirkel, hvori Intet er fast eller
uforanderligt, først eller sidst, Aarsag eller Virkning, Middel eller Øjemed.
Men det sande Uendelige kan ikke staa i udvortes, og altsaa endeligt Forhold
til Endeligheden, der breder sig i Tiden og i Rummet. Skal det Uendelige ikke
have sin Skranke i det Endelige, og altsaa begrænses ved dette, hvorved det jo
netop vilde ophøre at være uendeligt, --- saa maa det selv aabenbare sig i
Endeligheden, saa at denne er en nødvendig Form, gennem hvilken det naar til
fuld Bestemthed. Og det er jo netop, hvad Kristendommen lærer: gennem
Menneskeblivelsen træder Gud selv ind i Endeligheden; den anden Person i
Treenigheden betegner netop Differensen og Endeligheden, forsaavidt denne har
Plads i Gud; og i den tredie Person, Aanden, er da den fuldkomne Harmoni af det
Uendelige og det Endelige udtrykt. --- Ved denne Opfattelse mener Heiberg at
være ude over Striden mellem Rationalisterne og de Ortodoxe. Medens hine sætte
det Uendelige udenfor det Endelige, sætte disse det Endelige udenfor det
Uendelige. Begge Dele er selvmodsigende.

Ogsaa hos Hegel forekommer logiske Abstraktioner som Tydninger af religiøse
Antagelser. Men man mærker dog ofte hos Hegel en personlig Livserfaring eller
en historisk Fordybelse i den Kamp og Smerte, „Endeligheden“ volder, inden den
kan „ophæves“. Heiberg gaar for rask over fra Erfaring til ren Logik. Ikke for
Intet satte han Calderon langt over Shakespeare som poetisk Repræsentant for
Menneskeheden, fordi hin opløser Alt i Uendelighedens rene Form, medens denne
bliver stikkende i Karakterskildringer og Begivenheder, der vække Forundring og
Skræk, men ikke pege ud over det Endelige! --- Og for Heiberg selv var den rent
logiske Tydning afgørende for Muligheden af at hævde et religiøst Standpunkt.
Han skriver til H. C. Ørsted (1825), at han vilde være Ateist, hvis hans
Filosofi ikke lærte ham
% --- p. 136 ---
\setcounter{footnote}{0}%
\opage{136}den i Naturen og Historien aabenbarede Aand. Det har sikkert været
hans Alvor med denne Bekendelse, han har været overbevist om den Forsoning af
Tro og Viden, hans Filosofi lærte. En Tro, der ikke saaledes kunde omsættes til
begrebsmæssig Viden, vilde for ham være umulig. Han har nærmest hørt til, hvad
man i 40'erne kaldte det Hegelske Højre, ikke til det Venstre, der med Strauss
og Feuerbach drog radikale Konsekvenser af Hegels Religionsfilosofi.

Den særligt abstrakte Karakter, som Hegels Filosofi fik i Heibergs Hænder,
finder ikke blot sin Forklaring ved, at han ikke som Mesteren møjsommeligt
havde hentet Ideerne op fra Virkelighedens Gruber, men ogsaa deri, at han i sin
Digtning fandt det bedste Udtryk for den mere personlige og levende Side i sin
Aand. I Vaudevillerne var det Endeligheden, der --- forklaret i Spøg og lyrisk
Stemning --- hævder sin Plads som et Moment i det Uendelige; men i hans
Tankedigte er det en inderlig Forening af Ideer og Følelser, der faa et i vor
Literatur enestaaende Udtryk. Mange abstrakte Forflygtigelser kunne tilgives
den, hvis Livsanskuelse har ført ham til at skrive de skønne Linier:

Den Gud, du søger, er din egen Gud;

Hvad vil du mere?

Han træder ikke som en Genstand ud

Blandt andre, flere.

Du føler ham, hvergang dit Hjerte banker

Af hellig Lyst;

Du hører ham i dine bedste Tanker

Som evig Røst.

Idet du søger, du ham alt har fundet,

Idet du spørger, er alt Svaret vundet.

Da Heiberg i Aaret 1841 udgav sine „Prosaiske Skrifter“, mindede han i Fortalen
om, hvilken Modstand den Hegelske Filosofi havde mødt, idet de kaldede Lærere
først havde for%
% --- p. 137 ---
\setcounter{footnote}{0}%
\opage{137}sømt den og derpaa søgt at holde den under streng Karantæne [et Hib
til Sibbern og Poul Møller]; men han glædede sig over, at „det, man havde
udskreget som en Modesag og et tomt Formelvæsen, nu ved en yngre Arbejders
kraftige, geniale Virken er blevet causa victrix“. --- Denne yngre Arbejder var
\emph{Hans} \emph{Lassen} \emph{Martensen} (1808---1884), som just da under
stor Tilslutning holdt sine Forelæsninger over spekulativ Filosofi og Teologi
ved vort Universitet.

I sine Studieaar havde Martensen været stærkt betagen af Hegel. Især var han
begejstret over den Udsigt, Hegels Filosofi gav til en spekulativ Teologi, der
kunde staa som en højere Enhed overfor Rationalisme og Ortodoxi. Til
Rationalisterne regnede han ogsaa Schleiermacher, dels fordi denne søgte et i
den menneskelige Natur liggende Grundlag (Afhængighedsfølelsen) for sin
Teologi, dels fordi han sluttede sig til Kants Indskærpen af den videnskabelige
Erkendelses Grænse. Da Schleiermacher i Aaret 1833 besøgte København, opsøgte
den unge teologiske Kandidat ham og drøftede Forholdet mellem Dogmatik og
Filosofi med ham. „Engang,“ fortæller Martensen i sit Levned, „spurgte jeg
ganske naivt, om han antog, at en filosofisk Erkendelse var mulig af Guds Væsen
i sig selv, af den indre, evige Livsproces i Gud... Han svarede ganske roligt:
Ich halte es für eine Täuschung.“ Det advarede dog ikke den unge, begejstrede,
spekulative Teolog. Selvom man, mente han, gav Schleiermacher Ret i, at vi kun
kunne tænke i Modsætninger, udelukker dette dog ikke et videnskabeligt
Gudsbegreb, da Gud ikke vilde være den levende Gud, hvis der ikke var
\textit{indre} Modsætninger i hans Væsen --- hvad jo netop Treenighedsdogmet
hævder. Her sluttede han sig til Hegel. Men at han brød med Hegel kom af, at
denne ligesom Kant og Schleiermacher vilde gaa den rent menneskelige
Erkendelses Vej. I sin Disputats „om den menneskelige Selvbevidstheds Autonomi
i vor Tids dogmatiske Teologi“%
% footnote *) on p. 137
\footnote[1]{\textit{De autonomia conscientiæ sui humanæ in Theologiam nostri
temporis introducta.} Hauniæ 1837. --- En dansk Oversættelse udkom 1841.}, fra
hvilken man har regnet „én ny Æra“ i Teologien, hævder han, at der ikke gives
en rent selvstændig
% --- p. 138 ---
\setcounter{footnote}{0}%
\opage{138}Erkendelsesteori, men at Grundlaget for al Erkendelse er
Bevidstheden om Mennesket som et af Gud skabt Væsen. Dette have Kant og hans
Efterfølgere (Schleiermacher og Hegel) overset. Kun naar Bevidstheden i Kraft
af en oprindelig Samviden (eller Samvittighed) fra først af gaar ud fra al
Realitets Kilde, kan dens Viden blive mere end rent formel, og da bliver en
Erkendelse af det Uendelige og Absolute mulig, skønt Mennesket selv er et
endeligt Væsen.

Den selvstændige (autonome) menneskelige Viden kommer aldrig ud over sig selv,
griber aldrig den egentlige Realitet, men stanser ved en uovervindelig
Modsætning mellem Viden og Magt. Kun gennem Forholdet til Gud, i hvem Viden og
Magt ere evigt forbundne, komme vi ud over denne Modsætning. Gud er autonom,
Mennesket kan ikke være det. Ganske vist beror det paa Menneskets Villie, om
det vil anerkende hin Menneskeaanden iboende oprindelige Gudskundskab; men
ligesaa vist er det, at begynder man ikke med Gud, saa begynder man med at
negte ham. Tro er Forudsætning for Erkendelse. Credo ut intelligam (jeg tror
for at kunne forstaa), som Anselmus sagde.

Martensen indrømmer, at Teologien maa forudsætte en Erkendelsesteori, og gør
opmærksom paa, at enhver Tids Dogmatik er afhængig af vedkommende Tids
Opfattelse af den menneskelige Erkendelse. Men denne Indrømmelse kommer han
igen let over ved det simple Kunstgreb at gøre al Erkendelsesteori til Teologi.
Den filosofiske Erkendelsesteori bygger derimod paa en Analyse af den faktisk
givne Videnskab og er et Forsøg paa at fremstille dennes Forudsætninger i dens
Sammenhæng med den menneskelige Aands Natur. Martensen har hverken i sit
Begyndelsesarbejde (i flere Henseender det mest Gennemtænkte, han har ydet)
eller senere indladt sig paa at vise, hvorledes man ud fra hans dogmatiske
Forudsætninger kan komme til den faktisk givne Videnskabs Metoder og
Resultater, ej engang, hvorledes hans Forudsætninger kunne bestaa sammen med
disse. Og dog føler han selv en Brod her, idet han erklærer, at Vanskeligheden
ikke er at forstaa, at Gud er, men at forstaa, at der kan være
% --- p. 139 ---
\setcounter{footnote}{0}%
\opage{139}Andet \textit{foruden} Gud, da det synes at følge af Guds absolute
Natur, at der Intet kan være foruden ham. Denne Vanskelighed er paa den anden
Side allerede bortfalden ved det Skabelsesdogme, hvorpaa den teologiske
Erkendelsesteori bygger. --- Selvfølgeligt falde alle Gaader bort, naar man fra
først af gaar ud fra, at de ere løste.

Alligevel lægger Martensen stor Vægt paa, at Dogmet underkastes en nærmere
Analyse. Det er netop Teologiens Opgave at give en spekulativ Udfoldelse af
det, der \textit{ligger} i Dogmet%
% footnote *) on p. 139
\footnote[1]{Munus theoligiæ dogmaticæ in iis quæ dogmate implicite continentur
speculative explicandis versatur. \textit{De Autonomia.} p. 29. f.}. Udvikler
man ikke de i Religionen indeholdte spekulative Ideer, saa faar man paa den ene
Side en Videnskab, der staar udenfor Religionen, og paa den anden Side en rent
historisk Teologi, ja, kun en Bibliologi som i den gængse Tro
(supranaturalismus vulgaris). --- Martensen bebuder her en Opfattelse af
Teologien, ifølge hvilken den er alle Videnskabers fælles Moder og Midtpunkt.
Det var ogsaa, som han fortæller i sit Levned, hans Ungdomsideal, fremkaldt og
næret ved Steffens' teologiske Skrifter: „Der maatte kunne gives en Verdens- og
Livsanskuelse, i hvilken Alt hvad der i Tilværelsen har Betydning, Natur og
Aand, Natur og Historie, Poesi, Kunst, Filosofi føjer sig harmonisk sammen til
et Aandens Tempel, i hvilket Kristendommen er det altbeherskende og
altforklarende Midtpunkt.“

Til at rejse dette Tempel, paa hvilket han i Teori og Praxis arbejdede hele sit
Liv igennem, uanfægtet af selv den voldsomste Indsigelse, vilde han dog ogsaa
gøre Brug af filosofiske Hjælpemidler. Hin „Explication“ skulde finde Sted
gennem den af Hegel fundne dialektiske Metode. Derfor erklærer han, at
Teologien og Filosofien have den samme Metode. Han var overbevist om, at naar
Tanken fordyber sig i Dogmet, vil dettes Indhold udfolde sig efter dets Natur,
lade sine forskellige Momenter fremtræde i indre Sammenhæng. De specielle
Dogmers Række fremstillede han da saaledes tolv Aar efter i sin Dogmatik (1849)
--- og unegteligt mærkedes det her ikke, at der var store Vanskeligheder at
overvinde. (
% --- p. 140 ---
\setcounter{footnote}{0}%
\opage{140}Kun lige tilsidst, ved Dogmet om den evige Fordømmelse, finder han,
--- hvad der tjener mere til Ære for hans Humanitet end for hans Logik --- et
„Kors for Tanken“). Det er dog ikke som hos Hegel rent logiske Overgange
(gennem Paavisning og Løsning af Modsigelser), men det er den til historisk
Dogmekundskab støttede og af religiøs Fantasi vejledede Anskuen, der udretter
Arbejdet. Det er intet Tilfælde, naar Martensen paa sine ældre Dage hellere
taler om, at hans Dogmatik var en Reproduktion af Kirkens Lære, end at den var
en Explikation.

Martensens Forhold til Heiberg er her ganske ejendommeligt. Heiberg forstod
ikke ret, hvorfor Martensen ikke var tilfreds med, hvad Hegel gav ham. Hegel
forudsætter, hævder Heiberg, stedse det Umiddelbare; det er jo netop det, han
vil omsætte i Begrebets Form; men naar man begynder at filosofere, er det
Umiddelbare den Genstand, der filosoferes over, og det kommer da kun an paa, om
man i sin Filosoferen kan udtømme det. Et Barn har Forældre; men det skal selv
leve sit Liv senere hen, som om det stod ene i Verden og ikke tilbringe sit Liv
med at tænke paa sine Forældre. Man kan godt sige med Martensen, at Gudsfrygt
er Begyndelsen til Visdom; men naar vor Visdom formes til System, maa vi
begynde med Tvivlen og se til, hvorvidt den kan løses; Gudsfrygten skal man
have i sit Hjerte, Tvivlen i sit Øje. Da Martensen imidlertid vil beholde
Hegels Metode, vil Heiberg ikke regne ham (saaledes som Poul Møller) til
Desertørerne%
% footnote *) on p. 140
\footnote[1]{\textit{Perseus.} I. p. 33---41. --- Heibergs Udtalelser angaa en
Afhandling af Martensen i „Maanedsskrift for Litteratur“, der ligger forud for
hans Disputats.}.

Hvad Heiberg ikke saa, var, at „det Umiddelbare“ hos Martensen ikke blot
betyder Genstand; det betyder ogsaa Autoritet. Han begynder med et Dogme:
Skabelsesdogmet, og bliver indenfor det. Martensen har med fuld Ret hævdet, at
han ikke var Hegelianer, da han optraadte som Forfatter. Han havde derimod
sluttet sig til den religiøse Tænker \emph{Franz} \emph{Baader}, hos hvem da
ogsaa Hovedsætningerne i
% --- p. 141 ---
\setcounter{footnote}{0}%
\opage{141}hans Disputats allerede findes. For Baader fører Autonomien (Tankens
og Villiens Selvstændighed) til Anomi (Lovløshed) --- og tilsidst til aandelig
Død. Ti, siger han, „enhver Aand lever kun i og af Beundring og Tilbedelse;
idet den underordner sig under det, den beundrer og tilbeder, gaar den ud paa
at forherlige, udbrede og fremstille det. Ved at vende sig bort fra den sande
og eneste Genstand for sin Beundring, Tilbedelse og Underkastelse, volder den
sin egen Død, saavidt den formaar det.“ Baader gør udtrykkeligt opmærksom paa,
at Hegel sammenblander „det Umiddelbare, der \textit{hæver mig op,} og det,
\textit{jeg ophæver, ---} det livgivende Aandepust fra oven, og den jordiske Næring“%
% footnote *) on p. 141
\footnote[1]{\emph{Franz} v. \emph{Baader}: \textit{Fermenta Cognitionis.}
Berlin 1822. Erstes Heft. p. 39; 67. --- Martensens Standpunkt minder iøvrigt
ikke blot om Baader, men ogsaa om Jacobi, af hvem den af Martensen som religiøs
Genius beundrede Biskop Mynster var saa stærkt paavirket. Smlgn. ovenfor
Afsnittet om Tyge Rothe, hos hvem vi ogsaa have truffet Indflydelse af
Jacobi.}. Det vil netop sige: Hegel (og Heiberg med ham) skelne ikke mellem
Genstand og Autoritet. Og filosofisk ere de her sikkert i deres gode Ret. Det
var i Virkeligheden ikke en „højere“ Filosofi, Martensen fandt; det var en
Overgang fra Filosofien, der har Genstande, men ingen Autoriteter, til den ikke
blot paa Autoritet, men paa Forvexling af Genstand og Autoritet grundede
Teologi. Med saare liden Ret saa Martensen ned paa supranaturalismus vulgaris.
At han vilde beholde Hegels Metode, var en Inkonsekvens, og man mærker heller
ikke Meget til den i hans Dogmatik; han har stiltiende ladet den falde, skønt
han ynder at operere med „højere Enheder“, især naar han er i Forlegenhed.

Paa de Forelæsninger, han i Fyrrerne holdt ved Universitetet, og som dels hans
Aandfuldhed og Veltalenhed, dels de store Forjættelser om en højere Enhed af
Tro og Viden gjorde til nogle af de mest besøgte og beundrede i vort
Universitets Historie, var det nu netop hans Opgave at vise, at der gaves en
højere Spekulation end Hegels, en Spekulation, man kun fik Adgang til paa den
kristelige Tros Grundlag. Han siger
% --- p. 142 ---
\setcounter{footnote}{0}%
\opage{142}selv i sit Levned om disse Forelæsninger, at naar han skulde føre
sine Tilhørere ind til sine Problemer og sin Løsning, maatte det ske gennem en
Oversigt over Filosofien fra Kant til Hegel. Han maatte først begejstre dem for
Hegel --- og saa føre dem ud over Hegel. Han havde Brug for den Tænkningens
Skole, han ikke fandt hos Nogen saa omfattende som hos Hegel. Han benyttede de
Hegelske Formler til at udvikle en Anskuelse, der skulde være meget forskellig
fra den Hegelske. Han indrømmer selv, at det i mange Tilfælde ikke lykkedes ham
at besværge de Aander, han havde manet frem. Mange stansede, naar han havde
ført dem til Hegel og vilde ikke gaa videre. Han trøster sig med, at Mange, der
senere kunde regnes til den danske Kirkes ypperste Mænd, trofast fulgte ham
helt igennem.

En ung Teolog fra hine Dage, der ikke holdt fast ved den spekulative Teologi,
skønt han ogsaa paa sine ældre Dage omfattede Martensens Personlighed med
Beundring og Kærlighed, --- Digteren og Præsten \emph{Johannes} \emph{Fibiger},
--- har givet en Skildring af den studerende Ungdom paa den Tid og af
Martensens Forelæsninger.

„Teologien var dengang anset for en ligesaa hæderlig som naadig Frue at tjene,
nok istand til at blænde. Jeg kom da ogsaa snart højt op i fuld Flugt paa
hendes fantastiske Vinger, farlige nok. Det laa saa i Tiden. Man maa have levet
i den Tid for at kunne gøre sig en Forestilling om dens underlige Væsen, ja,
blot for at kunne tro paa Muligheden deraf. Under den tyske Filosofis Enevælde
ivrede alle Tanker for at bygge paa Fantasteriets Babelstaarn; hvad vi hørte
omkring os, var intet Mindre, end at enhver Stortaler satte sin nærmeste
Livsopgave i at bygge et endnu højere. Universet med alle dets store og smaa
Lønkamre var afsøgt og klaret i Begrebet, alle Gaader vare løste, Hegel og hans
berlinske Discipelskare havde fuldbragt Værket. Men derved var ogsaa
Kunstbegrebet afsløret og Kampprisen fremlagt. Det maatte kunne eftergøres og
bringes endnu videre; man skulde kun bygge sit eget System og gaa ud over
Hegel, og selv blive Aandeverdenens Stormand. Vi indaandede snart ingen anden
Luft; det Eneste,
% --- p. 143 ---
\setcounter{footnote}{0}%
\opage{143}vi ved Universitetet hørte, der klang som en Fremtidsrøst, var
Martensens glimrende Tale, og der var ikke saa lidt i af luftige Stene til den
Slags Bygning. Teologien... havde dengang sin Glansperiode. Alle Aander svor
til Teologien; den var Aandrighedens Helligaand og havde vundet sig en ny
Skikkelse, var hverken Skriftlærdom eller Kirkens Troslære eller en
Præsteskole, Sligt var aflagte Stadier; --- den var spekulativ. Dette mystiske
Ord, der optraadte med det „Absolutes“ Klang, glimted gennem enhver af vor
Mesters Kraftsætninger som Solens Blink i et Vandspring. Hvorfra det kom,
vidste vi ikke; det var forbi som et Trylleri. Men da det kom til nogen bedre
Forstaaelse, viste det sig at være den Kunst med et eklektisk Aggregat af tyske
Systemer at kunne bringe Kristendommens Mysterier ind under Begrebets Dagslys,
saa sikkert som havde man en Engel sat i Spiritus stillet frem til Analysering.
Credo, ut intelligam, var Formlen. Troen „poneres“, Englen sættes frem.
Derefter begynder Intelligensen at lade sine Straaler falde gennem alle dens
Hjerteslag, Nerver og Muskler. Alt blev gjort fuldkommen klart paa spekulativ
Maade; kun var det ikke ganske klart, hvorledes det vilde gaa, hvis den bly
Genius, ubehageligt berørt af Intelligensens kolde Instrumenter, skulde give
sig til at forsvinde“%
% footnote *) on p. 143
\footnote[1]{\emph{Johannes} \emph{Fibiger}: \textit{Mit Liv og Levned, som jeg
selv har forstaaet det.} (1898). p. 73 f. --- Smlgn. Skildringen af Martensen i
hans ældre Dage p. 323---327.}.

Denne sidste Ytring berører det afgørende Punkt. Martensen vilde filosofere paa
Grundlag af den positive Religions ophøjede Billedsprog og lade sine
Spekulationer nyde godt af den Stemning, de overleverede Billeder fremkalde.
Han gjorde ikke som Hegel Alvor af at lade Begrebet afløse Billedet. Han vilde
„ikke føre bort fra Forestillingen [Billedet] for at komme til Begrebet, men
tvertimod søge at klare Forestillingen gennem Begrebet, fremstille Enhed af
Begreb og Anskuelse“. Men er det muligt at naa en saadan „Enhed“? For Begrebet
er ethvert Billede, enhver Anskuelse kun en Analogi eller et Exempel; men de
religiøse Forestillinger skulle
% --- p. 144 ---
\setcounter{footnote}{0}%
\opage{144}jo være Aabenbaringer, uerstattelige ved andre Forestillinger.
Derved alene er aabenbaret Religion forskellig fra Poesi. Martensen staar da
ogsaa vaklende her. Han giver ingen højere Enhed af Tanke og Billede, men giver
snart en Tanke, snart et Billede. Og som Teolog er han enig med Grundtvig i, at
„med Billedsproget tør der ikke leges“; konsekvent, naar Billederne skulle være
Aabenbaringer; --- men hvad skal saa egentlig Begrebet? Paa sine gamle Dage
udtalte han, at om det, der ligger ud over Erfaringen, kan man kun tale i
Billeder. „Dersom vi om disse Ting ikke vilde tale i Billeder, da vilde vi kun
komme ud i det Tomme, vilde komme til nogle fattige Begreber, der intet Lys
vilde give os. Lader os derfor holde fast ved disse Billeder, som Guds Ord
giver os... Det er Billeder, og dog er det mere end Billeder, ti de betegne de
højeste Virkeligheder“%
% footnote *) on p. 144
\footnote[1]{\emph{Martensen}: Af mit Levnet. II p. 7; 54. --- III p. 240. Den
sidst anførte Udtalelser fra en Paaskeprædiken 1881.}. Hvorfra han ved dette
Sidste, faa vi ikke at vide; men det er klart, at den højere Enhed af Begreb og
Billede er opgiven. Martensen har, forat bruge Fibigers Lignelse, fastholdt
Englen og lagt Instrumentet bort.

Men heller ikke dette gør han altid. Aaret efter den sidst anførte Udtalelses
Fremkomst udgav han sin Bog om Jakob Böhme (1882), denne Mystiker, til hvem han
var bleven ført af Baader, og som han altid havde interesseret sig for.
Lejlighedsvis kommer han her til at tale om Jesu Himmelfart og dens
Forenelighed med den moderne Opfattelse, at „Himmelen“ ikke er et fysisk Sted i
Rummet. Han paastaar, at „Himmel“ betyder „den ydre Verdens indre Grund“, og af
den synlige Himmelfart kun var et Tegn for Disciplene paa, at deres Mester
havde forladt Rummets Verden! Her har han ladet Begrebet omtyde Billedet, og
det paa en Maade, der minder om den gamle Rationalisme, som jo dog skulde være
et overvundet Standpunkt%
% footnote **) on p. 144
\footnote[2]{Smlgn. nærmere om Forholdet mellem den videnskabelige Opfattelse
af Rummet og de bibelske Forestillinger min \textit{Religionsfilosofi} II A b.
--- Smlgn. ogsaa ovenfor om Heibergs og Grundtvigs Forhold til det
kopernikanske System. Hos Rasmus Nielsen ville vi støde paa det samme
Spørgsmaal igen. Ingen af dem har været optagen af Problemet paa den Maade,
Holberg var det. Romantiken havde idethele en Tendens til at slaa det hen.}.
--- Den dybsindige Skomager
% --- p. 145 ---
\setcounter{footnote}{0}%
\opage{145}fra Görlitz, som Martensen saa begejstret hylder i den nævnte Bog,
var en friere og konsekventere Tænker end Martensen, og han maa derfor paa
flere Steder høre, at nu tænker han ikke kristeligt, men hedensk. Martensen har
overset en Udtalelse af Böhme, i hvilken denne overfor en lignende Indvending
siger: „Ich schreibe nicht heidnisch, sondern philosophisch!“ Distinktionen
mellem kristelig og hedensk „Tænkning“ hører ikke hjemme i Filosofien, men i
Religionshistorien. ---

Det er ikke Opgaven her at følge Martensen gennem hans teologiske Værker, i
hvilke Filosofien kom til at spille en stedse mindre Rolle. Typen for hans
Tænkning er dog overalt den samme: Paakaldelse af en højere Enhed. Som han
havde villet finde en højere Enhed af Anskuelse og Begreb, saaledes vilde han
som Kirkemand, i Tilslutning til Mynster, hævde en Retning, der betegnede en
højere Enhed i Forhold til Grundtvigianismens „ensidige Objektivitet“ paa den
ene Side og den indre Missions „ensidigt subjektive Retning“ paa den anden
Side. Den kunde da med Rette kaldes en højkirkelig Retning, et Ordspil, som dog
ikke maa skrives paa Martensens Regning. Og endeligt vilde han i sin
\textit{kristelige Etik} (1871---1878) paavise Foreningen af det Kristelige og
det Humane, det store Ideal, der, som vi har set, stod for ham lige fra hans
første Ungdom. Men denne Enhed kan, hævder han, kun naas ved ubetinget
Underkastelse under Kristendommen. Dette synes dog ikke just at være en
Forening, i hvert Tilfælde ikke en højere Enhed, men en Underordning af det ene
Modsætningsled under det andet. Om en Harmoni af denne Art overhovedet er
forenelig med selve Kristendommen, var et Problem, Martensen aldrig for Alvor
havde stillet sig.%
% footnote *) on p. 145
\footnote[1]{I Ugeskriftet „Nær og Fjern“ skrev jeg i Aaret 1878 en længere
Kritik af Martensens Etik. Martensens „Levnet“ (III p. 208) siger, vistnok med
Henblik til denne Anmeldelse, at Bogen blev angreben „fra den rent negative
Side“. Jeg véd altsaa af egen Erfaring, hvorledes man er tilmode, naar man er
„rent negativ“. --- Ogsaa af Martensens Bog om Jacob Bøhme skrev jeg i sin Tid
en Anmeldelse („Ude og Hjemme“. 1882).} Des skarpere blev det stillet af den
Tænker, til hvem
% --- p. 146 ---
\setcounter{footnote}{0}%
\opage{146}vi nu gaa over, og hvis Opgave da ogsaa førte med sig, at han kom
til at træde i den skarpeste, mest lidenskabelige Modsætning til Martensen
baade som Tænker og som Kirkemand.

\begin{center}\rule{2cm}{0.4pt}\end{center}

% --- p. 147 ---
\setcounter{footnote}{0}%
\opage{147}

\begin{center}\large 14. SØREN KIERKEGAARD\end{center}

Kierkegaard indtager indenfor dansk Tænkning en Plads, der kan sammenlignes med
Pascals Stilling indenfor fransk Tænkning. Fælles for dem er den dybe Følelse
af Livets Modsætninger, den skarpe Fremhæven af de store Problemer og
Overbevisningen om, at saa stor en Horizont, Tanken end aabner for os, falde
vore højeste Opgaver dog indenfor en snever Kreds, tilsidst indenfor vor egen
Personlighed, i det inderligste Opgør af vort Forhold til Sandheden, den
Sandhed, vi have Mod og Ærlighed nok til at erkende.

Hans Livsopfattelse danner en afgørende Modsætning baade til Grundtvigs glade
Tryghed, der svælger i Historie, Billeder og Sang, og til Martensens
Spekulation, der lader Troen, som forudsættes given, udfolde sig i dialektisk
Bevægelse.

Filosofien har ikke Ret til at tilegne sig Alt i denne Personlighed og dens
Værk, der for en stor Del hører hjemme paa Religionens og Literaturens
Omraader. Men filosofisk Tænkning danner dog den Form, i hvilken han tilsidst
søgte at gøre sig Rede for sin Virksomhed som Helhed. Hans filosofiske
Hovedskrift, \textit{Afsluttende Efterskrift til Filosofiske Smuler} (1846),
har han selv erklæret for Midtpunktet i hele sin Produktion; paa begge Sider af
dette centrale Værk ligger den religiøse og den æstetiske Produktivitet.%
% footnote *) on p. 147
\footnote[1]{\textit{Af Søren Kierkegaards efterladte Papirer. 1849.} p. 69.
--- Dog er det netop i Dagbøgerne fra dette Aar, at Kierkegaard syslede med det
Spørgsmaal, om han ikke skulde have Ret til at paastaa, at han fra først af
havde havt et bevidst religiøst Formaal med hele sin Produktion, den æstetiske
og filosofiske saavel som den religiøse. Han kom til det Resultat, at han ikke
turde paastaa det. Produktionen i de forskellige Retninger var foregaaet
uvilkaarligt, og han turde i hvert Tilfælde ikke tilskrive sig selv Æren for,
at den i sin Helhed kunde tjene et religiøst Formaal. Se herom mine
\textit{Mindre Arbejder.} Anden Række. p. 183---186.}

% --- p. 148 ---
\setcounter{footnote}{0}%
\opage{148}Hvad han har ydet som Filosof, kan henføres under tre Synspunkter.
Han har givet et Udkast til en Erkendelsesteori, i hvilken han genoptager Kants
Problem, som Hegel havde ladet ligge. Han har, i sin Lære om „Stadierne“, givet
en sammenlignende Livsfilosofi paa Grundlag af en dyb Forstaaelse af
Personlighedens Væsen og Trang. Og han har stillet Problemet om Forholdet
mellem Kristendom og Humanitet i saa bestemt og principiel en Form, som ingen
anden Tænker i det nittende Aarhundrede. Paa alle tre Punkter faar hans Tanke
sin Lidenskab fra den gærende og dybt grebne Personlighed, den udspringer af.
Med voldsom Fart stiger den til de højeste Tinder, hvor det gælder Alt eller
Intet, og Modsætningerne spændes saa stærkt, at ikke blot Tanken, men
Personligheden selv tilsidst sprænges.

Den blivende Belæring, han kan yde, er overvejende af indirekte Art. Saaledes
forholdt det sig jo ogsaa med Sokrates, et af hans store Forbilleder. Men han
har gjort et Tankearbejde, som kan være af stor Betydning selv for den, der i
væsentlige Henseender vælger sig helt andre Udgangspunkter. Og direkte,
positivt har han hævdet Livsanskuelsens og Tænkningens indre Sammenhæng med
Personligheden og derved peget paa en stadig Kilde til Rensning og Fornyelse i
Aandens Verden. Det af Poul Møller opstillede Begreb personlig Sandhed har han
ført ud i dets Konsekvenser.

Hans Bane blev ikke lang (han er født 1813 og døde 1855), men hans
lidenskabelige Produktivitet gjorde, at han fik udtalt, hvad han efter sin
Personlighed og sit Standpunkt havde at sige, og Døden kom, da han netop havde
sagt det sidste Ord, der kunde siges fra hans Standpunkt.%
% footnote *) on p. 148
\footnote[1]{En udførligere Karakteristik af Kierkegaards Personlighed og Lære,
end jeg her kan give, har jeg givet i min Bog \textit{Søren Kierkegaard som
Filosof} (1892).}

Det er ikke blot hans Lære, det er selve hans Personlighed, der har filosofisk
Interesse ved de psykologiske Gaader, den indeholder. Efter min Opfattelse, som
jeg støtter paa
% --- p. 149 ---
\setcounter{footnote}{0}%
\opage{149}Studiet af hans efterladte Papirer, er Tungsindet, næret ved en
glædeløs Barndom og en pietistisk Opdragelse, den dybeste Strømning i hans
Væsen. Han følte sig stedse som en stækket Fugl, der istedetfor at kunne bruge
sine Vinger frit maatte bevæge sig i sin snevre Kreds, uden at kunne løse den
Gaade, om det var Sygdom eller Brøde, der stedse holdt den tilbage. Tungsindet
fjernede ham fra Livet og gjorde stadige og nære Forhold til andre Mennesker
til en Umulighed for ham. Det gav ham et andet Blik paa Livet, end de
sædvanlige, letlevende Mennesker havde; han saa Afgrunde, som de gik sorgløst
forbi, og han opdagede Elementer i Kristendommen, som de lod upaaagtede. Men
ved Siden af Tungsindet stod saa den geniale Produktionstrang, den store
Digter- og Tænkerevne. Den nærede Tungsindet ved at udspinde Grublerier og
endeløse Reflexioner, i stadig Kredsen om sig selv; men den tjente tillige til
Afløb og til Lindring, idet han kunde glemme det indre Mørke ved at fordybe sig
i Produktionen. Han holdt sig, som han udtrykker sig, i Live ligesom
Scheherezade („Tusind og en Nat“) ved stadigt at digte. --- Endnu en ny
Modsætning trængte sig dog her frem. Reflexionen og Fantasien og deres Trang
til at brede sig kom i stærk Kontrast til den af Tungsindet og den dybe Uro
affødte Trang til et absolut Tilhold, til et Nødanker paa Livets oprørte Sø,
--- til Underkastelse under en absolut Autoritet, en overnaturlig Myndighed.
Trods al Reaktion mod den glædeløse Barndom og den strænge Opdragelse formaaede
han ikke at løsrive sig fra den Livsopfattelse, de havde lagt Grunden til;
Vingen var nu engang stækket, og det især af indre, ikke af ydre Aarsager. Men
mod den faste Autoritet tørnede Tankens Bevægelser atter og atter for stedse
paany at kastes tilbage. Hans Tankeliv er karakteriseret ved, hvad han etsteds
udtaler: „Enhver Lidenskabs højeste Potens er altid at ville sin egen
Undergang, og saaledes er det ogsaa Forstandens højeste Lidenskab at ville
Anstødet, uagtet Anstødet paa en eller anden Maade maa blive dens Undergang.
Dette er da Tænkningens højeste Paradox, at ville opdage Noget, den ikke selv
kan tænke.“%
% footnote *) on p. 149
\footnote[1]{\textit{Filosofiske Smuler} (1884). p. 52. (Samlede Værker. IV p.
204). --- Udtrykket „Anstød“ antager jeg for en Reminiscens for Fichte, hos
hvem „Anstoss“ betyder den Skranke, Tænkningen stedse støder paa, hvor den
staar overfor noget blot Givet, noget Nyt, der ikke er en Konsekvens af selve
det foregaaende Tankearbejde.} --- Denne Trang til at gaa til Grænsen maa
% --- p. 150 ---
\setcounter{footnote}{0}%
\opage{150}enhver kritisk Filosof have, der vil følge det sokratiske Bud om at
skelne mellem hvad han véd, og hvad han ikke véd. Men den fik hos Kierkegaard
en ganske særlig Vælde, fordi hans Indre --- dels paa Grund af det oprindelige
Tungsind, dels paa Grund af den ved Opdragelsen rodfæstede religiøse Autoritet
--- frembød en Modsætning mellem Noget, der skulde staa absolut fast, og en
aldrig hvilende Stræben efter dog at drage Alt ind under Tankens Belysning.
Derved stod han fra først af anderledes overfor det religiøse Problem end Mænd
som Sibbern, Heiberg og Martensen. Sibbern tog det Kristelige som noget Givet;
han forbeholdt sig at prøve, hvorvidt det ogsaa virkeligt var \textit{givet,}
og i Løbet af sit lange Liv fandt han, at alt det Dogmatiske i Kristendommen
slet ikke var \textit{givet.} Heiberg tog Religionen som en Genstand, der
gennem dialektisk Udvikling blev omsat til ren Tænkning, netop hvad det
Dogmatiske angaar. Og for Martensen var Religionen baade Genstand og Autoritet,
uden at det var klart, naar den hørte op at være det Ene og begyndte at være
det Andet. For Kierkegaard var den kun Autoritet, en hellig, ærværdig
Autoritet, mild og trøsterig, men med et Medusaansigt overfor den, der vilde
behandle den som Genstand for Erkendelse, saavel som overfor den, der vilde
slaa af paa dens ideale Fordringer. Og dette sidste Synspunkt traadte for ham
mere og mere i Forgrunden, jo mindre Samtiden, og især den samtidige Kirke,
vilde anerkende dets Berettigelse. Saaledes blev han Skridt for Skridt ført til
at rejse den voldsomste aandelige Strid, der nogensinde har været ført i
Danmark. Da han havde sagt sit sidste Ord i denne Strid, lukkede Døden hans Øjne. ---

Saaledes var den psykologiske Jordbund, paa hvilken Kierkegaards Filosoferen
voxede frem. Naturligvis har han orienteret sig i Samtidens Literatur og
Filosofi. Hans Forhold til Poul Møller er allerede berørt. Efter sine
teologiske Studier
% --- p. 151 ---
\setcounter{footnote}{0}%
\opage{151}beskæftigede han sig indgaaende med græsk Filosofi og forfattede en
aandfuld Afhandling for Doktorgraden \textit{Om Begrebet Ironi, med særligt
Hensyn til Sokrates} (1841). Han havde havt en hegelsk Periode, men var,
sandsynligvis under Paavirkning af Poul Møller, kommen ud over den. Sokrates
blev hans Forbillede paa Grund af den nøje Forening, han frembød mellem
Tænkning og Personlighed. Naar han dog her fremstillede Sokrates som en
abstrakt og negativ Ironiker, var det en Opfattelse, han senere udtrykkeligt
berigtigede, ligesom ogsaa Begrebet Ironi undergik en Forandring hos ham under
Gennemarbejdelsen af Læren om Stadierne (Livsanskuelserne), som vi paa sit Sted
skulle paavise. Af sin Samtids Filosofer satte han, efter at have forladt
Hegel, den skarpsindige Adolph Trendelenburg meget højt, og udtaler flere
Steder med Varme, hvor Meget han skylder ham. --- Af ældre følte han sig
særligt tiltrukken af Lessing og Hamann.

\begin{center}\rule{2cm}{0.4pt}\end{center}

Kierkegaards Erkendelsesteori sammenfatter han selv i de to Sætninger: Der kan
gives et Tankesystem --- men et Tilværelsens System kan der ikke gives.

Der er et Indgreb af Tanker, som vi nødvendigvis maa benytte, naar vi ville
forstaa Tilværelsen. Men disse Tanker (Kategorierne) udspringe maaske selv af
Tilværelsen, ere en Art forkortede Udtryk for vore Erfaringer! Eller skulde de
kunne blive til i os uden at have Nogetsomhelst at gøre med den virkelige
Tilværelse? Kierkegaard giver intet Svar herpaa, men bebuder en nærmere
Undersøgelse, som han dog ikke kom til at foretage, da den religiøse
Forfattervirksomhed snart lagde helt Beslag paa ham. Men i hvert Tilfælde
forlanger han, at hine Grundtanker skulle udvikles i streng logisk Form, uden
at man optager Forudsætninger, for hvis Herkomst man ikke gør Rede. --- Og man
maa tillige gøre det klart, hvorledes Besyndelsen paa den Tankerække, man
udvikler, gøres. En absolut Begyndelse kan det ikke være; ti
% --- p. 152 ---
\setcounter{footnote}{0}%
\opage{152}af Intet kommer Intet. Og forudsætter man en Eftertanke, der
analyserer et givet Indhold, saa maa der spørges, hvorfor denne Eftertanke
stanser paa et vist Punkt for at begynde at bygge Tankesystemet op derfra.
Hvorfor fortsætter Eftertanken ikke sit analyserende Arbejde istedetfor at
stanse og begynde Konstruktionen? Kun ved en Beslutning, altsaa med et Ryk
eller ved et Spring kan den første Sten til Systemet lægges; ellers bliver man
ved at grave. Der gives altsaa ikke en nødvendig Overgang fra Erfaring og
Analyse til Tankesystemet. Hegel kunde [det synes at være Kierkegaards Tanke]
meget vel have tilbragt sit Liv med at fortsætte „Phänomenologie des Geistes“,
uden nogensinde at være kommen til at skrive sin „Logik als Wissenschaft“. ---
Endeligt maa det vel erindres, at selvom det lykkes at konstruere et
Tankesystem, i hvilket hvert af vore Grundbegreber finder sin Plads, bliver
endnu det Spørgsmaal at besvare, hvorledes dette System forhodler sig til
Filosofens egen Personlighed, som jo dog ikke helt kan gaa op i det.

Et Tankesystem lader altsaa, selvom det kan naas, forskellige Spørgsmaal
tilbage. Og et Tilværelsens System kan der ikke gives, da Tilværelsen for os,
der existere i Tiden, aldrig kan komme til at foreligge afsluttet. Fortiden er
ganske vist afsluttet, men vi, som existere og tænke over Fortiden, ere ikke
afsluttede; Tiden løber videre, alt imedens vi tænke, og stiller sine Krav til os.

I „Filosofiske Smuler“ drøftes udførligt det Spørgsmaal, om det Forbigangne
virkeligt er mere nødvendigt end det Tilkommende, og det hævdes, at jo mere man
gør sig \textit{samtidig} med det Forbigangne, des mere vil man blive klar
over, at det under Oplevelsen stedse kun har staaet som muligt, altsaa som
usikkert, --- som Genstand for Tro. End ikke overfor Fortiden, der dog
foreligger afsluttet, kan der da gives en absolut Viden, naar man gør sig den
ret nærværende.%
% footnote *) on p. 152
\footnote[1]{Man ser, hvorledes Kierkegaard her kommer tilbage til Humes
Problem, hvorvidt Overgangen fra et Øjeblik til et andet var nødvendigt eller
ikke. Ogsaa Hume afviste Appellen til Fortiden (dog særligt fra det Synspunkt,
at han nægtede Retten til at slutte fra Fortid til Fremtid, medens Kierkegaard
søger at blive samtidig med Fortidens Begivenheder, at gøre Fortid til Nutid).
--- Kierkegaard er enig med Hume i, at Nødvendighed ikke kan iagttages. Men han
begaar den Inkonsekvens at mene (i \textit{Begrebet Angst),} at „Springet“,
Sammenhængens absolute Afbrydelse skulde kunne iagttages. Se herom
\textit{Søren Kierkegaard som Filosof.} p. 78---81.} Lessing faar Ret: ikke
Besiddelsen af Sandheden, men den sta%
% --- p. 153 ---
\setcounter{footnote}{0}%
\opage{153}digt fortsatte Stræben efter Sandhed er det Højeste, et Menneske kan
naa til. I Existensen dække Tænken og Væren ikke hinanden. Vi leve forlængs,
men forstaa baglængs. --- Kun for et evigt Væsen udgør Tilværelsen et afsluttet
System, ikke for Væsener, der stikke midt i Tiden. ---

Man har i Filosofien for tidligt forladt Kants ærlige Vej. Hegel har ikke, som
han og hans Tilhængere mene, gendrevet Kant. Forskellen mellem Tænkning og
Virkelighed staar stedse i Kraft. --- Kierkegaard er En af de Første, der har
vist hen til en Genoptagelse af den kritiske Filosofis Undersøgelser.
Kriticismen var i den romantiske Periode ikke forsvunden; den holdt sig som en
Understrøm gennem saa dygtige Forskere som Herbart, Fries og Beneke, for ikke
at tale om Schleiermacher og Schopenhauer. Kierkegaard har neppe noget
egentligt Forhold til disse Tænkere. Han omtaler med Anerkendelse
Schleiermachers Troslære (i „Begrebet Angst“), og han læste i sine sidste Aar
(se Dagbøgerne fra 50'erne) Schopenhauer's Skrifter med Interesse; men i
erkendelsesteoretisk Retning er han neppe paavirket af dem. ---

Det er ikke blot mod den spekulative Filosofi, det er ogsaa mod Teologien i
dens forskellige Former, Kierkegaard vender sig.

Den intellektuelle Interesse, der gaar op i at vinde videnskabelig Forstaaelse
og paavise Sammenhæng mellem vore Oplevelser er for Kierkegaard ikke den
højeste. Idet han gaar ud fra, at Sjælelivet ytrer sig i en Sammenfatten, en
Syntese, finder han, at denne Syntese er des kraftigere, jo større Spændingen
og Modsætningen mellem Elementerne er. Og større er Spændingen hos den, der
lever i Længsel og Stræ%
% --- p. 154 ---
\setcounter{footnote}{0}%
\opage{154}ben mod et evigt Maal, og som altsaa skal sammenfatte Tiden med
Evigheden i urolig Forventning, end hos den Kontemplative, der arbejder med at
bringe Sammenhæng mellem forskellige foreliggende Data. Derfor staar Tro højere
end Viden, og derfor undrede Kierkegaard sig over, at Martensen kunde sidde
roligt og „arrangere en Dogmatik“, medens han antog Troen som en given Sag. I
Videnskaben tænke vi kun med en Del af vor Personlighed, men i enhver alvorlig
Livsanskuelse er det hele Personligheden, der tænker, og alle Sjælens Kræfter
virke sammen. Kierkegaards Værdimaaler, som vi ville komme tilbage til i hans
Lære om „Stadierne“, er den Fylde og den Grad af Modsætning, der kan forbindes
i Personlighedens Enhed. Han er her stegen ned til en Tanke, der ikke blot
siden Leibniz og Kant ligger til Grund for Erkendelsesteorien, men som ogsaa
maa gøre sig gældende i enhver energisk Maade at tage Livet paa. --- For mit
eget Vedkommende var det denne Tanke, der greb mig i min Ungdom og holdt mig
fast, saa jeg atter og atter vender tilbage til den. Jeg ser især paa dette
Punkt Kierkegaards blivende Betydning, saa meget min Livsanskuelse end har
fjernet sig fra ham, og saa lidet jeg særligt er enig i den Afstand, i hvilket
han sætter det videnskabelige Arbejde fra det øvrige personlige Liv. ---

Rent teoretisk kan efter Kierkegaard det religiøse Problem, særligt
Spørgsmaalet om Kristendommens Sandhed, ikke løses. Ti her er intet Givet, der
kan omsættes i Begreb; her er Tale om et fjernt Maal, om hvis Opnaaelse ingen
objektiv (videnskabelig) Vished kan haves, men til hvilket man kun kan forholde
sig i Tro, --- i den atter og atter gentagne Bekræftelse og Forvisning, i
hvilken Spændingens Uro overvindes. Troens Genstand existerer kun for den
Troende, ikke for den, der søger almindelige Begreber og Love.

Ligesaalidt nytter det at henvise til Bibelen som historisk Aktstykke. Ti naar
bliver den historisk-kritiske Undersøgelse færdig? og kan den --- selv i
heldigste Tilfælde --- nogensinde blive saa sikker, at man derpaa kan bygge den
mest personlige Afgørelse? --- Ej heller nytter det at henvise til Kirken
% --- p. 155 ---
\setcounter{footnote}{0}%
\opage{155}eller den i Kirken overleverede Trosbekendelse. Ti hvem garanterer
dens Ægthed og Oprindelighed? Her er en ny Tumleplads for historisk Granskning.
Men Afgørelsen kan ikke ligge der. Ingensomhelst objektiv Garanti er mulig.

Det Afgørende er stedse tilsidst den personlige Trang til et Tilhold, --- til
et Nødanker under Tidsexistensens stærke Bølgegang. Jo mere Spørgsmaalet staar
uafgjort for den objektive, videnskabelige (filosofiske og historiske)
Betragtning, des mere bliver det klart, at --- om den er mulig --- maa den
ligge i Personlighedens Inderlighed. Subjektiviteten er Sandheden. Det er
Individets Forhold til det, som tros, der er Hovedsagen: det kommer an paa, om
Individet selv er i Sandhed, selvom det, det tror, er usandt. Den, der beder i
Usandhed (upersonligt, uden Inderlighed) til den sande Gud, beder til en Afgud;
den, der beder i Sandhed (med hele sin Sjæls Inderlighed) til en Afgud, beder
til den sande Gud. Det kommer an paa \textit{et Hvorledes,} ikke paa et
\textit{Hvad.} Den højeste Sandhed, et Væsen, der existerer i Tiden, kan naa,
er karakteriseret ved objektiv Uvished, fastholdt i den mest lidenskabelige
Inderligheds Tilegnelse. ---

Kierkegaard har her ført Poul Møllers Ide om personlig Sandhed ud i dens
yderste Konsekvens. Han er tillige, mærkværdigt nok, kommen til at udtale en
lignende Tanke, som Heiberg ud fra helt andre Forudsætninger udtalte i de
ovenfor anførte Linjer.

Uden at gaa ind paa en speciellere Kritik kan man dog ikke undlade her at gøre
to Spørgsmaal, det ene fra Kierkegaards eget Standpunkt, det andet fra et
almenmenneskeligt Standpunkt. --- Hvorledes kan der være Tale om en
Aabenbaring, en virkelig Afsløring af Tilværelsens Hemmelighed, naar Sandheden
ikke er et Hvad, men et Hvorledes, og naar Hovedsagen er Personlighedens
inderlige Stræben, ikke dens bestemte Indhold?%
% footnote *) on p. 155
\footnote[1]{Smlgn. om dette Punkt den interessante Afhandling af den grundige
Kierkegaardkender Chr. \emph{Schrempf}: \textit{Kierkegaards Stellung zu Bibel
und Kirche.} (Zeitschrift für Theologie und Kirche. 1891).} Man forstaar paa
dette Punkt Kierkegaards dobbeltsidige Virkning: Nogle har han ført dybere
% --- p. 156 ---
\setcounter{footnote}{0}%
\opage{156}ind i Forholdet til Kristendommen; Andre ere netop under hans
Indflydelse blevne førte bort fra Kristendommen, fordi de kun derved kunde
bevare deres personlige Sandhed og Inderlighed. --- Det andet Spørgsmaal er
det, om det virkeligt skulde være saa, at Sjælens hele Energi maa forbruges til
at sikre sig et Nødværge? Til dette Spørgsmaal ville vi vende tilbage ved
Drøftelsen af „Kierkegaards „Stadier“.

\begin{center}\rule{2cm}{0.4pt}\end{center}

Personligheden (Subjektiviteten) fremtræder i forskellige Former, paa
forskelligt Grundlag. Allerede i det Foregaaende har der vist sig to Typer, den
ene karakteriseret ved rent intellektuel Interesse, den anden ved religiøs
Interesse. Idet Kierkegaard vil indskærpe Betydningen af den personlige
Sandhed, skildrer han en hel Række af Typer, eller som han med et mindre
heldigt Ord kalder dem, Stadier. Naar der ikke findes et særligt intellektuelt
eller intellektualistisk Stadium mellem dem, er det sikkert, fordi et saadant
Stadium for ham blot betød en Forflygtigelse af Personligheden. Han vil se
Individet i dets Forhold til det virkelige Liv og skildrer de forskellige
Maader at tage dette paa.

De forskellige Stadier eller Typer staa hos Kierkegaard skarpt afskaarne i
Forhold til hverandre. Man faar ikke at vide, hvorledes de blive til, eller
hvorledes en Overgang fra det ene til det andet er mulig. I den Hegelske
Filosofi skete Overgangen fra et Trin til et højere blot derved, at man hæver
det lavere Trin til fuld Bevidsthed; selve Reflexionen, den klare Bevidsthed
om, hvor man stod, skulde være nok til at virkeliggøre et højere Stade.
Kierkegaard hævder, at Overgangen forudsætter et Spring, fordi Standpunkterne
ere principielt forskellige, --- hvilket jo er det Samme som at sige, at man
ikke kan forstaa Overgangen.

Kierkegaards Fremstilling af „Stadierne“ er noget forskellig i de poetiske
Skrifter („Enten---Eller“ og „Stadier paa Livets Vej“) og i hans filosofiske
Hovedskrift („Uvidenskabelig Efterskrift“. Tillægget til andet Afsnits andet
Kapitel).
% --- p. 157 ---
\setcounter{footnote}{0}%
\opage{157}Hist har han kun tre Stadier, det æstetiske, det etiske og det
religiøse. Her faar han sex, idet Ironien skydes ind imellem det æstetiske og
det etiske Stadium, og Humoren mellem det etiske og det religiøse Stadium, og
idet det religiøse Stadium deles i to, hedensk-human og kristelig Religiøsitet.
--- Vi ville nu give en Karakteristik af disse sex Stadier.

\textit{Det æstetiske Stadium} skildres noget forskelligt i „Enten---Eller“
(1843) og i „Stadier paa Livets Vej“ (1845). Hist er det Karakteristiske en
virkelighedsfjern Dvælen i Tungsind og Fantaseren med Udmalelse af Stemninger,
raffinerede Reflexioner og mefistofeliske Tankeexperimenter. Det er den tænkte
Forfatter af første Del af „Enten---Eller“, der repræsenterer det æstetiske
Stadium, ikke blot de Skikkelser (f. Ex. Johannes Forføreren), han „i sin
famlende Existens“ fremmaner. Dog er et enkelt af Afsnittene i første Del, det,
der kaldes „Vexeldrift“, typisk for dette Stadium, ogsaa saaledes som det
skildres i første Afdeling af „Stadier paa Livets Vej“, hvor „Æstetikerne“
fremstilles existerende i Kød og Blod, ikke som blotte Muligheder.

Det Karakteristiske er idethele et skiftende Stemningsliv, der skyr faste
Former og Forhold, en kaad Legen med alle Livsforhold. Man tager overalt det
Nye, Friske, Pikante, Øjeblikkelige, uden at lade sig fastholde. Man nyder det
Æggende uden at hildes, det Sanselige uden at forfalde, og har Elasticitet til
at bryde af i rette Øjeblik. Den store Fjende er Gentagelsen og Pligten. Man
vil flygtig Berøring med Mennesker, og man prøver Beskæftigelser; men man tager
kun Blomsten af Alt, og undgaar Ægteskab og Venskab, Borgerskab og Videnskab.
Det er især det Erotiske, der benyttes som Exempel til at belyse den flygtige
Tangeren af Livets Cirkler, som er ejendommeligt for dette Standpunkt. Alt
gøres til Billede og nydes, som om det kun var et Billede. Det er paa Grund af
denne Reflexion, at dette Stadium faar Navn af det æstetiske. Det er ikke et
poetisk Stadium; der skildres ikke en Livspoesi, som ikke selv véd af, at den
er Poesi, fordi den udfolder sig som uvilkaarlig Følge af Fordybelsen i Livet.
I sit Haandexemplar af „Enten---Eller“ har Kierkegaard skre%
% --- p. 158 ---
\setcounter{footnote}{0}%
\opage{158}vet: „En egentlig Forelskelse lod sig ikke bruge i første Del...
Hvad jeg kunde bruge var en Mangfoldighed af erotiske Stemninger.“

\textit{Ironikeren} gaar ikke op i de enkelte Øjeblikke; hans Opgave er at
bevare sin Sammenhæng, at hævde sit Selv overfor de forskellige Forhold, han
stedes i, og som ikke ere ham Midler til Nydelse, men derimod Farer, der kunde
hindre ham i at være sig selv. Han gaar da med tilsyneladende Alvor ind paa den
enkelte Situation og synes at være optagen af den og at tillægge den Værdi; men
det er en Tribut, han yder, for at kunne bevare sig overfor Udvortesheden og
Spredtheden.

Ironikeren har intet positivt Indhold eller Formaal. Han forholder sig
væsentligt afværgende i sin Stræben efter at bevare sin Frihed. Men saaledes er
han en Mellemform mellem Æstetikeren, der forholder sig nydende, og Etikeren,
der har en positiv Opgave. I sin Disputats havde Kierkegaard behandlet Ironien
(og Sokrates som Repræsentant for den) som rent formel og negativ; men selv der
havde han i en af de vedføjede Teses udtalt, at ligesom Filosofien begynder med
Tvivlen, saaledes begynder et Liv, der skal kunne kaldes menneskeværdigt, med
Ironien. Ironien kan være Etikens Inkognito, dens ydre Form, da det etisk set
er den første Betingelse at have et Midtpunkt og ikke lade Livet opløses i
enkelte Øjeblikke. Der er i Ironien en Lidenskab, som vel væsentligt ytrer sig
i Abstraktionen fra alle endelige Forhold, men som dog fører til at hævde en
Modsætning mellem det Indre og det Ydre, der er af saa stor Betydning for et
etisk Standpunkt. ---

\textit{Det etiske Stadium} karakteriseres i de digteriske Fremstillinger som
Samling i Modsætning til Æstetikerens Spredthed, og som aaben Hengivelse i
Samliv i Modsætning til æstetisk Drømmen og Experimenteren. Det Etiske
forudsætter en Beslutning, der tager det op med Gentagelsen, den, Æstetikeren
frygtede. Derved bliver der Sammenhæng i Livet, og i Isolationen ophæves. Det
er nærmest Hegels Opfattelse af det
% --- p. 159 ---
\setcounter{footnote}{0}%
\opage{159}Etiske som Opgaaen i et Alment, i et Samfundsliv, Kierkegaard her
lægger til Grund, kun at det er Familien, Ægteskabet, han bruger som Type, ikke
Staten. Det borgerlige Liv kunde Kierkegaard kun se under Synspunktet af
Filisteri. --- I den filosofiske Fremstilling er Opfattelsen bleven en anden.
Begrebet den Enkelte spiller en Hovedrolle; det er fra en foreløbig Anvendelse
i Forordet til en af Kierkegaards religiøse Taler efterhaanden rykket frem i
Forgrunden for hans Tankegang, og det Etiske staar nu for ham som den Enkeltes
Sag, som beregnet paa Individualitet. Man har i sin Samvittighed kun med sig
selv at gøre. Og al Etik er religiøs. Uden Gudsforhold og Udødelighedstro er
det Etiske kun Skik og Brug. Det Sociale og Historiske har ingen Betydning, og
Ægteskabet er kun en Sinkelse. Saaledes gaar den religiøse Karakter af det
Etiske i Et med den individualistiske. Grunden til denne Forskydning i
Opfattelsen maa søges i en Tankegang, Kierkegaard havde udviklet i det lille
Skrift \textit{Frygt og Bæven} (1843), hvor der drøftedes Muligheden af, at et
religiøst Bud kunde stride mod Etikens almene Lov. Fortællingen om Abrahams
Lydighed overfor Buddet om at ofre sin Søn bruges som Exempel. I en saadan
Konflikt mellem det Religiøse og det Etiske staar Individet som den Enkelte,
isoleret fra alle Andre; Ingen kan forstaa ham og hans Færd. Gennem et voldsomt
Brud maa Overgangen fra det Etiske til det Religiøse ske, naar det Etiske
betyder det Almene, det for Alle Gældende. Uvilkaarligt har da Kierkegaard
senere lagt Begrebet den Enkelte til Grund allerede i Etiken; den religiøse
Karakter af Begrebet er fulgt med --- og dermed er da egentligt med det Samme
det Etiske ophævet som et særligt Stadium. Nu blive Poesi og Religiøsitet de to
eneste idelle Muligheder: „Imellem Poesi og Religiøsitet opfører den verdslige
Levevisdom sin Vaudeville. Ethvert Individ, der ikke lever netop poetisk eller
religiøst, er dumt.“%
% footnote *) on p. 159
\footnote[1]{\textit{Uvidenskabelig Efterskrift.} p. 349. (Værker VII p. 397).}
Jeg antager det ikke for tilfældigt, at Kierkegaard her taler om det Poetiske
og ikke om det Æstetiske; han tænker paa noget
% --- p. 160 ---
\setcounter{footnote}{0}%
\opage{160}mere Umiddelbart og Uvilkaarligt end den Type, han har
karakteriseret som „æstetisk“.%
% footnote *) on p. 160
\footnote[1]{I \textit{Stadier paa Livets Vej} (Værker VI p. 158) hedder det (i
Afsnittet om det Etiske): „Kvinden er i sin Umiddelbarhed væsentligt æstetisk,
men fordi hun er det væsentligt, derfor er ogsaa Overgangen til det Religiøse
ligeved. Den kvindelige Romantik er i næste Øjeblik det Religiøse.“ --- Jeg
antager, at „æstetisk“ her betyder poetisk, og ikke er taget i samme Betydning,
som naar Kierkegaard taler om „det æstetiske Stadium“. I hvert Tilfælde ser
man, at for Kvindens Vedkommende existerer det etiske Stadium slet ikke.} ---

Etikeren gaar op i sine Opgaver, der skulle løses ved hans egen Virksomhed.
Derfor har han sit Centrum i sig selv og søger paa positiv Maade at udfylde sin
Plads. For \textit{Humoristen} er hans egen Existens, hele den Kreds, i hvilken
han bevæger sig, og overhovedet alt Endeligt forsvindende i Forhold til det
Uendelige og Evige. Han bevæger sig med sit Centrum som en Epicykel paa en
større Kreds, hvis Midtpunkt han er uendelig fjernt fra. Derfor bliver
Endeligheden komisk, naar den overhovedet vil betyde Noget. Selv sin Lidelse og
sin Skyld ser Humoristen som forsvindende overfor den uendelige Horizont, der
aabner sig for ham. Humoren indeslutter Vemod ved og Sympati med Smerten i
Livet, og dens Smil og Spøg rummer en dyb Alvor. Ironikeren vil hævde sig selv
gennem den tilsyneladende Alvor overfor det Endelige; men Humoristens Alvor
hindrer ham ikke i en mild Spøg med alt Endeligt, der trods al Travlhed og
Vigtighed dog aldrig kan blive af afgørende Betydning. Humoren er det sidste
Stadium forud for det Religiøse. ---

Kierkegaard har her antydet en Livsanskuelse, der svarer til den Grundstemning,
Tankens Forstaaelse af Tilværelsen vækker. Det Enkelte ses i dets
Ejendommelighed og dets Værdi --- og dog som forsvindende i Forhold til den
uoverskuelige Tilværelse, som Tanken aabner for os. Modsætninger kunne
fremtræde i deres Skarphed og dog fortone sig i Forhold til Forvisningen om en
dyb Harmoni, en stor Enhedslov. Det er den Stemning, den enkelte modus i
Spinozas System maa stedes i overfor den uendelige Natur. Det er for mit
Vedkommende det af de Kierkegaardske Stadier, jeg nær%
% --- p. 161 ---
\setcounter{footnote}{0}%
\opage{161}mest vilde føle mit eget Standpunkt og min egen Stemning ved Livet
beslægtet med. Det er Skade, at Kierkegaard med de Evner, der stode til hans
Raadighed, ikke har udført Skildringen af det videre. Netop da hans Øjemed med
den hele Fremstilling af „Stadierne“ er af vise, „hvilket uhyre Existensomfang
der er muligt uden for Kristendommen“, det vil sige, hvilken høj og ædel
Livsopfattelse der er mulig paa rent menneskelig Grund, burde han have
karakteriseret dette sidste rent humane Stadium mere indtrængende. Han lader
flere af sine Pseudonymer%
% footnote *) on p. 161
\footnote[1]{Kierkegaard har fremstillet „Stadierne“ ved Hjælp af en Række
Pseudonymer, fordi han derved kunde opnaa „den psykologiske Konsekvens, hvilken
ingen faktisk virkelig Person i Virkelighedens sædelige Begrænsning tør tillade
sig eller kan ville tillade sig.“ („En første og sidste Forklaring“, efter
„Uvidenskabelig Efterskrift“).} tale i dets Navn --- men de have alle Retning
mod det femte og det sjette Stadium, hvis Overhøjhed de anerkende, skønt de
ikke selv tilhøre dem. Af samme Grund er det, at de tidligere Stadier ikke
udføres saa klart og fuldstændigt, som Opgaven krævede. Det er foruden det
„æstetiske“ egentligt kun de religiøse Stadier, der faa en fuldstændig
Udførelse --- og, som vi have set, blev Kierkegaard efterhaanden tilbøjelig til
at udskyde alle de mellemliggende Former.

\textit{Det religiøse Stadium} træder i Kraft, naar den egne Existens ikke blot
ses som forsvindende i Forhold til det Evige og Uendelige, men naar Individet
ser sin egentlige Opgave i at gøre sig til Et med dette Evige, have sit
Midtpunkt i Midtpunktet for den større Kreds, i hvis Periferi det lever sit
Liv. Det faar da to Centra, har sit Liv baade her og hisset, men skal egentligt
have det hisset. Der føles en dyb Spaltning, en smertelig Modsigelse. Smerten
kan ikke her som i Humoren, fortones i Vemod, men er stadigt forbunden med
Tidsexistensen paa Grund af Modsætningen mellem Tid og Evighed. Individet er
tilmode som en Fisk paa Landjorden. Det lever i et fremmed Element.

„Ikke er Fuglen i Buret, og ikke er Fisken paa Strandbredden, og ikke den Syge
paa Sygelejet, og ikke Fangen i det
% --- p. 162 ---
\setcounter{footnote}{0}%
\opage{162}snevreste Fængsel saaledes fangen, som den er det, der er fangen i
Gud; ti som Gud er, saa er den fængslende Forestilling [om Gud] allesteds og i
ethvert Øjeblik.“

Et Standpunkt som dette er muligt indenfor Hedenskabet og maa vel skelnes fra
det kristelige, skønt mange saakaldte Kristne ikke engang naa op til det. I
\textit{Kristendommen} skærpes Modsætningen mellem Tid og Evighed saaledes, at
den kun kan hæves ved, at det Evige selv gaar ind i Tiden, lever i
Tidsexistensens Form som endeligt Væsen, --- at Gud bliver Menneske. Her bliver
Sandheden et Paradox: Guden i Tiden! Mennesket kan paa sin Vej gennem Tiden
møde Guden selv: Alt beror paa, om det da kender ham og tror paa ham; --- men
selve Evnen til at kende og tro kan det ikke give sig selv; det maa modtage den
af Guden selv --- og det kommer da an paa, om det vil modtage den!%
% footnote *) on p. 162
\footnote[1]{I \textit{Filosofiske Smuler} (1844) var Paradoxet i det
kristelige Grunddogme om Gudmennesket fremstillet i skarp Modsætning til det
sokratiske (menneskelige) Standpunkt.} --- Hele Livsanskuelsen bestemmes
herved. Ti overfor det ene afgørende Punkt, Gudens Aabenbarelse i Tiden, taber
alt Andet sin Betydning. Kun paa dette ene Punkt berøres Tiden af Evigheden;
Kristus er --- som al Idealitet i Grunden er det --- en Tangent til Jorden, til
Virkeligheden.%
% footnote **) on p. 162
\footnote[2]{Det er karakteristisk for Kierkegaard, at baade det Stadium, han
sætter højest, og det, han sætter lavest, kan bestemmes ved samme Formel: at
være Tangent til Livets Cirkel. Som Kristus, ja al Idealitet (se Efterladte
Papirer“. 1849. p. 30; 242), saaledes kan ogsaa den æstetiske Maade at tage
Livet paa beskrives ved, at den gaar bort efter Tangenten! (Se især
„Vexeldrift“ i første Del af „Enten---Eller“). Og det er maaske især
karakteristisk for ham, at det er disse to Standpunkter, han bedst forstaar og
har mest Respekt for.} Hele den øvrige Tilværelse er gudsforladt, værdiløs. Al
Værdi er koncentreret paa et eneste Punkt. --- Deraf følger igen, at kun en
psykologisk, hverken en spekulativ eller en historisk Indledning til
Kristendommen er mulig. Det kommer an paa, hvor meget der er levet, --- om
Mennesket har lært at fortvivle om alt Andet for saa at finde sit Tilhold i
dette ene Punkt, det eneste, fra hvilket der falder Lys over Tilværelsen.

% --- p. 163 ---
\setcounter{footnote}{0}%
\opage{163}Denne sidste Tanke udviklede Kierkegaard nærmere i \textit{Sygdommen
til Døden (1849)}. Han undersøger i dette sit dybsindigste Skrift de
forskellige Arter af Fortvivlelse. Den umiddelbare Glæde og Lykke synes at være
fjernest fra Fortvivlelse, og dog er Trygheden og Roen illusorisk. Angsten
ligger bagved. Ti der stilles til Mennesket den store Fordring at være Aand
efter den absolute, guddommelige Maalestok. Kristendommen vil gøre Mennesket
til noget saa Overordentligt, at det ikke kan faa det i sit Hoved. Dette er
Forargelsens Grund. Svagheden, Indesluttetheden, Trodsen ere forskellige
Former, i hvilke Mennesket unddrager sig den høje Fordring; alle ere de Former
for Fortvivlelse, og der skjuler sig i dem alle en Villie, der vil unddrage sig
Fordringen, snart ved Selvopgivelse, snart ved Selvhævdelse. Ogsaa
Aandløsheden, der synes at være for pjattet til at kunne kaldes Synd, opstaar
dog kun ved Menneskets egen Skyld. Selve Syndsbevidstheden kan antage Form af
Fortvivlelse. Den rette Modsætning til Synd er ikke, som Hedningerne mente,
Dyd, men Tro. Og til Troen kommer man gennem Fortvivlelsen, der lærer at henfly
til Naaden. Fortvivlelse er derfor en Sygdom, som det er den største Ulykke
ikke at have havt. Men Kristenhedens Syndsbegreb gaar ikke dybere end
Hedningernes. Man har udjevnet de store Modsætninger Synd og Tro, Fortvivlelse
og Naade. Udjevningen er sket først spekulativt, saa pøbelagtigt.

Det var Kierkegaards Opgave i det nævnte Skrift at vise, hvor dybtgaaende indre
Oplevelser, hvor skarpe Modsætninger den maa kende, der med Sandhed skal kunne
anvende det nye Testamentes Psykologi paa sig selv. Ingen Gengivelse af
Skriftets Indhold kan give nogen Forestilling om, med hvilken Lidenskab og med
hvilket Skarpblik det er skrevet. Mesteren fører os langs Afgrunde og lærer os
aandelige Muligheder at kende, som han selv ikke turde give sig ud for at have
prøvet, men som efter hans Overbevisning den maatte kende paa første Haand, der
med Rette skulde kunne kalde sig en Kristen.

Medens „Sygdommen til Døden“ indskærper de psykolo%
% --- p. 164 ---
\setcounter{footnote}{0}%
\opage{164}giske Erfaringer, Kristendommen forudsætter, lægger
\textit{Indøvelse i Kristendom} (1850) Vægten paa Mødet med „Guden i Tiden“.
Man forholder sig kun til Kristus ved at være samtidig med ham. Af Historien
kan man intet lære. Det er stadigt den Fornedrede, ikke den Ophøjede, der taler
til os. Han er jo endnu ikke kommen tilbage, og de senere Slægter kunne ikke
have noget Fortrin for de tidligere; Opgaven for dem alle er at kendes ved
Kristus som den fornedrede, som det ringe Menneske, der dog sagde sig at være
Gud. Stedse er det i Samtidighedens Situation, at Prøven skal gøres. I Forhold
til det Højeste er der kun én Tid, den nærværende. At man har faaet at vide
gennem Dogmatik eller Katekismus, senere Tiders Visdom, hvem hint af sin Tids
Autoriteter ringeagtede Menneske var, er kun en Hindring for at blive samtidig,
altsaa komme i det afgørende Forhold til den Maade, hvorpaa nu engang det
Højeste er blevet aabenbaret paa Jorden. I denne Henseende er ingen Forandring,
intet Fremskridt muligt, og det er Tegn paa, at man har afskaffet
Kristendommen, naar det ser ud, somom det var lettere at blive Kristen nu end i
den med Kristus samtidige Generation. En historisk Kristendom er Galimatias og
ukristelig Forvirring. Forargelsens Tegn og Troens Genstand er bleven en
Æventyrskikkelse. Man priser ham for det, man mest vilde forarges over, hvis
man var samtidig med ham, og man bygger sin Tro paa Udfaldet og paa, hvad man
mener, er Historiens Vidnesbyrd. Man har faaet Facit at vide og umager sig
derfor ikke med selve Regnestykket. --- Den eneste Redning er nu at stanse
Forvanskningen og Forvirringen ved at erkende sin Afstand fra det oprindelige
Kristelige, --- indrømme, hvor man staar, og bøje sig for Idealitetens
Fordringer, selvom man ikke kan gøre dem Fyldest. ---

Skridt for Skridt er Kierkegaard saaledes i sin sammenlignende Livsfilosofi
stegen op fra „Æstetikerens“ kaade Leg med Livet til den skærende Alvor i den
højt spændte Idealitets strenge Fordring. Vort Spørgsmaal maa nu blive, hvilken
den Maalestok er, ifølge hvilken han ordner sine Stadier og kalder dem lavere
eller højere i Forhold til hverandre. Vi
% --- p. 165 ---
\setcounter{footnote}{0}%
\opage{165}vende herved tilbage til det Spørgsmaal, der syntes at skulle komme
til Forhandling mellem to af vore største Digtere, Baggesen og Grundtvig, men
som dengang blev ved Udfordringen.

Kierkegaards Maalestok er i og for sig den eneste, der kan være Tale om. Den er
hentet fra det Grundbegreb om Aand, Sjæl og Personlighed (hvad man nu vil kalde
det), til hvilket man stedse vil blive nødt til at komme tilbage: Begrebet Syntese.

Hvor der er Bevidsthed, forudsættes det stedse, at et mangfoldigt Indhold
sammenfattes til Enhed. Deraf følger, at et aandeligt Livstrin staar des
højere, jo større og jo mere forskelligartet et Indhold der er givet, og jo
fastere og inderligere det sammenarbejdes. Eller, for at bruge Kierkegaards
Udtryk, jo større Existensomfanget er, og jo skarpere Modsætningerne ere, naar
dog Enheden hævdes, des højere et Stadium have vi for os%
% footnote *) on p. 165
\footnote[1]{Smlgn. hermed min Bestemmelse af Begrebet psykisk Energi
\textit{(Psykologi.} Femte Udg. p. 67; 123; 182) og Pierre Janets dertil
svarende Begreb tension psychologique \textit{(Les obsessions et la
psychathénis. Paris} 1903. I p. 265 ff; 488 ff).}. Den psykologiske Spænding
maa være des større, jo større de Modsætninger ere, som skulle sammenholdes.
Saadanne Modsætninger ere ifølge Kierkegaard Øjeblikket og Livshelheden, Tiden
og Evigheden, Virkeligheden og Idealet, Verden og Gud. Men det er Tiden med
sine mange Øjeblikke i Modsætning til Evighedens stadige Nu, der betinger den
største Modstand, Syntesen har at overvinde. Det var jo ogsaa Tidsforholdet,
der betingede den stadige Skranke for vor Erkendelse.

Men foruden Modsætningen og Inderligheden har Kierkegaard endnu et Kendemærke
paa den fuldkomnere Tilstand, et Kendemærke, der vel er afledet, men dog
efterhaanden trænger sig frem til Hovedpladsen i Kierkegaards Vurdering. Det
var ikke blot den Maade, paa hvilken en stor Indholdsfylde formes til Enhed,
der blev det Afgørende. Modsætningerne og Modsigelserne indenfor det sjælelige
Indhold betonedes mere og mere, og Vægten blev lagt paa den Smerte, som
stridende Modsætninger fremkalde. Lidelsen blev da
% --- p. 166 ---
\setcounter{footnote}{0}%
\opage{166}efterhaanden den egenlige Maalestok. Nu kan der ogsaa i den Grad af
Lidelse, særligt aandelig Smerte, et Menneske kan bære, ligge et Vidnesbyrd om
dets Aandskraft. Lidelsen kan være Symptom paa store Opgaver og stor Livsfylde,
paa et Mod, der i store Formaals Tjeneste tager Kampen op selv mod en
forholdsmæssigt stærk Modstand. Men Lidelsen er --- ligesaalidt som
Lystfølelsen --- ikke i og for sig et absolut sikkert Kendemærke. Som
Lystfølelsen kan være Tegn paa, at man ikke har stillet tilstrækkeligt store
Fordringer til sig selv (det var derfor, Kant saa det som et uheldigt Symptom,
naar man var altfor glad ved at gøre sin Pligt), --- saaledes kan Lidelsen
skyldes en krampagtig Stræben efter det Umulige. Kierkegaard faar Ret, dersom
det virkeligt er saa, at der er to Midtpunkter for vort Liv, der trække hver
til sin Side, og dersom den Kreds, hvori vort Liv i Erfaringen bevæger sig, og
hvor alle vore positive Opgaver ligge, ikke er vort egentlige Element, saalidt
som Landjorden er Fiskens Element. Fisken, der ligger gispende paa Land, har to
Midtpunkter, et i Vandet, hvor den ikke er, et paa Landet, hvor den er.
Saaledes skal Mennesket efter Kierkegaard være stedt: det lever i Tiden, men
det Hinsidige er dets Hjem. Herved kommer en indre Strid op, der aldrig kan
løses, saalænge Livet varer. Harmonisering er umulig. Paa dette Punkt staa de
to sidste Stadier i afgørende Modsætning til de fire første.

Dersom Troen paa en fra det nærværende Liv principielt forskellig hinsidig
Tilværelse skal være det absolut afgørende Princip for vor Maade at tage Livet
paa, saa at vort Liv her skal bestemmes ved, hvad vi mene om Livet hisset, faar
Kierkegaard Ret. Og han har den store Fortjeneste at have gjort opmærksom paa
Vigtigheden af at blive klar over, hvor vi staa. Hvor vor Skat er, maa ogsaa
vore Hjerter være, og have vi Skatten baade her og hisset, maa vort Hjerte være
delt. Det var denne Delthedens Smerte, Kierkegaard fandt saa lidt af, naar han
saa sig omkring i Kristenheden. Man havde gennemgaaende baade i Teori og Praxis
lagt sig Tingene til Rette saaledes, at Modsætningen mellem de to Midtpunkter
slappedes. Man havde sat et Baade---Og istedetfor
% --- p. 167 ---
\setcounter{footnote}{0}%
\opage{167}et Enten---Eller. Og selve dette Baade---Og var ikke blevet til en
højere Enhed. Man levede faktisk „her“, og Udsigten til „hisset“ spillede ingen
praktisk afgørende Rolle, stod kun som en sidste Trøst.

Den humane Etik kan lære Meget af Kierkegaard; han staar i Etikens Historie som
Alkymisten i Kemiens Historie. Under sit Forsøg paa at løse en umulig Opgave
har han kastet betydningsfulde Strejflys ud over det indre, kæmpende Sjælelivs
Verden. Enhver etisk Opfattelse, der lægger Vægt paa Personlighedens Stræben
efter at bøje Livets forskellige Kræfter sammen til Enhed, vil kunne lære af
den danske Pascal, selvom den overfor hans asketisk-romantiske Ideal stiller
den Forvisning, at det højeste Menneskelige kan naas ved Uddybelse, Berigelse
og Koncentration af Livet, som det kan føres i Erfaringens Verden.

\begin{center}\rule{2cm}{0.4pt}\end{center}

I sine Skrifter og især i sine Dagbøger har Kierkegaard atter og atter udtalt,
at hans egen Tro og hans eget Liv ikke naar op i Højde med den strenge, ideale
Opfattelse af Kristendommen, han havde gjort gældende. Han følte sig her som en
ulykkelig Elsker og gennemkæmpede en ydmygende Strid overfor Tanken om at
optræde som kristeligt Sandhedsvidne overfor sin Tid. Han blev klar over, at
det Højeste, han kunde naa, var Selvforstaaelse og Ærlighed overfor det, der i
hans Øjne var Idealet. Han mente nu ogsaa, at Kirken og dens Repræsentanter
maatte gøre en saadan Indrømmelse, og hans Skrift „Indøvelse i Kristendom“ var
en Henstilling i denne Retning --- en Henstilling, som ikke blev fulgt%
% footnote *) on p. 167
\footnote[1]{Der var dog Præster, som gjorde den forlangte Indrømmelse. Se
\emph{Johannes} \emph{Fibigers} Levned. p. 285---295.}. Længe overvejede han da
Muligheden af en offentlig, agitatorisk Protest mod den bestaaende Kirkes
Forvanskning af det oprindelige Ideal. Hvad der en Stund holdt ham tilbage, var
dels Selvkritik, dels Pieteten overfor Biskop Mynster, dels, og ikke mindst,
Hensynet til de Svage og Fattige, for hvem
% --- p. 168 ---
\setcounter{footnote}{0}%
\opage{168}Kristendommen i dens milde Form var en Trøst, de ikke kunde undvære.
Men han skød alle Hensyn til Side, da Martensen i en Prædiken efter Mynsters
Død fremstillede den afdøde Biskop som et Sandhedsvidne, et Led i den hellige
Kæde af Sandhedens Vidner, der strækker sig gennem Tiderne fra Apostlenes Dage.
Han offentliggjorde en skarp Protest imod denne Anvendelse af Begrebet
Sandhedsvidne i Bladet „Fædrelandet“ den 18. December 1854. Og dermed begyndte
den store Sandhedsvidnestrid, der fra Kierkegaards Side førtes i Bladartikler i
„Fædrelandet“ fra December 1854 til Maj 1855 og derefter i Smaaskrifter
(„Øjeblikket“ Nr. 1---9) fra Maj til September 1855.%
% footnote *) on p. 168
\footnote[1]{Alle Kirkegaards Indlæg i Striden findes nu i Samlede Værker XIV.}

Denne ydre Sandhedsvidnestrid --- der, som ovenfor antydet, til Forgænger havde
havt en indre Sandhedsvidnestrid i Kierkegaards eget Sind, --- angik først
Spørgsmaalet om, hvad der maatte kræves af dem, der skulde anerkendes som
Kristendommens Vidner. Men den førte derfra naturligt over til Forholdet mellem
Urkristendommen („det nye Testamentes Kristendom“) og det nittende Aarhundredes
protestantiske Kirke. Og derved kom Striden egentligt til at dreje sig om den
nytestamentlige Livsanskuelse og om, hvorvidt vor Tids Kristne hylde denne
Livsanskuelse. Det er dette sidste Synspunkt, der giver Striden filosofisk
Interesse den Dag idag. Kun Hovedpunkterne i Kierkegaards Indlæg kan jeg her
dvæle ved. ---

Kun en eneste Tesis har jeg, siger Kierkegaard, medens Luther havde 95. Den
lyder saaledes: Det nye Testamentes Kristendom er slet ikke til. Den, der lader
være at deltage i den offentlige Gudsdyrkelse, som paastaar at være det nye
Testamentes Kristendom, har en stor Skyld mindre! --- Vil den officielle
Kristendom overhovedet stille sig i Forhold til, hvad det nye Testamente kalder
Kristendom, har den kun Et at gøre: at indrømme sin Afstand fra den.

Kierkegaard kæmper ikke for Kristendommen og tør end ikke kalde sig en Kristen.
Han kæmper for Redelighed, og
% --- p. 169 ---
\setcounter{footnote}{0}%
\opage{169}hvad han vover ved sin Optræden, vover han for Redeligheden.

„Hvad er nemlig ifølge det nye Testamente det at blive Kristen, hvortil den
idelige og idelige Formaning mod ikke at forarges, og hvorfra de frygtelige
Kollisioner (at hade Fader, Moder, Hustru, Barn o. s. v.), hvori det nye
Testamente aander? mon ikke begge Dele, fordi Kristendommen meget godt véd, at
det at blive Kristen er at blive menneskeligt talt ulykkelig for dette Liv, dog
saligt forventende en evig Salighed?“

Hvad der er sket, er dette: Kristendommen er bleven afskaffet ved at florere.
Lidelsernes Religion er ved at blive bestaaende Kirke bleven Livstilfredshedens
Religion. Modsætningsforholdet til Verden er forsvundet, idet Alle ere blevne
Kristne --- i det Intetsigendes Betydning. Hvis det havde sin Rigtighed med
alle disse Kristne, saa kunde i hvert Tilfælde det nye Testamente, der stadigt
gaar ud fra Modsætningsforholdet til Verden, ikke længere være Vejledning for
den Kristne, men vilde i denne Henseende være som en forældet Rejsebog. Da Gud
ikke kan have givet os en saadan meningsløs Bog, saa kan Forklaringen kun søges
i et Falskneri fra Menneskenes Side, i et Oprør --- ikke i Trods, men i
Hykleri. Allerede Apostlene vare noget for ivrige i at udbrede Kristendommen og
toge det ikke saa nøje med de Enkelte. Og siden har man forfalsket
Kristendommen; man vilde ikke opgive dens Trøst, men man vilde endnu mindre
fyldestgøre dens Fordring. Man tillempede den efter sit Behov. Af et Par lysere
Satser i den Sørgemarsch, den oprindelige Kristendom var, komponerede man sig
en fejende Galop. Som det er gaaet Kristendommen, gaar det med alle store
Ideer: de fremtræde kun i Begyndelsen i deres Renhed, men forvanskes i Løbet af
den historiske Proces og ende ofte i fuldstændig Modsætning til deres Begyndelse.

En væsentlig Aarsag til denne Forvanskning af Kristendommen var den, at det
Dogmatiske blev skudt i Forgrunden, og det Etiske skudt tilside. Kristus som
Forsoneren traadte iste%
% --- p. 170 ---
\setcounter{footnote}{0}%
\opage{170}detfor Kristus som Forbilledet; man holdt sig til hans Velgerninger
og glemte Efterfølgelsen.

Og saa fik man Søndagskristne, ligesom man har Søndagsjægere, --- Batailloner
af Kristne, som ikke engang Djævelen vil ulejlige sig med at hente. Og
Menneskene lærte idethele at stille ringere Fordringer til sig selv. De saa, at
Livet blev lettere, jo mere ubetydelige de gjorde sig selv, medens det bliver
vanskeligere og vanskeligere for dem, der se Menneskets Opgave i en Stræben
efter at være i Slægtskab med Guddommen. Vær Pjat, og du skal se, alle
Vanskeligheder forsvinde, og denne Verden bliver en herlig Verden! ---

Kort efter at have udtalt denne sidste Sætning (i det sidste Nummer, der udkom
af „Øjeblikket“) faldt Kierkegaard syg om paa Gaden og blev bragt til Frederiks
Hospital, hvor han døde den 11. November 1855.

\begin{center}\rule{2cm}{0.4pt}\end{center}

Det Angreb Kierkegaard rettede mod den bestaaende Kirke og Kristenhed, førtes
ud fra et lignende Standpunkt som det, hvorfra \emph{John} \emph{Henry}
\emph{Newman} nogle Aar tidligere (1836---1843) havde angrebet den engelske
Kirke, og fra hvilket han førtes over til den katolske Kirke. Hans
Hovedspørgsmaal var: Hvad have vi vovet for Kristus? Og det var hans
Overbevisning, at „hvis den nærværende Nød, om hvilken Paulus taler [I Kor. 7,
hvor han raader fra Ægteskab], ikke angiver den kristne Kirkes stadige
Stilling, saa er det nye Testamente ikke skrevet for os, men maa omarbejdes for
at kunne anvendes“. Denne Ytring minder ganske om Kierkegaards Sammenligning
mellem det nye Testamente og en forældet Rejsehaandbog; kun at ifølge ham
Menneskene allerede havde besørget Omarbejdelsen. Det var Newmans Paastand, at
det nye Testamente meget tydeligt beskriver den urkristelige Livstype, men at
vi have læst vedkommende Steder saa ofte, at vi have mistet Evnen til at tage
dem i deres rette Mening.%
% footnote *) on p. 170
\footnote[1]{Smlgn. \emph{Richard} H. \emph{Hutton}: \textit{Cardinal Newmann.}
London 1891. p. 110---119.} --- Naar Kierkegaard oftere udtalte, at Forvansk%
% --- p. 171 ---
\setcounter{footnote}{0}%
\opage{171}ningen var større i den protestantiske Kirke end i den katolske,
viser der sig en ny Overensstemmelse med Newman, hvem han ganske sikkert ikke
har kendt. Dermed skal ikke være sagt, at Kierkegaard, hvis han havde levet,
vilde være gaaet samme Vej som Newman. Dertil spillede Begrebet den Enkelte for
stor en Rolle hos ham, og desuden havde han udtrykkeligt indskærpet, at
Meningen i Kristendommen allerede brast, da man skelnede mellem to Slags
Kristne, nogle (Munke og Nonner), der opfyldte Fordringen, andre (Lægfolket),
der gik paa Akkord.

Han har selv fremhævet, at hans Kamp mod den bestaaende Kristenhed fandt en
Støtte i den Maade, paa hvilken „den sidste Formation af Fritænkere (Feuerbach
og hvad dertil hører) “ er optraadt. „De have egentligt paataget sig den Opgave
at forsvare Kristendommen mod de nulevende Kristne... Feuerbach angriber de
Kristne ved at vise, at deres Liv ikke svarer til Kristendommens Lære... At han
vistnok er en malitiøs Dæmon, kan nok være, men i taktisk Henseende er han en
brugbar Figur“.%
% footnote *) on p. 171
\footnote[1]{\textit{Efterladte Papirer.} 1849. p. 463.} --- Før Feuerbach
havde i det attende Aarhundrede Reimarus, i det nittende Schopenhauer gjort
opmærksom paa Modsætningen mellem Urkristendommens og den moderne Kristenheds Etik.

Hvorhen Kierkegaard havde vendt sig, hvis han havde levet, er det ørkesløst at
drøfte. Det Mærkelige var, at da han havde sagt det sidste Ord, han fra sit
Standpunkt kunde sige, havde han ogsaa forbrugt sin Livskraft.%
% footnote **) on p. 171
\footnote[2]{Ogsaa hans sidste Penge var da brugte. „Den Dag, han tog ud paa
Hospitalet, sendte han Giødvad 300 Rdlr., forat han kunde afgøre nogle mindre
Regninger for ham, og vel ogsaa betale Opholdet paa Hospitalet. Disse 300 Rdlr.
var den sidste Rest af hans Formue; af dem kunde netop Hosptalets Regning
betales, men næppe engang Begravelsesomkostningerne.“ (Brev fra Brøchner til
Chr. Molbech. 17. Febr. 1856).}

Derimod er det af religionshistorisk Interesse at undersøge, hvad der ligger
til Grund for den af Kierkegaard og Andre, fra saa forskellige Standpunkter
hævdede Modsætning mel%
% --- p. 172 ---
\setcounter{footnote}{0}%
\opage{172}lem den nytestamentlige og den moderne kristelige Livsanskuelse, en
Modsætning, der utvivlsomt er tilstede.

Kierkegaard har ikke fremdraget den Forudsætning, paa hvilken den
nytestamentlige Etik hviler, og uden hvilken dens negative eller indifferente
Forhold til menneskelig Kultur ikke forstaas. Denne Forudsætning er den levende
Forventning om Kristi nære Genkomst, om Dommen og om alle Tings Ende. „Guds
Rige“ betyder i det ny Testamente hverken noget Fjernt eller noget evigt
Nærværende, men Noget, som er i Frembrud og snart vil komme helt for Dagen.
Derfor advarer Apostelen mod at stifte Familie; derfor fik personlig Frihed,
Videnskab, Kunst og Statsliv ingen Betydning. Blikket var ene rettet mod det
Store og Herlige, der nærmede sig, og som den daværende Slægt ifølge sin
Mesters Ord skulde opleve. Trods „den nærværende Tids Lidelser“ var man vis
paa, at Sejren over „Verden“ var nær. Urkristendommen var ingen Sørgemarsch,
men en Sejrsmarsch. Den begejstrede Forventning var dens ledende Stemning.%
% footnote *) on p. 172
\footnote[1]{Smlgn. om dette hele Spørgsmaal min \textit{Religionsfilosofi} IV C.}

Mærkeligt nok bortforklarer Kierkegaard denne Kendsgerning. Hin Forventning var
efter ham en Illusion, der med psykologisk Nødvendighed --- hos Kristus og
Apostlene som hos de Kristne --- opstod paa Grund af den voldsomme Modsætning
og Spænding, i hvilken Livet skulde føres. „Det er i den Grad kvalfuldt at være
den sande Kristen, at det ikke var til at holde ud, hvis man ikke bestandigt
ventede Kristi Genkomst som forestaaende nu. Kvalen, Lidelsen føder en
nødvendig Illusion.“ Og deraf kan da fremdeles den Konsekvens drages, at naar
man først venter Kristi Genkomst om mange Aarhundreder, i en fjern Tid, saa er
man ikke længere sand Kristen --- føler ikke Spliden.%
% footnote **) on p. 172
\footnote[2]{\textit{Efterladte Papirer.} 1849. p. 247 f. --- Karakteristisk
for Kierke gaard er det, at han kalder Forventningen om Genkomsten „en
\textit{subjektivt sand} Replik.“ „Saaledes maa ikke blot Kristus, og dernæst,
hvad jo ogsaa er faktisk Apostlene, men saaledes maa enhver sand Kristen sige.“
Religionspsykologisk er denne hele Vending hos Kierkegaard af stor Interesse.}
Det skal altsaa
% --- p. 173 ---
\setcounter{footnote}{0}%
\opage{173}ikke være Forjættelsen og den af Forjættelsen vakte Forventning, der
har fremkaldt Spændingen, men omvendt: Spændingen har fremkaldt Forventningen
og --- Forjættelsen! --- Her optræder Kierkegaard for en Gangs Skyld som
Rationalist. Mod hans Forklaring strider det, at Forventningen om Messias'
Komme hørte hele den Tids Jødedom til, som det ses af en Række jødiske
Apokalypser. Det nye Testamentes Standpunkt er i denne Henseende kun forskellig
fra disse samtidige jødiske Apokalypser ved den Tro, at Messias var kommen én
Gang, saa at Forventningen altsaa gik ud paa en Genkomst. Men Kierkegaard var
ingen Ynder af apokalyptisk Literatur. Enhver Udmaling af himmelske Forhold, af
Saligheden, af den nye Himmel og nye Jord, der skulde komme, betragtede han som
„Æstetik“, der ikke hørte hjemme i Religionen. Det kom ham an paa Vejen, ikke
paa Maalet. Men naar nu Vejen netop bestemmes ved Maalet? Urkristendommen har
ikke frygtet „Æstetiken“, men har med Lidenskab og Inderlighed givet sig den
religiøse Fantasi i Vold. Her som paa mange andre Punkter viser sig
Kierkegaards Forskellighed fra Mænd som Grundtvig og Martensen, der netop følte
Trang til at udmale sig „Maalet“. Men han havde jo ogsaa sin egen Mening om,
hvorledes det gik dem med „Vejen“. ---

For den religionshistoriske Betragtning er det i og for sig naturligt, at
Kristendommen gik over til at blive en Kulturmagt, idet den oprindelige
Forventnings Energi omsattes til et Arbejde for det menneskelige Livs
Udvikling, medens den ideale Afslutning blot blev et fjernt Endemaal. Den
Strid, der da bliver at føre, er, hvorledes Kirken har fyldestgjort sin Opgave
som opdragende og trøstende Autoritet. Paa dette Punkt stillede, som vi have
set, Sibbern sig meget kritisk paa sine gamle Dage. Kierkegaard tog slet ikke
Sagen op fra denne Side.

\begin{center}\rule{2cm}{0.4pt}\end{center}

Kierkegaard staar som en tragisk Skikkelse i dansk Tænknings Historie. Med sit
store Princip om Sandheden som Personlighedens Sag, som liggende i
Subjektiviteten, tørnede han
% --- p. 174 ---
\setcounter{footnote}{0}%
\opage{174}mod, hvad der vedblev at staa for ham med den ubetingede Autoritets
Magt, den dogmatiske Formulering af Kristendommen. Han saa ikke, at de to
Principer vare uforenelige. Dogmet er jo et Facit, som er givet, og givet med
allerhøjeste Sanktion. Og man kan saa i Grunden ikke fortænke Menneskene i, at
de ikke gøre sig Ulejlighed med at foretage Regnestykket. Kierkegaards
Hovedsætning faar kun ret sin Betydning, naar man med den styrer ud paa et Hav,
hvor Dogmerne --- de Fyrtaarne, tidligere Slægter have rejst --- ikke længere
viser Vej. Ud over alt Givet søger den forskende Tanke, naar Grænseproblemerne
for Alvor gaa op for den. I det Arbejde, der her stedse er at gøre, kom
Kierkegaard ikke til at deltage, optagen, som han udelukkende blev, af „at læse
de individuelle, humane Existensforholds Urskrift, det Gamle, Bekendte og fra
Fædrene Overleverede, igennem endnu engang“.

Men dybt har han gravet i Tankens og Aandens Jordbund. Han har følt Spændingen
og Afstanden mellem Ideal og Virkelighed, baade paa Erkendelsens og paa Livets
Omraade, inderligere og lidenskabeligere end Nogen i hans Aarhundrede. Og
hvorledes man saa stiller sig til hans personlige Stade, --- Et har han lært
os: kun gennem personlig Tilegnelse og Gennemarbejden af Livstankerne (hvorfra
vi saa hente dem) og gennem fuld Oprigtighed og Klarhed over, hvor vi staa, kan
vort aandelige Liv blive sundt og sandt. Og derigennem har han i Sandhed øvet
en Filosofs Gerning.

\begin{center}\rule{2cm}{0.4pt}\end{center}

% --- p. 175 ---
\setcounter{footnote}{0}%
\opage{175}

\begin{center}\large 15. FREDERIK DREIER\end{center}

Medens Kierkegaard i sin Ensomhed følte sig i stedse skarpere Modsætning til
den bestaaende Kirke, medens Sibbern fordybede sig i sine Fremtidsfantasier, og
medens H. C. Ørsted arbejdede paa „Aanden i Naturen“, var der en ung medicinsk
Student, der i en Række af Smaaarbejder udkastede et Program for fri Tænken og
for Reform af Samfundet, et Program, der paa en Gang stod i Modsætning til hine
ældre Tænkeres Standpunkter og dannede en Fortsættelse af dem.

Det var Romantikens Ideer, de ældre Tænkere hver paa sin Maade gik ud fra,
selvom de søgte at bringe dem i mere kritisk og realistisk Form. Ørsted og
Sibbern lagde vel Vægt paa Naturerkendelsen, men det var dog den aandelige Side
af Tilværelsen, der laa til Grund ved deres Tydning af Verdensgaaden. Howitz
havde vel bragt det naturvidenskabelige Synspunkt mere i Forgrunden, men hans
Indflydelse havde ikke strakt sig vidt.%
% footnote *) on p. 175
\footnote[1]{Kierkegaard anede, at Naturvidenskaben vilde komme til at spille
en mere fremragende Rolle end hidtil. I Striden mellem Mennesket og Gud vil,
mener han, Mennesket (naar Spekulationens Tid er forbi) forskanse sig bag
Naturvidenskaben. Han morer sig over den nye Kvide, Teologien da vil komme i;
det er Nemesis, fordi den vilde være Videnskab! I hans egne Øjne er Sagen ikke
vanskelig: den kristelige Etik er og bliver den samme, selvom man forvexler et
Æble med en Pære, og hvad enten Solen gaar rundt om Jorden eller omvendt. Og
Gud blander maaske med Villie noget Galimatias ind i Verden for at prøve
Menneskene og forarge Naturforskerne og Selskabet for Naturlærens Udbredelse!
\textit{(Efterladte Skrifter.} 1851---53. p. 287---289).}

Nu forsøgte Frederik Dreier (f. 1827, d. 1853) at lægge Naturvidenskaben til
Grund, ja egentligt at gøre den til den
% --- p. 176 ---
\setcounter{footnote}{0}%
\opage{176}eneste Videnskab. Støttet til den vil han først og fremmest virke
for Bevidsthedens Reform, for Tankens Frigørelse ikke blot fra ydre
Autoriteter, men ogsaa fra Overtro, Mystik og Abstraktion, som endnu raade selv
hos Naturforskere og Filosofer. Er dette sket, vil Resten følge af sig selv;
frigøres Tanken fra Teologien, vil Arbejdet frigøres fra Kapitalen, Folket fra
Staten.

Der gaar en stor Begejstring gennem den unge Forfatters Skrifter.%
% footnote *) on p. 176
\footnote[1]{Hovedskriftet i filosofisk Henseende er \textit{Aandetroen og den
frie Tænkning.} 1852. --- Forud havde Dreier anonymt udgivet \textit{„Folkenes
Fremtid.} Af en Fritænker“ (1848) og \textit{„Fremtidens Folkeopdragelse.} Af
en Socialist.“ (1848), samt en skarp Kritik af Goldschmidt og et Indlæg i Klara
Raphael Fejden. --- En udførlig Biografi og Fremstilling af hans Teorier findes
i C. E. \emph{Jensen} og F. J. \emph{Borgbjerg}: \textit{Socialdemokratiets
Aarhundrede.} Andet Bind. p. 32---135.} Trods sin teoretiske Materialisme er
han en praktisk Idealist. Han var en overvejende intellektuel Natur. De andre
Glæder, Livet kan skænke, havde i hans Øjne store Fremskridt at gøre, før de
kunne naa Glæden ved videnskabelig Tænkning og Syslen. Og det videnskabelige
Fremskridt fører nødvendigvis til det sociale. Med den intellektuelle Glæde
forbinder sig derfor naturligt Iveren for Samfundets Reform. Resignationen har
ingen Plads i Dreiers Livsanskuelse; den viger Pladsen for en kraftig Virken
for Reform i alle Retninger. Hans højeste Ønske er at virke for Harmoni og
Skønhed i Livet. Den frie Videnskabs Mænd og Arbejderne skulle i Forening føre
Kampen mod al gammel Slendrian. Hvad Religionen var for tidligere Tider, det
kan i vor Tid Reformens Lidenskab være. Al Videnskab, al dygtig og ædel
praktisk Virksomhed understøtter den og staar i dens Tjeneste. Det berettes
ogsaa om ham, at han fra Barndommen af nærede stor Medfølelse med Lidende og
Ulykkelige.

H. C. Ørsted og Sibbern høre i hans Øjne til den mystiske og dualistiske
Retning i Tænkningen; han finder dog i dem Forløbere for Antagelsen af Arternes
Oprindelse gennem naturlig Udvikling, saa at ogsaa Mennesket er opstaaet som
Frugt af en lang Udviklingsgang. Ørsteds Naturfilosofi kriti%
% --- p. 177 ---
\setcounter{footnote}{0}%
\opage{177}serer han skarpt; Udkast til et specielt Skrift om dette Emne findes
blandt hans Papirer paa det kongelige Bibliotek. Kierkegaard tager han til
Indtægt for den skarpe Modsætning mellem Teologi og Videnskab, vor Tids
skarpeste Modsætning. Der er, netop ved den extreme Modsætning i Realiteten, en
vis Analogi mellem Kierkegaards og Dreiers Tænkning. Begge ville de bort fra
„Begrebsmystiken“ til Existensen. Men medens Kierkegaard mest eller udelukkende
har de existerende Subjekter for Øje („Subjektiviteten er Sandheden“), lægger
Dreier Hovedvægten paa de existerende Objekter („Alt udenfor os er Stof, og Alt
indeni os er Stof“). I begge Henseender kunde der da ogsaa efter Romantikens
overstrømmende Tankeudviklinger trænges til en ædruelig Reduktion.

Sin Hovedpaavirkning har Dreier faaet fra sin Tids Naturvidenskab, især som det
synes fra Dubois Reymond og Carl Vogt. Han beundrer desuden det attende
Aarhundredes Encyklopædister, især Holbach og Condorcet, i hvilke han ser sine
Forgængere, og han finder paa mange Punkter Vejen banet for sig af den
radikale, venstre Fløj af Hegels Disciple (Feuerbach, Stirner), saavel som af
de franske socialistiske Forfattere. Af engelske Forfattere nævner han Bentham,
Owen og især Stuart Mill, hvem han beundrer som Logiker, som „en skarp,
konsekvent Tænker, der ikke respekterer Begrebsmystiken“, medens han kritiserer
hans Nationaløkonomi.

Det er, som man ser, Strømninger, der ikke hidtil havde gjort sig gældende her
i Landet, som den unge Reformator er kommen i Berøring med. Vi skulle nu se,
hvorledes hans Hovedtanker have udformet sig under hans korte Løbebane. Det
blev kun til Spirer, men var hans Liv ikke saa tidligt blevet afbrudt, vilde de
sikkert have udfoldet sig i frodig Væxt.

\begin{center}\rule{2cm}{0.4pt}\end{center}

Det er den filosofiske Side af Dreiers Anskuelse, jeg her især holder mig til,
skønt det er hans socialreformatoriske Ideer, der i den nyeste Tid især have
genvakt Opmærksomheden for ham.

% --- p. 178 ---
\setcounter{footnote}{0}%
\opage{178}Dreier følte sig i en stærk Modsætning til Romantikens Filosofi, i
hvilken han kun kunde se et sidste Forsøg paa at forsvare „Aandetroen“. Han
inddeler al Filosofi i apologetisk og kritisk Filosofi. De apologetiske Skoler
forudsætte uden videre den nedarvede Tros Grundbegreber og have kun til Hensigt
at lempe dem sammen. De kritiske Skoler gaa ud paa at ophæve de Modsætninger
(Gud og Verden, det Gode og det Onde o. s. v.), som Teologien forudsætter og de
teologiserende Filosofer omforme. Naar dette kritiske Arbejde er gjort, er
Filosofien ikke længere en særegen Videnskab; den smelter sammen med hele den
øvrige Videnskab som dennes almindeligste begrebsadskillende og
begrebssammenfattende Metode.

Ved denne sidste Antydning gaar Dreiers Opfattelse i lignende Retning som
Kierkegaards Ide om et Tankesystem, der skulde omfatte de Grundbegreber, vor
Erkendelse arbejder med. Og det er en Opgave, Filosofien til alle Tider maa
stille sig, den mest centrale af alle dens Opgaver. Hvad den kan udrette til
Kritik af gængse og overleverede Forestillinger, beror tilsidst paa dens
Forstaaelse af selve vor Tænknings, vor Forestillingsevnes Natur og Virkemaade.
Det er, som det ligger i Dreiers Ord, de sidste Forskelle og de sidste Enheder,
Filosofien skal finde, saavidt det kan ske paa det Trin af Erkendelsens
Udvikling, hvor Filosofen foretager sit Arbejde.

Dreier tænker sig kun de teologiske Forestillinger og deres Eftervirkninger som
Genstand for denne Kritik. Han har ikke set, at selve Naturvidenskabens
Grundbegreber og Forudsætninger maa være Genstande for samme Art Kritik, og at
denne Kritik forlængst var i Gang. Derved kommer der noget Naivt i hans
Tankegang, naturligt fremkaldt ved hans Begejstring for sin specielle Videnskab
og dens Mission.

I sin Opfattelse af Erkendelsen bygger Dreier paa den Forudsætning, at Sansning
er Grundlaget for al Tænkning. Men idet han ved Sansning kun forstaar Modtagen
af Indtryk, Indtryk, der saa kunne efterlade Rester (Sansebilleder), som kunne
gentages, selvom ingen Indtryk foreligge, forudsætter han dog --- som
forskellig fra denne Modtagen og dens
% --- p. 179 ---
\setcounter{footnote}{0}%
\opage{179}Eftervirkninger --- den Proces, hvormed Sansebillederne forbindes og
sammenlignes. Det ser ud, som om de enkelte Sansebilleder vare uafhængige af
enhver Forbindelse og Sammenligning. Men saa maa „Tænkningen“ da have et andet
Grundlag end „Sansningen“ (i den Betydning, i hvilken Dreier tager dette Ord).
Paa lignende Maade forklarer han Bevægelse som Resultat af Tingenes
Vexelvirkning --- somom Tingene først vare der, saa kom i Vexelvirkning, og
Bevægelsen igen var Følge af denne Vexelvirkning, --- og somom Tingenes (de
materielle Tings) Vexelvirkning skyldtes Andet end Bevægelse.

Dreier klager over den Begrebsmystik, der fremkommer ved Tilbøjeligheden til at
gøre de abstrakte Begreber til lignende, for sig selv fremtrædende Enkeltheder
som selve Sanseforestillingerne. Men han er selv Mystiker, hvad de simple
Sanseforestillinger angaar, tillægger dem en Selvstændighed, de ikke have. Og
han, der hader al Dualisme, gaar ud fra en Dualisme mellem de enkelte
Forestillinger paa den ene Side, deres Forbindelse og Sammenligning paa den
anden Side.

Paa Naturopfattelsens Omraade er det især Modsætningen mellem Stof og Kraft,
han vil tillivs. Det er i hans Øjne den sidste Rest af Modsætningen mellem
Verden og Gud, Aandetroens sidste Udløbere. Man maa ikke betragte Stof og Kraft
som selvstændige Væsener, og han priser Système de la nature for at have ryddet
godt op her. Efter hans egen Mening maa Ophævelsen af hin Dualisme ske ved, at
Stoffet gøres til Alt: „Alt, hvad der er til udenfor os, er Stof; Alt, hvad der
er til indeni os, i vort Legeme eller i vor Bevidsthed, er ogsaa Stof, i visse
bestemte Forbindelser, Former og Virksomheder.“ --- Men er „Stof“ ikke ogsaa et
abstrakt Begreb? Existerer det afset fra de specielle Former og Virksomheder,
som er det, vi have givet, og som vor Tænkning skal undersøge?

Som det gaar med Modsætningen mellem Stof og Kraft, skal det ogsaa gaa med
Modsætningen mellem Legeme og Sjæl, der har samme teologiske Oprindelse.
Sjælebegrebet
% --- p. 180 ---
\setcounter{footnote}{0}%
\opage{180}skyldes en af den begyndende Tænknings sædvanlige Personifikationer.
At Sjælen ikke er materiel, vil sige, at den ikke er virkelig; ti vi kende
Intet, som ikke er legemligt. Hvad man mente med Ordet Sjæl, var, hvad vi nu
kalde Hjernevirksomhed. Der er ganske vist Forskel mellem Bevidsthed og Hjerne,
men det er den samme Forskel, som mellem Maven (med Tarmene) og Fordøjelsen.
Carl Vogt har ikke Ret i at sige, at Tanken forholder sig til Hjernen som
Galden til Leveren; ti Tanken er en Virksomhed%
% footnote *) on p. 180
\footnote[1]{Carl Vogt bruger iøvrigt en mere uartig Sammenligning. Se om hans
Udtalelse og om den materialistiske Teori idethele min \textit{Psykologi} II, 8
c.}. --- Her kommer altsaa den samme Forskel frem igen, som før, mellem Stoffet
og dets Virksomhed. Og desuden bestaar et Organs Virksomhed jo kun i Bevægelse.
Hvorfra kommer da Forskellen mellem en Tanke og en Bevægelse? ---

Den begejstrede unge Mediciner har ikke Tid til at besvare saadanne Spørgsmaal.
Han vil have Tankelivet saa simpelt og klart som muligt for at kunne gaa over
til sine sociale Reformplaner. Han har end ikke Tid til ret at tænke over sin
egen Begejstring. Hvad man kalder Følelse er i hans Øjne noget Underordnet,
ofte kun uklar Erkendelse. Analyserer man det nærmere, bestaar det dels af
Tanker, Forestillinger og Fantasibilleder, dels af Organfornemmelser, der
stamme fra Forandringer i Ernæring, Afsondring, Blodomløb o. s. v.%
% footnote **) on p. 180
\footnote[2]{Som man ser en Antydning i Retning af den senere af William James
og Carl Lange opstillede Teori.} Han vilde dog maaske ved nærmere Overvejelse
have fundet denne Teori noget for elementær.

At han ikke kan tillægge religiøs Følelse nogen Betydning, følger af, hvad der
er sagt i det Foregaaende. Det er dernæst hans Overbevisning, at det vil gaa
Poesien, som det gaar Religionen. Dens Tid er forbi. Hvor Poesien endnu raader,
er der et Hul i Videnskaben, ligesom hvor Religionen raader. Oprindeligt
svarede Poesien ganske til Folkenes intellektuelle Udvikling, idet den gav en
saadan Beskrivelse og Forklaring af Virkeligheden, som man paa et lavt Trin af Erken%
% --- p. 181 ---
\setcounter{footnote}{0}%
\opage{181}delse kunde have. Men nu maa man jo studere gamle Mytologier for at
forstaa Poesien! --- Fantasien er uundværlig for Videnskaben, men den skal være
tjenende. I Modsætning til Poesien beskriver Videnskaben Virkeligheden saa
nøjagtigt og med saa stor Sammenhæng som muligt; Poesien bruger Billeder og
Forestillinger, der høre en længst forsvunden Tid til. Billedkunstens Betydning
er at være Illustration, Landskabsmaleriet f. Ex. for Geografien. En lignende
Betydning vil den psykologiske og sociale Roman kunne have paa dens Omraader%
% footnote *) on p. 181
\footnote[1]{I et Manuskript med Overskrift \textit{Poesien,} der findes blandt
Dreiers Papirer paa det kgl. Bibliotek, forekommer denne Tankegang.}.

Det forholder sig paa lignende Maade med Nationalfølelsen. Det danske Folk er
Indbegrebet af alle Danske. Det forsvinder ikke, fordi hver enkelt Dansk taler
Tysk, ligesaalidt som det enkelte Individ forsvinder, fordi han faar et andet
Sprog. Hvortil saa alle de heftige nationale Stridigheder? Især da Antagelsen
af det fremmede Sprog kun vil være foreløbig, da et almenmenneskeligt Sprog vil
være den nødvendige Konsekvens af den almenmenneskelige Videnskab. Og allerede
nu er der ikke det Fællesskab indenfor et Folk, som der er mellem enkelte
Stænder i Folket og de tilsvarende Stænder i andre Folk.

Istedetfor de Følelser, hvis Tid er forbi, vil en ny Samfundsorden udvikle nye
Følelser. Nu er det Konkurrencen, der behersker Alt. Men det gælder om at
fremelske Kærlighed til selve Arbejdet og om at muliggøre den Erkendelse, at
Menneskenes Interesser væsentligt ere de samme, og at den Enes Lykke er knyttet
til den Andens. Dette vil være muligt ved en socialistisk Samfundsordning, hvor
Kapitalen er Fælleseje. Da vil det Væsentlige i Kristendommen kunne ske
Fyldest: den store Protest mod al Uretfærdighed, Kærlighedsbudet og Fordringen
om ej at gøre mod Andre, hvad man ej vil, de gøre mod os. De første Kristnes
Kommunisme var en Forudsætning for deres Broderlighed. Det var nu en mindre
heldig Organisationsform; men den moderne sociale Reform vil have en lignende
Virkning.

% --- p. 182 ---
\setcounter{footnote}{0}%
\opage{182}De virkeligt værdifrembringende Klasser i Samfundet ere Arbejderne,
Kunstnerne og Videnskabsmændene. Det er deres Sag at bekæmpe de gamle
Illusioner og fremvirke en ny Tingenes Orden. Reformerne ville ikke udgaa fra
Staten, men fra de enkelte Individer, som maa slutte sig sammen og dernæst
efterhaanden gennem Udvidelse af den kommunale og politiske Valgret opnaa
fuldstændigt Selvstyre. Det gælder da om at hindre privat Kapitaldannelse,
hvilket kan ske ved Udelukkelse af Mellemhandel, idet der oprettes Varehuse
under offentlig Bestyrelse, hvor Producenter kunne aflevere deres Varer, og
Forbrugerne modtage dem. En gensidig Kreditforening ordner Omsætningen og
fastsætter Arbejdsløn og Varepriser%
% footnote *) on p. 182
\footnote[1]{Se Dreiers Tidsskrift \textit{Samfundets Reform.} 1853. (Især
Artiklen „Kapital og Arbejde). --- Om den store Betydning, Brugsforeninger
senere have faaet, smlgn. min \textit{Etik.} 3. Udgave. p. 382 f, hvor deres
Betydning for Løsningen af det sociale Spørgsmaal paavises.}. --- Hvad man end
mener om Gennemførligheden af en saadan Ordning, betegner Ideen dog et
Fremskridt i Forhold til Sibberns naive Henvisning til, at Enhver kan tage,
hvad han vil, da Ingen har Grund til at tage Mere, end han behøver. Dreier
afviger ogsaa derved fra Sibbern, at han lader den økonomiske Ordning være
Betingelse for det broderlige Sindelag, medens dette hos Sibbern er
Forudsætningen for hin (saaledes som det vist ogsaa var Tilfældet ved den
urkristelige Kommunisme, skønt Dreier synes at mene det Modsatte). Dreier
minder her om Karl Marx, hvis første Skrifter han har kendt, og ifølge hvem al
Samfundsorden og al aandelig Udvikling bestemmes ved den økonomiske Produktion
og dens Ordning til hver given Tid. Der er dog her en stadig Vexelvirkning
mellem Tænkemaade og Samfundsordning; de paavirke stadigt hinanden, og det var
en spekulativ Eftervirkning hos Marx, naar han vilde konstruere al
Kulturudvikling ud fra det økonomiske Synspunkt%
% footnote **) on p. 182
\footnote[2]{Smlgn. om Marx mit Værk \textit{Den nyere Filosofis Historie.}
Anden Udgave. II p. 282, og min \textit{Etik.} Tredie Udgave p. 387---394.}.

Dreiers pædagogiske og sprogvidenskabelige Undersøgelser gaar jeg ikke nærmere
ind paa. Selv uden at tage Hensyn
% --- p. 183 ---
\setcounter{footnote}{0}%
\opage{183}til dem vil man have faaet Indtrykket af, hvor mange Ideer der optog
den unge Mands Interesse. Hvilken Skæbne de vilde have faaet i hans Sind og
udenfor ham, hvis hans Bane ikke var blevet afbrudt, er ikke godt at sige. Nu
kom han blot til at staa som en Bebuder af, at Romantikens Tid var forbi, og at
der skulde tages ædrueligt Stade til Virkeligheden baade i Praxis og Teori.

\begin{center}\rule{2cm}{0.4pt}\end{center}

% --- p. 184 ---
\setcounter{footnote}{0}%
\opage{184}

\begin{center}\large 16. RASMUS NIELSEN\end{center}

Dreiers Optræden var en lille Episode, der ikke blev meget bemærket, skønt den
nok kunde give Anledning til Eftertanke hos de Filosoferende. Der kom i hans
Smaaskrifter, ligesom hos Kierkegaard, Fordringen frem om en Filosofi, der
gennemførte en Undersøgelse af selve Videnskabens Grundbegreber og
Grundforudsætninger. Men Spekulationen øvede endnu sin Magt, og den
Undersøgelse, der var krævet baade fra et religiøst og et naturvidenskabeligt
Synspunkt, blev foreløbigt opsat.

Forat følge Filosofiens fortsatte Skæbne i Danmark maa vi nu gaa tilbage igen
til 1840, det Tidspunkt, da baade Martensen og Kierkegaard begyndte deres
Forfattervirksomhed, den Hegelske Periode.

Ligesom Martensen havde Rasmus Nielsen (f. 1809, d. 1884)%
% footnote *) on p. 184
\footnote[1]{Biografi og Fremstilling hos P. A. \emph{Rosenberg}:
\textit{Rasmus Nielsen.} Nordens Filosof. 1903. --- Jeg tror, at min gamle
Lærer, hvis Paavirkning jeg altid mindes med Taknemmelighed, vilde have været
bedre tjent med en mere kritisk Fremstilling end den Panegyrik, Hr. Rosenberg
her har givet. Jeg har i hvert Tilfælde forladt Rasmus Nielsens Skole med en
stærkere Følelse af Problemernes Brod, end der spores i den anførte Bog, hvor
der næsten „løses“ et Problem paa hveranden Side.} sat sig den Opgave at
anvende den Hegelske Filosofi til Bedste for Teologien, dog saaledes, at den
nævnte Filosofi tillige undergik en nødvendig Ændring. Nielsen disputerede 1840
for den teologiske Licentiatgrad med en Afhandling om \textit{den spekulative
Metodes Anvendelse paa den hellige Historie.} Den Indvending, der her rettes
mod Hegel, er omtrent
% --- p. 185 ---
\setcounter{footnote}{0}%
\opage{185}den samme, som Martensen nogle Aar forud havde fremført. Det er
Forskellen mellem Tanken og Virkeligheden, der indskærpes: ikke Menneskets
Tanke, kun den guddommelige Tanke, der tillige er skabende, er udtrykt i
Virkeligheden.%
% footnote *) on p. 185
\footnote[1]{\textit{De speculativa historiæ sacræ tractandæ methodo
commentatio.} 1840. --- Grundtanken er den samme som i Martensens Disputats
nogle Aar forud: Troen som Grundlag for Spekulationen, i Tillid til, at der i
Troen indeholdes en skjult Dialektik (p. 43 f).} Martensen havde udtrykt denne
Forskel som en Modsætning mellem Viden og Magt, og denne Modsætning kom til at
spille en stor Rolle i Nielsens Tænkning gennem alle dens forskellige Perioder.
Begrebet Magt havde Hegel indført i sin Logik i det Afsnit, i hvilket han gør
Overgangen fra Begrebet Substans som blot Iboen eller Grund til Begrebet
Kausalitet, og Martensen havde hævdet, at man ikke kunde foretage denne
Overgang uden Forudsætning af Skabelsesdogmet, da kun den guddommelige Tanke
kan være Aarsag til en fra den selv forskellig Realitet.

Det Arbejde, der var at gøre, blev nu delt mellem de to Mænd, der dog
foreløbigt arbejdede meget sammen, Martensen paa sin Dogmatik og Nielsen paa et
filosofisk System, som skulde lade Hegel undergaa samme Skæbne, denne selv
havde beredt saa mange tidligere Filosofer, nemlig at gøres til Element i en
„højere Enhed“. For Nielsens livfulde Fantasi, en Evne, der var stærkt udviklet
hos ham, formede det nye System sig allerede i sin Grundtanke, og han udgav en
Subskriptionsindbydelse til et spekulativt Værk, som dog ikke blev fuldendt,
saavel som til en Moralfilosofi, der ikke udkom. Kierkegaard spottede over de
mange Bebudelser af „Systemet“, der skulde løse alle Gaader, selv dem, Hegel
ikke havde kunnet magte.

Nielsens vigtigste Værk fra denne Periode er hans \textit{Propædeutisk Logik}
(1845). Han udvikler den allerede antydede Kritik af Hegel, hvem han bebrejder
en Formalisme, der stedse lader Forholdet mellem Tanke og Virkelighed staa
uopklaret. „Den formelle Begriben mangler Almagten!“ Hegels Dialektik, den
stadige gensidige Belysning, Begreberne un%
% --- p. 186 ---
\setcounter{footnote}{0}%
\opage{186}derkastes, anerkendes og beundres. Men ved den uklare Identificeren
af Begreb og Virkelighed bliver Hegel en logisk Mystiker. --- Den
Begrebsmystik, Dreier klagede saa meget over, kan han maaske være bleven
opmærksom paa gennem Nielsens Kritik af Hegel. ---

Det var efter denne Tankegang ikke vanskeligt at gøre Overgangen fra Filosofi
til Teologi; den var jo egentligt allerede gjort. Gudsbegrebet var allerede
indført som sidste eller første Aarsag (Almagt), og det gør tillige Tjeneste
som absolut Subjekt, hvorved den Vanskelighed skal være løst, at der skulde
existere Objekter, der ikke vare til for noget Subjekt, altsaa ikke
objektiveredes. De to Begreber Magt og Objektivering gaa i Et. Den
Sammenblanding af en logisk (erkendelsesteoretisk) og en metafysisk Tankegang,
som den tyske Spekulation hviler paa, gør sig altsaa stadigt gældende. Det var
den, Kierkegaard protesterede imod ved sin Distinktion mellem et logisk System
og et Tilværelsens System, og som iøvrigt allerede Fries havde kritiseret hos
Fichte og Schelling.

Forholdet mellem Filosofi og Teologi, som det stillede sig for Nielsen i denne
første Periode, er udtrykt i følgende Sætninger: „Det videnskabelige Medium, i
hvilket Filosofiens teoretiske og Religionens praktiske Idé gensidigt forklare
og eticere hinanden, er Teologien. Ti Teologien søger ikke blot en
Livserkendelse for Sandhedens Skyld, men ogsaa en Sandhedserkendelse for Livets
Skyld, og er desaarsag med Rette bleven betragtet som alle Videnskabers
væsentlige Midtpunkt, som den egentlige og sande Centralvidenskab“. Teologien
er her altsaa baade Udgang og Afslutning.

Det var den Harmoni mellem Teologi og Filosofi, som raadede ved vort
Universitet i Fyrrerne under Martensens og Nielsens Samarbejde, kun forstyrret
af Kierkegaards Ironi og Kritik og af nogle radikale Hegelianeres (Brøchners og
Becks) Indsigelser. Sibbern foretog samtidigt i Stilhed Overgangen til sit
senere Stade, som han udtalte i Programmerne for 1846 og de følgende Aar.

Da kom der en brat Ende paa Samarbejdet og Harmonien.

\begin{center}\rule{2cm}{0.4pt}\end{center}

% --- p. 187 ---
\setcounter{footnote}{0}%
\opage{187}Det gik nemlig op for Nielsen, at der i Kierkegaards Skrifter saa
energisk var peget paa den personlige, existentielle Side ved Kristendommen, at
Muligheden af at omsætte de kristelige Dogmer i Spekulationens Form maatte
falde bort. „Systemet“ sprængtes, en spekulativ Teologi blev en Umulighed. Han
traadte i Forbindelse med Kierkegaard, blev hans ivrige Tilhænger. Og ligesom
man i København tidligere havde set Sibbern spasere med Poul Møller, og Poul
Møller med Kierkegaard, saaledes kunde man nu (hver Torsdag) se Kierkegaard
spasere med Rasmus Nielsen. Alle disse Spasereture have havt deres Betydning
for dansk Aandsliv.

Da Martensens Dogmatik udkom, rettede Nielsen, støttet til Kierkegaards
„Uvidenskabelig Efterskrift“ et stærkt Angreb paa den, og der udspandt sig en
hel teologisk-filosofisk Fejde. Nielsen udgav desuden i de følgende Aar flere
populære Skrifter, af hvilke især maa nævnes \textit{Evangelietroen og den
moderne Bevidsthed} (1849).

I to Aar (1848---1850) bestod det personlige Forhold mellem Kierkegaard og den
Discipel, han saa uventet havde faaet, og som havde brudt med Venner og
Kolleger for at slutte sig til den ensomme Tænker. Det var ikke, som man har
sagt, af krænket Forfængelighed, Kierkegaard igen trak sig tilbage, idet han
nemlig skulde have taget Nielsen ilde op, at han i sine nye Skrifter ikke
tilstrækkeligt havde givet sin nye Mester Æren for sine Synsmaader. Ganske vist
fandt Kierkegaard, at Nielsen pyntede sig endel med laante Fjer. Men sin
Opfattelse af Forholdet har han udtalt i følgende Ord (i Slutningen af
Dagbøgerne fra 1849): „R. Nielsens Fortjeneste er, at han er blevet opmærksom
(ikke Lidet, naar man er en ældre Mand). Hans Fejl er, at han saa dog gør det,
der ingen Videnskab er, til Videnskab.... Og dernæst er det hans Fejl,
istedetfor at indskrænke sig til enfoldigt at tilstaa Noget om sig selv, strax
at ville benytte det Lærte til at angribe en Anden.“ --- Hvad Kierkegaard har
savnet hos Nielsen, har været den Inderlighed i det personlige Liv, der først
og fremmest gennemarbejder det Modtagne og søger at prøve det praktisk. Og det,
der her var det Modtag%
% --- p. 188 ---
\setcounter{footnote}{0}%
\opage{188}ne, var jo netop en praktisk Opfattelse, en personlig Holdning. Men
som forhen Hegels Indflydelse, saaledes ansporede nu Kierkegaards Indflydelse
først og fremmest den begavede Mands Fantasi og Spekulationsiver. Han var en
livlig og vækkende Taler, der forstod at fremstille Tankens og Livets
Modsætninger i store Billeder og stærke Træk.

I sin Hegelske Periode var han bunden til stedse at skulle finde den „højere
Enhed“, hvad enten han nu kunde finde den indenfor Filosofien eller tog
Teologien til Hjælp. Men nu, da ikke blot Filosofi og Teologi vare blevne til
uforsonlige Modsætninger, men det overhovedet saa galt ud med al „højere
Enhed“, maatte Evnen og Lysten til at virke ved Modsætninger komme til at
spille en saa meget des større Rolle. Og han øvede en ikke ringe Virkning
gennem de Antiteser, han ofte fremstillede med slaaende Kraft og med æggende
Skarphed.

\begin{center}\rule{2cm}{0.4pt}\end{center}

Men det nye Standpunkt blev dog ogsaa Indledning til en stor Arbejdsperiode.
Anerkendelsen af, at Forholdet mellem „Livserkendelse“ og „Sandhedserkendelse“
ikke var saa simpelt, som han i sin første Periode havde ment, maatte føre ham
til at underkaste Forholdet mellem Tanken og Virkeligheden overhovedet en ny
Prøvelse. Der kom nu en Række Arbejdsaar, i hvilke han fordybede sig i Studiet
af Matematik og Naturvidenskab for at faa et solidt Grundlag for Belysningen af
det nævnte Forhold. I denne Periode øvede han sikkert sin største Virkning,
fordi han selv var en Søgende.

Dog tog han ikke Spørgsmaalet op fra Grunden af. Han blev ikke
Erkendelsesteoretiker, men fastholdt trods Alt Hegels Grundtanke, skønt han nu
gav den en ny Modifikation. Vel raader der i de forskellige
Indledningsarbejder, han under Navn af Propædeutik udgav i Halvtredserne og
Begyndelsen af Tredserne, endnu en vis kritisk Forsigtighed. Fra Formens Side
er Propædeutiken fra 1857 hans mest fuldendte Arbejde; her raader den
Begrænsning, han ellers ikke var Mester i. Men det stærkere Præg af selvstændig
Realitet,
% --- p. 189 ---
\setcounter{footnote}{0}%
\opage{189}de enkelte Kendsgerninger nu efter Programmet skulde have, blev dog
ikke opnaaet. Enten optager han Stoffet fra de andre Videnskaber helt
uforarbejdet, eller ogsaa anvendes det blot som Exempler paa den Hegelske
Dialektik. Matematiske Funktioner, fysiske og kemiske Processer, organiske
Udviklingsformer, Alt optraadte som dialektiske Bevægelser --- kun at, som
sagt, de højere Enheder blev mere og mere usynlige. Den Sammenhæng, i hvilke
Naturfænomener stedse staa med hverandre, saaledes at f. Ex. to Stoffer kunne
ved Forbindelse frembringe et tredie, fortolker Dialektikeren uden videre som
en logisk Sammenhæng --- til Fordel for Godtgørelsen af „det objektive Begrebs“
Gyldighed. Og gennemgaaende gøres det, der væsentligt er Metode eller brugbar
Analogi, som Infinitesimalregningen og Atomteorien, til absolute Synspunkter.
Enheden i Tilværelsen blev under forskellige Former anskueliggjort ved
Analogien med menneskelige Begrebers Forhold til deres Indholdselementer
(Momenter), uden at det dog her, ligesaalidt som hos Hegel, blev indrømmet, at
det (kun) var Analogi, der raadede. Selve Berettigelsen og den definitive
Betydning af en saadan Analogi blev ikke drøftet.

Efter Afslutningen af sine Forstudier begyndte Nielsen Udgivelsen af et stort
Værk, hvis Navn var \textit{Grundideernes Logik.} Kun to Dele udkom (1864---66)
--- og de udgjorde kun to Trediedel af det paatænkte Værks første Niendedel!
Atter her viste det sig, at Nielsens Fantasi ilede forud for hans Tanke. Han
saa i Aanden, at der til at beskrive Vexelforholdet mellem det Subjektive og
det Objektive, eller som han stadigt kalder det, Viden og Magt, behøvedes en
uendelig Analyse, da ethvert Subjekt forudsætter et Objekt, og ethvert Objekt
igen et Subjekt. Naar man ikke vil slutte af paa spekulativ og teologisk Maade,
er det en Holmgang, der aldrig kan faa Ende --- og det var derfor intet
Tilfælde, at Værket ikke blev afsluttet. Videnskabeligt set maa man her lade
sig nøje med at paavise det stadige Forhold mellem Subjekt og Objekt som en
Lov: enhver Egenskab, man tillægger Subjektet, kan selv gøres til Genstand
(Objekt) og maa have sin
% --- p. 190 ---
\setcounter{footnote}{0}%
\opage{190}Grund i visse objektive Forhold; og ethvert Objekt, eller ethvert
objektivt Forhold opfattes paa en ved det opfattende Subjekts Natur bestemt
Maade --- og saaledes kan man blive ved i Kraft af Loven for Forholdet mellem
Subjekt og Objekt.

Nielsen opfattede nu ikke Sagen paa denne Maade. Hans Tankegang var, at da
ethvert Objekt maa objektiveres, d. e. forudsætter et Subjekt, og da det
menneskelige Subjekt ikke kan opfatte (objektivere) Alt, maa der, hvis
Objekternes Realitet skal hævdes, være et absolut („ontologisk“) Subjekt, for
hvilket den absolute Realitet er til. Han overser, at Spillet maa begynde igen
her; selv en Gud er bunden til Loven om Forholdet mellem Subjekt og Objekt.
Forsøget paa gennem „Grundideernes Logik“ at begrunde en abstrakt Teisme er
derfor ikke lykkedes. Vi kunne ikke komme længere end til at bestemme og
beskrive et Subjekt, i Forhold til hvilket visse Fænomener (Objekter) gælde,
ligesom Astronomen maa bestemme et Punkt (paa Jorden, paa Solen eller
hvorsomhelst), ud fra hvilket Himmellegemernes Stillinger og Bevægelser kunne
beskrives, somom de vare absolute.

Der findes mange fine Analyser i Nielsens ufuldendte Værk. Og hvad der i den
danske Filosofis Historie har særlig Interesse, er den stærke Vægt, der lægges
paa det Konkrete og Individuelle, som ikke gaar op i noget Begreb. Nielsen
bruger her endog Udtrykket „det Antilogiske“ for at betegne dette i almindelige
Begreber Uopløselige. Det er for ham ikke, som for de strenge Hegelianere, en
ligegyldig Rest. Det er for ham det slaaende Udtryk for en Magt, som vi ikke
kunne komme udenom. Dog er ved Udtrykket „det Antilogiske“ Modsætningen til
Erkendelsen ikke træffende karakteriseret. Det er det Uudtømmelige, det stadigt
Inkommensurable i ethvert konkret Fænomen, som vi her staa overfor; men
„antilogisk“ er ikke det rette Udtryk, da der stedse gennem nye Iagttagelser og
nye Sammenligninger er Mulighed for nye Bestemmelser. ---

Som den utrættelige Arbejder, Nielsen var, søgte han stedse at finde nye Former
for sine Ideer. Foruden en hel Række af populære Skrifter maa nævnes hans
\textit{Religionsfilosofi} (
% --- p. 191 ---
\setcounter{footnote}{0}%
\opage{191}1869), hans \textit{Naturfilosofi} (1873) og tilsidst
\textit{Filosofiske Grundproblemer} (i Universitetets Festskrift 1879); dette
sidste er den klareste og mest sammentrængte Fremstilling, Nielsen har givet af
sin Opfattelse, saaledes som den formede sig i hans senere Aar. Han tillægger
her Erkendelsesproblemet større Selvstændighed, end han har gjort i sine
tidligere Arbejder.

Jeg vil her blot dvæle ved den Maade, hvorpaa det religiøse Problem stillede
sig i Nielsens Filosofi efter Tilslutningen til Kierkegaard.

\begin{center}\rule{2cm}{0.4pt}\end{center}

Da jeg i Aaret 1861 kom til Universitetet, stod Nielsen i sin fulde Kraft, fuld
af Iver og Virkelyst, og det Kursus, jeg hørte hos ham i mit første
Studenteraar, fik stor Betydning for mig. Han læste væsentligt over
Kierkegaard, idet han gennemgik „Stadierne“ efter et lille Udvalg af
Hovedsteder i Kierkegaards Skrifter. Vel var atter her Anlægget saa stort, at
det ikke blev udført, men med stor Veltalenhed førte han dog sine unge
Tilhørere paa Vej til Forstaaelse af den store Skikkelse, der nyligt havde
forladt Kamppladsen, og hvis Optræden endnu mærkedes i sine Eftervirkninger.
Til denne Fremstilling sluttede sig saa et lille Omrids af Logiken. Jeg fik her
et levende Indtryk baade af det mest Personlige og af det mest Abstrakte,
Tænkningen maa komme til at staa overfor, de to Endepunkter, hvis indbyrdes
Forhold det gælder at bestemme. Nielsen var dengang endnu en Søgende, og netop
derfor viste han ud til store Horizonter. Da jeg nu havde min Nød med
Teologien, i hvilken jeg hverken ret kunde se det Personlige eller det Logiske
komme fuldt ud til deres Ret, var det dualistiske Standpunkt angaaende Tro og
Viden, som Nielsen i den Periode hævdede, mig en Beroligelse. Jeg deltog nogle
Aar senere i den Strid om Tro og Viden, Nielsens Udtalelser vakte, med et lille
(højst umodent) Skrift. Min Udvikling førte mig bort fra dette Standpunkt,%
% footnote *) on p. 191
\footnote[1]{Jeg tillader mig herom at henvise til den lille Afhandling, jeg
har kaldet „Forord og Efterskrift til min Religionsfilosofi“, og som findes i
\textit{Mindre Arbejder.} Anden Række. (Især p. 47---51).}
% --- p. 192 ---
\setcounter{footnote}{0}%
\opage{192}men hvad jeg skylder Nielsen, er den Interesse, han vakte for
Problemet, og som har holdt sig hos mig, selv om min Interesse for det er
bleven overvejende psykologisk.

Psykologi var netop Noget, man ikke kunde lære hos Nielsen. Enten talte han i
retoriske og opbyggelige Vendinger eller i den abstrakte Logiks Sprog. Jeg
søger deri Grunden til, at han, uden fra først af selv at mærke det,
efterhaanden gled bort fra den strenge Kierkegaardske Opfattelse af Problemet.
Som Brøchner har gjort opmærksom paa, var der fra først af en ejendommelig
Forskel mellem Kierkegaards og Nielsens Begrundelse af deres fælles Standpunkt.
For Kierkegaard var, som vi have set, Hovedsagen den, at det existerende
Individ lever midt i Tiden, uden nogen sikker Afgørelse af de mest brændende
Spørgsmaal, og mellem Modsætninger, der ikke kunne forsones; for Nielsen havde
Problemet sin Oprindelse fra Modsætningsforholdet mellem det Enkelte og det
Almene, der jo tilsteder saa mange Nuancer og Overgangsformer. Denne Forskel
hænger igen sammen med, at Kierkegaard stod den kritiske Filosofi nærmere, og
jo ogsaa fordrede den genoptagen, medens Nielsens spekulative Dialektik vel en
Stund kunde dæmpes, men snart efter udfoldede sig igen, saavidt det efter den
skete Stansning var muligt.

Det Kierkegaardske Paradox, den Grænse for Tænkningen, Dogmets absurde Indhold
sætter, naar det dog skal anerkendes, mildnedes efterhaanden for Nielsen. I
sine populære Skrifter, i sin Religionsfilosofi og i sin Afhandling om
Grundproblemerne sætter han „Mysteriet“ i al Ubestemthed istedetfor Paradoxet,
og det Usandsynlige istedetfor det Umulige. Og en endnu betydningsfuldere
Vending blev gjort ved den Maade, han stillede sig til Miraklet. Hovedvægten
--- saaledes er hans Tankegang --- er religiøst set ikke at lægge paa alle
Enkeltheder i den ydre Begivenhed; det Afgørende er dens aandelige Betydning;
den Virkelighed, der her er Tale om, er „Aandsvirkelighed“, „Virkelighed for
Aanden“. I sit sidste Skrift forklarer han dette Udtryk saaledes:
„Aandsvirkelighed betyder, at her virkeligt er et Trosindhold, en utvetydig
Aabenbarelse af den guddommelige Villie; Virkelig%
% --- p. 193 ---
\setcounter{footnote}{0}%
\opage{193}hed for Aanden betyder, at det Ydre ikke tør fastholdes i saa
sanseligt umiddelbar Udvorteshed, at her kan være Tale om Brud paa Naturens
Love.“ Her staar Nielsen lige ved Begrebet Myte. Naar Underet ikke behøver at
fremtræde „saa udvortes“, som virkelige Begivenheder ellers pleje, og naar dets
hele Betydning beror paa, hvad det betyder aandeligt set, saa have vi netop,
hvad man kalder en religiøs Myte. Underet er et Symbol, et Tegn, --- „et
udvortes Tegn paa et væsentligt Indvortes“. Vægten ligger da f. Ex. ikke paa,
hvad der virkeligt skete, da Jesus „for til Himmels“, men paa, hvad Apostlene
oplevede i deres Indre, da denne Begivenhed skete. Det blev en lignende
Forklaring af Fortællingen om Himmelfarten, som vi allerede have truffet hos
Martensen, --- og som forresten Nielsen allerede havde antydet i sin Disputats.
Der er, hed det der, altid et Element af Illusion i Aabenbaringsfænomenerne (in
phænomenis revelationis semper \textit{momentum illusionis} insit, necesse
est), fordi Aabenbaringen maa lempes efter den menneskelige Opfattelse. Det gør
da ingen Forskel, om Himmelfarten foregaar i den ydre Tingenes Orden (in
externo rerum ordine), eller om det kun var for Apostlenes Tanke (apostolis
meditantibus), at Fænomenet frembragtes%
% footnote *) on p. 193
\footnote[1]{\textit{De speculativo methodo.} p. 148.}. Det er en haarfin
Forskel, der her er mellem den spekulative Teologi og den mytiske Opfattelse.
Det forholder sig altsaa saaledes, at Gud altid har været Kopernikaner, eller
rettere Brunonianer, men har maattet lempe sig efter den menneskelige
Opfattelse. Da nu Filosofen ogsaa er Brunonianer, maa han ogsaa lempe sig; hvad
der for ham staar som et klart Begreb, er for den positive religiøse Bevidsthed
et Syn, et Billede. Dette var jo netop Grundtanken i Hegels Religionsfilosofi;
til Filosofiens højeste Begreber svarer der hos de Troende konkrete Anskuelser
eller Billeder, der udtrykke det samme aandelige Indhold i Forestillingens
eller Fantasiens Form.

Naar Nielsen saaledes uvilkaarligt var vendt tilbage til sin Ungdoms
Opfattelse, havde det sin Grund i, at han kom til at
% --- p. 194 ---
\setcounter{footnote}{0}%
\opage{194}staa overfor et Spørgsmaal, som Kierkegaard altid havde skudt
tilside, og som man ogsaa i Almindelighed gaar udenom, nemlig den Fortolkning,
Dogmer og Mirakler maa undergaa, naar de virkeligt skulle have personlig og
aandelig Betydning. Det er det religiøse Billedsprogs Problem, der allerede
havde sysselsat Grundtvig og Martensen. Man kan gerne udstede Forbud mod „at
lege med Billedsproget“, men det Spørgsmaal lader sig nu engang ikke afvise,
hvilken virkelig Værdi Billedsproget har for det aandelige Liv, og hvorfor
netop dette bestemte Billedsprog bruges, og ikke et andet%
% footnote *) on p. 194
\footnote[1]{Smlgn. nærmere om Forholdet mellem Tanke og Billede, særligt med
Hensyn til det religiøse Problem, min \textit{Religionsfilosofi} II C.}.

Det var intet Tilfælde, at Nielsen i sine senere Aar følte megen Sympati med
den Grundtvigske Retning, der netop lagde stor Vægt paa Billedsproget og dets
Tydning, men samtidigt lagde liden Vægt paa videnskabelig Dogmatik. Her kunde
han fastholde sin Lære om Tro og Viden som absolut uensartede Principer og sin
Protest imod videnskabelig Teologi, medens dog en fri aandelig Udlægning maatte
tilstedes. Sympatien var dog ikke ganske gensidig, da man fra Grundtvigsk Side
var noget betænkelig ved „Aandsvirkeligheden“, der ikke med Urette mindede om
mytisk Symbolisme, saasnart man skulde tænke noget Bestemt ved den. Nielsen
„legede med Billedsproget“ paa en mistænkelig Maade.

Men det egentlige Motiv til Nielsens Tilnærmelse til Grundtvigianismen laa dog
vistnok paa det etiske Omraade. For Kierkegaard var, som vi have set, den
kristelige Etik Hovedsagen. Naar han erklærede, at det nye Testamentes
Kristendom ikke var til, tænkte han først og fremmest paa den nytestamentlige
Etik. Over Naturvidenskabens Fremtrængen trøstede han sig jo med, at den
kristelige Etik maatte blive den samme, hvad enten man i Naturvidenskaben kom
til det ene eller det andet Resultat. Det var Modsætningen mellem det Religiøse
og det Humane, han vilde hævde. Netop denne Modsætning søgte Nielsen ud over,
ligesom Grundtvigianismen. Og han skiller sig definitivt fra Kierkegaard ved i
sin Bog \textit{Om Aandsdannelse} (1877) at skelne mellem Liv og
% --- p. 195 ---
\setcounter{footnote}{0}%
\opage{195}Existensform. „Liv og Existens er ikke enstydige Ord. Det samme Liv
kan give sig Udslag i forskellige Existenser.... Det evige Liv i Kristendommens
Væsen bliver ikke staaende ved en enkelt Tidsalders stereotype
Existensskema.... Apostelexistensen hørte op, da Gennembrudsperioden var
afsluttet.“ Det religiøse Liv kan, mener Nielsen, være det samme som i
Apostlenes Dage, idet der kan leves i samme Tro og Haab; men man kan ikke kaste
Alt bort, hvad det lange Kulturarbejde har tilvejebragt, og i sin Livsførelse
uden videre efterligne den Maade at tage Livet paa, som var ejendommelig for
den store Epoke, da Kristendommen traadte ind i Verden. Kierkegaard havde Ret
overfor sine kirkelige Modstandere, fordi disse selv højtideligt anerkendte, at
det nytestamentlige Existensskema stadigt gjaldt. Han havde ogsaa Ret i, at
Idealerne skulle forkyndes. Men et andet Spørgsmaal er, hvorledes de skulle
realiseres og forme sig til Existens. Han har ikke forstaaet Nutidsexistensen
og dens aandelige Vilkaar. Det gælder om at leve med i den moderne Kultur og
lade de store kristelige Idealer gennemtrænge den. --- Dette var nu netop den
Opfattelse, mod hvilken Kierkegaard i „Øjeblikkene“ havde rettet det voldsomste
Angreb, tømt hele sit Kogger af Spot og Harme. Og det var den Opfattelse, som
Nielsens tidligere Fælle og senere Fjende Aaret efter i anden Del af sin Etik
udviklede langt rigere, end Nielsen formaaede det. Mere og mere maatte det da
blive klart, at der trængtes til en Undersøgelse af Etikens Forudsætninger og
Resultater. Der var jo den Mulighed, at Nutidsexistensen, og den menneskelige
Existens idethele, ud af sine egne Kræfter og Livsvilkaar kunde vinde et
Grundlag for en etisk Livsførelse. Denne Mulighed havde Kierkegaard peget paa i
sine Stadier (det andet, tredie og fjerde), men ladet falde i sin Iver for at
naa til det „Existensskema“, det var hans Mission at opstille i dets
oprindelige Klarhed overfor alskens Tillempelser, der ikke vilde være ved, at
de vare Tillempelser. At Nielsen lod denne Mulighed ligge, var en Hovedgrund
til, at hans Efterfølgere maatte slaa ind paa helt andre Veje end hans.

\begin{center}\rule{2cm}{0.4pt}\end{center}

% --- p. 196 ---
\setcounter{footnote}{0}%
\opage{196}

\begin{center}\large 17. HANS BRØCHNER\end{center}

Efter Planen for dette Arbejde skal det slutte med de Filosofer, der have været
mine Lærere. Saa langt tror jeg, det vil være muligt at gennemføre en rent
historisk Opfattelse; mit eget personlige Standpunkt gør sig naturligvis
gældende i hele Fremstillingen, baade ved det i Tiden Nærmere og ved det
Fjernere, men behøver ikke at træde mere frem ved hint end ved dette.

Da jeg kom til Universitetet, stod Sibbern trods sin Stivnen i sære udvortes
Former dog som saa ærværdig og interessant ved, hvad jeg vidste om ham, at jeg
ikke kunde tænke mig Andet end, at jeg maatte høre ham. Det var dog først langt
senere, at hans Værd og Betydning gik op for mig, saa at jeg fra en væsentlig
Side ønskede og stræbte at fortsætte hans Værk. Foreløbigt var det Rasmus
Nielsen, hvis begejstrede Henvisning til Kierkegaard og hvis vækkende
Veltalenhed fik den største Indflydelse paa mig. Men jeg hørte dog i mine
Studenteraar stadigt Brøchners Forelæsninger over Filosofiens Historie, og
deltog ogsaa i Øvelser hos ham. Senere skulde han blive den af mine Lærere, jeg
kom i det nærmeste personlige Forhold til. Om Ingen, jeg har kendt, gælder det
som om ham, at han havde og søgte sin Styrke i sin Begrænsning. Han var et
Villiesmenneske fra først til sidst, saa meget han end færdedes i Tankens
abstrakte Verden (og endog gjorde denne mere abstrakt, end den behøvede at
være). I et Brev til sin Ven Chr. Molbech (24. Decbr. 1862) skriver han: „Jeg
bestræber mig redeligt for at lære min Styrke og min Begrænsning at kende for
at kunne stille
% --- p. 197 ---
\setcounter{footnote}{0}%
\opage{197}mit Maal saaledes, at jeg kan nærme mig til det; ti \textit{naa} et
selvstillet Maal skal man ikke. Og skulde jeg tage fejl af Maalet og Vejene, Et
beholder jeg saa dog, hvori man ikke kan tage fejl: det at have villet, og den
Idealitet, der ligger deri.“ Disse Ord fremkom i Anledning af et paatænkt
filosofisk Arbejde, men de ere karakteristiske for ham baade som Menneske og
som Tænker. Han følte sin Begrænsning som det Middel, der stod til hans
Raadighed, naar han skulde give sit Bidrag til den aandelige Udvikling og
derigennem udvikle sin egen Personlighed. Der var en Resignationens Stemning
over ham, næret ved Sygdom og Sorg, saavel som ved Bevidstheden om, hvor Lidet
han naaede af, hvad der stod for hans Tanke, og hvor ufuldkomment et Udtryk han
havde kunnet give det, han havde naat. Men hans kraftige Villie hindrede
Resignationen i at faa Bitterhedens eller Kuldens Klangfarve. Som den Tankens
Ridder han var, nedlagde han ikke Vaabnene saalænge det endnu var ham muligt at
kæmpe%
% footnote *) on p. 197
\footnote[1]{I Brøchners Breve til Chr. Molbech, og i den Maade, paa hvilken
hans Ven i sine Breve kommer til at skildre ham, haves et smukt Billede af ham,
bedre end nogen Biografi kunde give. I Indledningen til Udgaven af disse Breve
har jeg givet en kort Karakteristik af min Lærer og Ven. \textit{(Hans Brøchner
og Chr. K. F. Molbech.} En Brevvexling. 1902).}.

Brøchner (f. 1820, d. 1875) var i sin Ungdom, ligesom Nielsen og Kierkegaard,
stærkt paavirket af Hegel, og denne Paavirkning tabte sig aldrig ganske, heller
ikke i hans Sprog, der bevarede sin abstrakte Karakter, uden, som saa ofte hos
Hegel, at være gennemvævet med slaaende Billeder og mægtig Patos. Han forlod
Hegels Skole i en anden Retning end de andre Danske (Poul Møller, Martensen,
Kierkegaard og Nielsen), der havde tilhørt Skolen og derpaa banede sig egne
Veje. Han gav Strauss og Feuerbach Ret i deres Kritik af Hegels Forsøg paa at
løse Problemet om Tro og Viden. Strauss' historisk-kritiske Fremstilling af den
kristelige Troslære har han oversat paa Dansk. Den Vej, ad hvilken han søgte
videre, blev for en væsentlig Del bestemt ved Kierkegaards person%
% --- p. 198 ---
\setcounter{footnote}{0}%
\opage{198}lige Paavirkning. Det er af Interesse baade til Forstaaelse af
Brøchner og af Kierkegaard at se følgende Udtalelse af Brøchner faa Uger efter
Kierkegaards Død: „Paa mig har hans Død vel gjort et smerteligt, men ikke noget
nedbøjende Indtryk. Han har været meget for mig, baade ved sine Skrifter og ved
det personlige Forhold, i hvilket jeg har staaet til ham. Der er Ingen, hvis
Personlighed i den Grad har indvirket vækkende og tilskyndende paa mig, og den
venlige Velvillie, han altid viste mig, gav mig tit Mod, naar jeg var ved at
miste det. Det er med Glæde at jeg erindrer ham i den lange Række af Aar, jeg
har kendt ham, næsten 20 Aar, hvorledes det Milde og Kærlige i hans
Personlighed bestandigt mere fik Overvægten over det stærke ironiske og
polemiske Element, der af Naturen var i ham, og hvorledes hans Tanke bestandigt
blev rigere, sikrere og klarere, saa at tit et Ord af ham kunde virke
beroligende og forsonende, og opklare en Forvirring, som man med egne Kræfter
ikke kunde faa Bugt med. Jeg vil nok komme til at savne ham. Men naar jeg
tænker paa, hvor fuldstændigt han har faaet det udrettet, der var hans Livs
Opgave, --- hvor rigt og fyldigt dette Liv har været i dets korte Varighed, og
hvor Meget af ham der bliver tilbage, kan jeg ikke med nogen trykkende Følelse
tænke paa hans Død; tvertimod forekommer den mig smuk og lykkelig.“ (Brev fra
Brøchner til Molbech 2. Decbr. 1855)%
% footnote *) on p. 198
\footnote[1]{Mellem Brøchners Papirer fandtes en lille stemningsfuld Opsats:
„Erindringer om Søren Kierkegaard“, som jeg udgav i Tidsskriftet „Det nittende
Aarhundrede“ (1876).}. --- For at forstaa dette Forhold mellem Kierkegaard og
Brøchner trods deres saa højt forskellige Anskuelser maa man erindre, at for
Kierkegaard var personlig Sandhed og Oprigtighed det Første og det Afgørende.
Hos sin yngre Ven har han fundet den Stræben efter Sandhed, Lessing havde
erklæret for det Højeste, og for Kierkegaard fik den dogmatiske Formulering af
den fundne Sandhed stedse mindre og mindre Betydning. Det har været en rent
sokratisk Indflydelse, der her er øvet af den Ældre paa den Yngre. Men det har
naturligvis ogsaa været af stor Be%
% --- p. 199 ---
\setcounter{footnote}{0}%
\opage{199}tydning, at de begge stode i polemisk Forhold til Teologien og den
bestaaende Kirke. Som vi paa et tidligere Sted have set, hævdede Kierkegaard,
at „den sidste Formation af Fritænkere (Feuerbach og hvad dertil hører)“ i
Grunden forsvarede Kristendommen mod de moderne Kristne. --- Forholdet mellem
Kierkegaard og Brøchner har været af væsentligt anden Art end Forholdet mellem
Kierkegaard og Nielsen.

Det var først sent hen i sit Liv, at Brøchner skred til en Fremstilling af sin
filosofiske Anskuelse i dens Helhed. Alt imedens hans Livsanskuelse formedes
praktisk og i Stilhed under hans Livs Tilskikkelser fordybede han sig i Studiet
af Filosofiens Historie, som han egentligt har grundlagt herhjemme. I en
grundig Monografi over \textit{Spinoza} (1857) gav han en Fremstilling af det
filosofiske Standpunkt, han følte sig nærmest beslægtet med. I sit Skrift
\textit{Om Udviklingsgangen i Filosofiens Historie} (1869) gav han en stærkt
sammentrængt Fremstilling af de Grundtanker, der have raadet i Filosofien, af
deres Forhold til andre Sider af Kulturudviklingen og af den historiske
Sammenhæng, der bevidst eller ubevidst gør sig gældende mellem dem. Sin
Afslutning fandt dette historiske Studium i en Haandbog i \textit{Filosofiens
Historie} (1872---74), den eneste danske Fremstilling, der baade omfatter den
græske og den moderne Filosofi.

Hans vigtigste Skrifter ere de to Bøger \textit{Problemet om Tro og Viden
(1868),} der giver en Kritik af Martensen, Kierkegaard og Nielsen, og
\textit{Det Religiøse i dets Enhed med det Humane} (1869), der udvikler hans
egen religionsfilosofiske Anskuelse.

\begin{center}\rule{2cm}{0.4pt}\end{center}

Brøchners Udgangspunkt er psykologisk, og er beslægtet med det, der som
Maalestok ligger til Grund ved Kierkegaards Lære om Stadierne: Menneskeaanden
som Energi og Totalitet. Den Livsanskuelse, der skal kunne bestaa, maa være en
saadan, der muliggør en Udfoldelse af den menneskelige Natur, saaledes at de
forskellige Elementer i denne
% --- p. 200 ---
\setcounter{footnote}{0}%
\opage{200}kunne virke harmonisk sammen. I Begrebet Selv ligger baade Energi
(Virksomhed) og Helhed. I Selvbevidstheden kundgøre begge sig i Trangen til at
forme og harmonisere Tilskyndelser og Drifter, saavelsom til at ordne alle
Tanker i indre Sammenhæng og Helhed, den Trang, som er Menneskets Væsenstrang.
Mennesket adskiller sig fra Dyret ved denne Trang til at gaa ud over de enkelte
Tilskyndelser og de enkelte Sansninger og naa det Almengyldige, det Universelle.

Brøchner bruger Udtrykket Væsenstotalitet for at betegne sit Hovedsynspunkt. En
saadan Totalitet ligger efter hans Opfattelse bagved alle sjælelige
Livsytringer, selvom den ikke direkte og positivt kommer til Bevidsthed. Herved
faar hans Psykologi en dogmatisk Karakter. Væsenstotaliteten grundes ikke paa
Analyse af Erfaringen, men opstilles paa Forhaand som utvivlsom. Overfor de
Tænkere, med hvilke han nærmest forhandlede, kunde han ogsaa meget vel antage
Forudsætninger om den menneskelige Aands Væsen som Totalitet given, da de hver
paa sin Maade anerkendte den; men det vilde have tjent til Klarhed og
bestemtere Afgørelse, naar der var bleven givet en nærmere Begrundelse. Og
afset fra Mangelen af tilstrækkelig erfaringsmæssig Begrundelse er der ogsaa
den Betænkelighed ved Brøchners Princip, at han for meget bruger det til at
forklare enkelte psykologiske Fænomener. Saaledes lærer han, at Driften „med
indre Nødvendighed“ udvikler sig fra at være Villiens rudimentære Form til at
blive en sand Villie. Nu viser dog Erfaring, at ikke alle Drifter udvikle sig
saaledes, og at naar det sker, er det, fordi andre psykologiske Fænomener,
særligt andre Drifter optræde, og det er gennem Vexelvirkninger og Brydninger
med dem, at Udviklingen i den antydede Retning kan ske. En hel Mængde
psykologiske Forhold --- Kontrastvirkninger, Associationer og Forskydninger ---
kunne herved gøre sig gældende. For Brøchners spekulative Psykologi udspringer
derimod hin „indre Nødvendighed“ af „Væsenstotaliteten“, idet det under
Forudsætning af denne bliver selvmodsigende at stanse ved en enkelt Tilskyndelse i
% --- p. 201 ---
\setcounter{footnote}{0}%
\opage{201}dens Isolation. --- Paa tilsvarende Maade skal det Nye, der ved
Beslutningen kommer til den forudgaaende ufuldendte Erkendelse (Overvejelse),
være en fra hin Væsenstotalitet som „umiddelbar Kilde“ stammende Supplering%
% footnote *) on p. 201
\footnote[1]{Det Vexelforhold, der her viser sig mellem det Bevidste og det
Ubevidste, har Brøchner nærmere drøftet i en Afhandling i „Nyt dansk
Maanedsskrift“ (1871).}. I al praktisk Sans kundgør det Samme sig, f. Ex. naar
der handles i Modsætning til ensidige Teorier. I sympatiske Tilskyndelser ytrer
sig en Helhedstrang, der endnu ikke er kommen til fuld etisk Klarhed. Den sande
Handlen er den, i hvilken Mennesket udvikler og bekræfter sig selv gennem
Hengivelsen til Handlingens Genstand. Og dette er kun muligt, naar denne
Genstand ikke er noget Tilfældigt, Endeligt og Forgængeligt; det er en
Modsigelse at nære en uendelig Interesse for Sligt. Med Endeligheden i
Formaalet falder al Egoisme bort. Paa lignende Maade er Erkendelsen først
fuldendt, naar dens Indhold for den selv ikke længere staar blot som givet og
modtaget, men formes til en Enhed, der svarer til den Enhed, som kundgør sig i
Selvbevidstheden. Som en Selvhedsenergi kundgør Menneskets Væsen sig baade,
hvor den virker i Erkendelsens, og hvor den virker i Villiens Form. Af disse to
Former maa Erkendelsen stedse beholde Principatet, da de specielle
Aandsvirksomheder stedse ere ledsagede af Bevidsthed og kun derved ere
Aandsvirksomheder. Iøvrigt er der et stadigt Vexelforhold mellem Erkendelse og
Villen: der gives en Erkendelsens Lidenskab, og der gives en Villen af det
erkendte Almengyldige. Hvor der i Teori eller i Praxis stanses ved et Endeligt
og Begrænset med Tilbøjelighed til at gøre det til et Absolut, vil den indre
Væsenstotalitet tidligere eller senere føre ud over Begrænsningen. Selv i det
Onde, d. v. s. den Villen, der benytter al sin Energi til at gaa op i et
Endeligt, er der en uholdbar Modsætning mellem Villiens Form og dens Indhold.
Den Kraftudvikling og Øvelse, her anvendes, gaar ikke tabt, men træder gennem
Angeren i det Godes Tjeneste. Under Kraftens Forvandling
% --- p. 202 ---
\setcounter{footnote}{0}%
\opage{202}er her, ligesom paa det fysiske Omraade, Kraftsummen bestaaende. Og
formedelst den indre Modsigelse, der er i den onde Villie og ytrer sig i
hvileløs Uro og Stræben, enten som Omtumlen i Nydelse eller trodsig Kamp mod
det Gode, maa det Onde mere og mere udhule sig selv, saa at en Retningsændring
bliver den eneste Udvej.

„Væsenstotaliteten“ eller Menneskets Idé stræber under alle disse Former efter
at udfoldes, saa haard Kampen end kan blive. Det er Brøchners Grundtanke, paa
væsentlige Punkter en under Kierkegaards Indflydelse foregaaet Uddybelse af
Hegels Lære, ved hvilken denne førtes tilbage til de psykologiske og etiske
Synspunkter, fra hvilke den for en stor Del oprindeligt er udsprungen. Det er
formedelst denne Grundtanke, at han, skønt han ikke tror paa overnaturlige
Kræfters Indgriben, fastholder Troen paa etisk Udvikling og Fremskriden.

\begin{center}\rule{2cm}{0.4pt}\end{center}

Men naar Brøchner antager „en indre Nødvendighed“ i Udviklingen og antager en
uendelig Fylde af Muligheder, der kunne realiseres gennem denne Udvikling, er
han klar over, at han herved ikke kan støtte sig til streng Bevisførelse.
Stiller man Viden og Tro i Modsætning til hinanden, er han ikke i Tvivl om, at
han vil vælge Viden, fordi det er ved Erkendelsen, ved Tanken, at Menneskets
Væsen ene kan udfolde sig saaledes, at det fortjener Navn af Aand. Men kun hvis
man ved Tro forstaar Autoritetstro eller Tro paa overnaturlige Kræfter og
Aabenbaringer, existerer hin Modsætning for Brøchner.

Paa intet Udviklingstrin kan et endeligt Væsen fuldt ud realisere det uendelige
Indhold, den Væsenstotalitet, der er dets inderste Grund. Ofte maa det
foregribe, hvad der først senere kan godtgøres, og ofte maa det stanse ved
Sætninger, der opstilles, uden at deres Forudsætninger kunne drages frem; disse
Forudsætninger ligge skjulte i dets inderste Væsen. Erfaringen er stedse
ufuldstændig, og derfor kan Erken%
% --- p. 203 ---
\setcounter{footnote}{0}%
\opage{203}delsen kun afsluttes ved et Postulat, en Tro. Men da denne Tro er
Videns Selvsupplering i Kraft af det samme Princip, ifølge hvilket der
overhovedet existerer en Viden, bliver der her ingen nødvendig Strid mellem Tro
og Viden. Det Princip, ifølge hvilket der existerer en Viden, er et Enheds- og
Totalitetsprincip. Troen gaar ud paa, at et saadant Princip har Gyldighed i den
hele Tilværelse: „Den ikke af Erfaringen udtømte Tilværelses Enhedsprincip
gribes ikke som det, hvis Vished Erfaringen betrygger, men som det, der naas
ved et nødvendigt Postulat ud fra den individuelle Bevidsthed, som uden det
vilde tabe Tilholdet i sig selv, blive usikker paa sin egen Enhed, forvirres i
sig selv.“ Det er en rationel Tro, idet den bevarer Overensstemmelsen med den
rationelle Erkendelse. Der er en Hengivelse i Troen; men heller ikke i
Erkendelsen mangler dette Element: „I en sand Erkendelse er der en Tankens
Lutring, en Opgiven af individuel Hildethed, Menen og Fordom, med ét Ord, af
Tankens Selviskhed; der er i den en Selvhengivelse til Sandhedens absolute
Magt.“ I Troen træder dette Element kun stærkere frem.

Det er ad denne Vej, man naar til Begrebet om det Absolute, om „det, der hviler
i sig selv og har Vished i sig selv,“%
% footnote *) on p. 203
\footnote[1]{Det er Spinozas Definition af den absolute Substans, der her
ligger til Grund: „Ved Substans forstaar jeg det, der bestaar i sig selv og
forstaas ved sig selv, d. v. s. hvis Begreb ikke forudsætter nogen anden Tings
Begreb, af hvilket det skulde dannes.“ \textit{Spinoza:} Ethica. I. Def. 3.}
Menneskets højeste og sidste Tanke. Denne Tanke nærme vi os til ved Analyse af
Tilværelsen, saaledes som Erfaringen viser os den, idet Vexelvirkning og
Aarsagssammenhæng omfatter Alt, hvad vi kende, og forudsætter en Energi, der
rører sig i Alt, en uendelig Totalitet, indenfor hvilken de endelige
Totaliteter, Erfaringen udviser, finde deres Forklaring. En saadan Analyse
fører kun indirekte hen imod den sidste Forudsætning. Direkte nærme vi os den,
naar vi gaa ud fra de højeste Tilværelsesformer, fra Livet, der er uforklarligt
ud fra det blot Materielle, men især fra det Sjælelige
% --- p. 204 ---
\setcounter{footnote}{0}%
\opage{204}og det Aandelige. Disse Fænomener gøre det nødvendigt at antage en
absolut Subjektivitet (Aand), af hvilken de i Erfaringen fremtrædende
Subjektiviteter kun ere begrænsede Former. --- Brøchner foretrækker det mere
abstrakte Udtryk Subjektivitet (Aand) for det konkretere Udtryk Personlighed,
da dette Udtryk kun kan bruges om et endeligt Væsen, der staar i Forhold til
andre Væsener. Han er Panteist, ikke Teist. „Naar vi da ikke begrænse
Panteismens Begreb til de ufuldkomne og ensidige Former af den, saa maa vi
erkende den panteistiske Opfattelse som væsentlig og nødvendig for
Filosofien... Ingen Filosofi kan tænke det Uendelige som værende adskilt fra og
udenfor det Endelige... Al Filosofi søger Erkendelsens Enhed og Sammenhæng og
er kun Filosofi, forsaavidt den forfølger denne... Enhver højere Tanke, til
hvilken den stiger op, er netop kun den højere derved, at den omfatter de
lavere.“ --- Ligesaalidt som andre Hegelianere vil Brøchner indrømme, at han i
sine afsluttende Tanker støtter sig til en Analogi. Analogien er i
Virkeligheden alle metafysiske og religiøse Forestillingers Kilde.%
% footnote *) on p. 204
\footnote[1]{Smlgn. min \textit{Religionsfilosofi II} C og mine \textit{Mindre
Arbejder.} Anden Række. p. 41---46.} Det laa egentligt Brøchner nær at gøre
denne Indrømmelse, saasnart han skilte sig fra Hegel ved den relative Forskel
mellem Viden og Tro, en Forskel, Hegel ikke vilde anerkende i denne Form. Men
Indrømmelsen vilde have medført Anerkendelsen af en større Medvirken af andre
sjælelige Elementer end de intellektuelle ved hine Forestillingers Opstaaen og
Udvikling.

Overfor de positive Religioner er det Brøchners væsentligste Indvending, at de
stanse ved endelige Forestillinger. Det Guddommelige fremtræder i Form af
endelige Personligheder, og derfor med sanselige Egenskaber, det er begrænset i
Tid og Rum, og det virker (i Underet) udefra paa Naturen, uden Lov og indre
Sammenhæng. Dermed hænger det Egoistiske i de positive Religioner sammen;
Mennesket forholder sig i dem som et endeligt Væsen til et andet, et Jeg til et
Du. Egentlig Selvhengivelse uden Egoisme er kun mulig overfor det Uendelige, og
kun en saadan Selvhengivelse er Et med
% --- p. 205 ---
\setcounter{footnote}{0}%
\opage{205}den sande Selvbekræftelse, idet der udløses ny Kraft gennem dette
Forhold, i hvilket „Væsenets Kilder springe frem fra det Uendeliges Dyb.“ ---
Stedse ligger hos Brøchner den Forudsætning til Grund, at der er et uendeligt
Indhold, et Kraftfond, der skal og kan udfoldes eller udløses, og at den
inderligste Selvbesindelse for Menneskets Vedkommende er Midlet til at faa Del
deri. Hans Tankegang faar her Karakteren af en Mystik, og han vilde sikkert
ikke selv helt fornegte dette Slægtskab. I den inderlige Selvbesindelse opdage
vi, at vi ere Led af en uendelig Sammenhæng og opdage en Lov, der omfatter
saavel os selv som hele den øvrige Tilværelse.---

Overfor Feuerbach hævder Brøchner, at der i alle Religioner, særligt de højere,
findes rationelle Elementer, og at man ikke blot maa hefte sit Blik paa de
enkelte Religioner hver for sig, men se paa Udviklingsgangen i Religionens
Historie, fra Naturreligionens laveste Former til Aandsreligionernes højeste.
Det er en Kendsgerning, at det har været muligt for Menneskeaanden at arbejde
sig op fra lavere til højere Former af Religion. Og dette er ikke sket ved, at
det egentligt Religiøse er forsvundet, men netop ved, at det er traadt mere
frem, --- naar man ved Religion forstaar Menneskets absolute Forhold til det
Absolute som til et Princip, der begrunder og sammenfatter det Menneskelige
saavelsom hele den øvrige Tilværelse. Det er et Enhedsforhold, ikke et
Modsætningsforhold. Det Religiøse ytrer sig i enhver positiv Form for aandeligt
Liv. --- I en saadan Religion, Humanitetens Religion, er det Væsentlige i de
positive Religionsformer bevaret, medens de historiske Former opløses. Det er
Virkelighedens sande og evige Religion, der er tilstede, skønt ofte uklart og
ubevidst, i vor Tids humane Bevidsthed. Det er ikke Metafysik, der sættes i
Religionens Sted; de abstrakte Begreber skulle kun udrede, hvad der rører sig i
den humant-religiøse Livsvirkelighed.

Brøchner gør selv opmærksom paa, at det Religiøse efter hans Bestemmelse svarer
baade til Kierkegaards etiske Stadium og til det religiøse Stadium, der endnu
ikke er det kristelige. Han kunde her have taget det mellemliggende Stadi%
% --- p. 206 ---
\setcounter{footnote}{0}%
\opage{206}um, „Humoren“, med. Disse tre Stadier gaa for Brøchner sammen i Et;
navnligt vil han ikke stille det Religiøse i Modsætning til det Etiske, da det
for ham netop er Sjælen i dette. --- Det hænger sammen med hans stærkt
koncentrerede Fremstilling, at han ikke finder Anledning til at fastholde
Kierkegaards Distinktioner, eller dog selv opstille andre. Ogsaa her mangler
der en psykologisk Underbygning, og Forsøget paa at give en saadan vilde
sikkert have ført til at fremdrage Vanskeligheder og Spørgsmaal, der nu ikke
ret kunde komme frem.

Ved sit abstrakte Sprog har Brøchner afskaaret sig fra at blive tilgængelig for
større Kredse; men han har den Fortjeneste, at han paa sin Maade har hævdet
Sammenhængen i det personlige Liv, der forflygtigedes af den spekulative
Filosofi og sprængtes af Kierkegaard. Og i Modsætning til ethvert mere eller
mindre materialistisk Forsøg paa udelukkende at bygge den menneskelige
Livsanskuelse paa, hvad Erfaring lærer om den materielle Natur, har han hævdet
Tankens Realitet som den, hvis Kald det er at tyde alle Former for Tilværelse,
den materielle saavel som den aandelige. ---

Paa en Forelæsning, Brøchner holdt i det Aar, der skulde blive hans Dødsaar,
udtalte han, at han var en Tilhænger af „den idealistiske Filosofi, Tyskland
har Æren af at have frembragt.“ Vi have set, hvorledes denne filosofiske
Retning hos ham selv, saavelsom hos de andre Tænkere, der havde været
underkastet en lignende Paavirkning, brødes med den danske Interesse for
Psykologi, Etik og Kritik. Dersom i den filosofiske Generation, der fulgte
efter Brøchner og hans Samtidige, disse Interesser ere komne stærkere frem end
før, saa skyldes det ganske vist tildels den Paavirkning fra engelsk og fransk
Tænkning, der gjorde sig gældende herhjemme i den sidste Del af det nittende
Aarhundrede; men det tør ogsaa opfattes som en Fortsættelse af den særligt
danske Retning hos Forgængerne --- ligefra Holbergs Tid af.

\begin{center}\rule{2cm}{0.4pt}\end{center}

\end{document}
